\documentclass[11pt,a4paper]{article}
\usepackage[T1]{fontenc}
\usepackage{lmodern,amsmath,amssymb,amsthm,mathtools}
\usepackage[a4paper,margin=25mm]{geometry}
\usepackage{enumitem,microtype}
\usepackage[colorlinks=true,linkcolor=blue,urlcolor=blue,citecolor=blue]{hyperref}
\setlength{\emergencystretch}{3em}
\numberwithin{equation}{section}
\newtheorem{theorem}{Theorem}[section]
\newtheorem{lemma}[theorem]{Lemma}
\newtheorem{proposition}[theorem]{Proposition}
\newtheorem{corollary}[theorem]{Corollary}
\theoremstyle{definition}
\newtheorem{definition}[theorem]{Definition}
\newtheorem{convention}[theorem]{Convention}
\theoremstyle{remark}
\newtheorem{remark}[theorem]{Remark}
\newcommand{\D}{D}
\newcommand{\Dbc}{D^b_c}
\newcommand{\RHom}{R\!\operatorname{Hom}}
\newcommand{\cHom}{R\mathcal{H}\!om}
\newcommand{\Sh}{\operatorname{Sh}}
\newcommand{\Mod}{\operatorname{Mod}}
\newcommand{\Loc}{\operatorname{Loc}}
\newcommand{\Qbar}{\overline{\mathbf Q}_\ell}
\newcommand{\Ql}{\mathbf Q_\ell}
\newcommand{\Zl}{\mathbf Z_\ell}
\newcommand{\Fell}{\mathbf F_\ell}
\newcommand{\et}{\acute{e}t}
\newcommand{\proet}{\mathrm{pro}\acute{e}t}
\newcommand{\coker}{\operatorname{coker}}
\newcommand{\fib}{\operatorname{fib}}
\newcommand{\hocolim}{\operatorname{hocolim}}
\newcommand{\holim}{\operatorname{holim}}
\newcommand{\stacks}[1]{\href{https://stacks.math.columbia.edu/tag/#1}{\textnormal{Stacks, Tag~#1}}}

\newcommand{\Ab}{\operatorname{Ab}}
\newcommand{\Hom}{\operatorname{Hom}}
\newcommand{\Dc}{D_{\mathrm{cons}}}
\newcommand{\cO}{\mathcal O}
\newcommand{\Qlbar}{\overline{\mathbf Q}_\ell}
\newcommand{\lotimes}{\otimes^{\mathrm L}}
\newcommand{\Res}{\operatorname{Res}}
\newcommand{\Ext}{\operatorname{Ext}}
\newcommand{\Coext}{\operatorname{Coext}}
\newcommand{\Cons}{\operatorname{Cons}}
\newcommand{\IndCons}{\operatorname{IndCons}}
\newcommand{\Dic}{D_{\IndCons}}
\newcommand{\QbarX}{\overline{\mathbf Q}_{\ell,X}}
\usepackage{amscd}
\newcommand{\Z}{\mathbf Z}
\newcommand{\Q}{\mathbf Q}
\newcommand{\A}{\mathbf A}
\newcommand{\PP}{\mathbf P}
\newcommand{\Dctf}{D^b_{\mathrm{ctf}}}
\newcommand{\dt}{\mathbin{\otimes^{\mathbf L}}}
\DeclareMathOperator{\GL}{GL}
\newcommand{\Ob}{\mathcal O}
\newcommand{\Qb}{\overline{\mathbf Q}_{\ell}}
\newtheorem{example}[theorem]{Example}
\title{Condensed Coefficients, Proper Base Change\\
and Six Functors in Ordinary Derived Categories}
\author{Weil conjectures project}
\date{14 October 2026}
\begin{document}
\maketitle
\begin{abstract}
These notes set up the $\ell$-adic sheaf theory used in the proof of the
Weil conjectures. Let $k$ be a finite or algebraically closed field and
$\ell$ a prime invertible in $k$. For schemes of finite type over $k$ we
work with complexes of sheaves on the pro-\'etale site, with
coefficients in a finite commutative $\ell$-primary ring $A$, or in
$\Zl$, $\mathcal O_E$, $\Ql$, $E$ or $\overline{\mathbf Q}_\ell$, where
$E$ is a finite extension of $\mathbf Q_\ell$; taking $E=\Ql$ gives
$\Zl$ and $\Ql$. All these coefficient sheaves come from one
definition: a condensed ring $\Lambda$ on the point gives
$\Lambda_X=p_X^*\Lambda$ on every scheme $X$. Compatibility with
inverse image is then automatic, and the familiar completed descriptions
become comparison theorems. We first develop the pro-\'etale site,
its contractible local structure, derived inverse limits, and the
comparison with the \'etale site that permits torsion calculations.
We prove torsion proper base change through the projective line,
Gabber's affine henselian comparison and Chow's lemma, and pass from it
to the condensed coefficient rings. The main result is the existence of the
exceptional inverse image: for every separated morphism of finite type
$f$, the compactly supported direct image $Rf_!$ on bounded
constructible complexes of finite Tor amplitude has a right adjoint
$f^!$. Brown representability on the localizing
subcategory generated by these complexes constructs $f^!$ for all these
coefficients, and a geometric argument, on the projective line and on
closed immersions, shows that $f^!$ preserves bounded constructible
complexes of finite Tor amplitude. For integral and rational
coefficients the finite Tor condition is automatic. For smooth $f$ of relative dimension $d$ we prove smooth
purity, $f^!\simeq f^*(d)[2d]$, and construct the trace map.
Three appendices prove Nagata's compactification theorem and Chow's
lemma, construct the Picard scheme and the Jacobian of a curve, and
compute the cohomology of curves.
\end{abstract}
\tableofcontents

\section{Introduction}
\label{n:scope}
To formulate and prove the Weil conjectures one needs $\ell$-adic
cohomology with its full set of operations: for a morphism $f$ of
varieties over a finite field, the functors $f^*$, $Rf_*$, $Rf_!$ and
$f^!$, together with $\otimes^L$ and internal Hom, all acting on bounded
constructible complexes with coefficients in $\Qbar$. This chapter
constructs them on the pro-\'etale site of Bhatt and
Scholze~\cite{BS}. There the coefficients are sheaves of rings, and
constructible complexes form full subcategories of ordinary derived
categories of module sheaves.

The theory is uniform in the coefficient ring. With $A$ a finite
commutative ring killed by a power of $\ell$ and $E$ a finite extension
of $\Ql$ with ring of integers $\mathcal O_E$, it applies at once to
\[
 A,\quad\Zl,\quad\mathcal O_E,\quad\Ql,\quad E,\quad\Qbar.
\]
All of them enter the same localizing category and the same Brown
representability argument. The category on which the operations
act is the category of bounded constructible complexes of finite Tor
amplitude, written $\Dctf(X,\Lambda)$ and defined in
\S\ref{n:constructible}. For integral and rational coefficients it
equals the category $\Dbc(X,\Lambda)$ of all bounded constructible
complexes; for a finite ring such as $\mathbf Z/\ell^2$ the Tor
condition is needed.
Finite coefficients also serve as a tool: the comparison with finite
coefficients enters the proofs for every ring in the list.

The guiding example is Poincar\'e duality. Let $X$ be smooth of pure
dimension $d$ over an algebraically closed field, and let $\Lambda=E$ or
$\Qbar$. The trace should induce perfect pairings
\[
 H_c^i(X,\Lambda)\times H^{2d-i}(X,\Lambda(d))\longrightarrow\Lambda.
\]
Behind this pairing stands the exceptional inverse image of the
structure map $a\colon X\to\operatorname{Spec}k$: smooth purity
identifies $a^!\Lambda$ with $\Lambda_X(d)[2d]$ in $\Dctf(X,\Lambda)$. The twist $(d)$, the shift $[2d]$ and the trace
are what the construction has to produce.

Four of the operations come for free: they are derived functors of
sheaves of modules, and the only question is whether they preserve
constructibility. The exceptional inverse image $f^!$ is different. It
is defined as a right adjoint of $Rf_!$, and nothing guarantees a priori
that such an adjoint exists. Its existence, and the fact that it
preserves constructible complexes, is the main result of the chapter
(Theorem~\ref{n:existence}). The proof has two halves.
\begin{enumerate}[label=(\Alph*)]
\item Existence by Brown representability. Constructible complexes of finite
  Tor amplitude do not have arbitrary direct sums, so we pass to the
  smallest triangulated subcategory containing them and closed under
  direct sums, written $\Loc(\Dctf(-,\Lambda))$. For $f\colon X\to Y$,
  the functor
  $Rf_!\colon\Loc(\Dctf(X,\Lambda))\to\Loc(\Dctf(Y,\Lambda))$ preserves
  direct sums, and Brown representability produces a right adjoint
  $f^!\colon\Loc(\Dctf(Y,\Lambda))\to\Loc(\Dctf(X,\Lambda))$. The
  work lies in proving compactness and preservation of direct sums,
  which uses finite cohomological dimension; Brown representability
  then applies formally.
\item Preservation of $\Dctf$. Locally on
  source and target, every morphism is a closed immersion followed by a
  projection from affine space. For the projective line we write down
  the adjunction explicitly, and for closed immersions we reduce to
  finiteness of direct image along an open immersion. Together these
  show that $f^!$ carries $\Dctf(Y,\Lambda)$ into $\Dctf(X,\Lambda)$,
  that is, preserves both constructibility and finite Tor
  amplitude. For integral and rational coefficients this is precisely
  preservation of $\Dbc$.
\end{enumerate}
The last step is smooth purity, a canonical normalization: for a smooth
morphism $f\colon X\to Y$ of relative dimension $d$, $f^!$ is $f^*$
followed by the Tate twist $(d)$ and the shift $[2d]$, as functors
$\Dctf(Y,\Lambda)\to\Dctf(X,\Lambda)$. For a smooth curve $C$ over $k$
with structure map $a$, for instance, $a^!\Qbar\simeq\Qbar(1)[2]$ in
$\Dctf(C,\Qbar)$;
Poincar\'e duality and the trace map on $H^2_c$ come from this
isomorphism.

\subsection{Why this route}
The classical theory, in SGA~4 and SGA~5~\cite{SGA4,SGA5}, works with
torsion coefficients $\mathbf Z/\ell^n$ on the \'etale site, where all
six operations exist and are well understood. An $\ell$-adic sheaf is
then a compatible system $(F_n)$ of torsion sheaves, and
$\Qbar$-sheaves are obtained by inverting $\ell$ and passing to a
colimit over finite extensions. On the derived level, compatible
systems need a normalization condition and control of derived inverse
limits. One can form the derived category of all systems, but the adic
constructible theory has to be cut out of it by further work, as in
Ekedahl's construction~\cite{Ekedahl}.

The pro-\'etale site removes most of this. On $X_{\proet}$ the ring
$\Zl$ is a sheaf of rings, $\ell$-adic sheaves are
sheaves of modules, and their derived category is the derived
category of an abelian category, in which inverse limits are computed by $R\varprojlim$. The
operations $f^*$, $Rf_*$, $\otimes^L$ and internal Hom are the usual
derived functors of module sheaves, acting between the derived
categories of $\Lambda$-module sheaves on $X_{\proet}$ and on
$Y_{\proet}$ for $f\colon X\to Y$.

We also make the coefficients themselves uniform. Start with a condensed
ring $\Lambda$ on the point and put $\Lambda_X=p_X^*\Lambda$. The choice
of coefficients is then separate from the choice of scheme. The identity
$p_Y\circ f=p_X$ gives isomorphisms $f^{-1}\Lambda_Y\simeq\Lambda_X$ of
sheaves of rings on $X_{\proet}$, compatibly with composition, and a morphism of condensed rings gives
extension of scalars on every scheme at once. Finite, integral and
rational coefficient sheaves are instances of this one construction.
Their familiar descriptions, by completion, by inverting $\ell$ and by
a filtered colimit, become comparison theorems. These comparisons
connect the uniform definition to the classical finiteness theorems,
and they show exactly where a proof uses a property of the particular
coefficient ring.

The remaining operation $f^!$ accounts for most of the length of the
chapter. There are two natural ways to construct it. One can transport
the torsion $f^!$ of SGA~4, level by level, to $\ell$-adic systems;
this requires keeping track of compatibilities between infinitely many
levels, and the formalism of towers we use in
Section~\ref{sec:smooth-purity} shows how much bookkeeping that
involves. Or one can obtain $f^!$ abstractly, as a right adjoint given
by Brown representability, and then prove by geometry that it preserves
$\Dctf$. We take the second route: the adjoint argument then applies to
finite, integral and rational coefficients alike, and the geometry is
confined to the proof that $f^!$ preserves $\Dctf$. For smooth morphisms one more
thing is needed: a canonical form of the local formula. We construct it
with finite coefficients from the diagonal, carry the adjunction through
a category of coefficient towers, and identify the result with the
Brown adjoint by uniqueness of adjoints.

\subsection{What is essential and what is formal}
The essential inputs are few. On the coefficient side, every
constructible $\Qbar$-complex is defined over a finite extension $E$ of
$\Ql$, and every lisse $E$-sheaf contains a lattice after refining the
stratification (\S\ref{b:sec:coefficients}). On the categorical side,
constructible complexes of finite Tor amplitude are compact because of
continuity and uniform cohomological bounds for perfect local models on
strata; for integral and rational coefficients this can also be
verified through Ext groups between constructible sheaves
(\S\ref{b:sec:compactness}). On the geometric side, we need proper base
change (Section~\ref{sec:proper-base-change}), the explicit adjunction
for the projective line, built from the cycle class of the diagonal,
and Deligne's theorem that direct image along an open immersion
preserves constructibility (\S\ref{b:sec:finite-geometric}). Smooth
purity needs in addition the identification of the dualizing complex
with $\Lambda(d)[2d]$ through the deformation to the normal bundle of
the diagonal (\S\ref{sp:grp-finite}). The coefficient comparisons, and
the uniform bounds that control unbounded complexes, complete the list.

Given these inputs, representing the adjoint by Brown representability,
composing adjunctions, and extending comparison maps from constructible
generators are formal. Change of coefficients is not: each step uses a
specific comparison. Integral coefficients use derived completeness and
normalized towers, $E$ uses localization, and $\Qbar$ uses descent to a
finite extension of $\Ql$. The towers live in the ordinary derived
category of an abelian category of inverse systems, where a compatible
family of finite-level maps is a single morphism.

On a first reading one can take the following route. Read
Section~\ref{sec:proetale-foundations} for the site, weakly contractible
covers, and the scope of classical comparison, then
Section~\ref{sec:setting}. In Section~\ref{sec:proper-base-change},
read the exchange map, the statement of Theorem~\ref{pbc:torsion-theorem},
and the coefficient comparison in \S\ref{pbc:condensed}, and postpone
the geometric proof. In Section~\ref{b:sec:coefficients} read
the statements and the lattice argument. In Section~\ref{sec:formal-half},
accept compactness (Lemma~\ref{ctf:unified-compact}) and
Brown representability (Theorem~\ref{b:thm:brown-adjoint}) and skip
their proofs. Read Section~\ref{sec:restriction-chapter} carefully: it
contains the geometry. Skim Sections~\ref{i:heart} and~\ref{n:proof}.
In Section~\ref{sec:smooth-purity}, read the finite-coefficient part
\S\ref{sp:grp-finite} and the trace \S\ref{sp:grp-trace}, and accept the
passage through towers. The appendices can be read when their
results are first cited; the places are listed at the end of the
introduction.

\subsection{Conventions}
We fix a prime $\ell$ and a field $k$ that is finite or algebraically
closed, of characteristic different from $\ell$. Unless said otherwise,
schemes are separated and of finite type over $k$, and morphisms are
$k$-morphisms. The hypothesis on $k$ is used only to know that $k$ has
finite cohomological dimension (Lemma~\ref{b:lem:base-field-cd}); the
constructions of \S\S\ref{n:coefficients} and~\ref{n:ambient} make sense
for arbitrary schemes.

For the six operations, $\Lambda$ is an admissible
coefficient ring (\S\ref{n:coefficients}): a finite commutative $\ell$-primary ring $A$, or
$\Zl$, $\mathcal O_E$, $\Ql$, $E$ or $\Qbar$. The case $E=\Ql$ is
allowed, so $\Zl=\mathcal O_{\Ql}$ and $\Ql$ are the basic integral and
rational members of the family. All are regarded as condensed rings,
and in every case $\Lambda_X=p_X^*\Lambda$. For discrete finite $A$
this pullback is the classical constant sheaf; the completion,
localization and finite-extension descriptions of the other cases are
comparisons (\S\ref{n:coefficients}).

The operations act on $\Dctf(X,\Lambda)$, and $f^!$ is first
constructed on its localizing closure $\Loc(\Dctf(X,\Lambda))$; both are
defined in \S\ref{n:constructible}. We keep $\Dbc$ for bounded
constructible cohomology without a Tor condition: it equals $\Dctf$
for integral and rational coefficients, but can be larger for finite
rings. Results with a different coefficient scope state it: arguments
with cohomology sheaves in $\Dbc$, definitions for arbitrary condensed
rings, and torsion proper base change.

Derived categories are the classical ones: $D(X_{\proet},\Lambda)$ is
the derived category of the abelian category of sheaves of
$\Lambda_X$-modules. All functors are derived, and we write $f^*$,
$Rf_*$, $Rf_!$ and $f^!$; in particular $f^!$ always denotes a functor
on derived categories. For $f\colon X\to Y$, unless another category is
named, $f^*\colon D(Y_{\proet},\Lambda)\to D(X_{\proet},\Lambda)$ and
$Rf_*\colon D(X_{\proet},\Lambda)\to D(Y_{\proet},\Lambda)$; the
categories on which $Rf_!$ and $f^!$ act are fixed where these functors
are defined. Sizes are taken with respect to a fixed,
sufficiently large universe. The category $\Dbc(X,\Lambda)$ is a full
subcategory of $D(X_{\proet},\Lambda)$; we do not compare it with
$D^b(\Cons(X,\Lambda))$.

\subsection{Organization}
Section~\ref{sec:proetale-foundations} develops the pro-\'etale site,
its exactness properties, its classical part, and the adic cohomology
comparison. Section~\ref{sec:setting} introduces the coefficient rings
and the categories of constructible complexes, states the main theorem
and explains its proof. Section~\ref{sec:proper-base-change} proves
proper base change, which is needed before proper direct image can be
used; the proof goes through injective acyclicity, the projective line
and Chow's lemma. Section~\ref{b:sec:coefficients} collects the facts
that reduce questions about $E$- and $\Qbar$-coefficients to torsion
coefficients. Sections~\ref{sec:formal-half}
and~\ref{sec:restriction-chapter} carry out halves (A) and (B).
The existence theorem is uniform in the coefficients; its compactness
hypotheses are verified for finite rings, for fields, and, in
Section~\ref{i:heart}, for $\mathcal O_E$.
Section~\ref{n:proof} completes the proof and treats the remaining
operations and change of coefficients, and Section~\ref{sec:smooth-purity}
proves smooth purity.

Three appendices supply the algebraic geometry the main text uses; from
the main text they use only Section~\ref{sec:proetale-foundations}.
Appendix~\ref{app:compactification} proves Nagata's compactification
theorem, the existence of common refinements of compactifications, and
Chow's lemma. Nagata's theorem defines $Rf_!$
(\S\ref{sec:plan} and Section~\ref{sec:formal-half}), and Chow's lemma
reduces proper base change to the projective line
(\S\ref{pbc:sec-reduction}). Appendix~\ref{app:picard} constructs the
Picard scheme and the Jacobian of a curve, with its torsion points
prime to the characteristic. Appendix~\ref{app:curves} uses them to
compute the cohomology of curves with torsion coefficients; these
calculations give the case of the projective line in proper base change
(\S\ref{pbc:projective-line-section}) and the cohomology of $\PP^1$
used in Section~\ref{sec:restriction-chapter}.

All coefficient sheaves are defined on the pro-\'etale site. Finite
geometric calculations may be made on the \'etale site through the
comparison equivalence. The table records which coefficients and which
site each part uses.
\begin{center}
\small
\begin{tabular}{@{}p{0.2\textwidth}p{0.33\textwidth}p{0.39\textwidth}@{}}
\hline
Part & Coefficients and site & Purpose\\
\hline
\S\ref{sec:proetale-foundations}
 & module sheaves on $X_{\proet}$
 & local structure, derived limits, classical comparison\\
\S\ref{sec:setting}, \S\ref{b:sec:coefficients}
 & $A$, $\Zl$, $\mathcal O_E$, $\Ql$, $E$, $\Qbar$, pro-\'etale
 & coefficients and $\Dctf$; reduction to finite coefficients\\
\S\ref{sec:proper-base-change}
 & classical torsion on $X_{\proet}$
 & exchange, injective acyclicity, projective line, Chow reduction\\
\S\ref{sec:formal-half}
 & $A$, $\Zl$, $\mathcal O_E$, $\Ql$, $E$, $\Qbar$, pro-\'etale
 & Brown representability on $\Loc(\Dctf)$\\
\S\ref{b:sec:finite-geometric}
 & finite $\ell$-primary rings, pro-\'etale
 & $f^!$ on $\PP^1$ and open direct image, by \'etale calculations\\
rest of \S\ref{sec:restriction-chapter}
 & $A$, $\Ql$, $E$, $\Qbar$, pro-\'etale
 & $f^!$ preserves $\Dctf$\\
\S\ref{i:heart}
 & $\Zl$, $\mathcal O_E$, pro-\'etale
 & compactness and $f^!$ with integral coefficients\\
\S\ref{n:proof}
 & $A$, $\Zl$, $\mathcal O_E$, $\Ql$, $E$, $\Qbar$, pro-\'etale
 & six functors on $\Dctf$; change of coefficients\\
\S\ref{sec:smooth-purity}
 & $A$, $\Zl$, $\mathcal O_E$, $\Ql$, $E$, $\Qbar$
 & smooth purity, trace, Frobenius\\
App.~\ref{app:compactification}
 & schemes
 & Nagata compactification, refinements, Chow's lemma\\
App.~\ref{app:picard}
 & schemes, coherent cohomology
 & Picard scheme, Jacobian, its torsion\\
App.~\ref{app:curves}
 & classical torsion sheaves on $C_{\proet}$
 & cohomology of curves\\
\hline
\end{tabular}
\end{center}
The adic comparisons use the sequence
\[
 \mathcal O_E/\pi_E^n\ \rightsquigarrow\ \mathcal O_E
 \ \rightsquigarrow\ E\ \rightsquigarrow\ \Qbar .
\]
Each arrow is a coefficient comparison whose precise hypotheses are
stated in \S\ref{n:change}. The first needs compatible statements for
all $n$ and derived completeness, the second inverts $\ell$, and the
third uses that every object and morphism over $\Qbar$ is defined over
a finite $E$. The finite rings in this sequence are quotients of
$\mathcal O_E$; an arbitrary admissible finite ring $A$ is treated
directly by Theorem~\ref{ctf:unified-brown}.
To pass back from $E$ to finite coefficients, reducing an $E$-complex
modulo $\pi_E$ is useless, since it gives zero. Instead one chooses an
integral model, a lattice on each stratum (\S\ref{b:sec:coefficients}),
and reduces that.

\section{The pro-\'etale site and its comparison theorems}
\label{sec:proetale-foundations}

Every coefficient sheaf and every operation in these notes lives on
$X_{\proet}$, the site of Bhatt and Scholze~\cite{BS}. Three of its
properties are used. Weakly \'etale geometry supplies enough affine
objects with exact sections. These objects control products, sequential
limits and unbounded complexes. Finally, the classical part of the
pro-\'etale topos admits \'etale calculations, and continuous and adic
cohomology are recovered from it.

In this section schemes are arbitrary unless a result states a
finiteness hypothesis. Even when $X$ is a variety, objects of
$X_{\proet}$ need not be Noetherian or of finite type. The standing
assumption that schemes are of finite type over $k$ resumes in
Section~\ref{sec:setting}.

\subsection{Weakly \'etale geometry and the pro-\'etale site}
\label{pe:local-site}

The universe fixed in the conventions is taken large enough to contain
the covers constructed below; $X_{\proet}$ refers to that choice.

\begin{definition}[Ind-\'etale, pro-\'etale and weakly \'etale morphisms]
\label{pe:morphisms}
A ring map $R\to B$ is \emph{ind-\'etale} if
$B=\varinjlim_\alpha B_\alpha$ for a filtered system of finitely
presented \'etale $R$-algebras. For an affine base
$X=\operatorname{Spec}(R)$, the corresponding morphism
\[
 \operatorname{Spec}(B)=\varprojlim_\alpha
 \operatorname{Spec}(B_\alpha)\longrightarrow X
\]
is called \emph{affine pro-\'etale}. More generally, a
\emph{pro-\'etale affine over $X$} is an inverse limit of affine
schemes \'etale over $X$.

A scheme morphism $u:U\to X$ is \emph{weakly \'etale} if $u$ and its
diagonal $\Delta_u:U\to U\times_XU$ are flat. On affine charts this
means that both
\[
 R\longrightarrow B,
 \qquad B\otimes_RB\longrightarrow B
\]
are flat. Thus every \'etale morphism, and every affine pro-\'etale
morphism, is weakly \'etale. A weakly \'etale morphism locally of finite
presentation is \'etale.
\end{definition}

Flatness of the diagonal takes the place of finite presentation; flatness
of $u$ alone is too weak to define an object of the pro-\'etale site.
A weakly \'etale morphism need not have a global presentation as a
limit of \'etale morphisms. What holds instead is the covering
refinement of Theorem~\ref{pe:indetale-refinement}; see also
\stacks{097Z}.

\begin{lemma}[Stability]
\label{pe:weakly-etale-stability}
Weakly \'etale morphisms are stable under base change and composition,
and the property is local on source and target for the Zariski topology.
Every $X$-morphism between two weakly \'etale $X$-schemes is weakly
\'etale. Affine pro-\'etale morphisms are stable under affine base
change and composition.
\end{lemma}
\begin{proof}
Base change preserves flatness of a morphism and its diagonal. For a
composite $V\to U\to X$, the diagonal factors through
$V\times_UV\to V\times_XV$, which is a base change of
$\Delta_{U/X}$; this gives the assertion about composites.
For an $X$-map $h:U\to V$, the graph
$U\to U\times_XV$ is a base change of the flat diagonal of $V/X$, and
the projection $U\times_XV\to V$ is flat because $U/X$ is flat.
Their composite is $h$, while $\Delta_h$ is a base change of
$\Delta_{U/X}$. Zariski locality follows from locality of flatness.
For ind-\'etale rings, tensor products commute with filtered colimits,
and finite presentation lets one choose a finite stage for each
\'etale algebra; this gives base change and transitivity. The same
arguments show that a filtered colimit of \'etale algebras has flat
multiplication map; see \stacks{092A}.
\end{proof}

\begin{definition}[The small pro-\'etale site]
\label{pe:site}
The objects of $X_{\proet}$ are weakly \'etale $X$-schemes, with all
$X$-morphisms as arrows. A family $\{U_i\to U\}_{i\in I}$ is a
covering if it is an fpqc covering. Explicitly, for every affine open
$W\subset U$ there must be finitely many indices
$i_1,\ldots,i_m$ and affine opens
$V_j\subset U_{i_j}\times_UW$ with
\[
 W=\bigcup_{j=1}^m\operatorname{image}(V_j\to W).
\]
The indices may depend on $W$; no finiteness of $I$ is required.
Equivalently one may use finitely many quasi-compact opens above each
$W$. See \stacks{0989} and \stacks{098A}.
\end{definition}

The finite refinement condition is part of the topology; a jointly
surjective family without it is not a covering. For an affine
target every cover refines to a finite affine cover, whose disjoint
union is affine and faithfully flat over the target. Ordinary \'etale
and Zariski covers are pro-\'etale covers: flat maps locally of finite
presentation are open, so quasi-compactness gives the required finite
refinements; see \stacks{098B}.

\begin{proposition}[The topos and its functoriality]
\label{pe:topos-functoriality}
The above coverings define a subcanonical site. Its category of sheaves
$\Sh(X_{\proet})$ is a Grothendieck topos. Every scheme morphism
$f:X\to Y$ induces a geometric morphism
\[
 f_{\proet}:\Sh(X_{\proet})\longrightarrow\Sh(Y_{\proet}),
 \qquad f^{-1}\dashv f_*.
\]
Here $f^{-1}\colon\Sh(Y_{\proet})\to\Sh(X_{\proet})$ preserves finite
limits and all colimits, and is exact on
abelian sheaves. These morphisms are compatible with composition.
\end{proposition}
\begin{proof}
Identity covers, composites of covers, and base changes of covers satisfy
the axioms by Lemma~\ref{pe:weakly-etale-stability} and the fpqc covering
axioms. Representable presheaves satisfy descent for these covers by
fpqc descent for morphisms of schemes, so they are sheaves. Finally,
$V\mapsto V\times_YX$ preserves finite limits and covering families.
It therefore defines the displayed geometric morphism, with inverse
image left exact by the site construction; exactness on abelian sheaves
and compatibility with composition follow. These are the small-site
versions of the constructions in \stacks{0988}.
\end{proof}

\subsection{Local rings, profinite components, and contractible covers}

Every ring has a faithfully flat ind-\'etale cover by a ring over which
all ind-\'etale covers split (Theorem~\ref{pe:wc-cover}). Such rings are defined by a hierarchy of
local conditions.

\begin{definition}[The local hierarchy]
\label{pe:wlocal}
A spectral space $T$ is \emph{w-local} if its set $T^c$ of closed
points is closed and every point specializes to exactly one closed
point. A ring $R$ is w-local if $\operatorname{Spec}(R)$ is w-local.
It is \emph{w-strictly local} if it is w-local and
$R_{\mathfrak m}$ is strictly henselian for every maximal ideal
$\mathfrak m$. It is \emph{w-contractible} if it is w-strictly local
and the profinite space $\pi_0(\operatorname{Spec}(R))$ is
extremally disconnected.
\end{definition}

Recall that spectral spaces are quasi-compact and sober, with a basis
of quasi-compact opens closed under finite intersections. The space of
connected components of a spectral space is profinite. For a w-local
space the natural map
$T^c\to\pi_0(T)$ is a homeomorphism: each component consists of the
generalizations of its unique closed point, and $T^c$ is compact.
For $T=\operatorname{Spec}(R)$ this component is the image of
$\operatorname{Spec}(R_{\mathfrak m})$.
Thus w-strict locality gives strictly henselian geometry component by
component, and w-contractibility also controls how choices vary across
the profinite set of components. Closedness of $T^c$ is a separate
condition: it does not follow from uniqueness of closed
specializations. See \stacks{096A}, \stacks{0968}, and \stacks{0969}.

\begin{proposition}[Stonean projectivity]
\label{pe:stonean}
For a compact Hausdorff space $S$, extremal disconnectedness means that
the closure of every open subset is open. It is equivalent to
projectivity with respect to continuous surjections of compact
Hausdorff spaces. Every profinite space $T$ admits a continuous
surjection $S\to T$ with $S$ compact Hausdorff and extremally
disconnected; such $S$ is also profinite.
\end{proposition}
\begin{proof}[Topological input and construction]
The projectivity equivalence is the compact Hausdorff projectivity
theorem; see \stacks{08YN}. One proof chooses a minimal closed subspace
of a compact source that still surjects onto $S$; extremal
disconnectedness forces that irreducible surjection to be a
homeomorphism. For the cover, take the underlying set $T^{\mathrm{disc}}$
with the discrete topology, and let
$S=\beta(T^{\mathrm{disc}})$ be its Stone--\v{C}ech compactification.
For a discrete space this compactification is extremally disconnected,
and the map to $T$ extends uniquely and surjectively; see
\stacks{090D}. Only this discrete-space case is used here.
\end{proof}

\begin{proposition}[Splitting ind-\'etale covers]
\label{pe:wc-rings}
If $R$ is w-contractible, every faithfully flat ind-\'etale ring map
$R\to B$ has an $R$-algebra retraction $B\to R$.
\end{proposition}
\begin{proof}[Proof outline]
First refine $B$ by ind-Zariski operations so that its spectrum is
w-local and its closed points lie exactly above the closed points of
$\operatorname{Spec}(R)$. Here ind-Zariski means a filtered colimit of
local isomorphisms; these operations preserve local rings.
Over a closed point of $\operatorname{Spec}(R)$, strict henselianity
implies that the induced map of local rings for an ind-\'etale algebra
is an isomorphism, so what remains is a choice of component. Stonean
projectivity splits the surjection of profinite component spaces.
Restriction to that closed set of components is realized by a filtered
colimit of localizations at idempotents; it produces an $R$-algebra
isomorphic to $R$, hence a retraction. The algebraic realization of
these refinements is supplied by \stacks{097D} and \stacks{09AZ}; the
resulting argument is \stacks{0982}.
\end{proof}

\begin{theorem}[Existence of contractible covers]
\label{pe:wc-cover}
Every ring $R$ admits a faithfully flat ind-\'etale algebra $D$ that
is w-contractible.
\end{theorem}
\begin{proof}[Construction and its algebraic inputs]
The construction has two stages, following \cite[Theorem~2.4.9]{BS}
and \stacks{0983}. The first constructs a faithfully flat ind-\'etale
$R$-algebra $C$ that is w-strictly local: w-localize using finite
Zariski decompositions, adjoin solutions to \'etale lifting problems in
a filtered construction, and then refine the w-localization so that the
local rings at closed points are strictly henselian. This is the
simultaneous strict-localization construction of \stacks{097X}. The
finite stages are what make the result ind-\'etale.

For the second stage, choose a Stonean surjection
$S\to\pi_0(\operatorname{Spec}(C))$ using
Proposition~\ref{pe:stonean}, and realize this map by an ind-Zariski
$C$-algebra $D$ such that
\[
 |\operatorname{Spec}(D)|
 \simeq |\operatorname{Spec}(C)|
       \times_{\pi_0(\operatorname{Spec}(C))}S
\]
as topological spaces. Concretely, write $S$ as a limit of finite
quotients, use finite products of $C$, restrict to the relevant
closed sets of components by idempotent localizations, and take a
filtered colimit; see \stacks{097D}. This construction identifies local
rings, is surjective when $S$ is, and preserves w-locality. Hence
$C\to D$ is faithfully flat and the local rings at closed points of
$D$ remain strictly henselian. Its component space is $S$, so $D$ is
w-contractible, and the composite $R\to C\to D$ is ind-\'etale and
faithfully flat.
\end{proof}

Objects of $X_{\proet}$ are weakly \'etale, while the covers above are
ind-\'etale. The next theorem bridges the two: it is the reason a site
defined by weakly \'etale morphisms can be computed with inverse limits
of \'etale affines.

\begin{theorem}[Weakly \'etale maps admit pro-\'etale refinements]
\label{pe:indetale-refinement}
If $R\to B$ is weakly \'etale, there is a faithfully flat
ind-\'etale map $B\to B'$ such that $R\to B'$ is ind-\'etale.
\end{theorem}
\begin{proof}[The local rigidity argument]
The simultaneous strict-localization construction for the pair
$R\to B$ gives compatible ind-\'etale refinements $R\to C$
and $B\to D$, with both spectra w-local and with closed points of
$D$ mapping to closed points of $C$. The map $C\to D$ is still weakly
\'etale by Lemma~\ref{pe:weakly-etale-stability}, and the local rings
of $C$ at closed points are strictly henselian. Olivier's rigidity
theorem says that a weakly \'etale local map out of a strictly
henselian local ring is an isomorphism; see \stacks{092Z}. Applied at
closed points, and then at their generalizations, it shows that
$C\to D$ identifies local rings. The ind-Zariski refinement theorem for
maps identifying local rings now gives a faithfully flat ind-Zariski
$D\to D'$ such that $C\to D'$ is ind-Zariski. Taking $B'=D'$ proves
the assertion. See \stacks{097Z} for the compatible refinement and the
proof.
\end{proof}

\subsection{Affine bases and exact local computation}

\begin{proposition}[A cover-dense affine presentation]
\label{pe:affine-basis}
Every object of $X_{\proet}$ has a pro-\'etale covering by
pro-\'etale affines over $X$. For affine $X$ these form a full subsite
whose category of sheaves is equivalent to $\Sh(X_{\proet})$.
\end{proposition}
\begin{proof}
Cover $X$ by affine opens and cover their inverse images in $U$ by
affine opens. On such a chart the map is $R\to B$ weakly \'etale.
Theorem~\ref{pe:indetale-refinement} supplies
$\operatorname{Spec}(B')\to\operatorname{Spec}(B)$ as a single
faithfully flat cover, with $\operatorname{Spec}(B')$ a limit of
affine schemes \'etale over $\operatorname{Spec}(R)$, hence over
$X$. The site comparison lemma gives the equivalence of sheaf
categories. This is \cite[Lemma~4.2.4]{BS}.
\end{proof}

\begin{theorem}[Enough affine weakly contractible objects]
\label{pe:weakly-contractible-basis}
Every object of $X_{\proet}$ has a covering by affines
$U=\operatorname{Spec}(D)$ with $D$ w-contractible. For such $U$,
every epimorphism of sheaves $F\to h_U$ has a section. Equivalently,
$h_U$ is projective in $\Sh(X_{\proet})$.
\end{theorem}
\begin{proof}
Apply Theorem~\ref{pe:wc-cover} to each affine chart. For the
splitting assertion, a faithfully flat weakly \'etale $D$-algebra
refines, by Theorem~\ref{pe:indetale-refinement}, to a faithfully flat
ind-\'etale $D$-algebra. Proposition~\ref{pe:wc-rings} splits this
refinement and hence splits the original map. Every cover of the affine
$U$ refines to a finite affine cover; its disjoint union is a single
faithfully flat weakly \'etale affine cover and therefore splits. The
inverse images of its finitely many components give a finite clopen
partition of $U$ subordinate to the original cover.

Now an epimorphism $F\to h_U$ has lifts of $\mathrm{id}_U$ after
a covering of $U$. Refine that covering by the clopen partition just
constructed and glue the lifts on its disjoint pieces. This gives a
global lift of $\mathrm{id}_U$, hence a section. The same argument,
after pulling back an arbitrary epimorphism along a map from $h_U$,
proves projectivity. See \stacks{098F} and \stacks{0990}.
\end{proof}

\begin{corollary}[Exact sections on contractible affines]
\label{pe:wc-sections}
For an affine $U$ as in Theorem~\ref{pe:weakly-contractible-basis},
the functor $F\mapsto F(U)$ on abelian sheaves is exact. In particular
\[
 H^q(U,F)=0\quad(q>0),\qquad
 R\Gamma(U,K)\simeq K(U)
\]
for a complex $K$, with the last expression computed term by term.
\end{corollary}
\begin{proof}
Sections preserve kernels, and projectivity gives surjectivity on
sections for every epimorphism. The sections functor is therefore
exact, carries acyclic complexes to acyclic complexes, and agrees
with its derived functor, including for unbounded complexes.
\end{proof}

The affine objects in Theorem~\ref{pe:weakly-contractible-basis} form
a basis in the cover-dense sense. Their fibre products need not be
w-contractible, but they admit such covers again. Sheaf-theoretic
assertions can therefore be checked on these affines, while global
cohomology still records the gluing between them. This local exactness
is the input for the derived inverse-limit arguments below.

\subsection{Weakly contractible objects and module sheaves}
\label{pe:sec-homological}

On objects with exact sections, ordinary derived categories of module
sheaves can be computed sectionwise. In this subsection $X$ is an
arbitrary scheme and $\mathcal R$ is any sheaf of rings on $X_{\proet}$,
for instance the pullback of a condensed ring.

\begin{definition}[Weak contractibility]\label{pe:weak-contractible}
An object $W$ of $X_{\proet}$ is \emph{weakly contractible} if every
epimorphism of sheaves of sets $F\to h_W$ admits a section, where $h_W$
is the representable sheaf associated to $W$. Equivalently, $h_W$ is
projective in $\Sh(X_{\proet})$, or every epimorphism $F\to G$ induces a
surjection $F(W)\to G(W)$.
\end{definition}

For the equivalence, pull an epimorphism $F\to G$ back along a section
$h_W\to G$; the resulting epimorphism onto $h_W$ splits. Conversely,
surjectivity on sections lifts $\mathrm{id}_{h_W}$ along any
epimorphism onto $h_W$. This is the categorical definition
of~\cite[\S3.2]{BS}. The splitting condition concerns a cover as a
whole: a disconnected $W$ need not lift into a single member of a
covering family.

By Theorem~\ref{pe:weakly-contractible-basis} every object of
$X_{\proet}$ has an affine weakly contractible covering family, and by
Corollary~\ref{pe:wc-sections} each member has exact sections. The
consequence for module sheaves is used repeatedly.

\begin{corollary}[A conservative family preserving all module colimits]
\label{pe:sections}
Evaluation on all affine weakly contractible objects detects exactness
of complexes and isomorphisms of $\mathcal R$-module sheaves.
For each such object $W$, evaluation preserves all limits and all
colimits of $\mathcal R$-modules.
\end{corollary}
\begin{proof}
If a sequence becomes exact on all such $W$, a section in a kernel
lifts locally on a weakly contractible cover, so it lies in the
image sheaf. Applying this criterion to kernels and cokernels proves
conservativity, and applying it degreewise treats complexes.

Limits of sheaves are computed sectionwise. For colimits of
modules, first consider a direct sum $\bigoplus_i M_i$. A section is
locally a finite sum; quasi-compactness gives a finite covering on which
this description holds, so one finite set of summand indices
suffices on all covering members. The sheaf condition then identifies
the section with a section of that finite subsum. Evaluation therefore
preserves arbitrary direct sums (\stacks{0935}). An arbitrary
colimit of modules is a cokernel of a map between direct sums, and
exactness of evaluation gives the assertion. See \stacks{0945} for
the conservative criterion.
\end{proof}

The quasi-compactness used here is quasi-compactness as an object of the
site: coverings admit finite refinements. Affine objects of
$X_{\proet}$ have this property. Over a quasi-compact, quasi-separated
object, the quasi-compact, quasi-separated subsite with finite covers is
a coherent site. In topos terminology an object $P$ is coherent if it
is quasi-compact and $V\times_PV$ is quasi-compact for every
quasi-compact object $V$ mapping to $P$. Together with the affine
refinements, this provides the local coherent projective objects
required for $\Sh(X_{\proet})$ to be a \emph{locally weakly
contractible} topos in the sense of~\cite{ProetaleBlueprint}.

\begin{remark}[Evaluation is not generally a point]
\label{pe:evaluation-not-point}
For modules, sections over a quasi-compact weakly contractible $W$
preserve all colimits. The corresponding assertion for coproducts of
\emph{sheaves of sets} is generally false: on a disconnected $W$, a
section of $F\amalg G$ can choose different summands on different
components. Thus evaluation need not be the inverse image of a point
of the topos. It nevertheless gives the conservative exact module
functors needed below; see \stacks{0991}.
\end{remark}

\subsection{Repleteness, products, and sequential limits}

Inverse limits are needed to pass from torsion to adic coefficients.
Since surjectivity and exactness can be checked on sections over weakly
contractible objects, sequential limits along epimorphisms remain
epimorphic and products are exact.

\begin{definition}[Repleteness]\label{pe:replete-def}
A topos $\mathcal T$ is \emph{replete} if, for every inverse sequence
$\cdots\to F_2\to F_1\to F_0$ of sheaves of sets with epimorphic
transition maps, the projection $\varprojlim_nF_n\to F_m$ is an
epimorphism for every $m$.
\end{definition}

\begin{proposition}[Repleteness of the pro-\'etale topos]
\label{pe:replete}
For every scheme $X$, the topos $\Sh(X_{\proet})$ is replete.
\end{proposition}
\begin{proof}
Evaluate the sequence on a weakly contractible $W$; the transitions
$F_{n+1}(W)\to F_n(W)$ are then surjective. Given a section at stage
$m$, choose lifts successively at stages $m+1,m+2,\ldots$ (by countable
choice); the earlier stages are determined by the transition maps. This
gives a compatible section of $\varprojlim_nF_n(W)$, and since limits
are computed sectionwise, the projection is surjective on $W$.

An arbitrary $U$ is covered by affine weakly contractible objects, and
the construction on each member of the cover gives compatible lifts at
all stages. Hence the projection is an epimorphism of sheaves; see
\cite[\S\S3.1--3.2]{BS}.
\end{proof}

\begin{proposition}[Exact products]\label{pe:products}
Arbitrary set-indexed products are exact in
$\Mod(X_{\proet},\mathcal R)$. Products in
$D(X_{\proet},\mathcal R)$ are computed by taking termwise products
of representative complexes, and their cohomology sheaves are the
products of the cohomology sheaves.
\end{proposition}
\begin{proof}
On a weakly contractible $W$, a product of short exact sequences
evaluates to a product of short exact sequences of
$\mathcal R(W)$-modules, which is exact: by the axiom of choice one
lifts a tuple coordinatewise. Corollary~\ref{pe:sections} detects
exactness of sheaves on such $W$. (Repleteness alone gives the
countable case, as in \cite[\S3.1]{BS}; arbitrary products use the
weakly contractible basis.)

Exact products preserve quasi-isomorphisms. The derived product is
the termwise product of K-injective representatives, and by exactness
any quasi-isomorphic representatives may be used instead. Termwise
products therefore compute the product and commute with cohomology;
see \stacks{0947} and \stacks{08TB}.
\end{proof}

For a countable tower write $t_n:M_{n+1}\to M_n$ and set
$T((m_n)_n)=(t_n(m_{n+1}))_n$. For a tower of complexes, represented
by compatible complexes, the derived inverse limit is the homotopy
fibre
\begin{equation}\label{pe:rlim-cone}
 R\varprojlim_n K_n
 =
 \operatorname{Cone}\left(
  \prod_n K_n\xrightarrow{\,1-T\,}\prod_n K_n
 \right)[-1].
\end{equation}
The cone is formed on the level of complexes, which makes the
construction functorial in the tower; it agrees with the derived
inverse-limit functor on towers by \stacks{0940}.

\begin{proposition}[The Milnor sequence]\label{pe:milnor}
For a countable tower of $\mathcal R$-modules, $R^p\varprojlim=0$
for $p>1$, and
\[
 R^1\varprojlim_n M_n
 =\coker\left(\prod_n M_n\xrightarrow{\,1-T\,}\prod_n M_n\right).
\]
For a tower $(K_n)$ of complexes there is a short exact sequence of
sheaves
\begin{equation}\label{pe:milnor-sequence}
 0\longrightarrow R^1\varprojlim_n H^{q-1}(K_n)
 \longrightarrow H^q(R\varprojlim_nK_n)
 \longrightarrow \varprojlim_n H^q(K_n)
 \longrightarrow 0.
\end{equation}
If the transitions of a tower of modules are epimorphisms, its derived
limit is its ordinary limit in degree zero.
\end{proposition}
\begin{proof}
The formulas follow from the long exact cohomology sequence of
\eqref{pe:rlim-cone}, since by Proposition~\ref{pe:products} the
cohomology of each product term is the product of cohomology sheaves.
For epimorphic transitions, $1-T$ is surjective on every weakly
contractible $W$: given $(a_n)$, choose $m_0$ arbitrarily and then,
recursively, $m_{n+1}$ with $t_n(m_{n+1})=m_n-a_n$. By the
conservative criterion $1-T$ is an epimorphism of sheaves, so its
cokernel vanishes. See \stacks{0947}.
\end{proof}

More generally, the $R^1\varprojlim$ term vanishes if, on every affine
weakly contractible $W$, the tower of section modules satisfies the
Mittag--Leffler condition: evaluate the cokernel formula and apply the
Mittag--Leffler lemma for modules (\stacks{07KV}). The condition is
to be checked on this basis; over a general object, sections of an
epimorphism of sheaves need not be surjective.

\subsection{Left completeness and hypercover calculations}

\begin{proposition}[Left completeness]\label{pe:left-complete}
The standard $t$-structure on $D(X_{\proet},\mathcal R)$ is left complete:
for every complex $K$ the natural map
\[
 K\longrightarrow R\varprojlim_{n\geq0}\tau_{\geq-n}K
\]
is an isomorphism.
\end{proposition}
\begin{proof}
Represent $K$ by a complex and form its canonical truncation tower,
with maps $K\to\tau_{\geq-n}K$ and
$\tau_{\geq-(n+1)}K\to\tau_{\geq-n}K$. For fixed $q$, the sheaves
$H^q(\tau_{\geq-n}K)$ and $H^{q-1}(\tau_{\geq-n}K)$ are eventually
constant, with identity transition maps, equal to the corresponding
cohomology sheaves of $K$. In \eqref{pe:milnor-sequence} the $R^1$
term therefore vanishes and the limit term is $H^q(K)$. The natural
map induces this identification in every degree, so it is an
isomorphism. This is \stacks{0947}(2).
\end{proof}

K-injective resolutions of arbitrary complexes exist in the
Grothendieck category of $\mathcal R$-modules (\stacks{079P}), so
$R\Gamma(U,-)$ and $Rf_*$ are defined on the unbounded derived
categories: for $f\colon X\to Y$, $Rf_*$ goes from
$D(X_{\proet},\mathcal R)$ to the derived category of
$f_*\mathcal R$-module sheaves on $Y_{\proet}$. As derived right adjoints they preserve products, and
hence the derived inverse limits \eqref{pe:rlim-cone} (\stacks{0940}).
An unbounded calculation can therefore be carried out on bounded-below
truncations followed by a derived limit; the passage to individual
cohomology groups then involves the Milnor $R^1$ term.

For bounded-below complexes, weakly contractible objects give explicit
cochain complexes computing cohomology.

\begin{proposition}[Weakly contractible hypercovers]
\label{pe:hypercover}
Let $U$ be a quasi-compact, quasi-separated object of $X_{\proet}$.
There is a hypercover $W_\bullet\to U$ whose term in every simplicial
degree is one affine weakly contractible object. If
$K\in D^+(X_{\proet},\mathcal R)$ is represented by a complex
$K^\bullet$ vanishing below degree $a$, then
\begin{equation}\label{pe:hypercover-complex}
 R\Gamma(U,K)
 \simeq \operatorname{Tot}\bigl(K^b(W_r)\bigr)_{r\geq0,\ b\geq a}.
\end{equation}
The cosimplicial differential is the alternating sum of pullbacks
along face maps, combined with $(-1)^r$ times the differential of
$K^\bullet$. In total degree $q$ this complex has the finite sum
$\bigoplus_{0\leq r\leq q-a}K^{q-r}(W_r)$.
\end{proposition}
\begin{proof}
Start with a weakly contractible cover of $U$, and at each stage cover
the matching object by weakly contractible affines; a matching object
need not itself be weakly contractible. The matching objects are
finite limits of quasi-compact, quasi-separated schemes and stay in
that subsite, and finite disjoint unions give a single affine term at
each stage; see \stacks{09A1}.

For a sheaf $M$, the hypercover spectral sequence has
$E_1^{r,s}=H^s(W_r,M)$, and the rows with $s>0$ vanish by
Corollary~\ref{pe:wc-sections}. The remaining row is the
cosimplicial section complex. For bounded-below $K^\bullet$, choose
a bounded-below injective resolution $K^\bullet\to I^\bullet$. Since
evaluation on $W_r$ is exact, $K^\bullet(W_r)\to I^\bullet(W_r)$ is a
quasi-isomorphism. After a shift both double complexes lie in the
first quadrant, so their total complexes are quasi-isomorphic, and
hypercover descent for $I^\bullet$ identifies the second with
$R\Gamma(U,K)$. See \stacks{09A2}.
\end{proof}

For a single sheaf this is the cochain complex
$M(W_0)\to M(W_1)\to M(W_2)\to\cdots$; it gives no bound on the
cohomological dimension of arbitrary pro-\'etale sheaves. The
finiteness of the sum in each total degree uses the lower bound on
$K^\bullet$. For unbounded $K$ one uses instead
$R\Gamma(U,K)\simeq
R\varprojlim_n R\Gamma(U,\tau_{\geq-n}K)$, by
Proposition~\ref{pe:left-complete}; the sum totalization of an
unbounded double complex does not compute this.

\begin{corollary}[Bounded-below continuity]
\label{pe:bounded-continuity}
Suppose $U$ is quasi-compact and quasi-separated and
$(K_i^\bullet)$ is a filtered diagram of complexes of $\mathcal R$-modules
vanishing below one fixed degree $a$. Then for every $q$,
\[
 \varinjlim_i H^q(U,K_i^\bullet)
 \simeq H^q(U,\varinjlim_iK_i^\bullet).
\]
If $f:X\to Y$ is quasi-compact and quasi-separated, the analogous map
$\varinjlim_iR^qf_*K_i^\bullet\to
R^qf_*(\varinjlim_iK_i^\bullet)$ is an isomorphism of sheaves on
$Y_{\proet}$.
\end{corollary}
\begin{proof}
Use one hypercover $W_\bullet$ for all $i$. By
Corollary~\ref{pe:sections}, sections on each $W_r$ commute with
filtered colimits, hence so does the finite sum in each total degree
of \eqref{pe:hypercover-complex}. Filtered colimits of modules are
exact, which gives the first assertion. For direct images, evaluate the
cohomology presheaves on affine objects $V$ of $Y_{\proet}$: each
$X\times_YV$ is quasi-compact and quasi-separated, so the first
assertion applies, and sheafification commutes with colimits.
\end{proof}

Complexes in degree zero give continuity of $H^q(U,-)$ for filtered
colimits of module sheaves, and hence for direct sums, which are
filtered colimits of finite subsums. For general complexes the common
lower bound is needed. Direct sums of unbounded complexes
are handled later using uniform cohomological bounds for the classical
and constructible coefficient categories.

\subsection{Classical objects inside the pro-\'etale topos}
\label{pe:comparison-section}
Coefficients live on the pro-\'etale site throughout these notes. To
use the usual torsion arguments of \'etale cohomology, we identify the
part of the pro-\'etale topos that comes from the \'etale topos. This
part is not the whole topos: the integral and rational condensed
coefficients lie outside it (Example~\ref{pe:nonclassical-example}).
In this subsection $\Lambda$ denotes a discrete
commutative ring and $\Lambda_X$ its constant sheaf on $X_{\proet}$.

\begin{definition}[Classical pro-\'etale sheaves]
\label{pe:classical-definition}
A sheaf $F$ on $X_{\proet}$ is \emph{classical} if, for every
pro-\'etale affine presentation $U=\varprojlim_i U_i$ with $U_i$ affine
and \'etale over $X$, its restriction maps induce a bijection
\begin{equation}\label{pe:eq:classical-continuity}
 \varinjlim_i F(U_i)\xrightarrow{\ \sim\ }F(U).
\end{equation}
For sheaves of modules the map is required to be an isomorphism of
modules. A complex is classical if all its cohomology sheaves are
classical.
\end{definition}

\begin{proposition}[The comparison morphism and its inverse image]
\label{pe:comparison-morphism}
Inclusion of \'etale objects defines a morphism of topoi
\[
 \nu_X:X_{\proet}\longrightarrow X_{\et}.
\]
Its direct image $\nu_{X,*}\colon\Sh(X_{\proet})\to\Sh(X_{\et})$ is
restriction to the \'etale subsite. Its inverse image
$\nu_X^*\colon\Sh(X_{\et})\to\Sh(X_{\proet})$ is exact and has the
description
\begin{equation}\label{pe:eq:inverse-image-evaluation}
 (\nu_X^*F)(h:V\to X)
   =\Gamma(V_{\et},h_{\et}^*F),
\end{equation}
where $h_{\et}^*\colon\Sh(X_{\et})\to\Sh(V_{\et})$ is the \'etale inverse
image along the weakly \'etale map $h$. Both functors carry
$\Lambda$-module sheaves to $\Lambda$-module sheaves, so they restrict
to functors between $\Mod(X_{\et},\Lambda)$ and
$\Mod(X_{\proet},\Lambda_X)$.
For $U=\varprojlim_iU_i$ as above, this gives
\begin{equation}\label{pe:eq:affine-evaluation}
 (\nu_X^*F)(U)=\varinjlim_i F(U_i).
\end{equation}
\end{proposition}
\begin{proof}
An \'etale morphism is weakly \'etale and an \'etale covering is a
pro-\'etale covering, so the inclusion is a continuous functor
preserving finite limits and defines $\nu_X$. For
\eqref{pe:eq:inverse-image-evaluation}, define the right hand side for
every weakly \'etale $V$. Pro-\'etale coverings are fpqc coverings,
so by fpqc descent for \'etale sheaves this rule is a sheaf. Pulling
back the restriction of a pro-\'etale sheaf $G$ to $X_{\et}$ gives the
counit $\nu_X^*\nu_{X,*}G\to G$ in $\Sh(X_{\proet})$; on an \'etale
object $V$ the rule gives $F(V)$, which is the unit
$F\to\nu_{X,*}\nu_X^*F$ in $\Sh(X_{\et})$. The two adjunction identities follow by restriction and
pullback, so the rule is $\nu_X^*$ (\stacks{0GLZ}). In particular
\eqref{pe:eq:inverse-image-evaluation} describes the sheaf inverse
image directly.

For an affine limit $U$, sections of an \'etale sheaf are continuous
under affine inverse limits, which gives
\eqref{pe:eq:affine-evaluation} (\stacks{099S}).
\end{proof}

\begin{theorem}[Recognizing the classical part]
\label{pe:classical-sheaves}
The functor $\nu_X^*$ is fully faithful. Its essential image is exactly
the classical sheaves of Definition~\ref{pe:classical-definition}.
For modules it is an exact equivalence between
$\Mod(X_{\et},\Lambda)$ and the full subcategory of classical
$\Lambda_X$-module sheaves on $X_{\proet}$.
\end{theorem}
\begin{proof}
Evaluation on \'etale objects gives
$F\xrightarrow{\sim}\nu_{X,*}\nu_X^*F$ in $\Sh(X_{\et})$, hence full
faithfulness, and
every $\nu_X^*F$ is classical by \eqref{pe:eq:affine-evaluation}.
Conversely, let $G$ be classical. On a pro-\'etale affine the counit
$\nu_X^*\nu_{X,*}G\to G$ of sheaves on $X_{\proet}$ is
\[
 (\nu_X^*\nu_{X,*}G)(U)
  =\varinjlim_iG(U_i)\longrightarrow G(U),
\]
which is an isomorphism by \eqref{pe:eq:classical-continuity}. Since
pro-\'etale affines form a covering basis, the counit is an
isomorphism of sheaves.
\end{proof}

\begin{example}[A coefficient sheaf outside the classical part]
\label{pe:nonclassical-example}
Let $X$ be a geometric point and let $S=\mathbf Z_\ell$ be regarded as
a profinite object of $X_{\proet}$. The completed coefficient sheaf
$\mathbf Z_{\ell,X}$ has sections
$\operatorname{Cont}(S,\mathbf Z_\ell)$ on $S$, among them the
identity function. The pullback $\nu_X^*$ from $X_{\et}$ of the
constant sheaf with value the \emph{discrete} abelian group
$\mathbf Z_\ell$ has only locally
constant functions as sections, and these have finite image on the
compact space $S$. So the identity function does not come from any
finite stage of the profinite presentation of $S$, and
$\mathbf Z_{\ell,X}$ is not classical. Its reductions
$\mathbf Z/\ell^n$ are classical. This is what allows finite-level
\'etale arguments to be combined with a single pro-\'etale definition
of coefficients.
\end{example}

\subsection{Cohomology comparison}
\label{pe:derived-comparison-section}
For a bounded-below complex $K$ on $X_{\et}$, the pro-\'etale
cohomology of $\nu_X^*K$ is the \'etale cohomology of $K$, for every
scheme $X$; no bound on cohomological dimension enters. The proof rests
on the following acyclicity statement for the direct image
$\nu_{X,*}\colon\Mod(X_{\proet},\Lambda)\to\Mod(X_{\et},\Lambda)$.

\begin{lemma}[Acyclicity of pulled-back injectives]
\label{pe:injective-acyclicity}
If $I$ is an injective sheaf of $\Lambda$-modules on $X_{\et}$, then
$R^q\nu_{X,*}\nu_X^*I=0$ for $q>0$.
\end{lemma}
\begin{proof}
It suffices to show that positive cohomology vanishes on affine
\'etale objects. The restriction of $I$ to such an object is
injective, since restriction has an exact left adjoint, extension by
zero. On an affine base we use the basis of ind-\'etale affines: by
the compatible refinement of Theorem~\ref{pe:indetale-refinement},
the \v{C}ech-acyclicity criterion for a basis may be applied to finite
affine coverings by ind-\'etale maps.

Write each covering map as an inverse limit of affine \'etale maps
over the object being covered. By \eqref{pe:eq:affine-evaluation}, the
\v{C}ech complex of the covering is the filtered colimit of the
\v{C}ech complexes of these \'etale approximations. Each finitely
presented approximation descends to a common \'etale stage over the
base, is a covering after increasing that stage, and its nerve
descends by taking fibre products. Since $I$ is injective, each
augmented \v{C}ech complex at a finite stage has no positive
cohomology. Filtered colimits are exact, so the same holds for the
pro-\'etale covering, and the criterion gives the required vanishing.
This is the argument of \stacks{099U}; for module sheaves one uses
free $\Lambda$-module sheaves in the augmented resolution.
\end{proof}

\begin{theorem}[Bounded-below cohomology comparison]
\label{pe:cohomology-comparison}
For $K\in D^+(X_{\et},\Lambda)$ the unit of the adjunction between
$\nu_X^*\colon D^+(X_{\et},\Lambda)\to D^+(X_{\proet},\Lambda)$ and
$R\nu_{X,*}$ is an isomorphism
\begin{equation}\label{pe:eq:derived-unit}
 K\xrightarrow{\ \sim\ }R\nu_{X,*}\nu_X^*K\quad\text{in }D^+(X_{\et},\Lambda).
\end{equation}
Consequently $\nu_X^*\colon D^+(X_{\et},\Lambda)\to D^+(X_{\proet},\Lambda)$ is fully faithful,
and there is a canonical isomorphism in the derived category of
$\Lambda$-modules
\[
 R\Gamma(X_{\proet},\nu_X^*K)
   \simeq R\Gamma(X_{\et},K).
\]
More generally, for a pro-\'etale affine $U=\varprojlim_iU_i$ with
$U_i$ affine and \'etale over $X$,
\begin{equation}\label{pe:eq:derived-affine-continuity}
 R\Gamma(U_{\proet},(\nu_X^*K)|_U)
   \simeq \varinjlim_i R\Gamma((U_i)_{\et},K|_{U_i}).
\end{equation}
\end{theorem}
The colimit in \eqref{pe:eq:derived-affine-continuity} is a filtered
derived colimit, so its cohomology is the filtered colimit of the
cohomology groups.
\begin{proof}
Choose a bounded-below injective resolution $I^\bullet$ of $K$. By
Lemma~\ref{pe:injective-acyclicity}, $\nu_X^*I^\bullet$ is a
resolution by $\nu_{X,*}$-acyclic sheaves, so applying $\nu_{X,*}$
termwise computes $R\nu_{X,*}\nu_X^*K$. The underived unit is an
isomorphism on every term, so this complex is $I^\bullet$, which
proves \eqref{pe:eq:derived-unit} (\stacks{099V}). Composing derived
right adjoints gives the statement on global sections (\stacks{099W}).

For \eqref{pe:eq:derived-affine-continuity}, apply the global
sections statement to $U$ and then continuity of \'etale cohomology
under affine limits (\stacks{03Q6}). This gives the formula in each
degree for a sheaf. For a bounded-below complex only finitely many
terms of the hypercohomology spectral sequence contribute to each
total degree, and filtered colimits are exact.
\end{proof}

Write $D^+_{\mathrm{class}}(X_{\proet},\Lambda)$ for the full
subcategory of $D^+(X_{\proet},\Lambda)$ of classical complexes.

\begin{corollary}[The derived classical part]
\label{pe:classical-derived}
Classical module sheaves form a weak Serre subcategory, closed under
kernels, cokernels and extensions. There is an equivalence
\[
 D^+_{\mathrm{class}}(X_{\proet},\Lambda)
   \simeq D^+(X_{\et},\Lambda),
\]
with inverse functors $R\nu_{X,*}$ and $\nu_X^*$.
\end{corollary}
\begin{proof}
Kernels and cokernels are handled by exactness and full faithfulness
of $\nu_X^*$. Given an extension of two classical sheaves, apply
$\nu_{X,*}$; since $R^1\nu_{X,*}\nu_X^*F=0$, the result is a short
exact sequence whose pullback is the original extension.
By the derived unit every classical sheaf is $\nu_{X,*}$-acyclic, so
the bounded-below hypercohomology spectral sequence reduces the counit
$\nu_X^*R\nu_{X,*}K\to K$ in $D^+(X_{\proet},\Lambda)$ to the comparison for sheaves. Hence it is
an isomorphism for every bounded-below complex with classical
cohomology sheaves, whatever its terms (\stacks{09B1}).
\end{proof}

\subsection{Comparison of inverse and direct images}
\label{pe:functor-comparison-section}
Let $g:X\to Y$ be a morphism of schemes. Base change of site objects
gives the commutative square of topoi
\[
\begin{CD}
 X_{\proet}@>{\nu_X}>>X_{\et}\\
 @V{g_{\proet}}VV @VV{g_{\et}}V\\
 Y_{\proet}@>{\nu_Y}>>Y_{\et}.
\end{CD}
\]
It gives an isomorphism
$g_{\proet}^*\nu_Y^*\simeq\nu_X^*g_{\et}^*$ of functors
$\Mod(Y_{\et},\Lambda)\to\Mod(X_{\proet},\Lambda)$, compatible with
composition; since all four functors are exact, it is also an
isomorphism of functors $D(Y_{\et},\Lambda)\to D(X_{\proet},\Lambda)$.
Comparison of direct images needs a hypothesis on $g$.

\begin{theorem}[Direct-image comparison]
\label{pe:direct-image-comparison}
If $g$ is quasi-compact and quasi-separated, the canonical natural
transformation of functors $D^+(X_{\et},\Lambda)\to D^+(Y_{\proet},\Lambda)$
\begin{equation}\label{pe:eq:direct-image-comparison}
 \nu_Y^*Rg_{\et,*}K\longrightarrow
 Rg_{\proet,*}\nu_X^*K
\end{equation}
is an isomorphism for $K\in D^+(X_{\et},\Lambda)$. In particular,
direct image of a bounded-below classical pro-\'etale complex is
classical. The isomorphism is compatible with units, counits and
composition.
\end{theorem}
\begin{proof}
The map is adjoint to the composite of
$g_{\proet}^*\nu_Y^*\simeq\nu_X^*g_{\et}^*$ with the counit of the adjunction
$g_{\et}^*\dashv Rg_{\et,*}$ between $D^+(Y_{\et},\Lambda)$ and
$D^+(X_{\et},\Lambda)$; built from these adjunctions, it is
compatible with their units and counits. To show it is an
isomorphism, work over an affine open of $Y$ and evaluate on the
covering basis $V=\varprojlim_iV_i$, with $V_i$ affine and \'etale
over $Y$. For a sheaf $F$ the map in degree $q$ is the continuity map
\[
 \varinjlim_i H^q_{\et}(X\times_YV_i,F)
   \longrightarrow
 H^q_{\et}(X\times_YV,F).
\]
Since $g$ is quasi-compact and quasi-separated, so are its base
changes, and the map is an isomorphism by \stacks{03Q6}. Bounded-below
resolutions, or the convergent hypercohomology spectral sequences,
extend this from sheaves to $K$. This is the proof of \stacks{09A6}.
\end{proof}

For an open immersion $j:U\hookrightarrow X$, extension by zero
commutes with comparison:
\begin{equation}\label{pe:eq:jshriek-comparison}
 \nu_X^*j_{\et,!}\simeq j_{\proet,!}\nu_U^*\quad\text{as functors }\Mod(U_{\et},\Lambda)\to\Mod(X_{\proet},\Lambda).
\end{equation}
Indeed both functors are left adjoint to restriction followed by
$\nu_{U,*}$, which also identifies the adjunction maps (\stacks{09AD});
as the functors are exact, the identity also holds as functors
$D(U_{\et},\Lambda)\to D(X_{\proet},\Lambda)$.
For a closed immersion $i:Z\hookrightarrow X$, inverse-image comparison
gives $i^*\nu_X^*\simeq\nu_Z^*i^*$ as functors
$D(X_{\et},\Lambda)\to D(Z_{\proet},\Lambda)$, and
Theorem~\ref{pe:direct-image-comparison} gives the identity for
$Ri_*\colon D^+(Z_{\et},\Lambda)\to D^+(X_{\et},\Lambda)$ and its
pro-\'etale counterpart. These comparisons identify the open--closed
localization triangles on the two sites. For a separated morphism of
finite type $f:X\to Y$ with a compactification $f=\bar f\circ j$, where
$j$ is an open immersion and $\bar f$ is proper, they also identify
the functors $Rf_!\colon D^+(X_{\et},\Lambda)\to D^+(Y_{\et},\Lambda)$
and $Rf_!\colon D^+(X_{\proet},\Lambda)\to D^+(Y_{\proet},\Lambda)$,
both computed as $R\bar f_*j_!$.

\subsection{Finite local systems, constructibility and finite Tor amplitude}
\label{pe:constructibility-section}
The following notions are defined directly on $X_{\proet}$, as in
\S\ref{sec:setting}. A locally constant module sheaf is one
trivialized by a pro-\'etale covering; it is of finite presentation if
its constant models are finitely presented $\Lambda$-modules. For
Noetherian $\Lambda$, a sheaf is constructible if locally on $X$ a
finite partition into constructible locally closed strata makes it
locally constant of finite type on each stratum. We show that these
pro-\'etale objects are classical, so that they correspond to their
\'etale counterparts. The basic case is descent of torsors.

\begin{lemma}[Descent of classical-group torsors]
\label{pe:torsor-comparison}
Every pro-\'etale torsor under a classical sheaf of groups is
classical, and hence has a trivializing \'etale covering.
\end{lemma}
\begin{proof}
Write the group as $\nu_X^*G_0$ and work on an affine open of $X$.
Using Theorem~\ref{pe:indetale-refinement} and finite affine
refinements, refine a trivializing covering to a single faithfully
flat ind-\'etale affine $U\to X$, and write $U=\varprojlim_iU_i$ with
$U_i$ affine \'etale over $X$. A trivialization over $U$ gives a
transition element in $\nu_X^*G_0(U\times_XU)$. By affine continuity
it descends to $G_0(U_i\times_XU_i)$ for some $i$, and the cocycle
identity on the triple product holds after enlarging $i$. Since
$U\to X$ is surjective, each $U_i\to X$ is an \'etale covering, and
gluing the trivial torsor along this cocycle gives an \'etale torsor
with the prescribed pro-\'etale pullback. The same argument descends
trivializations (\stacks{099X}).
\end{proof}

\begin{proposition}[Locally constant sheaves of finite presentation]
\label{pe:finite-local-systems}
Every locally constant pro-\'etale $\Lambda_X$-module sheaf of finite
presentation is classical. Thus it has a trivializing \'etale
covering, and $\nu_X^*$ identifies the \'etale local systems of finite
presentation with these pro-\'etale ones. For a finite ring $A$ this includes every
finite local system of $A$-modules.
\end{proposition}
\begin{proof}
Decompose $X$ into the open pieces on which the constant model has a
fixed isomorphism type, and fix a finitely presented module $M$. By
finite presentation, the automorphism sheaf of $\underline M$ is the
constant discrete group $\operatorname{Aut}_\Lambda(M)$, so
$\mathit{Isom}(\underline M,F)$ is a torsor under this classical
group. By Lemma~\ref{pe:torsor-comparison} it has sections \'etale
locally; this is the descent argument of \stacks{099Y}. For a finite
model one can also descend its finite set of sections and transition
maps to finitely presented stages, and the fpqc sheaf condition then
descends the gluing.
\end{proof}

\begin{theorem}[Constructible comparison]
\label{pe:constructible-comparison}
Let $X$ be quasi-compact and quasi-separated. For a discrete
Noetherian $\Lambda$, every constructible
pro-\'etale $\Lambda_X$-module sheaf is classical. The functor
$\nu_X^*$ gives an equivalence from constructible \'etale sheaves of
$\Lambda$-modules to $\Cons(X,\Lambda)$, and an equivalence
$D^b_c(X_{\et},\Lambda)\to\Dbc(X,\Lambda)$ of bounded constructible
derived categories.
\end{theorem}
\begin{proof}
The preceding proposition handles a single stratum. For a finite
stratification argue locally, by induction on the number of strata.
At a given point, shrink to an affine open on which the stratum
containing it is closed, with inclusion $i$ and open complement $j$.
The exact sequence
\[
 0\longrightarrow j_!j^*F\longrightarrow F
   \longrightarrow i_*i^*F\longrightarrow0
\]
reduces the assertion to fewer strata, because $j_!$ and $i_*$ on
pro-\'etale module sheaves preserve classical sheaves by the functor comparisons and classical
sheaves are closed under extensions. The same argument with locally
closed constructible pieces covers the local qcqs definition. The
derived statement follows from Corollary~\ref{pe:classical-derived}
and bounded truncation. These comparisons are recorded in
\stacks{09AI}.
\end{proof}

\begin{corollary}[Finite Tor amplitude under comparison]
\label{pe:finite-tor-comparison}
For a qcqs scheme $X$ and a discrete Noetherian $\Lambda$, the preceding equivalence
$\nu_X^*\colon D^b_c(X_{\et},\Lambda)\to\Dbc(X,\Lambda)$ restricts to
\[
 \Dctf(X,\Lambda)\simeq
 D^b_{\mathrm{ctf}}(X_{\et},\Lambda).
\]
\end{corollary}
In particular the corollary applies to every finite commutative
$\ell$-primary ring $A$, regular or not.
\begin{proof}
The ring map for $\nu_X$ is an isomorphism, and inverse image is
exact and symmetric monoidal. Pullback of flat module sheaves is flat
(\stacks{05VD}), so the pullback of a bounded flat resolution
preserves the Tor bound. Conversely, let $K$ be an \'etale complex and
test $K\otimes_\Lambda^{\mathrm L}M$ against arbitrary \'etale module
sheaves $M$. Its inverse image is
$\nu_X^*K\otimes_{\Lambda_X}^{\mathrm L}\nu_X^*M$, and since $\nu_X^*$
is exact and fully faithful it detects vanishing of cohomology; so a
Tor bound on $\nu_X^*K$ gives the same bound on $K$. For constructible
complexes the bound can equivalently be checked on stalk complexes on
each stratum.
\end{proof}

\subsection{Unbounded complexes with finite coefficients}
\label{pe:unbounded-section}
Bounded-below comparison suffices for proper base change. Brown
representability works in the unbounded derived category and needs the
following extension. Over a field of finite cohomological dimension,
\'etale cohomology of finite type schemes is uniformly bounded, so
every computation in a fixed degree can be moved to a bounded-below
truncation.

\begin{definition}[Cohomological dimension]
\label{n:def-cd}
For a qcqs scheme $S$ and a qcqs morphism $g:X\to Y$, the
\emph{$\ell$-cohomological dimensions} are
\[
\begin{aligned}
 \operatorname{cd}_\ell(S)&=\sup\{q: H^q(S_{\proet},F)\neq0
     \text{ for some classical $\ell$-power torsion sheaf $F$}\},\\
 \operatorname{cd}_\ell(g)&=\sup\{q: R^qg_*F\neq0
     \text{ for some classical $\ell$-power torsion sheaf $F$ on }X_{\proet}\}.
\end{aligned}
\]
For a field $K$ we write $\operatorname{cd}_\ell(K)=\operatorname{cd}_\ell(\operatorname{Spec}K)$.
\end{definition}

By Theorems~\ref{pe:cohomology-comparison} and~\ref{pe:direct-image-comparison},
these are the usual $\ell$-cohomological dimensions computed on the
\'etale sites, so the classical bounds apply. If $K$ is a field with $\operatorname{cd}_\ell(K)<\infty$, then
$\operatorname{cd}_\ell(S)$ and $\operatorname{cd}_\ell(g)$ are finite for all
schemes $S$ and morphisms $g$ of finite type over $K$ (\stacks{0F10}).
For our field $k$ the hypothesis holds by the following lemma.

\begin{lemma}\label{b:lem:base-field-cd}
If $K$ is an algebraically closed field, then $\operatorname{cd}_\ell(K)=0$.
If $K$ is a finite field, then $\operatorname{cd}_\ell(K)\le1$.
\end{lemma}
\begin{proof}
\'Etale sheaves on $\operatorname{Spec}K$ are the same as discrete
modules over the absolute Galois group, and \'etale cohomology is Galois
cohomology (\stacks{03QQ}). For an algebraically closed field the Galois
group is trivial, so $\operatorname{cd}_\ell(K)=0$. Let $K$ be finite with $q$ elements. Its finite extensions are the
fields $\mathbf F_{q^n}$ of roots of $T^{q^n}-T$, with cyclic Galois group
generated by Frobenius of order exactly $n$ (a smaller order $d$ would
give $q^n$ roots of $T^{q^d}-T$). Hence the absolute Galois group is
$\widehat{\mathbf Z}$. We show $H^s(\widehat{\mathbf Z},M)=0$ for $s\geq2$
and every torsion discrete module $M$; this gives $\operatorname{cd}_\ell(K)\le1$.

\emph{Continuous cochains compute derived functors.} For a profinite
group $G$ and a discrete $G$-module $M$, continuous inhomogeneous
cochains form a $\delta$-functor: a continuous cochain has finite image,
so it lifts through a surjection of modules by lifting its finitely many
values. It is effaceable. Embed $M$ in the coinduced module
$I=\operatorname{Map}_{\mathrm{cont}}(G,M)$ by $m\mapsto(g\mapsto gm)$, with
$(hf)(g)=f(gh)$; this action is discrete because a locally constant
function on $G$ is invariant under an open subgroup. In homogeneous form,
$I$-valued equivariant cochains $F$ correspond to continuous functions
$f:G^{r+1}\to M$ via $F(g_0,\ldots,g_r)(t)=f(tg_0,\ldots,tg_r)$, with the
alternating-sum differential, and inserting $1$ as a first coordinate is
a contracting homotopy in positive degrees. By \stacks{010T}, continuous
cochain cohomology is therefore the right derived functor of invariants.

\emph{Reduction to finite cyclic groups.} A torsion discrete
$\widehat{\mathbf Z}$-module is the union of its finite stable submodules
(the orbit of an element is finite, and its elements have a common
finite order), and continuous cochains commute with this union since they
have finite image. So let $M$ be finite, with action through
$C_m=\mathbf Z/m$. A continuous function on a power of $\widehat{\mathbf Z}$
factors through a finite quotient, so the cochain complex of
$\widehat{\mathbf Z}$ is the colimit, under inflation, of those of
$C_{m'}$ for multiples $m'$ of $m$.

\emph{Inflation kills high degrees.} For $C_m=\langle\sigma\rangle$, the
periodic free resolution of $\mathbf Z$ over $\mathbf Z[C_m]$ has
differentials alternating between $\sigma-1$ and
$N_m=1+\sigma+\cdots+\sigma^{m-1}$; it is exact, because the kernel of
$\sigma-1$ is spanned by $N_m$ and the kernel of $N_m$ consists of the
elements with coefficient sum zero, which form the image of $\sigma-1$.
For $C_{rm}\to C_m$, the quotient map in degrees $0,1$ and $r^{\lfloor
s/2\rfloor}$ times the quotient map in degree $s$ is a chain map, because
$N_{rm}\mapsto rN_m$. It lifts the identity of $\mathbf Z$, so it
computes inflation (two lifts are homotopic, by lifting recursively
through the exact free resolution). On cohomology with values in $M$,
inflation in degree $s$ is therefore multiplication by $r^{\lfloor
s/2\rfloor}$, which is zero for $s\geq2$ when the exponent of $M$ divides
$r$. Hence the colimit vanishes in degrees $s\geq2$.
\end{proof}

\begin{proposition}[Unbounded comparison for finite coefficients]
\label{pe:unbounded-classical}
Let $X$ and $Y$ be of finite type over the standing field $k$, let
$A$ be a finite commutative $\ell$-primary ring, and let
$g:X\to Y$ be of finite type. Then
$\nu_X^*\colon D(X_{\et},A)\to D(X_{\proet},A)$ is fully faithful and
preserves arbitrary direct sums. The map
\eqref{pe:eq:direct-image-comparison}, now a natural transformation of
functors $D(X_{\et},A)\to D(Y_{\proet},A)$, is an isomorphism for every
$K\in D(X_{\et},A)$.
\end{proposition}
\begin{proof}
We first record the cohomological bounds. Every affine $U_i$ \'etale
over $X$ is affine of finite type over $k$ with
$\dim(U_i)\leq\dim(X)$, so \stacks{0F0V} gives
$\operatorname{cd}_\ell(U_i)\leq\dim(X)+\operatorname{cd}_\ell(k)$ for
$A$-module sheaves; this bound is independent of $i$. A bound of the
same kind holds for the fixed finite-type morphism $g$ and its
\'etale base changes (\stacks{0F10}). For the schemes $X\times_YV_i$
used below, choose a fixed finite affine open cover of $X$. Since
$X$ and $Y$ are separated over $k$, its intersections and their
pullbacks to the affine $V_i$ are affine; each pullback is \'etale
over the corresponding intersection and has dimension at most
$\dim(X)$. The finite \v{C}ech complex of this cover therefore adds a
bound independent of $i$ to the affine bound. On a pro-\'etale affine
limit the bound for classical sheaves follows from cohomology
comparison and continuity. All these bounds hold for arbitrary
$A$-module sheaves.

Next, the pro-\'etale truncation argument. Let $U$ be an affine object
of $X_{\proet}$, and let $L$ be a pro-\'etale complex with classical
cohomology and $H^q(L)=0$ for $q>b$. Left completeness gives
\[
 L\simeq R\varprojlim_n\tau_{\geq-n}L.
\]
For each bounded-below term, derived sections over $U$ have cohomology
only up to degree $b+\dim(X)+\operatorname{cd}_\ell(k)$. Derived sections
preserve the displayed inverse limit, so by the Milnor sequence
$R\Gamma(U,L)$ has no cohomology above $b+\dim(X)+\operatorname{cd}_\ell(k)+1$.
Apply this to $L=\tau_{<N}\nu_X^*K$: for any prescribed degree $m$,
choosing $N<m-\dim(X)-\operatorname{cd}_\ell(k)-2$ makes
\[
 H^mR\Gamma(U,\nu_X^*K)\longrightarrow
 H^mR\Gamma(U,\tau_{\geq N}\nu_X^*K)
\]
an isomorphism. The choice of $N$ depends only on $m$ and the bound,
so the argument applies to every $K\in D(X_{\et},A)$, bounded below or
not.

On the \'etale side, $\Gamma(U,-)$ on $A$-module sheaves has enough
injective acyclic objects and cohomological dimension at most
$\dim(X)+\operatorname{cd}_\ell(k)$, so
the derived-functor lemma for bounded cohomological dimension
(\stacks{07K7}) gives the same truncation for arbitrary complexes
directly. Thus in each degree the unit $K\to R\nu_{X,*}\nu_X^*K$ in
$D(X_{\et},A)$ is computed on a bounded-below truncation, where it is an isomorphism; hence the unbounded unit is an
isomorphism.

For direct images, evaluate on $V=\varprojlim_iV_i$ in the affine
pro-\'etale basis of $Y$. For a classical sheaf, cohomology on
$X\times_YV$ is the filtered colimit of cohomology on the
$X\times_YV_i$, and the uniform bounds allow the same fixed-degree
truncation on both sides. Explicitly, use the triangle
$\tau_{<N}K\to K\to\tau_{\geq N}K$. The relative bound for $g$ places
$Rg_{\et,*}\tau_{<N}K$ in degrees at most
$N-1+\operatorname{cd}_\ell(g)$, and the absolute bounds for the $V_i$
control its derived sections. On the pro-\'etale side, the absolute
bounds on $X\times_YV_i$ and continuity give the corresponding bound on
$X\times_YV$, up to the one extra Milnor degree found above. For $N$
sufficiently small both lower-tail contributions vanish in the
prescribed degree, and bounded-below direct-image comparison gives the
comparison for unbounded $K$. Finally, $\nu_X^*$ preserves sums
because it is an exact left adjoint.
\end{proof}

Classical sheaves are closed under arbitrary sums, since $\nu_X^*$ is
exact and preserves colimits (Theorem~\ref{pe:classical-sheaves}).
With the weak Serre property, it follows that cones and sums of
complexes with classical cohomology have classical cohomology, so the
localizing subcategory of $D(X_{\proet},A)$ generated by $\Dctf(X,A)$
lies in the classical part. That this localizing category is
identified with the whole \'etale derived category is the
compact-generation statement of Proposition~\ref{ctf:finite-comparison}.

\subsection{The geometric point and continuous cohomology}
\label{pe:coefficient-bridge}

The comparison theorems above concern discrete coefficients. Integral
and rational $\ell$-adic coefficients are defined directly on the
pro-\'etale site, and the corresponding discrete constant sheaves give
the wrong cohomology. Both facts are visible already over a separably
closed field.

\begin{proposition}[The pro-\'etale geometric point]
\label{pe:condensed-point}
Let $\Omega$ be a separably closed field. The affine pro-\'etale
basis of $\operatorname{Spec}\Omega$ is equivalent to the category
of profinite sets, with finite jointly surjective families as covers.
Its topos is consequently the topos of condensed sets. A discrete
scalar ring $A$ gives the classical constant sheaf, while a profinite
or locally profinite coefficient ring gives its sheaf of continuous
functions on these profinite test objects.
\end{proposition}
\begin{proof}
A finitely presented \'etale $\Omega$-algebra is a finite product of
copies of $\Omega$, so filtered colimits of such algebras correspond,
under $\operatorname{Spec}$, to inverse limits of finite sets. By the
affine pro-\'etale basis theorem sheaves are determined on this basis.
Faithfully flat covering maps become surjections of profinite sets,
and quasi-compactness gives the finite-cover condition. This proves
the equivalence of sites.
\end{proof}

Thus the \'etale topos of $\Omega$ is the topos of sets, while its
pro-\'etale topos retains profinite test objects. The final object is
weakly contractible, so global sections are exact, but evaluation on
that one object does not determine a condensed sheaf. Geometric stalks
suffice for classical sheaves; for general pro-\'etale module sheaves
we evaluate on the cover-dense basis of weakly contractible affines.
Example~\ref{pe:nonclassical-example} shows this for the completed
$\mathbf Z_\ell$ coefficient sheaf, whose sections on profinite tests
are all continuous functions; its finite reductions are classical,
and the adic theory uses exactly this combination.

\begin{proposition}[Continuous cohomology as pro-\'etale cohomology]
\label{pe:continuous-cohomology}
Let $X$ be a scheme and $(F_n)$ an inverse system of abelian sheaves
on $X_{\proet}$ which are classical, with epimorphic transition maps.
Put $F=\varprojlim_nF_n$ in the pro-\'etale topos. Then, in the
derived category of abelian groups,
\[
 R\Gamma(X_{\proet},F)
    \simeq R\varprojlim_n R\Gamma(X_{\proet},F_n).
\]
If $F_n=\nu_X^*F_{n,\et}$, this is the continuous \'etale cohomology
of the system $(F_{n,\et})$: the right derived functor of
$(G_n)\mapsto\Gamma(X_{\et},\varprojlim_nG_n)$, taken on the
abelian category of inverse systems.
\end{proposition}
\begin{proof}
By repleteness $F\simeq R\varprojlim_nF_n$, and derived global
sections, being a right adjoint, commute with derived inverse limits.
The bounded-below cohomology comparison identifies each
$R\Gamma(X_{\proet},F_n)$ with $R\Gamma(X_{\et},F_{n,\et})$. On the
category of inverse systems, the derived functor of sections followed
by inverse limit is $R\varprojlim_nR\Gamma(X_{\et},F_{n,\et})$; one
may compute it with an injective resolution in the diagram category,
or use the derived adjunctions to the constant-diagram functor. This
is \cite[Proposition~5.6.2]{BS}.
\end{proof}

For every $q$ there is a natural exact sequence
\begin{equation}\label{pe:eq:continuous-milnor}
0\longrightarrow\varprojlim
 {}^1_nH^{q-1}(X_{\proet},F_n)
\longrightarrow H^q(X_{\proet},F)
\longrightarrow\varprojlim_nH^q(X_{\proet},F_n)
\longrightarrow0.
\end{equation}
The first term need not vanish, since epimorphic transitions between
the sheaves $F_n$ do not make the transitions between their
cohomology groups surjective. It vanishes when the groups
$H^{q-1}(X_{\proet},F_n)$ satisfy the Mittag--Leffler condition, for
instance when they are finite, since descending chains of images then
stabilize at each level.

\begin{example}[The first cohomology of a finite field]
\label{pe:finite-field-adic}
For $X=\operatorname{Spec}\mathbb F_q$ and $\ell\nmid q$,
\[
 H^1(X_{\proet},\mathbf Z_{\ell,X})\simeq\mathbf Z_\ell,
 \qquad H^1(X_{\proet},\mathbf Q_{\ell,X})\simeq\mathbf Q_\ell.
\]
Indeed the finite-level first cohomology is
$\operatorname{Hom}_{\mathrm{cont}}(\widehat{\mathbf Z},
\mathbf Z/\ell^n)=\mathbf Z/\ell^n$, and the degree-zero groups are
finite, so their $\varprojlim^1$ vanishes. Localization commutes with
bounded-below cohomology on the qcqs scheme $X$ (\stacks{09A3}), which
gives the rational statement. For the discrete constant sheaf,
however, $H^1(X_{\et},\mathbf Q_\ell^{\mathrm{disc}})=0$: a
continuous homomorphism from a profinite group to the discrete
torsion-free additive group has finite, hence zero, image.
\end{example}
\subsection{Open--closed operations and derived completion}
\label{pe:localization-completion}

On the pro-\'etale site, inverse image along a closed immersion
preserves all small limits. The six-functor proof uses this feature
before any exceptional inverse image or proper base change theorem is
available.

\begin{proposition}[Open--closed operations on module sheaves]
\label{pe:open-closed}
Let $X$ be quasi-compact and quasi-separated, $j:U\hookrightarrow X$
a quasi-compact open immersion, and $i:Z\hookrightarrow X$ its
closed complement. For a sheaf of rings $B$ on $X_{\proet}$, use its
inverse images as the coefficient rings on $U$ and $Z$. Thus
$i^*$ goes from $B$-module sheaves on $X_{\proet}$ to $i^*B$-module
sheaves on $Z_{\proet}$, $i_*$ goes back, and $j^*$, $j_!$ relate
$B$-module sheaves on $X_{\proet}$ and $j^*B$-module sheaves on
$U_{\proet}$.
Then $i^*$ is exact and preserves all small limits and colimits;
$i_*$ is exact and fully faithful. The functors $j^*$ and $j_!$ are
exact, and every $B$-module sheaf $M$ has the exact sequence
\[
 0\longrightarrow j_!j^*M\longrightarrow M
   \longrightarrow i_*i^*M\longrightarrow0.
\]
Moreover $j_!$ preserves products and derived inverse limits of
module complexes. These operations commute with base change in
their usual open--closed squares.
\end{proposition}
\begin{proof}
The key limit assertion is local on $X$. A weakly contractible affine
test object $V$ in the affine ind-\'etale basis of $Z$ has a canonical
henselian lift $v(V)$ in the corresponding basis of $X$: the functor
$v$ is the fully faithful left adjoint to base change of these affine
bases, and $v(V)\times_XZ\simeq V$ (\stacks{09A9},
Lemmas~61.25.1--61.25.2). The pullback of a sheaf, evaluated on $V$,
is its evaluation on $v(V)$, because no sheafification is needed on a
weakly contractible object. Evaluation commutes with limits and these
test objects detect isomorphisms, so $i^*$ preserves limits. As a left
adjoint it preserves colimits, and it is exact. Lifting covers through
the henselian thickenings proves exactness of $i_*$ and the identities
$i^*i_*=\mathrm{id}$ and $j^*i_*=0$. The open--closed sequence follows
by the same cover descent, or from its kernel description.

On derived categories this gives, for $K\in D(U_{\proet},j^*B)$, the
triangle
\[
 j_!K\longrightarrow Rj_*K\longrightarrow i_*i^*Rj_*K
       \longrightarrow j_!K[1]\quad\text{in }D(X_{\proet},B).
\]
Its last two terms preserve products and derived inverse limits, hence
so does their fibre $j_!$. Exactness and exact products give the
corresponding limit statements on module sheaves. This is the
henselian-thickening argument of Bhatt--Scholze, \S6.1; see also
\stacks{09AA}, \stacks{09AH}, \stacks{09BM} and \stacks{09BP}.
\end{proof}

Closed-support arguments later use the local cohomology functor
\[
 R\underline\Gamma_ZK
   =i^*\operatorname{fib}(K\longrightarrow Rj_*j^*K)
     \in D(Z_{\proet},i^*B).
\]
As a functor $D(X_{\proet},B)\to D(Z_{\proet},i^*B)$ it is right
adjoint to the exact fully faithful functor
$i_*\colon D(Z_{\proet},i^*B)\to D(X_{\proet},B)$, since
$j^*i_*M=0$, and its definition needs only the open--closed operations
above. Once $i^!$ has been constructed by Brown representability, it
agrees with $R\underline\Gamma_Z$ on the appropriate localizing
category.

\begin{definition}[Derived adic completeness]
\label{pe:derived-completion}
Let $E/\Ql$ be finite, with uniformizer $\pi_E$, and let
$\mathcal O_{E,X}=p_X^*\underline{\mathcal O_E}$ and
$(\mathcal O_E/\pi_E^n)_X$ be the integral and finite coefficient
sheaves of \S\ref{n:coefficients}. A complex $K\in D(X_{\proet},\mathcal O_E)$ is
\emph{derived complete} if its canonical map
\[
 K\longrightarrow R\varprojlim_n
       (K\otimes^{\mathbf L}_{\mathcal O_{E,X}}(\mathcal O_E/\pi_E^{n})_X)
\]
is an isomorphism. For $K_n=K\otimes^{\mathbf L}_{\mathcal O_{E,X}}(\mathcal O_E/\pi_E^{n})_X$,
the canonical transitions satisfy the derived reduction identities
$K_{n+1}\otimes^{\mathbf L}_{(\mathcal O_E/\pi_E^{n+1})_X}(\mathcal O_E/\pi_E^{n})_X\simeq K_n$;
a compatible tower with these identities is called \emph{normalized}.
\end{definition}
Normalization is a property of the complexes $K_n$; their cohomology
sheaves $H^q(K_n)$ need not form a normalized tower under the ordinary
tensor product.

\begin{lemma}[Reduction and derived Nakayama]
\label{pe:derived-nakayama}
The finite quotient coefficient has the perfect resolution
$[\mathcal O_{E,X}\xrightarrow{\pi_E^n}\mathcal O_{E,X}]$, so $K_n$ is the cone of
$\pi_E^n:K\to K$. If a derived-complete complex has $K_1=0$, it is
zero. Consequently a morphism between derived-complete complexes is
an isomorphism whenever its first derived reduction is an isomorphism.
\end{lemma}
\begin{proof}
The coefficient sequence is exact on the profinite point, hence after
pullback, and it supplies the two-term resolution. The short exact
sequences between successive powers of $\pi_E$ show inductively that
$K_1=0$ implies $K_n=0$ for all $n$, and completeness then gives
$K=0$. Derived completeness is closed under cones, so applying this to
the cone proves the last assertion.
\end{proof}

For a complete $K$, comparison at finite levels gives
\[
 R\Gamma(X_{\proet},K)
    \simeq R\varprojlim_nR\Gamma(X_{\proet},K_n).
\]
When the $K_n$ are classical bounded-below complexes, each term on the
right can be computed by the \'etale comparison already proved, so
the normalized tower $(K_n)$ computes the cohomology of $K$. For Noetherian $X$, the integral category $\Dctf(X,\mathcal O_E)$
consists of the derived-complete complexes $K$ with $K_1$ constructible
of finite Tor amplitude (Proposition~\ref{n:integral-criterion};
\stacks{09C1}). The map from the constant sheaf $\mathcal O_E$ to $\mathcal O_{E,X}$ is
flat and an isomorphism modulo every $\pi_E^n$, so restriction of scalars
identifies their derived-complete categories (\stacks{099I}); the
identification concerns derived-complete objects only.

Inverting $\ell$ and passing to the colimit over $E$ give the rational
coefficient sheaves $\mathcal O_{E,X}[1/\ell]$ and $\varinjlim_E E_X$. Their
constructible objects have finite stratumwise perfect models.
Section~\ref{b:sec:coefficients} proves the lattice and
finite-extension descent arguments that compare them with the classical
adic categories. These comparisons concern constructible objects only.

\section{Coefficients, constructible complexes and the main theorem}
\label{sec:setting}

\subsection{The coefficient rings}
\label{n:coefficients}
We define each coefficient ring once, on a point, and pull it back to
every scheme. The point is the site of profinite sets described below;
Proposition~\ref{pe:condensed-point} explains its geometric
interpretation and how it differs from the ordinary \'etale point.
Pulling back gives $A_X$, $\mathcal O_{E,X}$, $E_X$ and $\QbarX$
by one construction. Here $A$ is a finite commutative $\ell$-primary
ring with its discrete topology; its pullback is the usual constant
coefficient sheaf.

\subsubsection{Condensed rings}
Let $*_{\proet}$ be the site of profinite sets, with finite jointly
surjective families as covers. A \emph{condensed ring} is a sheaf of
rings on $*_{\proet}$. A topological ring $V$ gives a condensed ring
$\underline V$, whose sections over $S$ are the continuous maps $S\to V$;
this is a sheaf because a surjection of profinite sets is a quotient map.

Fix an algebraic closure $\Qbar$ of $\Ql$. Throughout, $E$ denotes a
finite extension of $\Ql$ inside $\Qbar$, with ring of integers
$\mathcal O_E$ and uniformizer $\pi_E$. This gives condensed rings
$\underline{\mathcal O_E}$ and $\underline E$ for the $\ell$-adic
topologies. For $\Qbar$ we take the filtered colimit of the finite
coefficient fields. Explicitly, with $C(S,\Qbar)$ the continuous maps
for the valuation topology, put
\[
 \underline{\Qbar}(S)=
 \{a\in C(S,\Qbar):a(S)\subset E\text{ for some finite }E/\Ql\}
 =\varinjlim_E\underline E(S).
\]
This is already a sheaf: a cover of a profinite set can be refined by a
finite one, and finitely many finite extensions of $\Ql$ lie in a common
one. Since
filtered colimits of sheaves on $*_{\proet}$ are computed sectionwise,
$\underline{\Qbar}=\varinjlim_E\underline E$ as condensed rings.

This is smaller than $\underline V$ for $V=\Qbar$ with the valuation
topology. Algebraic numbers of unbounded degree can tend to $0$, and
such a sequence is a continuous map from a convergent sequence to $\Qbar$
with image in no finite extension. We exclude these maps because we want
a theory over finite coefficient fields: with our choice, every
constructible $\Qbar$-complex and every morphism between such complexes
descends to a finite $E$ (Lemma~\ref{b:lem:coefficient-finite-descent}).

\subsubsection{Pulling back to a scheme}
A profinite set $S$ defines a pro-\'etale scheme
$\operatorname{Spec}LC(S,\mathbf Z)$ over $\mathbf Z$, where
$LC(S,\mathbf Z)$ is the ring of locally constant functions. The functor
$S\mapsto\operatorname{Spec}LC(S,\mathbf Z)$ is continuous and preserves
finite limits, as one checks on finite quotients of $S$, so it induces a
morphism of topoi $\Sh((\operatorname{Spec}\mathbf Z)_{\proet})\to
\Sh(*_{\proet})$ \stacks{00X6}. Composing with the structure morphism of
$X$ we obtain
\[
 p_X:\Sh(X_{\proet})\longrightarrow\Sh(*_{\proet}),
\]
and for a condensed ring $\Lambda$ we write $\Lambda_X=p_X^*\Lambda$.
For a morphism $f:X\to Y$ we have $p_Y\circ f=p_X$, hence a canonical
isomorphism $f^{-1}\Lambda_Y\simeq\Lambda_X$, compatible with identities
and composition. Likewise a morphism $\Lambda\to\Lambda'$ of condensed
rings pulls back to $\Lambda_X\to\Lambda'_X$, which defines extension of
scalars $\Lambda'_X\otimes_{\Lambda_X}-$ from $\Lambda_X$-modules to
$\Lambda'_X$-modules on $X_{\proet}$. So one construction handles both pullback
along morphisms and change of coefficients.

The next lemma identifies these rings with the completed coefficient
rings of classical $\ell$-adic cohomology.

\begin{lemma}[Comparison with completed coefficients]
\label{n:coefficient-comparison}
Write $\mathcal O_{E,X}$, $E_X$ and $\QbarX$ for the pullbacks to
$X_{\proet}$ of $\underline{\mathcal O_E}$, $\underline E$ and
$\underline{\Qbar}$. There are canonical isomorphisms of sheaves of rings
\[
 \mathcal O_{E,X}\cong\varprojlim_n(\mathcal O_E/\pi_E^n)_X,
 \qquad E_X\cong\mathcal O_{E,X}[1/\ell],
 \qquad \QbarX\cong\varinjlim_E E_X,
\]
where $(\mathcal O_E/\pi_E^n)_X$ is the constant sheaf. They are
compatible with inverse image along every morphism of schemes.
\end{lemma}
\begin{proof}
On the point, $\underline{\mathcal O_E}$ is represented by the profinite
set $\mathcal O_E=\varprojlim_n\mathcal O_E/\pi_E^n$, so its pullback to
$X$ is represented by the pro-\'etale $X$-scheme
$\varprojlim_n\coprod_{\mathcal O_E/\pi_E^n}X$. By the Yoneda lemma the
sheaf this scheme represents is $\varprojlim_n(\mathcal O_E/\pi_E^n)_X$,
and the ring structures match. Working with the representing scheme is
what makes this easy: $p_X^*$ itself does not commute with inverse limits
in general.

A continuous map from a profinite set to $E$ has compact image, hence
lands in $\ell^{-m}\mathcal O_E$ for some $m$; thus
$\underline E=\underline{\mathcal O_E}[1/\ell]$ on the point. Since $p_X^*$
commutes with colimits, and inverting $\ell$ is a filtered colimit, we get
the second isomorphism, and the third in the same way. All three are built
from $p_X^*$, so they are compatible with pullback.
\end{proof}

The \emph{admissible coefficients} for the six-functor theorem are
$\Lambda\in\{A,\Zl,\mathcal O_E,\Ql,E,\Qbar\}$, where $A$ is a
finite commutative ring killed by a power of $\ell$ and $E/\Ql$ is
finite, including $E=\Ql$. All are regarded as condensed rings,
with coefficient sheaf $\Lambda_X=p_X^*\Lambda$ on $X_{\proet}$.
For $E=\Ql$ the preceding comparison shows that
$\mathbf Z_{\ell,X}=\mathcal O_{\Ql,X}$ and $\mathbf Q_{\ell,X}=(\Ql)_X$
satisfy
\[
 \mathbf Z_{\ell,X}\simeq\varprojlim_n(\mathbf Z/\ell^n)_X,
 \qquad
 \mathbf Q_{\ell,X}\simeq\mathbf Z_{\ell,X}[1/\ell].
\]
The main theorem, all six operations, the coefficient comparison,
proper base change and smooth purity are stated for this whole family.
Sections that treat only finite, integral or rational coefficients
state their scope. In the integral proofs $E=\Ql$ gives $\Zl$, and in the rational
proofs it gives $\Ql$. The categories of constructible complexes used
for every admissible $\Lambda$ are defined in
\S\ref{n:constructible}.
We write $\Mod(X_{\proet},\Lambda)$ for the category of
$\Lambda_X$-modules and $D(X_{\proet},\Lambda)$ for its derived category.

\begin{definition}[Tate twists]\label{def:sp-Tate-twist}
Let $\mu_{\ell^n,X}$ be the sheaf of $\ell^n$th roots of unity on $X_{\proet}$;
since $\ell$ is invertible on $X$, it is locally isomorphic to
$(\mathbf Z/\ell^n\mathbf Z)_X$. For a finite ring $A$ with $\ell^rA=0$ put
$A_X(1)=\mu_{\ell^r,X}\otimes_{\mathbf Z/\ell^r\mathbf Z}A_X$, which by the
Kummer maps does not depend on $r$. Put
$\mathcal O_{E,X}(1)=\varprojlim_n(\mathcal O_E/\ell^n)_X(1)$, and obtain
$E_X(1)$ and $\QbarX(1)$ by inverting $\ell$ and passing to the colimit
over $E$. Each $\Lambda_X(1)$ is locally free of rank one. For $d\ge0$,
$\Lambda_X(d)$ is the $d$-fold tensor power, $\Lambda_X(-d)$ its dual, and
$K(d)=K\otimes^{\mathbf L}_{\Lambda_X}\Lambda_X(d)$.
\end{definition}

\subsection{Derived functors of module sheaves}
\label{n:ambient}
For every condensed ring $\Lambda$ and $f:X\to Y$, the operations
$f^*\colon D(Y_{\proet},\Lambda)\to D(X_{\proet},\Lambda)$,
$Rf_*\colon D(X_{\proet},\Lambda)\to D(Y_{\proet},\Lambda)$, and
$\otimes^L$ and internal Hom on $D(X_{\proet},\Lambda)$, are derived
functors defined on the whole unbounded derived categories. The
remaining two operations need more: $Rf_!$ is built from
proper direct image and extension by zero
(\S\ref{b:sec:brown-construction}), and its right adjoint comes from
Brown representability.

Sheaves of $\Lambda_X$-modules can be viewed as module objects over the
ring object $\Lambda_X$ in $\Sh(X_{\proet},\Ab)$, with the tensor product
over $\Lambda_X$ given by the coequalizer of the two actions
$M\otimes\Lambda_X\otimes N\rightrightarrows M\otimes N$. The category of
$\Lambda_X$-modules is Grothendieck abelian, so complexes have K-injective
resolutions, which define $Rf_*$ and $\cHom$ (\stacks{07A5},
\stacks{08J7}), and K-flat resolutions, which define $\otimes^L$ with
$K\otimes^L-$ left adjoint to $\cHom(K,-)$ (\stacks{06YL}).

The pullback of a $\Lambda_Y$-module $M$ along $f:X\to Y$ is
$f^*M=\Lambda_X\otimes_{f^{-1}\Lambda_Y}f^{-1}M$. Since
$f^{-1}\Lambda_Y\simeq\Lambda_X$, this is just $f^{-1}M$. So $f^*\colon\Mod(Y_{\proet},\Lambda)\to
\Mod(X_{\proet},\Lambda)$ is exact and needs no deriving, and on derived
categories it is left adjoint to $Rf_*\colon D(X_{\proet},\Lambda)\to
D(Y_{\proet},\Lambda)$. For an open immersion $j:U\to X$, extension by
zero $j_!\colon\Mod(U_{\proet},\Lambda)\to\Mod(X_{\proet},\Lambda)$ is
exact and left adjoint to $j^*$.

Passing from $\mathcal O_E$ to $E$ means inverting $\ell$, and on derived
categories this is a Verdier localization:

\begin{proposition}[Rationalization]
\label{n:rationalization}
Extension of scalars $E_X\otimes_{\mathcal O_{E,X}}-\colon
D(X_{\proet},\mathcal O_E)\to D(X_{\proet},E)$ identifies
$D(X_{\proet},E)$ with the Verdier quotient of $D(X_{\proet},\mathcal O_E)$
by the complexes $K$ with $H^q(K)[1/\ell]=0$ for all $q$.
\end{proposition}
\begin{proof}
Inverting $\ell$ is exact, so $E_X$ is flat over $\mathcal O_{E,X}$ and
$H^q(E_X\otimes K)=H^q(K)[1/\ell]$. Restriction of scalars is right
adjoint to extension of scalars, and the counit is an isomorphism because
$\ell$ already acts invertibly on an $E_X$-module. So restriction is fully
faithful, and extension of scalars is the localization at its kernel.
\end{proof}

The kernel is larger than one might first guess:
$\bigoplus_{n\ge1}\mathcal O_{E,X}/\ell^n$ lies in it, although no single
power of $\ell$ kills it. On bounded constructible complexes this cannot
happen. There a complex killed by inverting $\ell$ is killed by a fixed
power of $\ell$, and extension of scalars induces an equivalence
\[
 \Dbc(X,\mathcal O_E)[\ell^{-1}]\xrightarrow{\ \sim\ }\Dbc(X,E),
\]
where the left side has the objects of $\Dbc(X,\mathcal O_E)$ and Hom
groups tensored with $E$ \cite[Proposition~6.8.14(2)]{BS}. The Hom part of this statement is
proved in \eqref{eq:sp-localized-Hom}. For essential surjectivity,
Lemma~\ref{b:lem:integral-models}(4) gives integral generators. By the
Hom statement, a connecting morphism between them lifts to an integral
one after multiplying by a power of $\ell$, so finite cones of the
generators also have integral models.

\subsection{Constructible complexes}
\label{n:constructible}
The definition of a constructible sheaf makes sense for any condensed
ring. Recall that for $g:Z\to X$ the restriction $M|_Z=g^*M$ of a
$\Lambda_X$-module is a $\Lambda_Z$-module, since
$g^{-1}\Lambda_X\simeq\Lambda_Z$; thus
$g^*\colon\Mod(X_{\proet},\Lambda)\to\Mod(Z_{\proet},\Lambda)$.

\begin{definition}[Local models and constructible sheaves]
\label{n:def-constructible}\label{n:def-field-constructible}
Let $\Lambda$ be a commutative condensed ring.
\begin{enumerate}
\item A $\Lambda_X$-module $M$ has a \emph{finite constant local model}
  if there is a pro-\'etale cover $\{U_i\to X\}$ and finitely presented
  modules $P_i$ over the scalar ring $\Lambda(*)$ such that
  \[
    M|_{U_i}\simeq\Lambda_{U_i}\otimes_{\Lambda(*)}P_i.
  \]
  Here $\Lambda(*)$ and $P_i$ are constant sheaves on $U_i$, acting
  through $\Lambda(*)\to\Lambda_{U_i}$. Equivalently, on each $U_i$ the
  module $M$ has a presentation by a finite matrix with entries in the
  scalar ring $\Lambda(*)$.
\item For qcqs $X$, a $\Lambda_X$-module $M$ is \emph{constructible} if
  there is a finite stratification $X=\coprod_aX_a$ by constructible
  locally closed subschemes such that each $M|_{X_a}$ has a finite
  constant local model. On an arbitrary scheme this condition is
  imposed locally on affine opens. The word constructible requires
  each stratum to be a constructible subset; for Noetherian schemes
  this is automatic for locally closed strata.
\end{enumerate}
We write $\Cons(X,\Lambda)\subset\Mod(X_{\proet},\Lambda)$ for the full
subcategory of constructible sheaves, and $\Dbc(X,\Lambda)\subset
D(X_{\proet},\Lambda)$ for the full subcategory of bounded complexes with
constructible cohomology sheaves. A complex $K$ has \emph{Tor amplitude
in $[a,b]$} if $K\otimes^{\mathbf L}_{\Lambda_X}N$ has cohomology only
in $[a,b]$ for every sheaf $N$ of $\Lambda_X$-modules. We write
$\Dctf(X,\Lambda)$ for the full subcategory of $\Dbc(X,\Lambda)$ of
complexes of finite Tor amplitude.
\label{def:sp-finite-category}\label{def:sp-tor-amplitude}
\end{definition}
For finite $A$ and qcqs $X$, $\Dctf(X,A)$ is the pullback of the usual
\'etale category $\Dctf(X_{\et},A)$ by
Corollary~\ref{pe:finite-tor-comparison}.

For a general condensed ring, kernels of maps between constructible
sheaves need not be constructible, and then $\Dbc(X,\Lambda)$ is not
triangulated. For the admissible coefficients the constructible sheaves form an
abelian category: for finite $A$ this is the classical finite-stalk
theory, and the other cases are Lemma~\ref{b:lem:coefficient-heart},
Corollary~\ref{b:lem:coefficient-heart-b} and Lemma~\ref{i:integral-heart}.
For the integral and rational rings the definition takes a familiar form.

For $\Lambda=E$ or $\Qbar$ the scalar ring $\Lambda(*)$ is a field, so
each $P_i$ is a finite-dimensional vector space and a finite constant
local model is a locally free module of finite rank. Thus a constructible
$E$- or $\Qbar$-sheaf is one that is \emph{lisse}, that is locally free
of finite rank, on the strata of a finite stratification
\cite[Definition~6.8.6]{BS}. Because the presentations have
entries in the scalar field, no rank stratification is needed to see
this.

For $\Lambda=\mathcal O_E$ torsion is allowed. Locally on a stratum, a
constructible $\mathcal O_{E,X}$-module is
$P\otimes_{\mathcal O_E}\mathcal O_{E,X}$ for a finitely generated
$\mathcal O_E$-module $P$; for example $\mathcal O_{E,X}/\pi_E$ is
constructible. In most proofs it is more convenient to use the
following criterion.

\begin{proposition}[The integral constructibility criterion]
\label{n:integral-criterion}
For $X$ Noetherian and $E/\Ql$ finite, a complex
$K\in D(X_{\proet},\mathcal O_E)$ lies in $\Dbc(X,\mathcal O_E)$ if and
only if it is derived $\pi_E$-complete and
\[
 K\otimes^{\mathbf L}_{\mathcal O_{E,X}}
             (\mathcal O_E/\pi_E)_X
\]
belongs to $\Dbc(X,\mathcal O_E/\pi_E)$ on $X_{\proet}$.
\end{proposition}
\begin{proof}
For $\mathcal O_E$ our definition of a constructible sheaf is
\cite[Definition~6.8.10]{BS}, and \cite[Proposition~6.8.11(3)]{BS}
identifies bounded complexes with such cohomology sheaves with the adic
constructible derived category. The Stacks Project describes this
category by completeness and reduction (\stacks{09C0}, \stacks{09C1}),
requiring the reduction to have finite Tor dimension; here that
condition is automatic, since $\mathcal O_E/\pi_E$ is a field. Two
identifications connect that description with the statement. First,
the Stacks Project uses an auxiliary \'etale description of the
reduction, and Theorem~\ref{pe:constructible-comparison} identifies it
with the pro-\'etale category. Second, it works over the constant sheaf
$\mathcal O_E$, whereas our coefficient sheaf is $\mathcal O_{E,X}$. The map between
them is flat and an isomorphism modulo every $\pi_E^n$, so restriction
of scalars identifies the derived complete objects on both sides
together with their reductions (\stacks{099I}). Under these
identifications the two criteria agree.
\end{proof}

Lemma~\ref{b:lem:integral-models} restates this comparison in the form
used later, and Section~\ref{i:heart} develops the integral
operations.

For $\Lambda=\Zl,\mathcal O_E,\Ql,E,\Qbar$ one has
\[
 \Dctf(X,\Lambda)=\Dbc(X,\Lambda).
\]
Indeed, a finite constant local model has a finite projective resolution
over the scalar ring: length at most one over the discrete valuation
ring $\mathcal O_E$, and length zero over a field. Pulling this
resolution back is exact, and finite stratification and boundedness
give a single Tor interval for a constructible complex
(see also Lemma~\ref{i:tensor}). For a finite nonregular ring the
condition is essential: $\mathbf F_\ell$ over $\mathbf Z/\ell^2$ is
constructible but has infinite projective dimension.

Finite Tor amplitude is a condition on the whole complex. The category
$\Dctf$ is thick but need not be stable under standard truncations:
for $A=\mathbf Z/\ell^2$ the complex $[A\xrightarrow{\ell}A]$ is
perfect, whereas its nonzero cohomology modules are $\mathbf F_\ell$,
which are not perfect. For this reason the compactness proofs for
finite $A$ work with bounded flat constructible models, and standard
truncations are taken in the ambient category $\Dbc$, which remains
the natural setting for cohomological arguments.

\begin{definition}[The localizing closure]
For admissible $\Lambda$, let $\Loc(\Dctf(X,\Lambda))$ be
the smallest full triangulated subcategory of $D(X_{\proet},\Lambda)$
containing $\Dctf(X,\Lambda)$ and closed under arbitrary direct sums.
For integral and rational coefficients this is also
$\Loc(\Dbc(X,\Lambda))$.
\end{definition}

Brown representability needs a category with arbitrary direct sums,
so it is applied to $\Loc(\Dctf(-,\Lambda))$. It needs two facts: that
the objects of $\Dctf(-,\Lambda)$ are compact generators of
$\Loc(\Dctf(-,\Lambda))$ and that proper direct image preserves direct
sums there, are proved in Section~\ref{sec:formal-half}. For $\mathcal O_E$ the generators are derived
complete, but their localizing closure also contains objects which
are not. Completeness of $f^!K$ for integral constructible $K$ is
part of the restriction theorem.

\subsection{The main theorem}
\label{n:brown}

\begin{theorem}[The exceptional inverse image]
\label{n:existence}
Let $\Lambda\in\{A,\Zl,\mathcal O_E,\Ql,E,\Qbar\}$ be an admissible
coefficient ring and $f:X\to Y$ a morphism of the standing schemes.
\begin{enumerate}
\item The objects of $\Dctf(X,\Lambda)$ are compact in
  $\Loc(\Dctf(X,\Lambda))$ and a small skeleton compactly generates it.
\item Compactifications define an exact functor
  $Rf_!:\Loc(\Dctf(X,\Lambda))\to\Loc(\Dctf(Y,\Lambda))$, preserving
  direct sums and sending $\Dctf(X,\Lambda)$ into $\Dctf(Y,\Lambda)$.
\item A single application of Brown representability gives an exact
  right adjoint $f^!:\Loc(\Dctf(Y,\Lambda))\to\Loc(\Dctf(X,\Lambda))$.
\item This right adjoint sends $\Dctf(Y,\Lambda)$ into
  $\Dctf(X,\Lambda)$, so its adjunction restricts to $\Dctf$.
\end{enumerate}
\end{theorem}
For $\Lambda=\Zl,\mathcal O_E,\Ql,E,\Qbar$ one has $\Dctf=\Dbc$, so
the theorem holds with $\Dbc$ in place of $\Dctf$. The finite ring $A$
is arbitrary; it need not be regular or Gorenstein. The $\Zl$ and
$\Ql$ cases are the integral and rational cases with $E=\Ql$.

Since $Rp_!=Rp_*$ as functors $\Loc(\Dctf(X,\Lambda))\to
\Loc(\Dctf(Y,\Lambda))$ for proper $p$, we get in particular:

\begin{corollary}
\label{n:existence-b}
For proper $p:X\to Y$ and admissible $\Lambda$, the functor
$Rp_*:\Dctf(X,\Lambda)\to\Dctf(Y,\Lambda)$ has a right adjoint $p^!$.
\end{corollary}

\subsection{Plan of the proof}
\label{sec:plan}
The proof has two parts: a formal existence theorem, giving
(1)--(3), and a geometric restriction theorem, giving (4).

Compactness is verified using perfect local models on strata. For
finite $A$, Proposition~\ref{ctf:finite-compact} works with bounded
flat constructible complexes. For fields the Ext argument of
Theorem~\ref{b:thm:constructible-compactness} applies, and the
integral version is Proposition~\ref{i:compactness}. These arguments
use local models, continuity and cohomological bounds; for $Rj_*$ they
use its finite cohomological dimension and continuity.
A compactification $f=pj$ exists by Nagata's theorem
(Theorem~\ref{app:nagata-theorem}). Put $Rf_!=Rp_*j_!$, a functor
$\Loc(\Dctf(X,\Lambda))\to\Loc(\Dctf(Y,\Lambda))$. Proper finiteness,
finite Tor control and the finite cohomological dimension of $p$ make
this a coproduct-preserving functor between the localizing categories,
and the comparisons between compactifications are built from proper
base change. Theorem~\ref{ctf:unified-brown} then applies Brown
representability once, for every admissible coefficient ring,
including finite rings.

The restriction to $\Dctf$ is geometric. For $\mathbf P^1$ the
diagonal class and trace give the unit and counit of the adjunction
explicitly; this computes $f^!$ for the projections of affine space. Closed immersions are treated by the
localization triangle and finiteness of $Rj_*$. This constructibility
of $Rj_*$ is proved after Brown, using the computed smooth adjunction;
for finite coefficients its finite Tor bound is checked by the scalar
projection formula. Local factorization as a closed immersion followed
by a smooth projection proves the general case. The finite-coefficient
restriction is Theorem~\ref{ctf:finite-restriction}; the integral and
rational cases are Theorems~\ref{i:restriction}
and~\ref{b:thm:restriction}.

Finally, smooth purity gives the canonical global normalization
$f^!\simeq f^*(d)[2d]$, compatible with trace, composition and
Frobenius. At finite levels these compatibilities are recorded by
towers, which model the adjoint $f^!$ of Theorem~\ref{n:existence}.

\section{Proper base change}
\label{sec:proper-base-change}

Proper base change identifies the stalks of a proper direct image with
the cohomology of the geometric fibres. More generally, it says that
pulling a proper family back to a new base and taking its cohomology can
be done in either order. This is the geometric input behind the base
change formula for compactly supported cohomology, and hence behind the
six-functor formalism developed below.

We first prove the torsion theorem for classical pro-\'etale complexes,
by calculations on the \'etale site, and then pass to the adic forms by
completion and change of coefficients; the comparison theorems of
Section~\ref{sec:proetale-foundations} relate the two sites. The torsion
proof follows~\cite{BlueprintPBC}. After constructing the exchange
morphism and proving the degree-zero case, we reduce the higher
vanishing to an acyclicity criterion for injective sheaves, prove the
case of the projective line by calculations on curves and henselian
comparison, and reduce arbitrary proper morphisms to that case. Since
inverse image is exact but does not in general preserve injectives,
the degree-zero isomorphism does not by itself give the derived
statement; this is why the acyclicity criterion is needed.

\begin{convention}\label{pbc:conventions}
In this section, schemes are arbitrary unless an additional hypothesis
is stated. The theorem concerns the classical part of
$D^+(T_{\proet},\Lambda)$, for a discrete commutative ring $\Lambda$.
The intermediate acyclicity, curve and Chow calculations are made on
$T_{\et}$, with abelian sheaves or modules over its discrete constant
ring. We write $\D^+(T_{\et},\Lambda)$ for the bounded below derived
category. Inverse image on these sheaves is denoted by $u^{-1}$; for
$u:S\to T$ it is a functor $\Sh(T_{\et})\to\Sh(S_{\et})$, and likewise
on module sheaves and $\D^+(T_{\et},\Lambda)\to\D^+(S_{\et},\Lambda)$. It is
exact and agrees with $u^*$ for the constant coefficient rings under
consideration. The derived direct image is $Ru_*$. Torsion always
means torsion of the underlying abelian sheaves.

The proper base change theorem imposes no invertibility condition on
the torsion coefficients, so a prime $r$ appearing in this section may
equal a residue characteristic. The standing condition on $\ell$ from
the introduction returns with the $\ell$-adic constructible theory.
\end{convention}

Cohomology of a sheaf of $\Lambda$-modules is canonically the
cohomology of its underlying abelian sheaf (\stacks{03FD}), because an
injective module sheaf is acyclic for the underlying section functor on
every object of the site. Applying this on the \'etale objects over the
target gives the same compatibility for $R^qu_*$ and, by bounded below
resolutions, for $Ru_*$. Forgetting the module structure need not
preserve injectives; only acyclicity is used.

\subsection{The exchange map}

Consider a Cartesian square
\begin{equation}\label{pbc:square}
\begin{CD}
X' @>{g'}>> X\\
@V{f'}VV @VV{f}V\\
Y' @>{g}>> Y.
\end{CD}
\end{equation}
The commutativity of the square supplies the canonical composition
isomorphism of functors $\Sh(Y_{\et})\to\Sh(X'_{\et})$
\[
 \theta_{\square}:f'^{-1}g^{-1}\xrightarrow{\sim}g'^{-1}f^{-1}.
\]
Since inverse images are exact, $\theta_\square$ also denotes the
induced isomorphism of functors $\D^+(Y_{\et},\Lambda)\to\D^+(X'_{\et},\Lambda)$.
For $F\in\Sh(X_{\et})$, the \emph{underived exchange map}, a morphism in
$\Sh(Y'_{\et})$,
\begin{equation}\label{pbc:underived-exchange}
 \operatorname{Ex}^0_{\square,F}:
 g^{-1}f_*F\longrightarrow f'_*g'^{-1}F
\end{equation}
is the adjoint, under $f'^{-1}\dashv f'_*$, of
\[
 f'^{-1}g^{-1}f_*F
 \xrightarrow{\theta_{\square}}g'^{-1}f^{-1}f_*F
 \xrightarrow{g'^{-1}\varepsilon_f}g'^{-1}F\quad\text{in }\Sh(X'_{\et}).
\]
Here $\varepsilon_f:f^{-1}f_*\to\mathrm{id}$ is the counit, a natural
transformation of functors $\Sh(X_{\et})\to\Sh(X_{\et})$.
Equivalently, first apply the unit
$g^{-1}f_*F\to f'_*f'^{-1}g^{-1}f_*F$, and then apply $f'_*$ to
the displayed composite. This construction is natural in $F$, so
$\operatorname{Ex}^0_\square:g^{-1}f_*\to f'_*g'^{-1}$ is a natural
transformation of functors $\Sh(X_{\et})\to\Sh(Y'_{\et})$, and
likewise for sheaves of $\Lambda$-modules.

On sections the map has a simple description, which we use to compute
it on stalks. A representative of a section of $g^{-1}f_*F$ over an
\'etale $V\to Y'$ is given by an \'etale $W\to Y$, a map $V\to
Y'\times_YW$, and a section over $X\times_YW$. Exchange restricts that
section along
\[
 X'\times_{Y'}V\longrightarrow X\times_YW,
\]
and passes to sheafification.

\begin{definition}[Derived exchange]\label{pbc:exchange}
For $K\in\D^+(X_{\et},\Lambda)$, the \emph{exchange morphism}
\begin{equation}\label{pbc:derived-exchange}
 \operatorname{Ex}_{\square,K}:
 g^{-1}Rf_*K\longrightarrow Rf'_*g'^{-1}K\quad\text{in }\D^+(Y'_{\et},\Lambda)
\end{equation}
is the adjoint, under $f'^{-1}\dashv Rf'_*$, of
\[
 f'^{-1}g^{-1}Rf_*K
 \xrightarrow{\theta_{\square}}g'^{-1}f^{-1}Rf_*K
 \xrightarrow{g'^{-1}\varepsilon_f}g'^{-1}K\quad\text{in }\D^+(X'_{\et},\Lambda),
\]
where $\varepsilon_f:f^{-1}Rf_*\to\mathrm{id}$ now denotes the derived
counit on $\D^+(X_{\et},\Lambda)$. The morphism is natural in $K$, so
$\operatorname{Ex}_\square:g^{-1}Rf_*\to Rf'_*g'^{-1}$ is a natural
transformation of functors $\D^+(X_{\et},\Lambda)\to\D^+(Y'_{\et},\Lambda)$.
\end{definition}

The derived adjunction exists because inverse image is exact and its
right adjoint preserves injectives. The construction of exchange needs
only a commutative square; the Cartesian hypothesis matters for the
question whether it is an isomorphism.

To compute \eqref{pbc:derived-exchange}, choose a bounded below
injective resolution $K\to I^\bullet$ and an injective replacement
$g'^{-1}I^\bullet\to J^\bullet$. Exchange is represented by the
following morphism of complexes of sheaves on $Y'_{\et}$
\begin{equation}\label{pbc:exchange-resolution}
 g^{-1}f_*I^\bullet
 \xrightarrow{\operatorname{Ex}^0_{\square,I^\bullet}}
 f'_*g'^{-1}I^\bullet\longrightarrow f'_*J^\bullet.
\end{equation}
The first arrow is the degreewise underived exchange map. Exactness of
$g^{-1}$ identifies the source with $g^{-1}Rf_*K$, and the last complex
computes $Rf'_*g'^{-1}K$. Comparison maps between injective resolutions
are unique up to homotopy, so this formula is independent of the
choices. It is the adjoint map of Definition~\ref{pbc:exchange}:
adjunction sends the first arrow to the square isomorphism followed by
the counit, and the replacement map gives the canonical passage to the
derived direct image.

\begin{lemma}[Identity and pasting]\label{pbc:exchange-pasting}
Let $K\in\D^+(X_{\et},\Lambda)$. Exchange is the identity for an
identity square, and is compatible with horizontal and vertical pasting
of Cartesian squares. More explicitly, if $h:Y''\to Y'$ is a second base
change, with induced maps $h':X''\to X'$ and $f'':X''\to Y''$, the
exchange map for $gh$ is the composite in $\D^+(Y''_{\et},\Lambda)$
\[
 h^{-1}g^{-1}Rf_*K
 \xrightarrow{h^{-1}\operatorname{Ex}_{\square,K}}
 h^{-1}Rf'_*g'^{-1}K
 \xrightarrow{\operatorname{Ex}_{h,g'^{-1}K}}
 Rf''_*h'^{-1}g'^{-1}K,
\]
under the canonical composition isomorphisms.

For composable maps $X\xrightarrow{f}Y\xrightarrow{e}Z$ and a base
change $h:Z'\to Z$, put $Y'=Y\times_ZZ'$ and $X'=X\times_ZZ'$.
Writing $e':Y'\to Z'$, $f':X'\to Y'$ and $g:Y'\to Y$,
$g':X'\to X$, the exchange map for $ef$ is the composite in
$\D^+(Z'_{\et},\Lambda)$, natural in $K$,
\[
 h^{-1}Re_*Rf_*K
 \xrightarrow{\operatorname{Ex}_{e,Rf_*K}}
 Re'_*g^{-1}Rf_*K
 \xrightarrow{Re'_*\operatorname{Ex}_{f,K}}
 Re'_*Rf'_*g'^{-1}K.
\]
\end{lemma}
\begin{proof}
For an identity square the square isomorphism is the identity modulo
the unit isomorphisms of inverse image. Its adjoint is the identity by
the triangle identity for an adjunction.

For pasting, expand each exchange map as unit, square isomorphism, and
counit. The square isomorphism for the outer rectangle is the composite
of the square isomorphisms for the inner squares. The adjacent unit and
counit cancel by the triangle identities; naturality moves the remaining
units and counits past the composition isomorphisms. The result is the
unit--square-isomorphism--counit formula for the outer rectangle. The
same calculation applies in the bounded below derived categories of the
\'etale sites. In the vertical
case the identifications $R(ef)_*\simeq Re_*Rf_*$ and
$R(e'f')_*\simeq Re'_*Rf'_*$ are valid because direct image, being
right adjoint to exact inverse image, preserves injectives. Alternatively,
one can compute the outer exchange on an injective resolution as in
\eqref{pbc:exchange-resolution}; the underived pasting calculation gives
the same map after replacement.
\end{proof}

\subsection{The base change property and the degree-zero case}

\begin{definition}[Torsion and the PBC property]\label{pbc:property}
A sheaf $F$ of abelian groups, or of $\Lambda$-modules, is
\emph{torsion} if every geometric stalk is a torsion abelian group.
Equivalently,
\[
 F=\varinjlim_{n\geq1}F[n],\qquad F[n]=\ker(n:F\to F),
\]
where the positive integers are ordered by divisibility. These are
sheaf colimits. The equivalence follows by taking geometric stalks,
which commute with colimits and finite limits and detect isomorphisms.

For a morphism $f:X\to Y$, write $\mathrm{PBC}_{\Lambda}(f)$ if the exchange morphism
\eqref{pbc:derived-exchange} is an isomorphism in $\D^+(Y'_{\et},\Lambda)$
for every base change
$g:Y'\to Y$ and every $K\in\D^+(X_{\et},\Lambda)$ with torsion
cohomology sheaves. Put $\mathrm{PBC}(f)=\mathrm{PBC}_{\Z}(f)$.
The definition makes sense for any $f$; for proper $f$ the property is
proved in Propositions~\ref{pbc:p1} and~\ref{pbc:reduction}.
\end{definition}

Torsion is a local condition on sections: every section is killed by an
integer after an \'etale covering, with the integer allowed to vary
between members of the covering. On a non-quasi-compact object a
section need not be killed by a single integer. For example, the
constant sheaf $\Q/\Z$ on a nonempty scheme is torsion, but on the
disjoint union of countably many copies of the scheme its sections
include $(1/2^n)_{n\geq1}$, which has infinite order.

In degree zero, proper base change holds for every \'etale sheaf of
sets, with no torsion hypothesis. It rests on the following comparison
of sections over a henselian pair.

\begin{theorem}[Henselian-pair comparison in degree zero]
\label{pbc:henselian-h0}
Let $(A,J)$ be a henselian pair, let $Z\to\operatorname{Spec}A$ be
proper, and let $i:Z_0=Z\times_A\operatorname{Spec}(A/J)\hookrightarrow Z$.
For every sheaf of sets $F$ on $Z_{\et}$, restriction is a bijection
\[
 \Gamma(Z,F)\xrightarrow{\sim}\Gamma(Z_0,i^{-1}F).
\]
The comparison is natural in $F$.
\end{theorem}

In particular the theorem applies to sheaves of abelian groups and of
modules. It is \stacks{0A0C}. Its geometric content is a connectedness
statement: every irreducible closed subset of a scheme proper over a
henselian pair meets the closed fibre in a nonempty connected subset.
Stein factorization and henselian connectedness prove this
(\stacks{0A0B}), and spectral-space restriction then gives equality of
sections for Zariski sheaves. The connectedness statement also holds
after every finite cover, since that cover remains proper over $A$.
A finite surjection whose Zariski-local pieces factor through the
\'etale covering, followed by descent on its self-intersection,
upgrades the Zariski comparison to \'etale sheaves (\stacks{09ZF}).
Over a general affine base the integral form of Stein factorization,
whose algebra need not be finite, suffices for this argument.

\begin{proposition}[Proper base change in degree zero]
\label{pbc:degree-zero}
For \eqref{pbc:square} with $f$ proper, and for every \'etale sheaf of
sets $F$ on $X$, the map
\[
 \operatorname{Ex}^0_{\square,F}:
 g^{-1}f_*F\xrightarrow{\sim} f'_*g'^{-1}F
\]
is an isomorphism of sheaves of sets on $Y'_{\et}$. Moreover, for a
geometric point $\bar y:\operatorname{Spec}\Omega\to Y$, restriction identifies
\begin{equation}\label{pbc:h0-fibre}
 (f_*F)_{\bar y}\xrightarrow{\sim}
 \Gamma(X_{\bar y},F|_{X_{\bar y}}).
\end{equation}
\end{proposition}
\begin{proof}
Let $Y_{(\bar y)}=\operatorname{Spec}\mathcal O^{\mathrm{sh}}_{Y,\bar y}$
and $X_{(\bar y)}=X\times_YY_{(\bar y)}$. The degree-zero strict-local
stalk formula gives
\[
 (f_*F)_{\bar y}\simeq\Gamma(X_{(\bar y)},F|_{X_{(\bar y)}}).
\]
It holds for quasi-compact quasi-separated $f$ by \'etale
continuity (\stacks{03Q9}; the degree-zero continuity proof there also
applies to sheaves of sets). Recall the mechanism: the left side is the
filtered colimit of $\Gamma(X\times_YV,F|)$ over pointed affine
\'etale neighbourhoods $V$ of $\bar y$. Their inverse limit is
$Y_{(\bar y)}$. A section on the limit is represented on a finite
\'etale covering of a quasi-compact scheme; the covering objects,
section representatives, and their finitely many agreement conditions
descend to one neighbourhood. This gives surjectivity of the
continuity map; descending the finite agreement conditions for two
representatives gives injectivity.

The local ring of $Y_{(\bar y)}$ is strictly henselian. Its residue
field $\kappa(y)^{\mathrm{sep}}\subseteq\Omega$ is separably closed.
Theorem~\ref{pbc:henselian-h0} identifies the last section set with
sections over $X_{\kappa(y)^{\mathrm{sep}}}$. Extending a separably
closed ground field preserves sections of a pulled-back \'etale sheaf
(\stacks{0A3I}), so these sections equal
$\Gamma(X_{\bar y},F|_{X_{\bar y}})$; here $F|$ denotes inverse
image of the given sheaf, also when $\Omega$ is transcendental over the
residue field. This proves \eqref{pbc:h0-fibre}.

Now choose $\bar y':\operatorname{Spec}\Omega\to Y'$ and use its
composite $\bar y=g\bar y'$ as the geometric point of $Y$. Properness
is preserved by base change, so \eqref{pbc:h0-fibre} applies to $f'$ as
well as $f$. The Cartesian square identifies
\[
 X'_{\bar y'}=X\times_Y\operatorname{Spec}\Omega=X_{\bar y}
\]
and identifies the two inverse-image sheaves there. Thus the stalks of
both sides of the desired map are the same section set. A section
represented over $X\times_YV$ maps under underived exchange to its
restriction over the base-changed neighbourhood, and hence to its
restriction on this same fibre. The stalk map is therefore the
identity under the displayed identifications. Geometric stalks detect
isomorphisms, proving the assertion.
\end{proof}
\subsection{Injective acyclicity and the passage to torsion complexes}

\begin{definition}[Injective acyclicity]\label{pbc:ia}
For a ring $\Lambda$, write $\mathrm{IA}_{\Lambda}(f)$ if for every
Cartesian square \eqref{pbc:square} and every injective sheaf $I$ of
$\Lambda$-modules on $X_{\et}$ which is torsion, one has
\[
 R^qf'_*g'^{-1}I=0\qquad(q>0).
\]
For a prime $r$, abbreviate
$\mathrm{IA}_r(f)=\mathrm{IA}_{\Z/r\Z}(f)$; every
$\Z/r\Z$-module sheaf is torsion. Write $\mathrm{IA}(f)$ if
$\mathrm{IA}_r(f)$ holds for every prime $r$.
\end{definition}

The pullback $g'^{-1}I$ of an injective sheaf is in general not
injective. Injective acyclicity asks only that it be $f'_*$-acyclic.
By the next lemma this is enough to compute $Rf'_*$ of the pullback of
an injective resolution term by term.

\begin{lemma}[Acyclic resolutions]\label{pbc:acyclicity}
Let $v:\mathcal A\to\mathcal B$ be a left exact additive functor
between abelian categories with enough injectives. If
$F\to A^\bullet$ is a resolution and each $A^m$ is $v$-acyclic, then
\[
 \mathcal H^q(vA^\bullet)\simeq R^qvF.
\]
More generally, for a bounded below complex $A^\bullet$ of
$v$-acyclic objects, the canonical comparison
$vA^\bullet\to Rv(A^\bullet)$ is an isomorphism in $D^+(\mathcal B)$.
\end{lemma}
\begin{proof}
For the first assertion put $Z^0=F$ and let $Z^{m+1}$ be the image of
$A^m\to A^{m+1}$. The exact sequences
\[
 0\longrightarrow Z^m\longrightarrow A^m\longrightarrow
 Z^{m+1}\longrightarrow0
\]
give
\[
 \operatorname{coker}(vA^m\to vZ^{m+1})\simeq R^1vZ^m,
 \qquad R^avZ^{m+1}\simeq R^{a+1}vZ^m\quad(a\geq1).
\]
Left exactness identifies the kernel of
$vA^q\to vA^{q+1}$ with $vZ^q$. Consequently, for $q\geq1$,
\[
 \mathcal H^q(vA^\bullet)
 =\operatorname{coker}(vA^{q-1}\to vZ^q)
 \simeq R^1vZ^{q-1}\simeq\cdots\simeq R^qvF,
\]
and degree zero is $vF$.

For the complex assertion first observe that a bounded below exact
complex of $v$-acyclic objects remains exact after applying $v$.
Indeed, its cycle sequences have the same form, starting with a zero
cycle at the first nonzero degree. The long exact sequences show
inductively that all cycles are $v$-acyclic, so each cycle sequence
remains short exact after $v$. Now choose a quasi-isomorphism
$A^\bullet\to J^\bullet$ to a bounded below complex of injectives.
Its cone is bounded below, exact, and termwise $v$-acyclic, since its
terms are finite direct sums of terms of $A^\bullet$ and of
$J^\bullet$. Applying $v$ to the cone shows that
$vA^\bullet\to vJ^\bullet$ is a quasi-isomorphism. The target
computes $Rv(A^\bullet)$, proving the assertion.
\end{proof}

The passage from sheaves killed by an integer to all torsion sheaves
uses continuity of \'etale cohomology. If $u:T\to S$ is quasi-compact
and quasi-separated, then
\begin{equation}\label{pbc:filtered-colimits}
 \varinjlim_i R^qu_*F_i\xrightarrow{\sim}
 R^qu_*\bigl(\varinjlim_i F_i\bigr)\qquad(q\geq0)
\end{equation}
for every filtered system of abelian \'etale sheaves, and the same
holds after any base change of $u$ (\stacks{09Z1}). It suffices to
check this on affine \'etale $V\to S$. There $R^qu_*F$ is the
sheafification of $V\mapsto H^q(T\times_SV,F|)$, and the schemes
$T\times_SV$ are quasi-compact and quasi-separated, so their
cohomology commutes with filtered colimits by the usual argument with
finite covers and finite hypercovers: in degree $q$, a class and a
witness that it dies involve only a finite $(q+1)$-skeleton, so the
cocycle and the finitely many equalities come from a single index.

\begin{theorem}[PBC and injective acyclicity]\label{pbc:ia-equivalence}
Let $f:X\to Y$ be proper. The following are equivalent:
\begin{enumerate}[label=(\roman*)]
\item $\mathrm{PBC}(f)$ holds for bounded below complexes with torsion
  cohomology sheaves.
\item For every torsion abelian sheaf $F$ and every base change
  \eqref{pbc:square}, the maps of abelian sheaves on $Y'_{\et}$
  induced by \eqref{pbc:derived-exchange} on cohomology sheaves
  \[
   g^{-1}R^qf_*F\longrightarrow R^qf'_*g'^{-1}F
  \]
  are isomorphisms for all $q\geq0$.
\item $\mathrm{IA}_r(f)$ holds for every prime $r$.
\end{enumerate}
For a commutative ring $\Lambda$ annihilated by a positive integer,
the corresponding fixed-coefficient statement is
$\mathrm{PBC}_{\Lambda}(f)\Longleftrightarrow
\mathrm{IA}_{\Lambda}(f)$.
\end{theorem}
\begin{proof}
Condition (i) implies (ii) by placing $F$ in degree zero. If (ii)
holds and $I$ is injective among $\Z/r\Z$-module sheaves, then
$Rf_*I=f_*I$ in $\D^+(Y_{\et},\Z/r\Z)$. Compatibility with
underlying abelian cohomology, explained after
Convention~\ref{pbc:conventions}, gives the same equality of underlying
derived direct images. Applying (ii) to $I$ makes
$Rf'_*g'^{-1}I$ isomorphic in $\D^+(Y'_{\et},\Z)$ to the degree-zero
complex $g^{-1}f_*I$.
This proves (iii).

Assume (iii). We first prove (ii) for sheaves killed by a prime $r$.
Choose an injective resolution $F\to I^\bullet$ among
$\Z/r\Z$-module sheaves. Exactness of $g'^{-1}$ makes
$g'^{-1}I^\bullet$ a resolution of $g'^{-1}F$, and (iii) makes its
terms $f'_*$-acyclic. Lemma~\ref{pbc:acyclicity} therefore identifies
$f'_*g'^{-1}I^\bullet$ with $Rf'_*g'^{-1}F$. The complex
$g^{-1}f_*I^\bullet$ computes $g^{-1}Rf_*F$ because inverse image
is exact. Proposition~\ref{pbc:degree-zero} gives an isomorphism
\[
 g^{-1}f_*I^m\xrightarrow{\sim} f'_*g'^{-1}I^m
\]
of sheaves of $\Z/r\Z$-modules on $Y'_{\et}$ in each degree.
Naturality makes these into an isomorphism of complexes.
Formula~\eqref{pbc:exchange-resolution} identifies that isomorphism with
the derived exchange map $\operatorname{Ex}_{\square,F}$ in
$\D^+(Y'_{\et},\Z/r\Z)$, and hence, forgetting the module structure, in
$\D^+(Y'_{\et},\Z)$. This proves all degrees
of (ii) for prime-torsion sheaves.

Next suppose $nF=0$. Induct on the number of prime factors of $n$,
counted with multiplicity. The case $n=1$ is trivial. Choose a prime
$r\mid n$ and use
\[
 0\longrightarrow F[r]\longrightarrow F
 \xrightarrow{r}rF\longrightarrow0.
\]
The kernel is killed by $r$, and the image $rF$ is killed by $n/r$.
Applying the two functors
$g^{-1}Rf_*,\ Rf'_*g'^{-1}\colon\D^+(X_{\et},\Z)\to\D^+(Y'_{\et},\Z)$
gives a morphism of distinguished triangles in $\D^+(Y'_{\et},\Z)$ whose outer exchange maps are isomorphisms by the preceding
case and induction. The middle exchange is an isomorphism, by the
long exact cohomology sequences and the five lemma. Thus (ii) holds
for all sheaves killed by an integer.

For a general torsion sheaf, write $F=\varinjlim_n F[n]$ as in
Definition~\ref{pbc:property}. Inverse images commute with colimits,
and both $f$ and $f'$ are quasi-compact and quasi-separated. Equation
\eqref{pbc:filtered-colimits} identifies the exchange map in degree
$q$ for $F$ with the filtered colimit of those for $F[n]$. Hence it
is an isomorphism. This proves (ii).

Finally let $K\in\D^+(X_{\et},\Z)$ have torsion cohomology. If $K$
is bounded, induction on its cohomological length using the canonical
truncation triangles and (ii) proves that its exchange map is an
isomorphism. In general fix a degree $m$ and choose $N\geq m$.
The complex $\tau_{\leq N}K$ is bounded and has the desired
comparison. Both functors $g^{-1}Rf_*$ and $Rf'_*g'^{-1}$ send
$\D^{\geq N+1}(X_{\et},\Z)$ into $\D^{\geq N+1}(Y'_{\et},\Z)$, since inverse images are exact
and derived direct images are left $t$-exact. Applying them to
\[
 \tau_{\leq N}K\longrightarrow K\longrightarrow
 \tau_{\geq N+1}K\xrightarrow{+1}
\]
shows that comparison for $K$ in degree $m$ is the comparison for
$\tau_{\leq N}K$. It is an isomorphism. Since $m$ was arbitrary,
(i) follows; no bound on the cohomological dimension of $f$ enters.

For the fixed ring $\Lambda$ annihilated by an integer, every module
sheaf is torsion. The same injective-resolution argument proves the
sheaf assertion directly from $\mathrm{IA}_{\Lambda}(f)$, and the
same truncation argument proves it for $\D^+$. The reverse implication
follows by applying exchange to an injective $\Lambda$-module sheaf.
\end{proof}

\begin{remark}\label{pbc:proof-division}
Theorem~\ref{pbc:ia-equivalence} reduces higher proper base change
to a vanishing statement, and the only geometric fact used in its
proof is the degree-zero comparison of Proposition~\ref{pbc:degree-zero}. It remains
to prove $\mathrm{IA}_r(f)$ for every proper $f$ and every prime $r$.
We prove it first for $\PP^1_S\to S$ and then reduce arbitrary proper
morphisms to that case. The curve calculations for primes invertible
on the base and for the residue characteristic are different, but
both are used only through this criterion.
\end{remark}
\subsection{Proper base change for the projective line}
\label{pbc:projective-line-section}

Let $p:\PP^1_S\to S$ be the projection. The proof of proper base
change for $p$ has three stages. First, the cohomology of a projective
line over a henselian local ring is compared with that of its closed
fibre. Second, this identifies the geometric stalks of $R^qp_*F$ with
the cohomology of projective lines over separably closed fields.
Third, these groups are compared by the invariance of curve cohomology
under extension of separably closed fields. Throughout, $r$ denotes a
prime; it need not be invertible on $S$.

\subsubsection{Henselian pairs and curves}

We use two results on henselian pairs, both valid with no Noetherian
or characteristic hypotheses.
\begin{enumerate}[label=(\roman*)]
\item \emph{Gabber's affine henselian comparison.}
If $(B,J)$ is a henselian pair and $G$ is a torsion abelian sheaf on
$\operatorname{Spec}(B)_{\et}$, restriction induces
\[
 H^q(\operatorname{Spec}B,G)
 \xrightarrow{\sim}
 H^q(\operatorname{Spec}(B/J),G|)
 \qquad(q\geq0).
\]
This is \stacks{09ZI}; its proof is an affine effacement argument.
\item \emph{Lifting finite \'etale covers over a proper henselian pair.}
If $Y$ is proper over a henselian pair $(B,J)$, and
$Y_0=Y\times_B B/J$, restriction gives an equivalence
\[
 \mathsf{F\acute Et}(Y)\xrightarrow{\sim}
 \mathsf{F\acute Et}(Y_0).
\]
This is \stacks{0GS2}, proved by algebraization, approximation, and
descent of finite presentation.  In particular, for a finite abelian
group $M$, it gives
\begin{equation}
 H^1(Y,\underline M)\xrightarrow{\sim}H^1(Y_0,\underline M).
 \label{pbc:finite-torsors}
\end{equation}
Indeed, these groups classify $M$-torsors.  Finite \'etale covers, their
group actions, and the equations expressing simple transitivity are
transported by the equivalence.
\end{enumerate}
We also use the degree-zero comparison of
Theorem~\ref{pbc:henselian-h0}. It concerns sections; (ii) adds the
lifting of finite \'etale covers.

The comparison of closed fibres rests on the following facts about
curves over an algebraically closed field. In the residue
characteristic, proper and affine curves behave differently, which is
why the list distinguishes them.

\begin{theorem}[Cohomology of curves]
\label{pbc:curve-input}
Let $k$ be algebraically closed, let $C$ be a separated scheme of finite
type over $k$ with $\dim C\leq1$, and let $F$ be a classical torsion
abelian sheaf on $C_{\proet}$. Cohomology is cohomology on
$C_{\proet}$. The following assertions hold:
\begin{enumerate}[label=$(C_{\arabic*})$,leftmargin=3.5em]
\item $H^q(C,F)=0$ for $q>2$.
\item If $C$ is affine, $H^q(C,F)=0$ for $q>1$.
\item If $\operatorname{char}k=p>0$ and $F$ is $p$-primary,
  $H^q(C,F)=0$ for $q>1$.
\item If $F$ is constructible with finite stalks and prime to
  $\operatorname{char}k$, all $H^q(C,F)$ are finite.
\item If $C$ is proper and $F$ is constructible with finite stalks,
  all $H^q(C,F)$ are finite.
\item If $k'/k$ is an extension of algebraically closed fields and $F$
  is prime to $\operatorname{char}k$, the pullback maps
  $H^q(C,F)\to H^q(C_{k'},F|)$ are isomorphisms.
\item If $C$ is proper, the maps in $(C_6)$ are isomorphisms for every
  torsion $F$.
\item For any open $U\subseteq C$, restriction
  $H^2(C,F)\to H^2(U,F|_U)$ is surjective.
\end{enumerate}
\end{theorem}

In characteristic zero every torsion sheaf is prime to the
characteristic, so $(C_4)$ and $(C_6)$ apply to all torsion $F$.

\begin{proof}
This is Theorem~\ref{app:curve-package}, proved in
Appendix~\ref{app:curves}. For a smooth proper connected curve and $n$
invertible in $k$, the Kummer sequence and the Jacobian
(Appendix~\ref{app:picard}) compute $H^q(C,\mu_n)$; for $p$-primary
coefficients the Artin--Schreier sequence is used instead; and general
torsion sheaves reduce to these cases by normalization, removal of
finitely many points and a finite \'etale cover of degree prime to the
torsion order.
\end{proof}

\begin{lemma}[Extension of a separably closed ground field]
\label{pbc:curve-field-extension}
Let $k\subseteq k'$ be separably closed fields, let $C$ be proper over $k$
with $\dim C\leq1$, and let $F$ be a classical torsion abelian sheaf
on $C_{\proet}$. Then the natural map
\[
 H^q(C,F)\xrightarrow{\sim}H^q(C_{k'},F|)
 \qquad(q\geq0)
\]
is an isomorphism.
\end{lemma}
\begin{proof}
This is Corollary~\ref{app:curve-invariance}. By the classical
comparison one may compute on the \'etale site. Choose compatible
algebraic closures $\bar k\subseteq\bar k'$; since $\bar k/k$ and
$\bar k'/k'$ are purely inseparable, $C_{\bar k}\to C$ and
$C_{\bar k'}\to C_{k'}$ are universal homeomorphisms and induce
equivalences of small \'etale topoi (\stacks{03SI}). This reduces the
claim to the pullback from $C_{\bar k}$ to $C_{\bar k'}$, which is an
isomorphism by $(C_7)$ of Theorem~\ref{pbc:curve-input}.
\end{proof}

\begin{remark}
Properness in $(C_7)$ is essential when the torsion includes the
characteristic.  For example, in characteristic $p$,
\[
 H^1(\A^1_k,\mathbf F_p)
 =k[t]/\{h^p-h:h\in k[t]\},
\]
and extending an algebraically closed ground field can add new classes.
This is why the field-extension step below needs properness: its
coefficients are arbitrary torsion sheaves.
\end{remark}

\subsubsection{Restriction of injectives to the closed fibre}

Let $A$ be a henselian local ring with maximal ideal $\mathfrak m$ and
residue field $k$.  Set
\[
 X=\PP^1_A,\qquad X_0=\PP^1_k,\qquad i:X_0\hookrightarrow X.
\]
For a sheaf $G$ on $X_{\et}$ write $G_0=i^{-1}G$.

We also use two standard finiteness facts about small \'etale topoi of
quasi-compact quasi-separated schemes.  A sheaf of
$\mathbf F_r$-modules is a filtered colimit of constructible ones
(\stacks{03SA}), and cohomology commutes with filtered colimits
(\stacks{03Q5}).  Moreover, a constructible sheaf $G$ of
$\mathbf F_r$-modules admits a monomorphism
\begin{equation}
 G\lhook\joinrel\longrightarrow
 \bigoplus_{a=1}^N(h_a)_*\underline{M_a},
 \qquad h_a:Y_a\to X\ \text{finite and finitely presented},
 \label{pbc:constructible-finite-embedding}
\end{equation}
with $M_a$ finite $\mathbf F_r$-modules (\stacks{09Z7}). These facts
hold on every quasi-compact quasi-separated scheme, so $A$ may be
non-Noetherian.

The sheaves on the right of \eqref{pbc:constructible-finite-embedding}
are controlled by the case of finite morphisms, Lemma~\ref{pbc:finite}
below, whose proof uses only the strict-local structure of finite
morphisms:
\begin{equation}
 R^qh_*G=0\quad(q>0),\qquad
 i^{-1}h_*G\xrightarrow{\sim}(h_0)_*i_Y^{-1}G
 \quad(h\text{ finite}).
 \label{pbc:finite-direct-image-used}
\end{equation}
Here $h_0\colon Y_0=Y\times_XX_0\to X_0$ and $i_Y\colon Y_0\to Y$ are
the base changes of $h$ and $i$, and the second map is the underived
exchange map \eqref{pbc:underived-exchange}, an isomorphism of sheaves on
$X_{0,\et}$. Both assertions follow from the decomposition of a finite
scheme over a strict localization into strictly henselian local schemes
(\stacks{03QP}, \stacks{0959}).

\begin{lemma}[Acyclicity on the closed projective line]
\label{pbc:p1-injective-fibre}
If $I$ is an injective sheaf of $\mathbf F_r$-modules on $X_{\et}$,
then $H^q(X_0,I_0)=0$ for every $q>0$.
\end{lemma}
\begin{proof}
First consider an affine open $W=\operatorname{Spec}B$ of $X$, and set
$J=\mathfrak mB$. The pair $(B,J)$ need not be henselian, so we pass
to its henselization $(B^h,JB^h)$.  It has quotient $B/J$, and $B^h$ is the filtered colimit
of the \'etale $B$-algebras $B'$ equipped with
$B'/JB'\simeq B/J$.  Gabber's affine comparison and \'etale cohomology
continuity give
\begin{equation}
 H^q(W_0,I_0|_{W_0})
 \simeq H^q(\operatorname{Spec}B^h,I|)
 \simeq\varinjlim_{B'}H^q(\operatorname{Spec}B',I|).
 \label{pbc:affine-neighbourhood-vanishing}
\end{equation}
Restriction of an injective to an \'etale object is injective: the
restriction functor has the exact left adjoint given by extension by
zero in the localized site.  Thus every group on the right vanishes for
$q>0$.

The two standard affine charts $U,V\subset X$ have affine intersection.
Consequently $U_0,V_0$, and $(U\cap V)_0$ are acyclic for the
restriction of $I$ in positive degrees, by
\eqref{pbc:affine-neighbourhood-vanishing}.  Mayer--Vietoris on $X_0$
then gives $H^q(X_0,I_0)=0$ for $q\geq2$.  This is the two-chart
case of \stacks{0A51}.

It remains to deal with $\xi\in H^1(X_0,I_0)$.  Write $I$ as a filtered
colimit of constructible $\mathbf F_r$-sheaves.  Pullback commutes with
this colimit, as does $H^1$ on $X_0$.  Hence $\xi$ is the image of a
class $\zeta\in H^1(X_0,G_0)$ for some constructible $G$ with a map
$G\to I$.  Choose \eqref{pbc:constructible-finite-embedding} and put
$P=\bigoplus_a(h_a)_*\underline{M_a}$.  Injectivity of $I$ extends
$G\to I$ to $P\to I$; let $\zeta'$ be the image of $\zeta$ in
$H^1(X_0,P_0)$.

By \eqref{pbc:finite-direct-image-used} and Leray,
\[
 H^1(X_0,P_0)
 \simeq\bigoplus_a H^1((Y_a)_0,\underline{M_a}),\qquad
 H^1(X,P)\simeq\bigoplus_a H^1(Y_a,\underline{M_a}).
\]
Each $Y_a$ is proper over $A$, since it is finite over $X$.  The
finite-torsor comparison \eqref{pbc:finite-torsors} therefore lifts
$\zeta'$ to a class $\widetilde\zeta\in H^1(X,P)$.
Its image in $H^1(X,I)$ is zero because $I$ is injective.  The natural
restriction square
\[
 \begin{CD}
 H^1(X,P) @>>> H^1(X,I)\\
 @VVV @VVV\\
 H^1(X_0,P_0) @>>> H^1(X_0,I_0)
 \end{CD}
\]
shows that the image $\xi$ of $\zeta'$ is zero.  This proves the remaining
degree.  The argument is the one in \stacks{0A5H}.
\end{proof}

\begin{proposition}[Henselian comparison for the projective line]
\label{pbc:p1-henselian}
For a henselian local ring $A$ and any torsion abelian sheaf $F$ on
$\PP^1_A$, restriction induces natural isomorphisms
\[
 H^q(\PP^1_A,F)\xrightarrow{\sim}
 H^q(\PP^1_k,F_0)\qquad(q\geq0).
\]
\end{proposition}
\begin{proof}
Suppose first that $rF=0$ for a prime $r$.  Take an injective resolution
$F\to I^\bullet$ in sheaves of $\mathbf F_r$-modules on $X_{\et}$.
Since $i^{-1}$ is exact, $F_0\to I_0^\bullet$ remains a resolution;
Lemma~\ref{pbc:p1-injective-fibre} makes it a resolution by
$\Gamma(X_0,-)$-acyclic sheaves.  Theorem~\ref{pbc:henselian-h0}
identifies the two complexes of sections term by term:
\[
 \Gamma(X,I^\bullet)\xrightarrow{\sim}
 \Gamma(X_0,I_0^\bullet).
\]
The left complex computes $R\Gamma(X,F)$, and the right computes
$R\Gamma(X_0,F_0)$.  Thus their natural restriction map is a
quasi-isomorphism.

The class of sheaves for which this holds is closed under extensions,
by the natural long exact cohomology sequences.  A sheaf killed by
$r^a$ has a finite filtration with quotients killed by $r$; a sheaf
killed by an integer $n$ is the finite direct sum of its primary parts.
This gives the result for sheaves of bounded torsion.  Finally,
$F=\varinjlim_n F[n]$, with positive integers ordered by divisibility.
Both $X$ and $X_0$ are quasi-compact and quasi-separated, so
\stacks{03Q5} passes the comparison to this filtered colimit.  All maps
used are restriction maps, which also proves naturality in $F$.
\end{proof}

\subsubsection{Geometric stalks and arbitrary base change}

The henselian comparison enters through the stalks of higher direct
images. For a quasi-compact quasi-separated $f:Z\to S$, a geometric
point $\bar s$, and an abelian \'etale sheaf $G$, it says
\begin{equation}
 (R^qf_*G)_{\bar s}
 \simeq H^q(Z\times_SS_{(\bar s)},G|),\qquad
 S_{(\bar s)}=\operatorname{Spec}\cO^{\mathrm{sh}}_{S,\bar s}.
 \label{pbc:strict-local-stalk}
\end{equation}
Indeed, the stalk is the filtered
colimit of $H^q(Z\times_SV,G|)$ over pointed \'etale neighbourhoods
$(V,\bar v)$ of $\bar s$, and \'etale cohomology continuity identifies
that colimit with the cohomology of the indicated limit scheme
(\stacks{03Q9}).

\begin{proposition}[Proper base change for the projective line]
\label{pbc:p1}
Let $S$ be any scheme and $\Lambda$ any commutative ring.  In the
cartesian square
\[
 \begin{CD}
 \PP^1_T @>{g'}>> \PP^1_S\\
 @V{p'}VV @VV{p}V\\
 T @>{g}>> S,
 \end{CD}
\]
the exchange map
\[
 g^{-1}Rp_*K\longrightarrow Rp'_*g'^{-1}K
\]
is an isomorphism in $D^+(T_{\et},\Lambda)$ for every
$K\in D^+(\PP^1_{S,\et},\Lambda)$ with torsion cohomology sheaves.  In particular, the projection satisfies
the injective-acyclicity criterion of Theorem~\ref{pbc:ia-equivalence}.
\end{proposition}
\begin{proof}
We first take a torsion sheaf $F$ in degree zero.  Choose a geometric
point $\bar t:\operatorname{Spec}\Omega\to T$ and put
$\bar s=g\bar t$.  The strict localizations fit into a natural map
$T_{(\bar t)}\to S_{(\bar s)}$ of strictly henselian local schemes.
Write $k_s,k_t$ for their residue fields; these are separably closed,
and the map induces an inclusion $k_s\hookrightarrow k_t$.
For each $q\geq0$, \eqref{pbc:strict-local-stalk} and
Proposition~\ref{pbc:p1-henselian} give the commutative diagram
\[
 \begin{CD}
 (g^{-1}R^qp_*F)_{\bar t}
 @>{\operatorname{Ex}_{\bar t}}>>
 (R^qp'_*g'^{-1}F)_{\bar t}\\
 @V{\sim}VV @VV{\sim}V\\
 H^q(\PP^1_{S_{(\bar s)}},F|)
 @>>> H^q(\PP^1_{T_{(\bar t)}},F|)\\
 @V{\sim}VV @VV{\sim}V\\
 H^q(\PP^1_{k_s},F|)
 @>>> H^q(\PP^1_{k_t},F|).
 \end{CD}
\]
The top arrow $\operatorname{Ex}_{\bar t}$ is the stalk at $\bar t$ of
the map of abelian sheaves on $T_{\et}$ induced by
\eqref{pbc:derived-exchange} on $q$th cohomology sheaves. The middle and
bottom arrows are the natural pullback maps.  The
top square commutes by the definition of exchange: a section over
$\PP^1_V$, for a pointed \'etale neighbourhood $V\to S$, is pulled
back to $\PP^1_{V\times_ST}$.  Passing to the two strict localization
limits gives exactly the middle arrow.  The bottom square is naturality
of restriction to the closed fibres.

The bottom arrow is an isomorphism by
Lemma~\ref{pbc:curve-field-extension}.  Hence every geometric stalk
of $g^{-1}R^qp_*F\to R^qp'_*g'^{-1}F$ is an isomorphism.
Geometric stalks are conservative, so exchange for $F$ is an
isomorphism in $D^+(T_{\et},\Z)$.  Cohomology and higher direct image of
$\Lambda$-module sheaves agree with those of their underlying abelian
sheaves, so this also proves the assertion with the stated coefficient
ring and its module structure.

Now let $K\in D^+(\PP^1_{S,\et},\Lambda)$ have torsion cohomology and
fix $m\in\Z$. The two functors
$g^{-1}Rp_*,\ Rp'_*g'^{-1}\colon D^+(\PP^1_{S,\et},\Lambda)\to
D^+(T_{\et},\Lambda)$ send $D^{\geq m+1}$ into $D^{\geq m+1}$, since inverse image is exact and direct image is
left exact.  The truncation triangle for $\tau_{\leq m}K$ therefore
reduces the degree-$m$ assertion to $\tau_{\leq m}K$, which is bounded.
Successive truncation triangles reduce this bounded case to the
shifted torsion sheaves $H^a(K)[-a]$.  Exchange is compatible with
these triangles and has already been proved for those sheaves.
It follows that exchange is an isomorphism in every degree for $K$.

Finally, if $I$ is a torsion injective sheaf of $\Lambda$-modules, then
$R^qp_*I=0$ for $q>0$.  The sheaf case of the theorem gives
$R^qp'_*g'^{-1}I\simeq g^{-1}R^qp_*I=0$.  For
$\Lambda=\mathbf F_r$ this is $\mathrm{IA}_r(p)$.
\end{proof}

\begin{remark}
In this proof the henselian comparison rests on affine effacement and
the lifting of finite torsors, and the passage between geometric stalks
rests on the cohomology of curves. The explicit cohomology of $\PP^1$
over a base, its trace, and smooth purity are treated later, using
Proposition~\ref{pbc:p1}.
\end{remark}
\subsection{Arbitrary proper morphisms}
\label{pbc:sec-reduction}

The projective-line case contains the cohomological difficulty of the
theorem. The reduction of an arbitrary proper morphism to it uses finite
morphisms, a finite cover of $\mathbf P^n$ by $(\mathbf P^1)^n$, and
Chow's lemma. Its one nonformal point is a direct-summand argument that
allows a proper surjection to be cancelled.

\begin{lemma}[Finite morphisms]\label{pbc:finite}
Let $u:V\to W$ be finite. For an abelian sheaf $F$ and a geometric point
$\bar w:\operatorname{Spec}(\Omega)\to W$, with $\Omega$ algebraically
closed, one has
\[
 (u_*F)_{\bar w}
 \simeq \prod_{\substack{\bar v:\operatorname{Spec}(\Omega)\to V\\
                         u\bar v=\bar w}}F_{\bar v}.
\]
The product is finite. The functor $u_*$ on abelian sheaves is exact,
and its formation commutes with arbitrary base change. Consequently, for
every base change $W'\to W$ and every constant coefficient ring
$\Lambda$, the derived exchange
\eqref{pbc:derived-exchange} is an isomorphism in
$\D^+(W'_{\et},\Lambda)$ for every $K\in\D^+(V_{\et},\Lambda)$.
\end{lemma}

\begin{proof}
After passage to the strict henselization of $W$ at $\bar w$, the finite
scheme $V$ is a finite disjoint union of strictly henselian local schemes.
On each such local scheme, sections of an \'etale sheaf are its stalk at
the closed geometric point. Indeed, every \'etale neighbourhood of that
point admits a section, so a germ extends to a global section, and two
extensions agreeing as germs agree after such a section. This also makes
sections exact on abelian sheaves. The strict-local description of direct
image therefore gives the displayed formula; see \stacks{03QP}.
Geometric stalks are conservative and finite products of abelian groups
are exact, so $u_*$ is exact.

In a cartesian base-change square, the two geometric fibres are
canonically the same after using the same geometric field. The
underived exchange is the identity on the corresponding factors of the
displayed formula. It is therefore an isomorphism of sheaves on
$W'_{\et}$. Both direct images in
the square are exact, so their derived functors and the exchange can be
computed term by term. The same argument applies to sheaves of modules.
\end{proof}

\begin{lemma}[Composition and cancellation]\label{pbc:composition-cancellation}
Let $X\xrightarrow{u}Y\xrightarrow{v}Z$ be proper morphisms.
\begin{enumerate}[label=(\roman*)]
\item If $\mathrm{PBC}(u)$ and $\mathrm{PBC}(v)$ hold, then
      $\mathrm{PBC}(vu)$ holds.
\item If $u$ is surjective and $\mathrm{PBC}(u)$ and
      $\mathrm{PBC}(vu)$ hold, then $\mathrm{PBC}(v)$ holds.
\end{enumerate}
\end{lemma}

\begin{proof}
We use Theorem~\ref{pbc:ia-equivalence}. Fix a prime $r$ and an
arbitrary base change $c:Z'\to Z$. Write $a:X'\to X$ and $b:Y'\to Y$
for the projections, and $u':X'\to Y'$, $v':Y'\to Z'$ for the induced
maps. Since inverse image of constant-coefficient module sheaves is
exact, its right adjoint $u_*$ preserves injectives. Thus the
Grothendieck spectral sequence for direct images
(\stacks{0734}) is available:
\begin{equation}\label{pbc:leray-reduction}
 R^pv'_*R^qu'_*a^{-1}J
 \ \Longrightarrow\ R^{p+q}(v'u')_*a^{-1}J.
\end{equation}
It is a first-quadrant spectral sequence, natural in the sheaf $J$.

For (i), let $J$ be injective among $\mathbf F_r$-module sheaves on $X$.
The condition $\mathrm{IA}_r(u)$ gives
$R^qu'_*a^{-1}J=0$ for $q>0$. Degree-zero proper base change
(Proposition~\ref{pbc:degree-zero}) identifies, as sheaves of
$\mathbf F_r$-modules on $Y'_{\et}$,
\[
 u'_*a^{-1}J\simeq b^{-1}u_*J.
\]
The sheaf $u_*J$ is injective, so $\mathrm{IA}_r(v)$ kills
$R^pv'_*b^{-1}u_*J$ for $p>0$. Thus
\eqref{pbc:leray-reduction} gives
$R^m(v'u')_*a^{-1}J=0$ for $m>0$. This proves
$\mathrm{IA}_r(vu)$ for every $r$, and hence (i).

For (ii), let $I$ be injective among $\mathbf F_r$-module sheaves on $Y$.
Choose a monomorphism of sheaves
\[
             u^{-1}I\lhook\joinrel\longrightarrow J
\]
with $J$ injective on $X$, and let $\iota:I\to u_*J$ be its adjoint.
This is a global sheaf morphism. It is monic: for a geometric point
$\bar y$ choose a geometric point $\bar x$ above it, which exists by
surjectivity of $u$. The composite
\[
 I_{\bar y}\xrightarrow{\iota_{\bar y}}(u_*J)_{\bar y}
 \longrightarrow J_{\bar x}
\]
is the stalk of the chosen monomorphism, under
$(u^{-1}I)_{\bar x}=I_{\bar y}$, so it is injective.
Injectivity of $I$ now supplies a sheaf map $s:u_*J\to I$ satisfying
$s\iota=\mathrm{id}_I$. Thus $I$ is a direct summand of $u_*J$.

The assumptions on $u$ and $vu$ give
\[
 R^qu'_*a^{-1}J=0\quad(q>0),\qquad
 R^m(v'u')_*a^{-1}J=0\quad(m>0).
\]
The spectral sequence \eqref{pbc:leray-reduction} collapses to its
$q=0$ row and gives
\[
 R^pv'_*b^{-1}u_*J
 \simeq R^pv'_*u'_*a^{-1}J=0\qquad(p>0).
\]
The splitting of $I\to u_*J$ remains a splitting after $b^{-1}$.
Higher direct image preserves finite direct sums and their direct
summands, so $R^pv'_*b^{-1}I=0$ for $p>0$. This proves
$\mathrm{IA}_r(v)$, and (ii) follows for all primes $r$.
\end{proof}

\begin{lemma}[A finite cover of projective space]\label{pbc:symmetric-product}
For every scheme $S$ and every $n\geq1$ there is a finite surjective
$S$-morphism
\[
 \sigma_n:(\mathbf P^1_S)^n\longrightarrow\mathbf P^n_S.
\]
If $\mathrm{PBC}(\mathbf P^1_T\to T)$ holds for every scheme $T$, then
$\mathrm{PBC}(\mathbf P^n_S\to S)$ holds for every $S$ and $n\geq0$.
\end{lemma}

\begin{proof}
Use the coefficients of a product of binary linear forms:
\[
 \prod_{i=1}^n(x_iT+y_iU)
       =\sum_{j=0}^nF_jT^{n-j}U^j,\qquad
 \sigma_n\bigl(([x_i:y_i])_i\bigr)=[F_0:\cdots:F_n].
\]
Every $F_j$ is homogeneous of degree one in each pair $(x_i,y_i)$.
The corresponding sections of $\mathcal O(1,\ldots,1)$ have no common
zero: at a geometric point every linear factor is nonzero, so their
product is nonzero. They therefore define the stated morphism over any
$S$, and the construction commutes with base change.

The elementary symmetric-polynomial description is visible on the
target chart $F_0\ne0$. Its preimage is the chart where every $x_i$ is
invertible. Over an affine base $\operatorname{Spec}(A)$ put
$w_i=y_i/x_i$ and $c_j=F_j/F_0$. The coordinate map is
\[
 A[c_1,\ldots,c_n]\longrightarrow A[w_1,\ldots,w_n],
       \qquad c_j\longmapsto e_j(w_1,\ldots,w_n).
\]
Each $w_i$ satisfies the monic equation
\[
 W^n-c_1W^{n-1}+c_2W^{n-2}-\cdots+(-1)^nc_n=0.
\]
This makes the map finite on that chart.

Global finiteness follows from properness and finite fibres. The source
is proper over $S$ and the target is separated over $S$, so $\sigma_n$
is proper, by the closed-graph argument. Over an algebraically closed
field a nonzero binary form of degree $n$ has a unique multiset of
linear factors up to nonzero scalars. Its fibre therefore has at most
$n!$ geometric points, given by their orderings. The fibres are of
finite type over fields; having finitely many geometric points makes
them zero dimensional. Thus $\sigma_n$ is quasi-finite, and a proper
quasi-finite morphism is finite (\stacks{02LS}). Factorization of binary
forms also shows that every geometric fibre is nonempty, so $\sigma_n$
is surjective. This argument includes forms with a root at infinity.

The projection $(\mathbf P^1_S)^n\to S$ is a composite of $n$
projective-line projections over successive bases. It satisfies PBC by
the hypothesis and Lemma~\ref{pbc:composition-cancellation}(i).
Lemma~\ref{pbc:finite} gives PBC for $\sigma_n$.
Lemma~\ref{pbc:composition-cancellation}(ii) cancels $\sigma_n$ and gives
PBC for $\mathbf P^n_S\to S$. The case $n=0$ is the identity.
\end{proof}

\begin{lemma}[Chow's lemma]\label{pbc:chow}
Let $f:X\to Y$ be proper with $Y$ affine. There exist a proper
surjection $\pi:\widetilde X\to X$ and a closed immersion
$i_Y:\widetilde X\hookrightarrow\mathbf P^n_Y$ such that
$f\pi=q_Yi_Y$, where $q_Y:\mathbf P^n_Y\to Y$ is the projection.
Moreover
\[
 i_X=(\pi,i_Y):\widetilde X\lhook\joinrel\longrightarrow
       X\times_Y\mathbf P^n_Y=\mathbf P^n_X
\]
is a closed immersion. In particular both $\pi$ and $f\pi$ are
projective morphisms.
\end{lemma}

\begin{proof}
This is Theorem~\ref{app:chow-projective}, whose closed-graph argument
gives both closed immersions.
\end{proof}

\begin{remark}[Chow's lemma over arbitrary bases]\label{pbc:chow-scope}
For the standing schemes of finite type over $k$, the Noetherian case
of Appendix~\ref{app:compactification} suffices. The qcqs version of
Theorem~\ref{app:chow-projective} gives the statement over the
arbitrary bases allowed in this section, for proper morphisms of finite
type.
\end{remark}

\begin{proposition}[Reduction to the projective line]\label{pbc:reduction}
If $\mathrm{PBC}(\mathbf P^1_S\to S)$ holds for every scheme $S$,
then $\mathrm{PBC}(f)$ holds for every proper morphism of schemes $f$.
\end{proposition}

\begin{proof}
The assertion is local on the target. Indeed, higher direct images
restrict to the corresponding higher direct images over an \'etale
open of the target, compatibly with the canonical exchange. For an
arbitrary base change $Y'\to Y$, the inverse images of an affine open
cover of $Y$ cover $Y'$, and geometric stalks detect an isomorphism in
$\D^+(Y'_{\et},\Z)$.
We may therefore suppose $Y$ affine.

Apply Lemma~\ref{pbc:chow}. The two maps $\pi$ and $f\pi$ factor as
closed immersions $i_X$ and $i_Y$ followed by projective-space
projections. Closed immersions are finite, so they satisfy PBC by
Lemma~\ref{pbc:finite}. Projective-space projections satisfy PBC by
Lemma~\ref{pbc:symmetric-product}. Composition therefore gives PBC for
both $\pi$ and $f\pi$. As $\pi$ is surjective, cancellation gives PBC
for $f$.
\end{proof}

\subsection{Proper base change and geometric fibres}
\label{pbc:sec-torsion-theorem}

\begin{theorem}[Proper base change on classical pro-\'etale complexes]
\label{pbc:torsion-theorem}
Let $\Lambda$ be a commutative ring, let $f:X\to Y$ be a proper
morphism of schemes, and let $g:Y'\to Y$ be arbitrary. Set
$X'=X\times_YY'$, with projections $f':X'\to Y'$ and $g':X'\to X$.
If $K\in D^+_{\mathrm{class}}(X_{\proet},\Lambda)$ has torsion
underlying abelian cohomology sheaves, the canonical exchange,
defined as in Definition~\ref{pbc:exchange} from the pro-\'etale
inverse and direct images $g^*,g'^*$ and $Rf_*,Rf'_*$ on
$D^+(\,\cdot_{\proet},\Lambda)$,
\[
 g^*Rf_*K\longrightarrow Rf'_*g'^*K
\]
is an isomorphism in $D^+_{\mathrm{class}}(Y'_{\proet},\Lambda)$.
\end{theorem}

The complex $K$ need not be constructible, and the torsion orders need
not be invertible on $Y$. Classicality of $K$ is required: it is what
brings the statement back to the \'etale site.

\begin{proof}
By Corollary~\ref{pe:classical-derived}, write $K=\nu_X^*K_{\et}$
with $K_{\et}\in D^+(X_{\et},\Lambda)$. On the \'etale site,
Proposition~\ref{pbc:reduction} proves proper base change for torsion
abelian sheaves, with the projective-line case
Proposition~\ref{pbc:p1} as its input. Three passages lead from there
to the theorem: to bounded below complexes, to $\Lambda$-modules, and
to the pro-\'etale site.

For the first two, write $K$ for the \'etale complex $K_{\et}$, and
suppose first that it is a bounded below complex of abelian sheaves
with torsion cohomology. Fix an integer $m$. Both functors
$g^{-1}Rf_*,\ Rf'_*g'^{-1}\colon D^+(X_{\et},\Z)\to D^+(Y'_{\et},\Z)$
send $D^{\geq m+1}$ to $D^{\geq m+1}$, so their $m$th
cohomology on $K$ is unchanged upon replacing $K$ by the bounded
truncation $\tau_{\leq m}K$. Using truncation triangles, its finite
Postnikov filtration reduces the exchange isomorphism in degree $m$ to
torsion sheaves and their shifts.

Next forget the $\Lambda$-action. The forgetful functor to abelian
sheaves is exact and conservative, and commutes with inverse image
and direct image. It also commutes with bounded below derived direct
image. Indeed, an injective sheaf of $\Lambda$-modules is acyclic, as an
abelian sheaf, for sections over every object of the site
(\stacks{03FD}); the same assertion on the \'etale slices gives
acyclicity for direct image, so a module-injective resolution
computes the underlying abelian derived direct image.
These identifications preserve the units and counits defining the
exchange. The underlying abelian map is therefore the exchange in
$D^+(Y'_{\et},\Z)$ already shown to be an isomorphism, and hence so is
the exchange in $D^+(Y'_{\et},\Lambda)$.

Finally apply $\nu_{Y'}^*$. Since $f$ and $f'$ are proper, hence qcqs,
the direct-image comparison of
Theorem~\ref{pe:direct-image-comparison} applies. Together with the
inverse-image comparison $g_{\proet}^*\nu_Y^*\simeq\nu_X^*g_{\et}^*$ it
identifies $\nu_{Y'}^*g^{-1}Rf_*K_{\et}$ and
$\nu_{Y'}^*Rf'_*g'^{-1}K_{\et}$ with $g^*Rf_*K$ and $Rf'_*g'^*K$ in
$D^+(Y'_{\proet},\Lambda)$. These comparisons respect the units and counits, so the
resulting isomorphism is the canonical pro-\'etale exchange map.
\end{proof}

\begin{corollary}[Geometric-fibre formula]\label{pbc:torsion-stalk}
Under the hypotheses of Theorem~\ref{pbc:torsion-theorem}, for every
geometric point $\bar y:\operatorname{Spec}(\Omega)\to Y$ there is a
canonical isomorphism in the derived category of $\Lambda$-modules
\[
 (Rf_*K)_{\bar y}\simeq
 R\Gamma\bigl(X_{\bar y,\proet},K|_{X_{\bar y}}\bigr).
\]
For a classical torsion sheaf $F$ of $\Lambda$-modules this says
\[
 (R^qf_*F)_{\bar y}\simeq
 H^q_{\proet}(X_{\bar y},F|_{X_{\bar y}})\qquad(q\geq0).
\]
\end{corollary}

\begin{proof}
Take $Y'=\operatorname{Spec}(\Omega)$ in the theorem and apply
global sections of $Y'_{\proet}$, which are exact on this pro-\'etale
geometric point by Proposition~\ref{pe:condensed-point}. Composition
gives the fibre cohomology on the right. On the left,
$(Rf_*K)_{\bar y}$ means $\Gamma(Y'_{\proet},\bar y^*Rf_*K)$. Since
the topos of $Y'_{\proet}$ is the topos of condensed sets, it is
classicality that identifies this complex with the
usual geometric stalk. Through the comparison of bounded below
complexes, the formula
also recovers the \'etale fibre formula used inside the proof.
\end{proof}

\begin{remark}[The role of properness]\label{pbc:reduction-inputs}
The reduction in Proposition~\ref{pbc:reduction} uses Chow's lemma, the
stalk formula for finite morphisms, the proper quasi-finite criterion,
and the Grothendieck spectral sequence. Its cohomological content is
the projective-line argument, together with degree-zero proper base
change and the injective-acyclicity criterion. Properness enters three
times: in the degree-zero comparison, in making the Chow cover
projective, and in the cancellation after base change.
\end{remark}

\subsection{Adic and condensed coefficient forms}
\label{pbc:condensed}

Recall from Section~\ref{n:coefficients} that a discrete finite
commutative ring $A$ gives a condensed ring whose pullback $A_X$ is the
usual constant coefficient sheaf. Finite rings are thus coefficient
rings in the same sense as $\mathcal O_E$, $E$ and $\Qbar$, and
Theorem~\ref{pbc:torsion-theorem} gives proper base change at each
finite level for every admissible coefficient ring, including $\Zl$
and $\Ql$.

\begin{corollary}[Classical finite coefficients on the pro-\'etale site]
\label{pbc:proet-finite}
Let $A$ be a finite commutative ring, $p:X\to Y$ proper, and $g:Y'\to Y$ any
morphism. Let $p':X'\to Y'$ and $g':X'\to X$ be the projections from
$X'=X\times_Y Y'$. If $K\in D^+_{\mathrm{class}}(X_{\proet},A)$, then the
pro-\'etale exchange map of Theorem~\ref{pbc:torsion-theorem}
\[
 g^*Rp_*K\longrightarrow Rp'_*g'^*K
\]
is an isomorphism in $D^+(Y'_{\proet},A)$.
\end{corollary}
\begin{proof}
This is the finite-coefficient instance of
Theorem~\ref{pbc:torsion-theorem}. By
Theorem~\ref{pe:constructible-comparison}, its classicality hypothesis
is automatic for bounded constructible finite-coefficient complexes on
qcqs schemes, in particular in our standing finite-type range, and
hence for every $K\in\Dctf(X,A)$ in that range.
\end{proof}

To pass from finite coefficients to $\mathcal O_E$ one has to deal with
completeness: the inverse image
$g^*\colon D(Y_{\proet},\mathcal O_E)\to D(Y'_{\proet},\mathcal O_E)$,
exact on module sheaves, need not preserve derived completeness. For
$M\in D(Y_{\proet},\mathcal O_E)$ define its \emph{completed inverse
image} $g^*_{\mathrm{comp}}M\in D(Y'_{\proet},\mathcal O_E)$ by
\[
 g^*_{\mathrm{comp}}M
   =R\varprojlim_n g^*(M\otimes^{\mathbf L}_{\mathcal O_{E,Y}}
                                  (\mathcal O_E/\pi_E^n)_Y).
\]
The corollary above and the proposition below hold for arbitrary
schemes. Statements about the constructible categories $\Dctf$ or
$\Dbc$ in this subsection are made under the finite-type $k$-scheme
assumptions of the introduction. In that setting the completed inverse
image of a constructible integral complex agrees with $g^*$, by the
integral constructibility criterion and the finite constant local
models of Section~\ref{n:constructible}. The completed form is the one
available before finiteness of proper direct image has been established.

\begin{proposition}[Proper base change for complete integral complexes]
\label{pbc:complete}
In the square above, fix $E/\Ql$ finite. Suppose that
$K\in D(X_{\proet},\mathcal O_E)$ is derived $\pi_E$-complete and that
each reduction
\[
 K_n=K\otimes^{\mathbf L}_{\mathcal O_{E,X}}
                           (\mathcal O_E/\pi_E^n)_X
\]
belongs to $D^+_{\mathrm{class}}(X_{\proet},\mathcal O_E/\pi_E^n)$.
Then the
completed exchange map, constructed in the proof,
\[
 g^*_{\mathrm{comp}}Rp_*K
       \xrightarrow{\ \sim\ }Rp'_*g'^*_{\mathrm{comp}}K
\]
is an isomorphism in $D(Y'_{\proet},\mathcal O_E)$.
\end{proposition}
\begin{proof}
Derived completeness identifies $K$ with $R\varprojlim_nK_n$.
The quotient $\mathcal O_E/\pi_E^n$ has the two-term finite free
resolution $[\mathcal O_E\xrightarrow{\pi_E^n}\mathcal O_E]$, so
reduction is computed by a cone and commutes with the exact derived
functor $Rp_*\colon D(X_{\proet},\mathcal O_E)\to
D(Y_{\proet},\mathcal O_E)$, giving isomorphisms in
$D(Y_{\proet},\mathcal O_E)$
\[
 (Rp_*K)\otimes^{\mathbf L}(\mathcal O_E/\pi_E^n)_Y
                         \simeq Rp_*K_n.
\]
These maps are compatible along the coefficient tower, and taking
derived limits gives
\[
 \begin{aligned}
 g^*_{\mathrm{comp}}Rp_*K
    &\simeq R\varprojlim_n g^*Rp_*K_n,\\
 Rp'_*g'^*_{\mathrm{comp}}K
    &\simeq R\varprojlim_n Rp'_*g'^*K_n,
 \end{aligned}
\]
where the second line uses that derived direct image commutes with
derived inverse limits. The completed exchange map is the derived
limit of the exchange maps $g^*Rp_*K_n\to Rp'_*g'^*K_n$, which are
isomorphisms in $D^+(Y'_{\proet},\mathcal O_E/\pi_E^n)$ by
Corollary~\ref{pbc:proet-finite}. The comparison of towers is
therefore an isomorphism in the derived category of towers of
$\mathcal O_E$-module sheaves on $Y'_{\proet}$,
since its cone has zero cohomology at every level, and applying
$R\varprojlim$ proves the claim.

Equivalently, one may work over the constant sheaf $\mathcal O_E$ and
apply \stacks{09C9}, whose perfect-quotient hypothesis is the two-term
resolution used above. The constant and completed coefficient rings
have equivalent derived-complete categories by \stacks{099I}, so this
yields the same statement for the pulled-back condensed ring
$\mathcal O_{E,X}$.
\end{proof}

For a constructible $K$, proper finiteness will show that $Rp_*K$ is
again constructible. All inverse images in the preceding proposition
are then already complete, and it becomes the ordinary exchange
isomorphism, between objects of $\Dbc(Y',\mathcal O_E)$,
\[
 g^*Rp_*K\xrightarrow{\sim}Rp'_*g'^*K
 \qquad(K\in\Dbc(X,\mathcal O_E)).
\]
Proposition~\ref{b:prop:proper-coefficients} proves this finiteness,
using Theorem~\ref{pbc:torsion-theorem} to identify the finite-level
exchange maps. Its proof then treats $E$ and $\Qbar$: stratified
lattices reduce to the integral case, inversion of $\ell$ commutes
with bounded-below proper direct image, and descent to a finite
extension of $\Ql$ handles $\Qbar$. For these coefficient rings
$\Dctf=\Dbc$, and for a finite admissible ring $A$ the torsion theorem
applies in particular to $\Dctf(X,A)$.

On the localizing categories $\Loc(\Dctf(X,\Lambda))$ the exchange
isomorphism is proved later. For finite $A$, the classical comparison
and the finite cohomological bound give the exchange isomorphism
$g^*Rp_*\simeq Rp'_*g'^*$ of functors
$\Loc(\Dctf(X,A))\to\Loc(\Dctf(Y',A))$ in
Proposition~\ref{ctf:finite-proper}. For the integral
and rational coefficient rings, once proper direct image is known to
preserve direct sums, its exchange isomorphism on the constructible
generators extends to an isomorphism of functors
$\Loc(\Dctf(X,\Lambda))\to\Loc(\Dctf(Y',\Lambda))$ by
Proposition~\ref{b:prop:proper-base-change-unbounded}. The bounded-below
torsion theorem needs no Tor-amplitude condition; finite Tor amplitude
enters through the generators of $\Loc(\Dctf(-,\Lambda))$ used in the Brown
construction.

\begin{example}[Fibrewise computation and the need for properness]
For a proper $p:X\to Y$ and a geometric point $\bar y\to Y$, the
torsion theorem gives
\[
 (R^qp_*F)_{\bar y}\simeq H^q(X_{\bar y},F|_{X_{\bar y}}).
\]
The same formula for constructible $\mathcal O_E$, $E$ and $\Qbar$
coefficients follows from the comparisons just described. Properness
is essential: if
$j:\mathbf G_m\hookrightarrow\mathbf A^1$ and $i:\{0\}\hookrightarrow
\mathbf A^1$ over an algebraically closed field, then
$i^*j_*A$ is nonzero for a nonzero finite constant ring $A$, whereas
the base change of $\mathbf G_m$ to $\{0\}$ is empty. The degree-zero
exchange map $i^{-1}j_*A\to0$ of sheaves on $\{0\}_{\et}$ already fails to be an isomorphism in this nonproper
example.
\end{example}

\section{Constructible sheaves and reduction to finite coefficients}
\label{b:sec:coefficients}

The common six-functor category is $\Dctf(X,\Lambda)$ for every
admissible $\Lambda$. For finite commutative $\ell$-primary rings,
including $\mathcal O_E/\pi_E^n$, its finiteness properties are
classical \'etale results, used with finite Tor amplitude
(Proposition~\ref{ctf:finite-proper}). For $\mathcal O_E$, $E$ and
$\Qbar$ we obtain the $\ell$-adic counterparts by comparison with finite
levels. For these integral and rational coefficients $\Dctf=\Dbc$, so
we use $\Dbc$ when discussing cohomology sheaves, truncations and
lattices. Three facts make the comparison possible. An object of
$\Dbc(X,\Qbar)$ is defined over some finite extension $E$ of $\Ql$; an
object of $\Dbc(X,E)$ is built by finitely many cones from sheaves
$a_!L$, with $L$ lisse on a locally closed stratum; and after refining
the stratum, $L$ contains a lattice $L_0$, so that $a_!L_0$ is an
integral object whose reductions modulo $\pi_E^n$ are classical torsion
sheaves. The first statement lifted in this way is that proper direct
image preserves bounded constructible complexes.

\subsection{Inverse image of the coefficient sheaves}

Because all coefficient sheaves are pulled back from the point
(\S\ref{n:coefficients}), they are compatible with inverse image. The
same holds for their Tate twists, and a finite extension of coefficient
fields gives a finite free extension of coefficient sheaves; both facts
are used repeatedly below.

\begin{lemma}[Inverse image of the coefficient sheaves]
\label{b:lem:coefficient-pullback}
Let $f:X\to Y$ be a morphism of schemes.
\begin{enumerate}
\item For every admissible $\Lambda$, the canonical map
  \[
   f^{-1}\Lambda_Y\longrightarrow\Lambda_X
  \]
  is an isomorphism of sheaves of rings on $X_{\proet}$. The same holds for its Tate
  twists, including the twists of a finite ring $A$.
\item If $E\subset F$ is a finite extension, then $\mathcal O_{F,X}$ is a
  finite free $\mathcal O_{E,X}$-module and $F_X$ is a finite free
  $E_X$-module.
\end{enumerate}
\end{lemma}

By (1), the pullback $f^*\colon\Mod(Y_{\proet},\Lambda)\to\Mod(X_{\proet},\Lambda)$
of modules over these sheaves is the exact functor $f^{-1}$, as already noted in \S\ref{n:ambient}.

\begin{proof}
(1) For every condensed ring $\Lambda$, the identity $p_Y\circ f=p_X$
gives
\[
 f^{-1}\Lambda_Y=f^{-1}p_Y^*\Lambda\simeq p_X^*\Lambda=\Lambda_X,
\]
which gives the assertion uniformly, including finite $A$.

The torsion twists $\mu_{\ell^n}$ are represented by finite \'etale group
schemes, which commute with base change. Their inverse limit is
represented by the inverse limit of these schemes, so its inverse image
is represented by the base change (\stacks{09A5}). This gives the
assertion for $\mathcal O_E(1)$. For a finite $A$ killed by
$\ell^n$, its twist is obtained by extending scalars from
$\mathbf Z/\ell^n\mathbf Z$ in the finite-level twist, so inverse image
commutes with it as well. Inverting $\ell$, extending scalars, and
taking tensor powers and duals gives the other twists.

(2) An $\mathcal O_E$-basis of $\mathcal O_F$ is an isomorphism of
topological $\mathcal O_E$-modules $\mathcal O_F\cong\mathcal O_E^{[F:E]}$,
hence of condensed modules
$\underline{\mathcal O_F}\cong\underline{\mathcal O_E}^{[F:E]}$. Pull it
back along $p_X$, and invert $\ell$ for the fields.
\end{proof}

Thus the condensed definition gives functoriality, while the explicit
representing scheme $\varprojlim_n\coprod_{a\in\mathcal O_E/\pi_E^n}X$
of Lemma~\ref{n:coefficient-comparison} gives the comparison with the
usual $\ell$-adic coefficients.

\subsection{Constructible sheaves with field coefficients}

In this subsection $\Lambda=E$ or $\Qbar$, so $\Dctf=\Dbc$ and we can
argue with constructible cohomology sheaves and standard truncations.
For a general finite $A$, finite Tor amplitude is a condition on the whole complex, and
the finite flat models of Lemma~\ref{ctf:finite-models} replace this
d\'evissage through the heart. As explained after
Definition~\ref{n:def-constructible}, a $\Lambda_X$-module is
constructible exactly when it is \emph{lisse}, that is locally free of
finite rank for the pro-\'etale topology, on each stratum of a finite
stratification of $X$. D\'evissage arguments reduce statements about
$\Dbc(X,\Lambda)$ to the simplest constructible sheaves, lisse sheaves
extended by zero from a stratum. For a locally closed immersion $a:S\hookrightarrow X$,
factor $a=i\circ u$ with $u:S\hookrightarrow\overline S$ the open
immersion into the closure and $i:\overline S\hookrightarrow X$ the closed
immersion, and put
$a_!=i_*u_!\colon\Mod(S_{\proet},\Lambda)\to\Mod(X_{\proet},\Lambda)$.
Both $u_!$ and $i_*$ are exact, so $a_!$ is exact, and it induces
$a_!\colon\D(S_{\proet},\Lambda)\to\D(X_{\proet},\Lambda)$ without deriving.

\begin{lemma}[Constructible sheaves]
\label{b:lem:coefficient-heart}
Let $\Lambda=E$ or $\Qbar$.
\begin{enumerate}
\item $\Cons(X,\Lambda)$ is closed under kernels, cokernels and
  extensions in $\Mod(X_{\proet},\Lambda)$. In particular it is an abelian
  category, and its inclusion into $\Mod(X_{\proet},\Lambda)$ is exact.
\item $\Cons(X,\Lambda)$ is essentially small.
\item Every $F\in\Cons(X,\Lambda)$ has a finite filtration
  $0=F_0\subset F_1\subset\cdots\subset F_m=F$ whose subquotients
  $F_k/F_{k-1}$ are of the form $a_!L$, with $a:S\hookrightarrow X$ a
  locally closed immersion and $L$ a lisse $\Lambda_S$-module.
\end{enumerate}
\end{lemma}
For a curve $X$, part (3) says that a constructible sheaf $F$ sits in an
exact sequence $0\to u_!L\to F\to i_*G\to0$, where $u$ is the inclusion
of a dense open subset on which $F$ is lisse, $L=u^*F$, and $i_*G$ is
supported on the finitely many closed points of the complement.

\begin{example}[A running example]\label{b:ex:running}
We use the same few schemes throughout:
\[
 \{0\}\xrightarrow{\ i\ }\mathbf A^1_k\xleftarrow{\ u\ }\mathbf G_{m,k},
 \qquad \mathbf A^1_k\xrightarrow{\ j\ }\mathbf P^1_k .
\]
The sheaf $u_!\Lambda_{\mathbf G_m}$ is constructible for the
stratification $\mathbf A^1=\mathbf G_m\sqcup\{0\}$. Its stalk is
$\Lambda$ at a geometric point of $\mathbf G_m$ and $0$ at the origin.
It is not lisse on $\mathbf A^1$, and the sequence of part (3) is the
localization sequence
\[
 0\to u_!\Lambda_{\mathbf G_m}\to\Lambda_{\mathbf A^1}\to i_*\Lambda\to0 .
\]
The example returns for compact supports (after
Lemma~\ref{b:lem:geo-p1-cohomology}), for $i^!$ (Example~\ref{b:ex:closed-support}) and
for the trace (\S\ref{sp:grp-trace}).
\end{example}

\paragraph*{Exercises.} (1) Compute the stalks of $a_!L$ for a locally
closed immersion $a$ and a lisse $L$. (2) For $K=\Lambda_X[1]$, in
which degree is the nonzero cohomology sheaf? (3) Show that
$u_!\Lambda_{\mathbf G_m}$ has no nonzero global sections on
$\mathbf A^1_k$.

\begin{proof}[Proof of Lemma~\ref{b:lem:coefficient-heart}]
(1) Given finitely many constructible sheaves, refine their
stratifications to a common one; all three closure properties can then
be checked on one stratum $S$.

\emph{Extensions.} Let $0\to F'\to F\to F''\to0$ be exact with $F'$ and
$F''$ lisse on $S$. After passing to a cover, $F''$ is free, and a basis
of $F''$ lifts to sections of $F$. This splits the sequence locally, so
$F$ is locally free of finite rank.

\emph{Kernels and cokernels.} On the topologically Noetherian scheme
$S$, the kernel and cokernel of a morphism of lisse $\Lambda_S$-modules
are lisse (Bhatt--Scholze, \emph{The pro-\'etale topology for schemes},
Corollary~6.8.5); in particular its rank is locally constant on $S$, and
no further stratification is needed. This applies to $E$ and to $\Qbar$,
since our coefficient sheaves are those of loc.\ cit., by its
Lemma~6.8.2 and Lemma~\ref{n:coefficient-comparison}. Restriction to
strata is exact, so the kernel and cokernel on $X$ are constructible.

(2) Choose the sites in a fixed universe. A constructible sheaf is
determined up to isomorphism by a finite stratification together with,
on each stratum, finite-rank descent data on a covering, and these form a
set.

(3) Refine the stratification of $F$, by Noetherian induction, so that
there is an open stratum $u:S\hookrightarrow X$ with closed complement
$i:Z\hookrightarrow X$ that is again a union of strata. The exact sequence
$0\to u_!u^*F\to F\to i_*i^*F\to0$ splits off the stratum $S$, and
$u^*F$ is lisse. Applying the same step to $i^*F$ on $Z$ and iterating
gives a finite filtration of $F$ whose subquotients are $a_!(F|_S)$ for
the strata $a:S\hookrightarrow X$, and each $F|_S$ is lisse.
\end{proof}

\begin{corollary}\label{b:lem:coefficient-heart-b}
Let $\Lambda=E$ or $\Qbar$. Then $\Dbc(X,\Lambda)$ is a triangulated
subcategory of $\D(X_{\proet},\Lambda)$. Every object of $\Dbc(X,\Lambda)$
is obtained by finitely many cones from objects $a_!L[n]$, with $a$ and
$L$ as in Lemma~\ref{b:lem:coefficient-heart}(3) and $n\in\mathbf Z$.
\end{corollary}
\begin{proof}
Let $K\to K'\to K''\to K[1]$ be a distinguished triangle with $K,K'$ in
$\Dbc(X,\Lambda)$. By the long exact cohomology sequence, each
$\mathcal H^q(K'')$ is an extension of
$\ker(\mathcal H^{q+1}K\to\mathcal H^{q+1}K')$ by
$\coker(\mathcal H^qK\to\mathcal H^qK')$. Both are constructible by
Lemma~\ref{b:lem:coefficient-heart}(1), hence so is $\mathcal H^q(K'')$.
Thus $\Dbc(X,\Lambda)$ is closed under cones, and it is clearly closed under
shifts.

For the second assertion, induct on the number of nonzero cohomology
sheaves of $K$. If $q$ is the largest degree with $\mathcal H^q(K)\ne0$,
the truncation triangle $\tau_{<q}K\to K\to\mathcal H^q(K)[-q]\to$ reduces
the claim to a single shifted constructible sheaf. Its filtration from
Lemma~\ref{b:lem:coefficient-heart}(3) gives finitely many short exact
sequences, that is, distinguished triangles, with terms $a_!L$.
\end{proof}

\subsection{Continuity of cohomology and of Ext}
\label{b:sec:ext-continuity}

The descent argument below, the uniform Ext bound and the compactness
argument for $E$ and $\Qbar$ (\S\ref{b:sec:compactness}) all use one
fact: cohomology, and Hom out of a constructible complex, commute with
filtered colimits of complexes with a common lower bound. It is needed
for $\Lambda=E$, $\Qbar$, and for $\mathcal O_E$, since the uniform Ext
bound is proved first for integral lattices and then transported to $E$
by inverting $\pi_E$.

We start with continuity of pro-\'etale cohomology, and of higher
direct images, for filtered colimits of complexes of sheaves that are
bounded below uniformly. It will be applied both to $R\Gamma$ on a
stratum and to $Rj_*$ for an open immersion $j$. Here a filtered colimit
of complexes means the termwise colimit of a filtered diagram of
complexes of sheaves.

\begin{lemma}[Continuity of cohomology]\label{b:lem:bounded-continuity}
Let $Y$ be a quasi-compact, quasi-separated scheme, let $A$ be a sheaf of
rings on $Y_{\proet}$, and let $(K_i^\bullet)$ be a filtered diagram of
complexes of $A$-modules that all vanish in degrees below a fixed integer
$a$. Then for every $q$ the natural map
\[
 \varinjlim_i H^q(Y,K_i^\bullet)
 \longrightarrow
 H^q(Y,\varinjlim_i K_i^\bullet)
\]
is an isomorphism. If $g:Y\to Z$ is quasi-compact and quasi-separated,
then likewise $\varinjlim_i R^qg_*K_i^\bullet\to R^qg_*(\varinjlim_iK_i^\bullet)$
is an isomorphism of sheaves on $Z_{\proet}$ for every $q$.
\end{lemma}
\begin{proof}
Choose a hypercovering $W_\bullet\to Y$ by affine weakly contractible
objects; by Proposition~\ref{pe:hypercover}, $R\Gamma(Y,K_i^\bullet)$ is
computed by the total complex of the double complex
$K_i^m(W_n)$. The point is that this total complex commutes with filtered
colimits.

First, the functor $M\mapsto M(W_n)$ is exact and commutes with direct
sums (Corollary~\ref{pe:sections}). An exact additive functor that
commutes with direct sums commutes with all colimits, since a colimit is
the cokernel of a map between two direct sums. Second, the term of total
degree $q$ is $\bigoplus_{n+m=q}K_i^m(W_n)$ with $n\geq0$ and $m\geq a$,
a finite direct sum with $0\leq n\leq q-a$. Hence the total complex of
$\varinjlim_iK_i^\bullet$ is the colimit of the total complexes, and
since filtered colimits of modules are exact, taking cohomology proves
the first assertion.

For the second, $R^qg_*K$ is the sheafification of
$V\mapsto H^q(Y\times_ZV,K)$ on affine objects $V$ of $Z_{\proet}$. Each
$Y\times_ZV$ is quasi-compact and quasi-separated, so the first assertion
applies, and sheafification commutes with colimits.
\end{proof}

Lemma~\ref{b:lem:coefficient-heart}(3)
gives a constructible sheaf with coefficients $E$ or $\Qbar$ a finite filtration whose subquotients are of the form $a_!L$, with
$a:S\hookrightarrow X$ locally closed and $L$ lisse on $S$. To reduce
statements about $\RHom(a_!L,-)$ to cohomology on $S$ we use the functor
$R\underline\Gamma_Z$ of sections with support in a closed subscheme $Z$.

\begin{lemma}[Sections with support in a closed subscheme]
\label{b:lem:formal-closed-support}
Let $\Lambda$ be admissible, and let $i:Z\hookrightarrow X$ be a
closed immersion with open complement $j:V\hookrightarrow X$.
For a sheaf $M$ of $\Lambda_X$-modules on $X_{\proet}$ let
$\underline\Gamma_Z(M)=i^*\ker(M\to j_*j^*M)$, the sheaf on $Z$ of
sections of $M$ with support in $Z$, and let
$R\underline\Gamma_Z\colon\D(X_{\proet},\Lambda)\to\D(Z_{\proet},\Lambda)$ be its
right derived functor.
\begin{enumerate}
\item $R\underline\Gamma_Z$ is right adjoint to the direct image
  $i_*$ from $\D(Z_{\proet},\Lambda)$ to $\D(X_{\proet},\Lambda)$, also in
  the form $\RHom_X(i_*A,K)=\RHom_Z(A,R\underline\Gamma_ZK)$, and
  \[
   R\underline\Gamma_Z K=i^*\fib(K\longrightarrow Rj_*j^*K)
   \quad\text{in }\D(Z_{\proet},\Lambda).
  \]
\item Let $a=i\circ u:S\hookrightarrow X$ with $u:S\hookrightarrow Z$ an
  open immersion, and let $L$ be a finite locally projective
  $\Lambda_S$-module with dual $L^\vee$. Then naturally in
  $K\in\D(X_{\proet},\Lambda)$
  \begin{equation}\label{b:eq:stratum-hom}
   \RHom_X(a_!L,K)
   =R\Gamma\bigl(S,L^\vee\otimes u^*R\underline\Gamma_ZK\bigr).
  \end{equation}
\end{enumerate}
\end{lemma}
\begin{proof}
We use three facts about the closed immersion $i$: $i^*$ is exact,
preserves limits and satisfies $i^*i_*=\mathrm{id}$ (\stacks{09AA});
$i_*$ is exact (\stacks{09BM}); and every sheaf $M$ sits in the
localization sequence $0\to j_!j^*M\to M\to i_*i^*M\to0$ (\stacks{09AH}).
By the last fact, a sheaf whose restriction to $V$ vanishes is of the form
$i_*$ of a sheaf on $Z$; this applies to $\ker(M\to j_*j^*M)$.

\emph{Adjunction.} A map $i_*A\to M$ composes to zero with
$M\to j_*j^*M$, because $j^*i_*A=0$. So it factors through
$\ker(M\to j_*j^*M)=i_*\underline\Gamma_Z(M)$, and by full faithfulness of
$i_*$ it corresponds to a map $A\to\underline\Gamma_Z(M)$. Thus
$\underline\Gamma_Z$ is right adjoint to
$i_*\colon\Mod(Z_{\proet},\Lambda)\to\Mod(X_{\proet},\Lambda)$. Since $i_*$ is
exact, $\underline\Gamma_Z$ preserves K-injective complexes, and K-injective
resolutions exist (\stacks{079P}). Hence $R\underline\Gamma_Z$ is computed
on K-injective complexes, and the adjunction passes to $\RHom$.

\emph{Fiber formula.} For an injective sheaf $I$ the unit
$I\to j_*j^*I$ of the adjunction $j^*\dashv j_*$ on
$\Mod(X_{\proet},\Lambda)$ is surjective. Indeed, by injectivity the inclusion
$j_!j^*I\to I$ extends along $j_!j^*I\hookrightarrow j_*j^*I$ to a map
$s:j_*j^*I\to I$. The composite of $s$ with the unit is the identity
after applying $j^*$, hence is the identity, because $j_*$ is fully
faithful. Now represent $K$ by a K-injective complex $I^\bullet$ with
injective terms (\stacks{079P}). Then $j^*I^\bullet$ is K-injective,
since the left adjoint $j_!$ of $j^*$ is exact, so $j_*j^*I^\bullet$
represents $Rj_*j^*K$. The termwise exact sequences
\[
 0\to i_*\underline\Gamma_Z(I^n)\to I^n\to j_*j^*I^n\to0
\]
give a distinguished triangle
$i_*R\underline\Gamma_ZK\to K\to Rj_*j^*K\to i_*R\underline\Gamma_ZK[1]$
in $\D(X_{\proet},\Lambda)$.
Applying $i^*$ and $i^*i_*=\mathrm{id}$ gives the fiber formula.

\emph{Formula \eqref{b:eq:stratum-hom}.} Since $a_!=i_*u_!$, part~(1),
the adjunction $u_!\dashv u^*$ and the identity
$\cHom(L,-)=L^\vee\otimes-$ for the finite locally projective sheaf $L$ give
$\RHom_X(i_*u_!L,K)=\RHom_Z(u_!L,R\underline\Gamma_ZK)
=R\Gamma(S,L^\vee\otimes u^*R\underline\Gamma_ZK)$.
\end{proof}

\begin{lemma}[Continuity of Ext out of a constructible complex]
\label{b:lem:constructible-ext-continuity}\label{b:lem:constructible-ext-continuity-b}
Let $\Lambda=E$, $\Qbar$ or $\mathcal O_E$. Let $C\in\D(X_{\proet},\Lambda)$
be obtained by finitely many extensions and shifts from sheaves $a_!L$,
with $a:S\hookrightarrow X$ a locally closed immersion and $L$ a finite
locally free $\Lambda_S$-module. Let $(K_i^\bullet)$ be a filtered diagram
of complexes of $\Lambda_X$-modules that all vanish below a fixed degree.
Then for every integer $q$ the natural map of $\Hom$ groups in
$\D(X_{\proet},\Lambda)$
\[
 \varinjlim_i\Hom(C,K_i^\bullet[q])
 \longrightarrow
 \Hom\bigl(C,(\varinjlim_i K_i^\bullet)[q]\bigr)
\]
is an isomorphism. The hypothesis on $C$ holds:
\begin{enumerate}
\item for every $C\in\Dbc(X,\Lambda)$ if $\Lambda=E$ or $\Qbar$, by
  Corollary~\ref{b:lem:coefficient-heart-b}; in particular, for
  $C\in\Cons(X,\Lambda)$ the functors $\Ext^q_\Lambda(C,-)$ commute with
  filtered colimits of $\Lambda_X$-modules;
\item for every finite extension of sheaves $a_!L_0$ with $L_0$ a finite
  locally free $\mathcal O_{E,S}$-module, if $\Lambda=\mathcal O_E$.
\end{enumerate}
\end{lemma}
\begin{proof}
Both sides are cohomological functors of $C$, so by the five lemma the
class of $C$ for which the claim holds is closed under extensions and
shifts. It therefore suffices to treat $C=a_!L$, with $a=i\circ
u:S\hookrightarrow Z\hookrightarrow X$ as in
Lemma~\ref{b:lem:formal-closed-support}. The only property of $L$ used
below is that $L^\vee\otimes-$ is exact.

By \eqref{b:eq:stratum-hom} we must show that the cohomology of
$R\Gamma(S,L^\vee\otimes u^*R\underline\Gamma_ZK)$ commutes with the
colimit. Apply the exact functor
$L^\vee\otimes u^*i^*(-)\colon\D(X_{\proet},\Lambda)\to\D(S_{\proet},\Lambda)$
to the triangle $K\to Rj_*j^*K$ whose fiber defines $R\underline\Gamma_Z$. By the
long exact sequence and the five lemma, it suffices to treat the two
terms
\[
 R\Gamma\bigl(S,L^\vee\otimes u^*i^*K\bigr)
 \quad\text{and}\quad
 R\Gamma\bigl(S,L^\vee\otimes u^*i^*Rj_*j^*K\bigr).
\]
The functor $L^\vee\otimes u^*i^*(-)$ is exact, commutes with colimits and
preserves the lower bound, so Lemma~\ref{b:lem:bounded-continuity} on $S$
handles the first term. For the second, the hypercohomology spectral
sequence
\[
 E_2^{p,r}=H^p\bigl(S,L^\vee\otimes u^*i^*R^rj_*(j^*K)\bigr)
 \Longrightarrow
 H^{p+r}\bigl(S,L^\vee\otimes u^*i^*Rj_*j^*K\bigr)
\]
is functorial in $K$. The sheaves $R^rj_*(j^*K_i)$ commute with the
colimit by the second part of Lemma~\ref{b:lem:bounded-continuity}, and
then $H^p(S,-)$ does by its first part. So the map of spectral sequences
is an isomorphism on $E_2$, after passing to the colimit. Because of the
common lower bound, only finitely many $E_2^{p,r}$ contribute to a given
total degree, so the abutments agree as well.
\end{proof}

\subsection{Descent to a finite coefficient field}

The coefficient ring $\Qbar$ is a filtered union of the fields $E$, but
integral lattices exist only over a single $E$. We therefore show that
every object and every morphism of $\Dbc(X,\Qbar)$ is defined over some
finite $E$. For $C\in\D(X_{\proet},E)$ and a finite extension $F/E$ we
write $C_F=F_X\otimes_{E_X}C$ and $C_{\Qbar}=\QbarX\otimes_{E_X}C$ for the
extensions of scalars, given by functors
$\D(X_{\proet},E)\to\D(X_{\proet},F)$ and
$\D(X_{\proet},E)\to\D(X_{\proet},\Qbar)$; since $F_X$ and $\QbarX$ are flat over $E_X$, these
need no deriving.

\begin{lemma}[Descent to a finite coefficient field]
\label{b:lem:coefficient-finite-descent}
\leavevmode
\begin{enumerate}
\item Every lisse $\QbarX$-module is isomorphic to $L_{\Qbar}$ for a lisse
  $E_X$-module $L$ and some finite $E/\Ql$.
\item For $C,D\in\Dbc(X,E)$ and every integer $r$, extension of scalars
  induces an isomorphism
  \[
   \varinjlim_{F\supset E}
   \Hom_{\D(X_{\proet},F)}(C_F,D_F[r])
   \;\xrightarrow{\ \sim\ }\;
   \Hom_{\D(X_{\proet},\Qbar)}(C_{\Qbar},D_{\Qbar}[r]),
  \]
  the colimit running over finite extensions $F$ of $E$ in $\Qbar$.
\item Every object of $\Dbc(X,\Qbar)$ is isomorphic to $K_{\Qbar}$ for
  some $K\in\Dbc(X,E)$ and some finite $E/\Ql$.
\end{enumerate}
\end{lemma}

To descend a lisse $\QbarX$-module we work with its transition matrices on a
pro-\'etale cover. Two facts describe the entries of such matrices of the entries of such matrices
on a quasi-compact object $V$ of $X_{\proet}$.
\begin{itemize}
\item Sections of a filtered colimit of sheaves over $V$ are the filtered
  colimit of the sections. Indeed, every covering of $V$ has a finite
  refinement by quasi-compact objects, and the sheaf condition for a
  finite covering is a finite limit, which commutes with filtered
  colimits (this is the degree-zero case of \stacks{0739}). Since
  $\QbarX=\varinjlim_EE_X$, finitely many sections of $\QbarX$ over $V$
  all lie in $E_X(V)$ for one finite $E/\Ql$.
\item $E_X(V)$ is the ring of continuous functions $|V|\to E$. A
  continuous function on a quasi-compact space has compact image, which
  lies in $\ell^{-m}\mathcal O_E$ for some $m$, so this follows from the
  description of $\mathcal O_{E,X}$ in
  Lemma~\ref{n:coefficient-comparison}.
\end{itemize}
Affine, hence quasi-compact, pro-\'etale coverings and quasi-compact
refinements exist by \stacks{0988}.

\begin{proof}
(1) Let $L$ be a lisse $\QbarX$-module. Since $X$ is quasi-compact, $L$
is free on a finite covering of $X$ by quasi-compact (we may take affine)
pro-\'etale objects $V_k$. Since $X$ is quasi-separated, the
intersections $V_k\times_XV_{k'}$ are quasi-compact. On them, $L$ is
described by finitely many transition matrices with entries in
$\QbarX$, and so are their inverses. By the two facts above, all these entries lie in $E_X$ for
one finite $E$. The cocycle identities hold over $E_X$ because
$E_X\to\QbarX$ is injective. Descent of sheaves along the covering then
produces a lisse $E_X$-module $L_E$ with $(L_E)_{\Qbar}\cong L$. The same
argument descends morphisms between lisse sheaves.

(2) Extension of scalars $F_X\otimes_{E_X}-$ is left adjoint to
restriction of scalars $\D(X_{\proet},F)\to\D(X_{\proet},E)$, so
$\Hom_{\D(X_{\proet},F)}(C_F,D_F[r])=\Hom_{\D(X_{\proet},E)}(C,D_F[r])$,
and similarly for $\Qbar$. Represent $D$ by a bounded complex $D^\bullet$.
The complexes $D^\bullet\otimes_{E_X}F_X$ form a filtered system of
complexes of $E_X$-modules, all vanishing below the lower bound of
$D^\bullet$, and their colimit is $D^\bullet\otimes_{E_X}\QbarX$, which
represents $D_{\Qbar}$. Lemma~\ref{b:lem:constructible-ext-continuity},
applied to $C\in\Dbc(X,E)$, gives
\[
 \varinjlim_F\Hom_{\D(X_{\proet},E)}(C,D_F[r])
 \xrightarrow{\ \sim\ }\Hom_{\D(X_{\proet},E)}(C,D_{\Qbar}[r]),
\]
which is the claim after the adjunction.

(3) First let $K$ be a constructible $\QbarX$-module, with filtration
$F_1\subset\cdots\subset F_m=K$ as in
Lemma~\ref{b:lem:coefficient-heart}(3). We descend $F_k$ by induction on
$k$. The subquotient $F_k/F_{k-1}=a_!L$ descends because $L$ does, by
(1), and $a_!$ commutes with extension of scalars, as functors
$\Mod(S_{\proet},E)\to\Mod(X_{\proet},\Qbar)$. Suppose $F_{k-1}$ and
$F_k/F_{k-1}$ descend to $E$. The extension
$0\to F_{k-1}\to F_k\to F_k/F_{k-1}\to0$ is a distinguished triangle,
classified by a morphism $F_k/F_{k-1}\to F_{k-1}[1]$. By (2) this morphism
descends to some finite $F\supset E$, and its cone over $F$ is a model
of $F_k$.

For a bounded complex $K$, induct in the same way along the truncation
triangles $\tau_{<q}K\to K\to\mathcal H^q(K)[-q]\to$. Once
$\tau_{<q}K$ and $\mathcal H^q(K)$ have models over a finite $E$, the
connecting morphism $\mathcal H^q(K)[-q-1]\to\tau_{<q}K$ descends after a
finite extension by (2), and the cone of the descended morphism is a
model of $K$. Each step enlarges $E$ finitely often, so after finitely
many steps $K$ is defined over one finite $E$.
\end{proof}

\subsection{Lattices}

Statements with coefficients $\mathcal O_E/\pi_E^n$ reach $E$-coefficients
only through integral lattices. A lisse $E_X$-module need not have a
lattice on all of $X$, but it has one on each stratum of a finite
stratification, and for d\'evissage arguments this suffices.
A global lattice can fail to exist when $X$ is not normal. Let $X$ be a
projective line with $0$ and $\infty$ glued to a node, and let $Y\to X$
be the infinite chain of projective lines covering it, an \'etale
$\mathbf Z$-torsor. Descending $E_Y$ along $Y\to X$, with the generator
of $\mathbf Z$ acting by multiplication by $\ell$, gives a lisse
$E_X$-module $L$ of rank one. Since $Y$ is connected, a lattice in
$L|_Y\cong E_Y$ is a constant lattice $\pi_E^m\mathcal O_E$, and no such
lattice is stable under multiplication by $\ell$; so $L$ has no lattice
on $X$. On the open stratum $X\smallsetminus\{\text{node}\}\cong\mathbf G_m$
the torsor is trivial and $L$ is constant.

\begin{lemma}[Lattices after a finite stratification]
\label{b:lem:stratified-lattice}
Let $L$ be a lisse $E_X$-module. There is a finite stratification of $X$
by locally closed subschemes $S$ such that each $L|_S$ contains a lisse
$\mathcal O_{E,S}$-submodule $L_0$ with $L_0[1/\ell]=L|_S$.
\end{lemma}

The idea is to find a lattice after an \'etale surjective base change and
then push it down by the trace. The first step uses that the lattices in
$E^r$ form a discrete set on which $\GL_r(E)$ acts continuously, so that
the sheaf of lattices in $L$ is an \'etale $X$-scheme. The second step
needs the base change to be finite \'etale, which holds on the strata of
a finite stratification.

\begin{proof}
The rank of $L$ is a locally constant function on $|X|$: on a
trivializing covering the subsets of constant rank are open and closed
and saturated, and they descend by fpqc descent of open subsets, so we
may assume that $L$ has constant rank $r$.

\emph{The sheaf of lattices is an \'etale $X$-scheme.} Consider the
sheaf on $X_{\proet}$ whose sections over $U$ are the lisse
$\mathcal O_{E,U}$-submodules $M\subset L|_U$ with $M[1/\ell]=L|_U$. We
show that it is represented by a separated surjective \'etale $X$-scheme.
On an affine pro-\'etale object $V$ where $L$ is free of rank $r$, this
sheaf is the constant sheaf on the set
\[
 \GL_r(E)/\GL_r(\mathcal O_E)
\]
of lattices in $E^r$. This set is discrete because $\GL_r(\mathcal O_E)$
is an open subgroup of $\GL_r(E)$. A change of basis on $V\times_XV$ is a
continuous function to $\GL_r(E)$, and since the stabilizer of a lattice
is open, its action on each fixed lattice is locally constant. Hence the
induced map from the disjoint union of copies of $V\times_XV$ indexed by
the lattices in $E^r$ to itself is a morphism of schemes: on each
component, the preimages of the target components are open and closed and
form a disjoint open cover. The cocycle of the transition matrices
therefore gives a descent datum of schemes for these disjoint unions of
copies of $V$.

For affine $V=\operatorname{Spec}B$, the disjoint union of copies of $V$
indexed by a set $I$ is quasi-affine: it is the union, inside
$\operatorname{Spec}\prod_IB$, of the principal open subsets defined by
the coordinate idempotents. Descent of quasi-affine schemes along fpqc
coverings is effective by \stacks{0APK}. Taking a finite affine
trivializing covering of each affine open of $X$ and gluing, we obtain a
scheme representing the sheaf of lattices. It is separated and \'etale
over $X$ by \stacks{02KU} and \stacks{02VN}, and it is surjective because
every free module has a lattice. The universal lattice exists on it by
descent of module sheaves.

\emph{Pushing down by the trace.} Choose finitely many quasi-compact
affine opens of this scheme whose images cover $X$, and let
$p:V\to X$ be their disjoint union. Then $p$ is quasi-compact, separated,
surjective and \'etale, and the universal lattice gives a lattice
$M_0\subset p^*L$. By \stacks{07RY} and \stacks{07RZ}, $X$ has a finite
stratification by locally closed subschemes over each of which $p$ is
finite \'etale and surjective. On such a stratum $S$ let
$\operatorname{Tr}:p_*p^*L\to L$, a morphism in $\Mod(S_{\proet},E)$, be
the trace of the finite \'etale map $p$ over $S$, and define $L_0\subset L|_S$ as the image of the composite
\[
 p_*M_0\longrightarrow p_*p^*L\xrightarrow{\operatorname{Tr}}L.
\]
To see that $L_0$ is a lisse lattice, split the finite \'etale cover
locally on $S$ into $d$ copies of $S$. Then $p_*M_0$ is a direct sum of
$d$ lattices $M_1,\dots,M_d$ in $L$, the trace is their sum, and
$L_0=M_1+\cdots+M_d$. Trivialize $L$ further pro-\'etale locally. Each
$M_k$ is then a locally constant function to the discrete set of lattices
in $E^r$, so after a common refinement all $M_k$ are constant lattices
$\Lambda_k\subset E^r$. Then $L_0$ is the constant sheaf attached to
$\Lambda_1+\cdots+\Lambda_d$, which is a finite free $\mathcal O_E$-lattice
spanning $E^r$. Hence $L_0$ is lisse and $L_0[1/\ell]=L|_S$.
\end{proof}

\subsection{Integral models}

For integral coefficients $\Dctf(X,\mathcal O_E)=\Dbc(X,\mathcal O_E)$:
the finite constant local models have finite projective dimension over
the discrete valuation ring. We retain the $\Dbc$ notation here to
compare the cohomological constructibility criterion with the classical
adic category. By Proposition~\ref{n:integral-criterion}, an object lies in
$\Dbc(X,\mathcal O_E)$ precisely when it is derived complete and its
reduction modulo $\pi_E$ is a bounded constructible complex pulled back
from $X_{\et}$.
The Stacks Project works with the category $\Dc(X,\mathcal O_E)$
of \stacks{09C1}, built from modules over the constant sheaf
$\mathcal O_E$ instead of $\mathcal O_{E,X}$. The next lemma identifies
the two, so that the results of the Stacks Project apply. Combined with
lattices, it also shows that the rational categories are generated by
scalar extensions of integral objects with classical torsion reductions.

\begin{lemma}[Integral models]
\label{b:lem:integral-models}
\leavevmode
\begin{enumerate}
\item Restriction of scalars along the map from the constant sheaf
  $\mathcal O_E$ to $\mathcal O_{E,X}$ is an equivalence between the
  derived complete objects of $\D(X_{\proet},\mathcal O_E)$ and the derived
  complete complexes of modules over the constant sheaf $\mathcal O_E$.
\item An object $K$ of $\D(X_{\proet},\mathcal O_E)$ lies in
  $\Dbc(X,\mathcal O_E)$ if and only if its restriction of scalars lies in
  $\Dc(X,\mathcal O_E)$.
\item $\Dbc(-,\mathcal O_E)$ is preserved by
  $f^*\colon\D(Y_{\proet},\mathcal O_E)\to\D(X_{\proet},\mathcal O_E)$
  for every morphism $f:X\to Y$, by
  $\otimes^{\mathbf L}_{\mathcal O_{E,X}}$ on $\D(X_{\proet},\mathcal O_E)$,
  and by
  $a_!\colon\D(S_{\proet},\mathcal O_E)\to\D(X_{\proet},\mathcal O_E)$ for
  every locally closed immersion $a:S\hookrightarrow X$.
\item Every object of $\Dbc(X,E)$ is obtained by finitely many cones and
  shifts from objects $(a_!L_0)[1/\ell]$, where $a:S\hookrightarrow X$ is
  a locally closed immersion and $L_0$ is a lisse finite free
  $\mathcal O_{E,S}$-module. Every object of $\Dbc(X,\Qbar)$ is obtained in
  the same way from the objects $\QbarX\otimes_{\mathcal O_{E,X}}a_!L_0$,
  for some finite $E/\Ql$. The reductions of $a_!L_0$ modulo $\pi_E^n$ are
  bounded constructible complexes on $X_{\proet}$ with coefficients
  in $\mathcal O_E/\pi_E^n$.
\end{enumerate}
\end{lemma}
\begin{proof}
(1) and (2). The map from the constant sheaf $\mathcal O_E$ to
$\mathcal O_{E,X}$ is flat and induces isomorphisms modulo every
$\pi_E^n$ (\stacks{099L}). Hence restriction of scalars is an
equivalence on derived complete objects, by \stacks{099I}; this is (1).
By \stacks{09C1}, an object of $\Dc(X,\mathcal O_E)$ is a derived complete
complex whose reduction modulo $\pi_E$ is constructible of finite Tor
dimension over $\mathcal O_E/\pi_E$. Since $\mathcal O_E/\pi_E$ is a field,
every bounded complex has finite Tor dimension over it, so this is the
description of $\Dbc(X,\mathcal O_E)$ in
Proposition~\ref{n:integral-criterion}; together with (1) this gives
(2).

(3) By \stacks{09C2} and \stacks{09C7}, every $K\in\Dbc(X,\mathcal O_E)$ is
bounded, and on each stratum of a finite stratification it is
pro-\'etale locally a bounded complex of finite free
$\mathcal O_{E,X}$-modules.

\emph{Inverse image.} By Lemma~\ref{b:lem:coefficient-pullback}, $f^*$
sends finite free $\mathcal O_{E,Y}$-modules to finite free
$\mathcal O_{E,X}$-modules, so the local description is preserved.

\emph{Tensor product.} Tensoring with a complex that is locally a bounded
complex of finite free modules preserves derived completeness, and its
reduction modulo $\pi_E$ is the tensor product of the reductions, which
is constructible. To globalize, use the finite open--closed
d\'evissage along the strata, and the identity
\[
 (a_!L)\otimes^{\mathbf L}M\simeq a_!(L\otimes^{\mathbf L}a^*M)
\]
in $\D(X_{\proet},\mathcal O_E)$, for a locally closed immersion $a$ and
$L\in\D(S_{\proet},\mathcal O_E)$, $M\in\D(X_{\proet},\mathcal O_E)$. To
prove it, compare restrictions:
on the stratum both sides restrict to $L\otimes^{\mathbf L}a^*M$, and on
the complementary strata both restrict to $0$, by the exactness of
open and closed restriction. These restrictions are jointly conservative
by \stacks{09AD}.

\emph{Extension by zero.} The functor $a_!=i_*u_!$ preserves derived
completeness: $u_!$ for a quasi-compact open immersion $u$ preserves
limits by \stacks{09BP}, and $i_*$ is exact and preserves limits by
\stacks{09AA} and \stacks{09BM}. Both functors commute with reduction
modulo $\pi_E$ and preserve constructibility of the reductions.

(4) For $\Qbar$, first descend to a finite $E$ by
Lemma~\ref{b:lem:coefficient-finite-descent}. By
Corollary~\ref{b:lem:coefficient-heart-b}, every object of $\Dbc(X,E)$ is
obtained by finitely many cones and shifts from sheaves $a_!L$ with $L$
lisse. By Lemma~\ref{b:lem:stratified-lattice}, after refining the
stratum $S$ by a finite stratification (which replaces $a_!L$ by finitely
many extensions of such sheaves), $L$ contains a lattice $L_0$ with
$L_0[1/\ell]=L$, so $a_!L=(a_!L_0)[1/\ell]$. The reduction of $L_0$
modulo $\pi_E^n$ is a locally constant sheaf of finitely presented
$\mathcal O_E/\pi_E^n$-modules, so it comes from the \'etale site by
\stacks{099Y}. Hence $L_0\in\Dbc(S,\mathcal O_E)$, so
$a_!L_0\in\Dbc(X,\mathcal O_E)$ by (3), and its reductions are classical.
\end{proof}

\subsection{Proper direct image}

The common construction of $Rf_!$ in
\S\ref{b:sec:brown-construction} starts from proper direct image on
$\Dctf$, so we need preservation of this category, proper base change
and the projection formula for every admissible coefficient ring.
The finite case is classical, with the Tor bound checked explicitly in
Proposition~\ref{ctf:finite-proper}. The integral case reduces modulo
$\pi_E$ to the same torsion theorems, and the rational cases follow
from Lemma~\ref{b:lem:integral-models}(4). For $Y=\operatorname{Spec}k$ with
$k$ algebraically closed, the proposition says that the cohomology groups
$H^q(X,\Qbar)$ of a proper $X$ are finite-dimensional and vanish outside
a bounded range of degrees.

\begin{proposition}[Proper direct image of constructible complexes]
\label{b:prop:proper-coefficients}
Let $p:X\to Y$ be proper and let $\Lambda$ be admissible. Then
$Rp_*\colon\D(X_{\proet},\Lambda)\to\D(Y_{\proet},\Lambda)$ maps
$\Dctf(X,\Lambda)$ into $\Dctf(Y,\Lambda)$. On these categories the
canonical proper base change and projection formula maps are
isomorphisms: for a cartesian square with $g:Y'\to Y$, $g':X'\to X$ and
$p':X'\to Y'$, the base change map $g^*Rp_*\to Rp'_*g'^*$ is an
isomorphism of functors $\Dctf(X,\Lambda)\to\Dctf(Y',\Lambda)$, and the
projection formula map
$Rp_*K\otimes^{\mathbf L}M\to Rp_*(K\otimes^{\mathbf L}p^*M)$ is an
isomorphism in $\D(Y_{\proet},\Lambda)$ for $K\in\Dctf(X,\Lambda)$ and
$M\in\Dctf(Y,\Lambda)$.
\end{proposition}
For $\Lambda=\mathcal O_E,E,\Qbar$ the category $\Dctf$ is $\Dbc$, and
the proof works with $\Dbc$.
\begin{proof}
\emph{Finite coefficients.} For finite $A$, proper constructible
finiteness and a uniform cohomological bound give bounded
constructibility; the constant-module projection formula gives a
finite Tor bound. The details are in Proposition~\ref{ctf:finite-proper}.
Proper base change is Theorem~\ref{pbc:torsion-theorem} on the \'etale
sites, transported to an isomorphism in $\D(Y'_{\proet},A)$ by its finite
instance Corollary~\ref{pbc:proet-finite}; the projection formula is the
classical torsion formula (\stacks{0F0G}) in $\D(Y_{\et},A)$, transported
to $\D(Y_{\proet},A)$ under the same comparison.

\emph{Integral coefficients.} Let $K\in\Dbc(X,\mathcal O_E)$. The complex
$Rp_*K$ is derived complete by \stacks{099L}. Reduction modulo $\pi_E$
commutes with $Rp_*$:
\begin{equation}
 (Rp_*K)\otimes^{\mathbf L}\mathcal O_{E,Y}/\pi_E
       \simeq Rp_*(K\otimes^{\mathbf L}\mathcal O_{E,X}/\pi_E)
       \quad\text{in }\D(Y_{\proet},\mathcal O_E).
\label{b:eq:proper-reduction}
\end{equation}
Indeed, $K\otimes^{\mathbf L}\mathcal O_{E,X}/\pi_E$ is the cone of
multiplication by $\pi_E$ on $K$, and the exact functor $Rp_*$ commutes
with this cone. The right side of \eqref{b:eq:proper-reduction} is the
direct image of a complex with torsion coefficients, which is pulled back
from $X_{\et}$ (\stacks{09B1}); its direct image in
$\D(Y_{\proet},\mathcal O_E/\pi_E)$ is the pullback to $Y_{\proet}$ of its
direct image in $\D(Y_{\et},\mathcal O_E/\pi_E)$ (\stacks{09A6}). The \'etale direct
image is constructible by the finiteness theorem for proper direct image
with torsion coefficients (\stacks{0GL0}), and bounded because $p$ has
finite cohomological dimension (\stacks{0F10}). By
Lemma~\ref{b:lem:integral-models}(2), $Rp_*K\in\Dbc(Y,\mathcal O_E)$.

For constructible integral inputs, the canonical projection formula map
in $\D(Y_{\proet},\mathcal O_E)$ and the base change map in
$\D(Y'_{\proet},\mathcal O_E)$, built from the adjunctions and the usual maps between
tensor products and inverse images, have derived complete source and
target, by the preceding paragraph and
Lemma~\ref{b:lem:integral-models}(3). Their reductions modulo $\pi_E$ are
the torsion projection formula and proper base change maps. The first is
an isomorphism by the projection formula for torsion coefficients
(\stacks{0F0G}). The second is an isomorphism by
Theorem~\ref{pbc:torsion-theorem} and
Corollary~\ref{pbc:proet-finite}: the finite reductions are bounded
below classical complexes, so their canonical exchange maps are
isomorphisms in $\D(Y'_{\proet},\mathcal O_E/\pi_E)$. (Proposition~\ref{pbc:complete} gives a completed
integral version; here the inverse images of constructible complexes
are already complete.) Hence the cone of each integral map is derived complete
and has zero reduction modulo $\pi_E$. Such a complex is zero (derived
Nakayama): using the filtration of $\mathcal O_E/\pi_E^n$ with
subquotients $\mathcal O_E/\pi_E$, its reductions modulo every $\pi_E^n$
vanish, and a derived complete complex is the derived inverse limit of
these reductions. Hence the integral maps are isomorphisms. (Proper base
change also follows directly from \stacks{09C9}, since the inverse image
of a constructible integral complex is already derived complete.)

\emph{Rational coefficients.} For $K\in\Dbc(X,\mathcal O_E)$, the complex
$K[1/\ell]$ is the colimit of $K\xrightarrow{\ell}K\xrightarrow{\ell}\cdots$,
a system with a common lower bound, so
Lemma~\ref{b:lem:bounded-continuity} gives
\[
 Rp_*(K[1/\ell])\simeq(Rp_*K)[1/\ell]\quad\text{in }\D(Y_{\proet},E).
\]
In the same way, $Rp_*$ commutes with extension of scalars to a finite
$F/E$, as functors $\D(X_{\proet},E)\to\D(Y_{\proet},F)$ (which is a finite direct sum, by
Lemma~\ref{b:lem:coefficient-pullback}(2)) and, by
Lemma~\ref{b:lem:bounded-continuity} again, with the filtered colimit
$\QbarX\otimes_{E_X}-=\varinjlim_F F_X\otimes_{E_X}-$. Thus $Rp_*$
commutes with inverting $\ell$ and with extension of scalars to $\Qbar$.
By the integral case, $Rp_*$ therefore maps the generators of
Lemma~\ref{b:lem:integral-models}(4) into $\Dbc(Y,E)$ and $\Dbc(Y,\Qbar)$,
and since $\Dbc(Y,E)$ and $\Dbc(Y,\Qbar)$ are triangulated
(Corollary~\ref{b:lem:coefficient-heart-b}), $Rp_*$ preserves
$\Dbc(-,E)$ and $\Dbc(-,\Qbar)$. For the projection formula and proper
base change, rationalize the integral isomorphisms and extend to all
inputs by d\'evissage in each variable along the same generators. The
maps involved are the canonical ones, so the resulting isomorphisms
satisfy the usual compatibilities with composition.
\end{proof}
\paragraph*{What is now available.} For every admissible $\Lambda$,
proper direct image preserves $\Dctf$ and satisfies proper base change
and the projection formula there. For finite rings this rests directly
on the torsion theorems; for the other coefficients it goes through
integral models with classical finite-level reductions and descent from
finite extensions. For $E$, $\Qbar$ and integral lattices, in addition,
Hom out of a constructible complex commutes with bounded-below filtered
colimits.
This is used in the compactness proof for $E$ and $\Qbar$; compactness
for finite rings is proved from bounded flat constructible models.

\section{Construction of the exceptional inverse image}
\label{sec:formal-half}
Brown representability will give
$f^!\colon\Loc(\Dctf(Y,\Lambda))\to\Loc(\Dctf(X,\Lambda))$, as a right adjoint
of $Rf_!\colon\Loc(\Dctf(X,\Lambda))\to\Loc(\Dctf(Y,\Lambda))$, provided
$\Loc(\Dctf(X,\Lambda))$ is compactly generated by $\Dctf(X,\Lambda)$
and $Rf_!=Rp_*\circ j_!$, for a compactification $f=p\circ j$, preserves
direct sums and $\Dctf(-,\Lambda)$.
Compact generation of the ambient category $D(X_{\proet},\Lambda)$ would
give neither of the last two properties.

These estimates depend on the coefficients. For $\Lambda=E,\Qbar$ we
have $\Dctf(-,\Lambda)=\Dbc$, so d\'evissage along cohomology sheaves
is available; this case occupies the subsections below, through
Lemma~\ref{b:lem:proper-unbounded-window-b}. For finite rings the
estimates are proved in \S\ref{ctf:finite-foundations} from flat
constructible models, and for $\mathcal O_E$ they are
Proposition~\ref{i:compactness} and Corollary~\ref{i:compactness-b}.
The construction of $Rf_!$ and the comparison of compactifications that
follow are formal and hold for every admissible $\Lambda$; with Brown
representability they give $f^!$ on $\Loc(\Dctf(-,\Lambda))$
(Theorem~\ref{ctf:unified-brown}) and parts (1)--(3) of
Theorem~\ref{n:existence}. That $f^!$ preserves $\Dctf$ is a geometric
statement, proved in \S\ref{sec:restriction-chapter}.

\subsection{Compactness of constructible complexes}
\label{b:sec:compactness}

In this subsection $\Lambda=E$ or $\Qbar$. For these coefficients
$\Dctf=\Dbc$, and bounded constructible complexes are the proposed
compact generators. The equality $\Dctf=\Dbc$ is what allows us to use
the constructible heart and truncations below; for a finite ring it can
fail. Recall that $\Cons(X,\Lambda)$ is an abelian category for these
$\Lambda$ (Lemma~\ref{b:lem:coefficient-heart}). Unadorned
$\Hom$ and $\Ext$ are taken in $\D(X_{\proet},\Lambda)$, also when both
arguments lie in a smaller category; a subscript $\Lambda$ or
$\mathcal O_E$ records the coefficient ring when it needs emphasis.

For a constructible sheaf $C$ and a family $(K_\lambda)$ of objects of
$\Loc(\Dbc(X,\Lambda))$ we must show that
$\bigoplus_\lambda\Hom(C,K_\lambda)\to\Hom(C,\bigoplus_\lambda K_\lambda)$
is bijective; bounded constructible complexes then follow by d\'evissage.
Two facts do the work: $\Ext^q(C,-)$ commutes with filtered colimits
(Lemma~\ref{b:lem:constructible-ext-continuity}), and $\Ext^q(C,N)=0$ for
$q>d_C$, with $d_C$ independent of the constructible sheaf $N$
(Lemma~\ref{b:lem:uniform-constructible-ext}). The second controls
unbounded targets, whose summands $K_\lambda$ need not have a common
bound. The first fact was proved in \S\ref{b:sec:ext-continuity}; the
second is proved next.

\subsubsection{A uniform bound for Ext between constructible sheaves}

The bound $d_C$ on $\Ext^q(C,-)$ is a sum of the cohomological
dimensions $\operatorname{cd}_\ell$ of Definition~\ref{n:def-cd}.



The proof of the uniform bound computes $\RHom(a_!L_0,K)$, for a derived
complete target $K\in\D(X_{\proet},\mathcal O_E)$, as a derived inverse limit of its reductions modulo
$\pi_E^n$. For this, the functors in \eqref{b:eq:stratum-hom} must
commute with derived inverse limits.

\begin{lemma}\label{b:lem:formal-closed-support-b}
In the situation of Lemma~\ref{b:lem:formal-closed-support}, the functors
$R\underline\Gamma_Z\colon\D(X_{\proet},\Lambda)\to\D(Z_{\proet},\Lambda)$ and
$u^*R\underline\Gamma_Z\colon\D(X_{\proet},\Lambda)\to\D(S_{\proet},\Lambda)$
are left $t$-exact and commute with derived inverse limits of sequences.
So does tensoring with a finite locally free sheaf, as a functor
$\D(S_{\proet},\Lambda)\to\D(S_{\proet},\Lambda)$.
\end{lemma}
\begin{proof}
$R\underline\Gamma_Z$ is left $t$-exact because it is right adjoint to the
exact functor $i_*$. A derived inverse limit of a sequence is the fiber of
$1-\mathrm{shift}$ on a product, so it suffices that the functors preserve
products and fibers. This holds for $j^*$ and $u^*$ (open restriction is
exact and preserves limits), for $i^*$ (\stacks{09AA}), for $Rj_*$ (a right
adjoint), hence for $R\underline\Gamma_Z=i^*\fib(\mathrm{id}\to Rj_*j^*)$.
Tensoring with a finite locally free sheaf $L$ is exact and is right
adjoint to tensoring with $L^\vee$, so it preserves products. The lemma
is used below with $\Lambda=\mathcal O_E$.
\end{proof}

\begin{lemma}[Uniform Ext bound]
\label{b:lem:uniform-constructible-ext}
For every $C\in\Cons(X,\Lambda)$ there is an integer $d_C\geq0$ such that
\[
 \Ext^q_{\Lambda}(C,N)=0
 \qquad\text{for all }q>d_C\text{ and all }N\in\Cons(X,\Lambda).
\]
\end{lemma}
For $C=a_!L$, with $a$ factored as $S\hookrightarrow Z\hookrightarrow X$
and $j$ the open complement of $Z$, the proof gives
$d_C=\operatorname{cd}_\ell(j)+\operatorname{cd}_\ell(S)+2$. One of the two extra
degrees comes from the cone defining $R\underline\Gamma_Z$, the other from
a derived inverse limit over the reductions modulo $\pi_E^n$.
\begin{proof}
By the long exact Ext sequences, a bound that holds for every subquotient
of a finite filtration of the source holds for the source (take the
maximum), and a bound that holds for every subquotient of a finite
filtration of the target holds for the target, with the same value. By
Lemma~\ref{b:lem:coefficient-heart}(3) we may therefore assume
$C=a_!L$, with $a=i\circ u:S\hookrightarrow Z\hookrightarrow X$ as in
Lemma~\ref{b:lem:formal-closed-support} and $L$ lisse, and it suffices
to prove
\[
 \Ext^q(a_!L,b_!M)=0\qquad(q>\operatorname{cd}_\ell(j)+\operatorname{cd}_\ell(S)+2)
\]
for all locally closed $b:T\hookrightarrow X$ and lisse $M$ on $T$, where
$j:X\setminus Z\hookrightarrow X$ is the open complement of $Z$. Both
cohomological dimensions are finite by \stacks{0F10} and
Lemma~\ref{b:lem:base-field-cd}, and the bound depends only on $a$.

\emph{Lattices.} Take first $\Lambda=E$. After refining the
stratifications (Lemma~\ref{b:lem:stratified-lattice}) and enlarging $E$,
$L$ and $M$ contain lisse $\mathcal O_E$-lattices $L_0$ and $M_0$.
Enlarging $E$ is harmless: if $E'/E$ is finite, then $E'_X$ is finite free
over $E_X$, so $\Ext^q_{E'}$ of the extended sheaves is a finite direct sum
of copies of $\Ext^q_E$, and a bound over $E'$ gives the same bound over
$E$. The functor $b_!$ is the
composite of an extension by zero, which is exact and preserves limits
(\stacks{09BP}), and a closed direct image, which is exact
(\stacks{09BM}) and preserves limits as a right adjoint. Hence $b_!M_0$
is a $\pi_E$-torsion free, derived complete sheaf, and its reductions are
\[
 b_!M_0\otimes^{\mathbf L}_{\mathcal O_{E,X}}\mathcal O_{E,X}/\pi_E^n
 =b_!(M_0/\pi_E^n).
\]
These are constructible torsion sheaves pulled back from $X_{\et}$:
lisse torsion sheaves descend (\stacks{099Y}), and extension by zero and
closed direct image commute with this pullback (\stacks{09AF},
\stacks{09A6}).

\emph{Torsion targets.} By \stacks{09A6}, $Rj_*$ in
$\D(X_{\proet},\mathcal O_E/\pi_E^n)$ of a sheaf pulled back from the
\'etale site is the pullback of its direct image in
$\D(X_{\et},\mathcal O_E/\pi_E^n)$, so
$Rj_*j^*b_!(M_0/\pi_E^n)$ has cohomology in degrees
$[0,\operatorname{cd}_\ell(j)]$. The fiber formula of
Lemma~\ref{b:lem:formal-closed-support}(1) then gives
\[
 R\underline\Gamma_Z\,b_!(M_0/\pi_E^n)\in\D^{[0,\operatorname{cd}_\ell(j)+1]}(Z_{\proet},\mathcal O_E/\pi_E^n),
\]
with cohomology sheaves that are torsion sheaves pulled back from the
\'etale site, though not necessarily constructible. Tensoring with the
lisse $L_0^\vee/\pi_E^n$ and using the hypercohomology spectral sequence
on $S$ gives
\begin{equation}\label{b:eq:torsion-stratum-bound}
 H^q R\Gamma\bigl(S,(L_0^\vee/\pi_E^n)\otimes u^*R\underline\Gamma_Z\,b_!(M_0/\pi_E^n)\bigr)=0
 \quad(q>\operatorname{cd}_\ell(j)+\operatorname{cd}_\ell(S)+1),
\end{equation}
for every $n$, $b$ and $M_0$.

\emph{Integral targets.} Write $F=L_0^\vee\otimes u^*R\underline\Gamma_Z\,b_!M_0$.
By Lemma~\ref{b:lem:formal-closed-support-b} the functors forming $F$
commute with derived inverse limits of sequences; they also commute with
reduction modulo $\pi_E^n$, which is the cone of multiplication by
$\pi_E^n$. Since $b_!M_0$ is derived complete, so is $F$, and the
reductions of $F$ are the complexes in \eqref{b:eq:torsion-stratum-bound}.
Derived global sections preserve products and cones, hence derived
inverse limits. Using \eqref{b:eq:stratum-hom} with $\mathcal O_E$
coefficients, that is, for $\RHom$ in $\D(X_{\proet},\mathcal O_E)$,
\[
 \RHom_X(a_!L_0,b_!M_0)
 =R\Gamma(S,F)
 =R\varprojlim_n R\Gamma\bigl(S,F\otimes^{\mathbf L}\mathcal O_E/\pi_E^n\bigr).
\]
A derived inverse limit of a sequence raises an upper bound on cohomology
by at most one: it is the fiber of $1-\mathrm{shift}$ on the product, and
products of abelian groups are exact (this is the Milnor sequence,
\stacks{08TB}). Therefore
\begin{equation}\label{b:eq:integral-stratum-bound}
 \Ext^q_{\mathcal O_E}(a_!L_0,b_!M_0)=0
 \quad(q>\operatorname{cd}_\ell(j)+\operatorname{cd}_\ell(S)+2).
\end{equation}

\emph{Inverting $\pi_E$.} We have $a_!L=a_!L_0[1/\pi_E]$ and
$b_!M=(b_!M_0)[1/\pi_E]=\varinjlim(b_!M_0\xrightarrow{\pi_E}b_!M_0\xrightarrow{\pi_E}\cdots)$.
The adjunction between extension and restriction of scalars gives
\[
 \RHom_E\bigl(a_!L,b_!M\bigr)
 =\RHom_{\mathcal O_E}\bigl(a_!L_0,(b_!M_0)[1/\pi_E]\bigr),
\]
and by Lemma~\ref{b:lem:constructible-ext-continuity}(2) the cohomology of
the right side is the colimit of the groups
$\Ext^q_{\mathcal O_E}(a_!L_0,b_!M_0)$. So
\eqref{b:eq:integral-stratum-bound} gives the bound for $E$.

\emph{The case $\Lambda=\Qbar$.} By
Lemma~\ref{b:lem:coefficient-finite-descent}, $a_!L$ and $b_!M$ are
extensions of scalars of sheaves defined over a common finite $E$. Since
$\QbarX=\varinjlim_{E'}E'_X$ is flat over $E_X$, extension and restriction
of scalars and Lemma~\ref{b:lem:constructible-ext-continuity} identify
$\Ext^q_{\Qbar}(a_!L,b_!M)$ with the colimit of the corresponding groups
over finite extensions $E'$ of $E$, which vanish in the same range.
\end{proof}

\subsubsection{Ind-constructible sheaves}

The targets in the compactness statement are direct sums of objects of
$\Loc(\Dbc(X,\Lambda))$, whose cohomology sheaves are no longer
constructible. They are, however, filtered colimits of constructible
sheaves, and continuity and the uniform Ext bound extend to such
colimits.

\begin{proposition}[Ind-constructible sheaves]\label{b:prop:ind-heart}
Let $\IndCons(X,\Lambda)\subset\Mod(X_{\proet},\Lambda)$ be the full
subcategory of filtered colimits, computed in $\Mod(X_{\proet},\Lambda)$,
of constructible sheaves.
\begin{enumerate}
\item The functor $\operatorname{Ind}(\Cons(X,\Lambda))\to\Mod(X_{\proet},\Lambda)$
  sending a formal filtered colimit to the actual colimit is fully
  faithful, with essential image $\IndCons(X,\Lambda)$.
\item $\IndCons(X,\Lambda)$ is closed in $\Mod(X_{\proet},\Lambda)$ under
  filtered colimits, direct sums, kernels, cokernels and extensions. In
  particular it is a Grothendieck abelian category and its inclusion is
  exact.
\end{enumerate}
\end{proposition}
\begin{proof}
The key point is that constructible sheaves are finitely
presented in $\Mod(X_{\proet},\Lambda)$: by
Lemma~\ref{b:lem:constructible-ext-continuity} in degree zero,
$\Hom(C,-)$ commutes with filtered colimits for $C$ constructible. The
category $\Cons(X,\Lambda)$ is essentially small and abelian, and closed
under extensions (Lemma~\ref{b:lem:coefficient-heart}(1),(2)); fix a small
skeleton.

\emph{Full faithfulness.} For filtered diagrams $(C_\alpha)$ and
$(D_\beta)$ of constructible sheaves, finite presentation gives
\[
 \Hom_{\Lambda}(\varinjlim_\alpha C_\alpha,\varinjlim_\beta D_\beta)
 =\varprojlim_\alpha\varinjlim_\beta\Hom_{\Cons(X,\Lambda)}(C_\alpha,D_\beta),
\]
which is the definition of morphisms in the ind-category.

\emph{Canonical presentations.} For $B\in\IndCons(X,\Lambda)$ let $J(B)$ be
the category of constructible sheaves $C$ with a map $C\to B$. It is
filtered: two objects map to their direct sum, and two parallel arrows
are equalized by a cokernel in $\Cons(X,\Lambda)$. If
$B=\varinjlim_\alpha C_\alpha$, finite presentation says that every
$C\to B$ factors through some $C_\alpha$, and two factorizations agree at
a later stage; so the given diagram is cofinal in $J(B)$ and
$B=\varinjlim_{J(B)}C$.

\emph{Filtered colimits and direct sums.} Let $B=\varinjlim_\lambda B_\lambda$
be a filtered colimit with $B_\lambda\in\IndCons(X,\Lambda)$. The category
of pairs $(\lambda,C\to B_\lambda)$, $C$ constructible, is filtered (move
finitely many indices to a common upper bound, then use direct sums and
cokernels in $J(B_\lambda)$), and the colimit of the $C$ over it is
$\varinjlim_\lambda\varinjlim_{J(B_\lambda)}C=B$. A direct sum is the
filtered colimit of its finite subsums, so it lies in $\IndCons(X,\Lambda)$
too.

\emph{Kernels and cokernels.} Let $f:B\to B'$ be a map in
$\IndCons(X,\Lambda)$. We write $f$ as a filtered colimit of maps of
constructible sheaves. Let $P_f$ be the category of commutative squares
\[
 \begin{matrix} C&\xrightarrow{v}&D\\
 \downarrow&&\downarrow\\ B&\xrightarrow{f}&B',\end{matrix}
 \qquad C,D\in\Cons(X,\Lambda).
\]
It is filtered, because finite colimits of such squares are formed by
finite colimits of the arrows $v$ in the abelian category $\Cons(X,\Lambda)$.
The functor $P_f\to J(B)$, $(v)\mapsto(C\to B)$, is cofinal: given
$C\to B$, the composite $C\to B'$ factors through some $D\in J(B')$ by
finite presentation; two such choices are dominated by their direct sum;
and a compatibility along a map $C_0\to C$ that holds over $B'$ holds at
a finite stage, after passing to a cokernel. (Equivalently, two choices
over $C_0$ are dominated by the square with source $C_1\amalg_{C_0}C_2$
and target the quotient of $D_1\oplus D_2$ by the two images of $C_0$;
parallel choices are dominated by their cokernel.) The same operations,
starting from the square $0\to D$, show that $P_f\to J(B')$,
$(v)\mapsto(D\to B')$, is cofinal. Hence $f=\varinjlim_{P_f}v$, and
since filtered colimits are exact,
\[
 \ker f=\varinjlim_{P_f}\ker v,\qquad
 \coker f=\varinjlim_{P_f}\coker v,
\]
both of which lie in $\IndCons(X,\Lambda)$.

\emph{Extensions.} Let $0\to N\to M\to B\to0$ be exact in
$\Mod(X_{\proet},\Lambda)$, with $N=\varinjlim_\beta D_\beta$ and
$B=\varinjlim_\alpha C_\alpha$ ind-constructible. Pull the extension back
along $C_\alpha\to B$. By Lemma~\ref{b:lem:constructible-ext-continuity},
\[
 \Ext^1_{\Lambda}(C_\alpha,N)
 =\varinjlim_\beta\Ext^1_{\Lambda}(C_\alpha,D_\beta),
\]
so the pulled-back extension is pushed out from an extension $M_{\alpha\beta}$
of $C_\alpha$ by some $D_\beta$. Each $M_{\alpha\beta}$, and each of its
pushouts along $D_\beta\to D_{\beta'}$, is constructible because
$\Cons(X,\Lambda)$ is closed under extensions; their colimit over
$\beta'$ is the pullback of $M$ to $C_\alpha$, which is therefore
ind-constructible. Since filtered colimits are exact, $M$ is the colimit
over $\alpha$ of these pullbacks, so $M\in\IndCons(X,\Lambda)$ by the
closure under filtered colimits.

\emph{Grothendieck.} By the above, colimits in $\IndCons(X,\Lambda)$
exist and are computed in $\Mod(X_{\proet},\Lambda)$, so filtered colimits
are exact; and the direct sum of the objects of the skeleton of
$\Cons(X,\Lambda)$ is a generator. This is the definition
(\stacks{079A}).
\end{proof}

\begin{corollary}\label{b:prop:ind-heart-b}
Let $C\in\Cons(X,\Lambda)$. Then $\Ext^q_{\Lambda}(C,B)=0$ for all $q>d_C$
and all $B\in\IndCons(X,\Lambda)$, and each functor $\Ext^q_\Lambda(C,-)$
commutes with direct sums of objects of $\IndCons(X,\Lambda)$.
\end{corollary}
\begin{proof}
Write $B$ as a filtered colimit of constructible sheaves. By
Lemma~\ref{b:lem:constructible-ext-continuity}, $\Ext^q(C,B)$ is the
colimit of the corresponding groups for the constructible terms, which
vanish for $q>d_C$ by Lemma~\ref{b:lem:uniform-constructible-ext}. A
direct sum is the filtered colimit of its finite subsums, and $\Ext^q(C,-)$
commutes with finite direct sums; the same lemma gives the second claim.
\end{proof}

Objects of $\Loc(\Dbc(X,\Lambda))$ need not be bounded, and their
truncations need not lie in $\Loc(\Dbc(X,\Lambda))$. Truncation arguments,
here and in \S\ref{b:sec:brown-construction}, therefore take place in the
following larger category.

\begin{corollary}\label{b:cor:ambient-ind-derived}
The full subcategory
\[
 \Dic(X,\Lambda)
 =\{K\in\D(X_{\proet},\Lambda):H^n(K)\in\IndCons(X,\Lambda)\text{ for every }n\}
\]
of $\D(X_{\proet},\Lambda)$ is triangulated, closed under direct sums, and
stable under the truncations $\tau_{\leq n}$, $\tau_{\geq n}$. In
particular $\Loc(\Dbc(X,\Lambda))\subset\Dic(X,\Lambda)$.
\end{corollary}
\begin{proof}
In the long exact cohomology sequence of a triangle, each cohomology
sheaf of the third term is an extension of a kernel by a cokernel of maps
between cohomology sheaves of the other two; so
Proposition~\ref{b:prop:ind-heart} gives closure under cones. Direct sums
of sheaves are exact, so $H^n$ commutes with direct sums. Truncation only
removes cohomology sheaves. Since $\Dbc(X,\Lambda)\subset\Dic(X,\Lambda)$,
the last claim follows from the definition of $\Loc(\Dbc(X,\Lambda))$.
\end{proof}

\subsubsection{Compactness}

Because of the uniform bound $d_C$, maps out of a constructible sheaf
$C$ into an unbounded complex see only finitely many of its cohomology
sheaves.

\begin{lemma}[Reduction to finitely many degrees]\label{b:lem:finite-window}
Let $C\in\Cons(X,\Lambda)$, $K\in\Dic(X,\Lambda)$ and $m\in\mathbf Z$. The
truncation maps induce a natural isomorphism
\begin{equation}\label{b:eq:finite-window}
 \Hom(C,K[m])
 \cong\Hom\bigl(C,(\tau_{\geq m-d_C}\tau_{\leq m}K)[m]\bigr).
\end{equation}
\end{lemma}
\begin{proof}
Replacing $K$ by $K[m]$ we may assume $m=0$. The truncation
$\tau_{>0}K$ lies in $\D^{\geq1}$, so $\Hom(C,\tau_{>0}K[s])=0$ for
$s=0,-1$, and the truncation triangle for $\tau_{\leq0}$ shows that
$\Hom(C,K)=\Hom(C,\tau_{\leq0}K)$. It remains to remove the cohomology
in degrees $\leq-d_C-1$ from $K'=\tau_{\leq0}K$. By the triangle
$\tau_{\leq-d_C-1}K'\to K'\to\tau_{\geq-d_C}K'\to$, it suffices to show
\begin{equation}\label{b:eq:lower-tail-vanishing}
 \Hom(C,L)=0\qquad
 \text{for every }L\in\Dic(X,\Lambda)\text{ with }H^q(L)=0\text{ for }q>-d_C-1,
\end{equation}
applied to $L=\tau_{\leq-d_C-1}K'$ and to $L=\tau_{\leq-d_C-1}K'[1]$.

If $L$ is bounded, \eqref{b:eq:lower-tail-vanishing} follows from the
finite filtration of $L$ by its cohomology sheaves: a layer $H^q(L)[-q]$
with $q\leq-d_C-1$ contributes $\Ext^{-q}(C,H^qL)$, which vanishes by
Corollary~\ref{b:prop:ind-heart-b} since $-q>d_C$. For $L$ unbounded
below we use that the derived category is left complete.

The pro-\'etale site has a basis of weakly contractible objects
(Theorem~\ref{pe:weakly-contractible-basis}), so
Proposition~\ref{pe:products} shows that products of $\Lambda_X$-modules are
exact, products in $\D(X_{\proet},\Lambda)$ are computed termwise, and
\[
 L=R\varprojlim_n\tau_{\geq-n}L
\]
for every $L$. Concretely, choose a complex $L^\bullet$ representing $L$;
the limit is the derived inverse limit of the tower of good truncations
$\tau_{\geq-n}L^\bullet$ with their canonical transition maps. Since $\Hom(C,-)$ takes
products to products, the triangle defining $R\varprojlim$ gives the
Milnor exact sequence
\[
 0\longrightarrow\varprojlim_n{}^1
       \Hom(C,\tau_{\geq-n}L[-1])
 \longrightarrow\Hom(C,L)
 \longrightarrow\varprojlim_n\Hom(C,\tau_{\geq-n}L)
 \longrightarrow0.
\]
Each $\tau_{\geq-n}L$ is bounded with cohomology in degrees $\leq-d_C-1$,
so the right-hand term vanishes by the bounded case. For the left-hand
term, the triangle
\[
 H^{-n-1}(L)[n+1]\longrightarrow \tau_{\geq-n-1}L
 \longrightarrow \tau_{\geq-n}L\longrightarrow H^{-n-1}(L)[n+2]
\]
shows that the transition map
$\Hom(C,\tau_{\geq-n-1}L[-1])\to\Hom(C,\tau_{\geq-n}L[-1])$ sits between
$\Ext^{n}(C,H^{-n-1}L)$ and $\Ext^{n+1}(C,H^{-n-1}L)$. Both vanish for
$n>d_C$ by Corollary~\ref{b:prop:ind-heart-b}. So the tower is eventually
constant, and the $\varprojlim^1$ of an eventually constant tower is zero:
in the description of $\varprojlim^1$ as the cokernel of $1-\mathrm{shift}$
on the product, one solves the equation successively along the constant
tail and then along the finitely many earlier terms. This proves
\eqref{b:eq:lower-tail-vanishing}, and hence the lemma.
\end{proof}

\begin{theorem}[Compactness with rational coefficients]\label{b:thm:constructible-compactness}
For $\Lambda=E$ or $\Qbar$, every object of $\Dbc(X,\Lambda)$ is compact in $\Loc(\Dbc(X,\Lambda))$.
More precisely, for $C\in\Dbc(X,\Lambda)$ and every family $(K_\lambda)$
of objects of $\Dic(X,\Lambda)$, the natural map
\[
 \bigoplus_\lambda\Hom(C,K_\lambda)
 \longrightarrow\Hom\Bigl(C,\bigoplus_\lambda K_\lambda\Bigr)
\]
is an isomorphism.
\end{theorem}
\begin{proof}
The first assertion follows from the second, since
$\Loc(\Dbc(X,\Lambda))\subset\Dic(X,\Lambda)$
(Corollary~\ref{b:cor:ambient-ind-derived}) and direct sums in
$\Loc(\Dbc(X,\Lambda))$ are those of $\D(X_{\proet},\Lambda)$.

Let first $C\in\Cons(X,\Lambda)$, and suppose that all $K_\lambda$ have
cohomology in a fixed interval $[a,b]$. We prove the claim, for all
shifts of the $K_\lambda$ at once, by induction on $b-a$. If $a=b$, then
$K_\lambda=B_\lambda[-a]$ with $B_\lambda\in\IndCons(X,\Lambda)$, and the
claim is that $\Ext^q(C,-)$ commutes with direct sums of objects of
$\IndCons(X,\Lambda)$, which is Corollary~\ref{b:prop:ind-heart-b}. In
general, the triangles
$\tau_{\leq b-1}K_\lambda\to K_\lambda\to H^b(K_\lambda)[-b]\to$ are
compatible with direct sums, because $H^n$ commutes with direct sums. The
long exact Hom sequences, exactness of direct sums of abelian groups and
the five lemma give the induction step.

For arbitrary $K_\lambda$ apply Lemma~\ref{b:lem:finite-window} with
$m=0$. It replaces each $K_\lambda$, and $\bigoplus_\lambda K_\lambda$, by
its truncation to degrees $[-d_C,0]$; the interval is the same for the
whole family, and truncation commutes with direct sums. The bounded case
therefore applies.

Finally, for fixed $(K_\lambda)$ the objects $C$ for which the map is an
isomorphism are closed under shifts and cones (five lemma) and under
direct summands (\stacks{09QH}). A bounded constructible complex is a
finite extension of shifts of its constructible cohomology sheaves, so
it is among them.
\end{proof}

In particular we now have a compactly generated category, as Brown
representability requires.

\begin{corollary}\label{b:thm:constructible-compactness-b}
For $\Lambda=E$ or $\Qbar$, $\Loc(\Dbc(X,\Lambda))$ is compactly generated, by a set of objects of
$\Dbc(X,\Lambda)$.
\end{corollary}
\begin{proof}
Let $G$ be a small skeleton of $\Dbc(X,\Lambda)$; its objects are compact
by Theorem~\ref{b:thm:constructible-compactness}. It remains to see that
they detect zero objects. Suppose $\Hom(C[n],K)=0$ for all $C\in G$ and
all $n$. The objects $A$ with $\Hom(A[n],K)=0$ for all $n$ form a
localizing subcategory containing $G$, hence all of $\Loc(\Dbc(X,\Lambda))$.
Taking $A=K$ gives $\mathrm{id}_K=0$, so $K=0$.
\end{proof}

\subsection{The compactly supported direct image and Brown representability}
\label{b:sec:brown-construction}

Morphisms are separated of finite type. For a compactification
$f=p\circ j$, with $j$ an open immersion and $p$ proper, we define
$Rf_!=Rp_*\circ j_!$. By
Lemma~\ref{b:lem:coefficient-pullback}, for a morphism $g\colon X\to Y$
the inverse image $g^*\colon D(Y_{\proet},\Lambda)\to D(X_{\proet},\Lambda)$
is the exact inverse image of sheaves, so $g^*\dashv Rg_*$ is the usual
adjunction between the unbounded derived categories
$D(Y_{\proet},\Lambda)$ and $D(X_{\proet},\Lambda)$; the base change maps
below are natural transformations between functors on these categories,
built from it. Until the construction of $Rf_!$ we keep $\Lambda=E$ or
$\Qbar$, so that $\Loc(\Dctf(X,\Lambda))=\Loc(\Dbc(X,\Lambda))$. We write
$\Dic^{\leq b}(X,\Lambda)$ and $\Dic^{\geq a}(X,\Lambda)$ for the objects
of $\Dic(X,\Lambda)$ (Corollary~\ref{b:cor:ambient-ind-derived}) with
cohomology in degrees $\leq b$, respectively $\geq a$.

\subsubsection{Proper direct image of unbounded complexes}

Brown representability needs $Rf_!$ to preserve arbitrary direct sums.
For $j_!$ this is automatic, since $j_!$ is a left adjoint, so the
question is whether $Rp_*$ preserves direct sums for proper $p$. For
complexes with a common lower and upper bound this would follow from the
finiteness of $\operatorname{cd}_\ell(p)$, but the summands of a direct
sum in $\Loc(\Dbc(X,\Lambda))$ need not have any common bound. As for
compactness, the remedy is a uniform bound: $R^qp_*$ vanishes on
$\IndCons(X,\Lambda)$ for $q>2d+1$, where $d$ bounds the fibre dimensions
of $p$, so that $\mathcal H^q(Rp_*K)$ depends only on the cohomology of
$K$ in degrees $[q-2d-1,q]$.

\begin{lemma}[Cohomology commutes with filtered colimits]
\label{b:lem:proper-heart-continuity}
Let $Z$ be a quasi-compact and quasi-separated scheme and let $A$ be a sheaf
of rings on $Z_{\proet}$. For each integer $q$, the functor
$H^q(Z_{\proet},-)$ from sheaves of $A$-modules to abelian groups commutes with filtered colimits and direct
sums. If $p:X\to Y$ is quasi-compact and quasi-separated, then each
$R^qp_*$ on sheaves of modules on $X_{\proet}$ has the same two properties.
\end{lemma}

\begin{proof}
Let $W$ be an affine weakly contractible object. The functor
$\Gamma(W,-)$ is exact (Corollary~\ref{pe:wc-sections}) and preserves direct sums
(\stacks{0935}). An exact additive functor that preserves direct sums
preserves all colimits, because every colimit is the cokernel of a map
between two direct sums. So $\Gamma(W,-)$ also preserves filtered
colimits.

By Proposition~\ref{pe:hypercover} there is a hypercovering $Z_\bullet\to Z$ with every
$Z_n$ affine weakly contractible, and $R\Gamma(Z,M)$ is computed by the
cochain complex $M(Z_\bullet)$. Each
term of this complex commutes with filtered colimits and direct sums in
$M$, and so does taking cohomology of complexes of abelian groups. This
proves the first assertion.

For the second, let $V\to Y$ be affine weakly contractible and put
$X_V=X\times_YV$, a quasi-compact and quasi-separated scheme. Since
$\Gamma(V,-)$ is exact,
\[
 (R^qp_*M)(V)=H^q((X_V)_{\proet},M|_{X_V}).
\]
Restriction to $X_V$ preserves colimits, so the first assertion for
$Z=X_V$ shows that the comparison map for $R^qp_*$ is an isomorphism on
sections over $V$. Such $V$ form a basis of the site, so evaluation on them
detects isomorphisms of sheaves (Corollary~\ref{pe:sections}).
\end{proof}

\begin{lemma}[A uniform bound for proper direct image]
\label{b:lem:proper-heart-bound}
Let $p:X\to Y$ be proper, and let $d$ be an integer bounding the
dimensions of all geometric fibres of $p$ (for instance $d=\dim X$). For
$M\in\IndCons(X,\Lambda)$,
\[
 R^qp_*M=0\quad(q>2d+1),\qquad R^qp_*M\in\IndCons(Y,\Lambda)\quad(q\geq0),
\]
and $R^qp_*:\IndCons(X,\Lambda)\to\IndCons(Y,\Lambda)$ preserves arbitrary
direct sums.
\end{lemma}

The proof computes with an integral lattice. Its reductions modulo
$\pi_E^n$ satisfy the classical torsion bound $2d$, and the derived
inverse limit over $n$ costs one more degree, whence $2d+1$.

\begin{proof}
\emph{Reduction to a lattice.} The vanishing of $R^qp_*$ for $q>2d+1$ is
stable under extensions. By Lemma~\ref{b:lem:coefficient-heart}(3), every
$M\in\Cons(X,\Lambda)$ is a finite extension of sheaves $a_!L$, with
$a:S\hookrightarrow X$ a locally closed immersion and $L$ lisse. For
$\Lambda=\Qbar$, $L$ is the scalar extension $\QbarX\otimes_{E}L_E$ of a
lisse $E_S$-module $L_E$ for some finite $E/\Ql$
(Lemma~\ref{b:lem:coefficient-finite-descent}); for $\Lambda=E$ put
$L_E=L$. By Lemma~\ref{b:lem:stratified-lattice}, after refining $S$ by a
finite stratification (which replaces $a_!L$ by a finite extension of
sheaves of the same form), $L_E$ contains a lisse finite free
$\mathcal O_{E,S}$-lattice $L_0$ with $L_0[1/\ell]=L_E$. It is therefore
enough to prove the bound for $a_!L_0$, and then to pass to
$a_!L_E=(a_!L_0)[1/\ell]$ and, if $\Lambda=\Qbar$, to its scalar
extension $a_!L$.

The sheaf $a_!L_0$ is derived $\pi_E$-complete. Indeed $a$ is a closed
immersion composed with a quasi-compact open immersion; closed pushforward
is exact and preserves products (\stacks{09BM}), and extension by zero is
exact and preserves limits (\stacks{09BP}). So both preserve the derived
inverse limit that expresses the completeness of $L_0$.

\emph{The torsion bound.} Fix an affine weakly contractible $V\to Y$ and
let $p_V:X_V\to V$ be the base change of $p$; it is proper, with the same
bound $d$. Since $L_0$ is flat, $a_!L_0\otimes^{\mathbf L}_{\mathcal O_E}
\mathcal O_E/\pi_E^n$ is the sheaf $a_!(L_0/\pi_E^n)$, which is a classical
constructible \'etale sheaf of $\mathcal O_E/\pi_E^n$-modules pulled back
to the pro-\'etale site. The proper vanishing theorem \stacks{095U} says
that for a proper morphism whose fibres have dimension at most $d$ the
higher direct images $R^qp_*$ of torsion sheaves vanish for $q>2d$; it
applies over the non-Noetherian scheme $V$. Together with the comparison
\stacks{09A6} of \'etale and pro-\'etale direct images, it shows that
$R^qp_{V,*}$ of this sheaf vanishes for $q>2d$. As $V$ is weakly
contractible, $\Gamma(V,-)$ is exact on all pro-\'etale sheaves
(\stacks{0991}), so $H^q(X_V,-)=\Gamma(V,R^qp_{V,*}(-))$ and
\[
 H^q\bigl(X_V,a_!(L_0/\pi_E^n)\bigr)=0\quad(q>2d).
\]

\emph{The limit.} Restriction to $X_V$ is exact and commutes with
products, hence with derived inverse limits, so completeness gives
$(a_!L_0)|_{X_V}\simeq R\varprojlim_n a_!(L_0/\pi_E^n)|_{X_V}$.
Derived global sections also commute with derived inverse limits: a
product of K-injective complexes is K-injective and sections preserve
products. The Milnor exact sequence (Proposition~\ref{pe:milnor})
\[
 0\to{\textstyle\varprojlim^1_n}H^{q-1}\bigl(X_V,a_!(L_0/\pi_E^n)\bigr)
 \to H^q(X_V,a_!L_0)
 \to\varprojlim_nH^q\bigl(X_V,a_!(L_0/\pi_E^n)\bigr)\to0
\]
then gives $H^q(X_V,a_!L_0)=0$ for $q>2d+1$.

\emph{Back to $\Lambda$.} The sheaf $(a_!L_0)[1/\ell]$ is the filtered
colimit of $a_!L_0$ along multiplication by $\ell$, so
Lemma~\ref{b:lem:proper-heart-continuity} gives the same vanishing for
it. For $\Lambda=\Qbar$, scalar extension from $E$ to $\Qbar$ is the
filtered colimit of the extensions to finite $F\supset E$; as sheaves of
abelian groups each of these is a finite direct sum of copies of the
original sheaf, because $F$ is finite free over $E$. The continuity lemma
again gives the bound. Since this holds for every affine weakly
contractible $V$, the formula $(R^qp_*M)(V)=H^q(X_V,M|_{X_V})$ from the
proof of Lemma~\ref{b:lem:proper-heart-continuity} gives $R^qp_*(a_!L)=0$ for $q>2d+1$, and by the first
step $R^qp_*M=0$ for $q>2d+1$ and every $M\in\Cons(X,\Lambda)$.

\emph{Passage to $\IndCons(X,\Lambda)$.} By
Proposition~\ref{b:prop:proper-coefficients}, $Rp_*$ maps
$\Dbc(X,\Lambda)$ into $\Dbc(Y,\Lambda)$, so $R^qp_*M\in\Cons(Y,\Lambda)$
for constructible $M$. Every object of $\IndCons(X,\Lambda)$ is a filtered
colimit of constructible sheaves, and $R^qp_*$ commutes with filtered
colimits and direct sums by Lemma~\ref{b:lem:proper-heart-continuity}.
This gives all three assertions.
\end{proof}

\begin{lemma}[Proper direct image of unbounded complexes]
\label{b:lem:proper-unbounded-window}
Let $p$ and $d$ be as in Lemma~\ref{b:lem:proper-heart-bound}. Then
$Rp_*\colon D(X_{\proet},\Lambda)\to D(Y_{\proet},\Lambda)$ maps $\Dic(X,\Lambda)$ into $\Dic(Y,\Lambda)$, and
\[
 Rp_*\bigl(\Dic^{\leq b}(X,\Lambda)\bigr)
       \subset\Dic^{\leq b+2d+1}(Y,\Lambda),
 \qquad
 Rp_*\bigl(\Dic^{\geq a}(X,\Lambda)\bigr)
       \subset\Dic^{\geq a}(Y,\Lambda).
\]
Moreover, for $K\in\Dic(X,\Lambda)$ and every integer $q$, the canonical
truncation maps $K\leftarrow\tau_{\leq q}K\to\tau_{\geq q-2d-1}\tau_{\leq q}K$
induce an isomorphism in $\IndCons(Y,\Lambda)$
\begin{equation}
 \mathcal H^q(Rp_*K)\simeq
 \mathcal H^q\bigl(Rp_*\tau_{\geq q-2d-1}\tau_{\leq q}K\bigr).
 \label{b:eq:proper-window}
\end{equation}
\end{lemma}

\begin{proof}
\emph{Bounded below $K$.} The hypercohomology spectral sequence
\[
 E_2^{u,v}=R^up_*\mathcal H^v(K)\Longrightarrow\mathcal H^{u+v}(Rp_*K)
\]
converges. By Lemma~\ref{b:lem:proper-heart-bound} its only nonzero
columns are $0\leq u\leq2d+1$, and their terms lie in
$\IndCons(Y,\Lambda)$. So each $\mathcal H^n(Rp_*K)$ has a finite
filtration whose graded pieces are subquotients of objects of
$\IndCons(Y,\Lambda)$; since this category is closed under kernels,
cokernels and extensions (Proposition~\ref{b:prop:ind-heart}),
$\mathcal H^n(Rp_*K)\in\IndCons(Y,\Lambda)$. The two amplitude bounds can
be read off from the columns.

\emph{Arbitrary $K$.} By Proposition~\ref{pe:left-complete}, $K\simeq R\varprojlim_n
\tau_{\geq-n}K$. Being a right adjoint, $Rp_*$ preserves products and
hence derived inverse limits, so
$Rp_*K\simeq R\varprojlim_nRp_*\tau_{\geq-n}K$. The fibre of
$\tau_{\geq-n-1}K\to\tau_{\geq-n}K$ is $\mathcal H^{-n-1}(K)[n+1]$, whose
direct image has cohomology only in degrees $[-n-1,2d-n]$. Hence, for
fixed $q$, the towers $\mathcal H^q(Rp_*\tau_{\geq-n}K)$ and
$\mathcal H^{q-1}(Rp_*\tau_{\geq-n}K)$ have isomorphisms as transition
maps for $n$ large. An eventually constant tower has vanishing
$\varprojlim^1$, so the Milnor sequence of Proposition~\ref{pe:milnor} gives
\[
 \mathcal H^q(Rp_*K)\simeq
 \mathcal H^q(Rp_*\tau_{\geq-n}K)\quad\text{for all sufficiently large }n.
\]
The right side is covered by the bounded-below case. This proves
$\mathcal H^q(Rp_*K)\in\IndCons(Y,\Lambda)$, and if $K\in
\Dic^{\leq b}(X,\Lambda)$ it gives the bound $b+2d+1$. The lower bound
needs no argument for unbounded $K$, since an object of
$\Dic^{\geq a}$ is bounded below.

\emph{The isomorphism \eqref{b:eq:proper-window}.} Apply the bounds to
the truncation triangles
\[
 \tau_{\leq q}K\to K\to\tau_{\geq q+1}K,
 \qquad
 \tau_{\leq q-2d-2}K\to\tau_{\leq q}K\to\tau_{\geq q-2d-1}\tau_{\leq q}K.
\]
The direct image of $\tau_{\geq q+1}K$ lives in degrees $\geq q+1$, and
that of $\tau_{\leq q-2d-2}K$ in degrees $\leq q-1$. The long exact
cohomology sequences of the two triangles therefore give isomorphisms
$\mathcal H^q(Rp_*\tau_{\leq q}K)\simeq\mathcal H^q(Rp_*K)$ and
$\mathcal H^q(Rp_*\tau_{\leq q}K)\simeq
\mathcal H^q(Rp_*\tau_{\geq q-2d-1}\tau_{\leq q}K)$.
\end{proof}

\begin{lemma}[Proper direct image preserves direct sums]
\label{b:lem:proper-unbounded-window-b}
For $p$ proper, $Rp_*:\Dic(X,\Lambda)\to\Dic(Y,\Lambda)$ preserves
arbitrary direct sums.
\end{lemma}

\begin{proof}
First let $(M_\alpha)$ be a family with cohomology in a common finite
interval of degrees. Then the map
$\bigoplus_\alpha Rp_*M_\alpha\to Rp_*(\bigoplus_\alpha M_\alpha)$ is an
isomorphism in $D(Y_{\proet},\Lambda)$, by induction on the length of the interval. For length zero,
$M_\alpha$ is a shifted object of $\IndCons(X,\Lambda)$, and this is the
direct-sum assertion of Lemma~\ref{b:lem:proper-heart-bound} in each
degree. For the inductive step, sum the truncation triangles that split
off the top cohomology sheaf and apply the five lemma to the long exact
cohomology sequences.

Now let $(K_\alpha)$ be arbitrary and fix $q$. Direct sums of sheaves are
exact, so cohomology and truncations commute with direct sums. The
objects $\tau_{\geq q-2d-1}\tau_{\leq q}K_\alpha$ have cohomology in the
common interval $[q-2d-1,q]$. Applying \eqref{b:eq:proper-window} to each
$K_\alpha$ and to $\bigoplus_\alpha K_\alpha$, and the first paragraph to
the truncated family, shows that
$\bigoplus_\alpha Rp_*K_\alpha\to Rp_*(\bigoplus_\alpha K_\alpha)$ is an
isomorphism on $\mathcal H^q$. As $q$ is arbitrary, it is an isomorphism.
\end{proof}

\subsubsection{Construction of \texorpdfstring{$Rf_!$}{Rf!}}

From now on $\Lambda$ is any admissible coefficient ring. For finite
rings the role of Lemmas~\ref{b:lem:proper-heart-bound}--\ref{b:lem:proper-unbounded-window-b}
is played by Proposition~\ref{ctf:finite-proper}, proved in
\S\ref{ctf:finite-foundations} from flat constructible models, and for
$\mathcal O_E$ by Proposition~\ref{i:compactness} with
Corollary~\ref{i:compactness-b}.

Both factors of $Rp_*\circ j_!$ preserve the localizing categories. This
follows from the definition of $\Loc(\Dctf(X,\Lambda))$ in the following
form: if the objects of $\Loc(\Dctf(X,\Lambda))$ with some property form
a triangulated subcategory that is closed under direct sums and contains
$\Dctf(X,\Lambda)$, then every object of $\Loc(\Dctf(X,\Lambda))$ has the
property.

\begin{proposition}[Proper direct image and extension by zero]
\label{b:prop:proper-image-localizing}
For every admissible $\Lambda$, if $p:X\to Y$ is proper,
$Rp_*\colon D(X_{\proet},\Lambda)\to D(Y_{\proet},\Lambda)$ restricts to an exact functor
$\Loc(\Dctf(X,\Lambda))\to\Loc(\Dctf(Y,\Lambda))$ that preserves direct
sums. If $j:U\to X$ is an open immersion,
$j_!\colon D(U_{\proet},\Lambda)\to D(X_{\proet},\Lambda)$ restricts to an exact
functor $\Loc(\Dctf(U,\Lambda))\to\Loc(\Dctf(X,\Lambda))$ that preserves
direct sums and maps $\Dctf(U,\Lambda)$ into $\Dctf(X,\Lambda)$.
\end{proposition}

\begin{proof}
For $E$ and $\Qbar$, Lemma~\ref{b:lem:proper-unbounded-window-b}
shows that $Rp_*$ preserves direct sums on
$\Dic(X,\Lambda)\supset\Loc(\Dctf(X,\Lambda))$, and
Proposition~\ref{b:prop:proper-coefficients} gives preservation of
$\Dctf=\Dbc$. For finite $A$, the corresponding two facts are
Proposition~\ref{ctf:finite-proper}; for $\mathcal O_E$, they are
Proposition~\ref{i:compactness}, using $\Dctf=\Dbc$. In each case the
objects $K\in\Loc(\Dctf(X,\Lambda))$ with
$Rp_*K\in\Loc(\Dctf(Y,\Lambda))$ form a triangulated subcategory closed
under direct sums and containing $\Dctf(X,\Lambda)$, hence all of
$\Loc(\Dctf(X,\Lambda))$.

Extension by zero is exact, is left adjoint to $j^*$ and preserves
direct sums. It preserves constructibility by adding the closed
complement as a zero stratum. For finite $A$ it also preserves finite
Tor amplitude, as follows stalkwise from extension by zero; for
$\mathcal O_E$, $E$ and $\Qbar$ this follows from the local models of
constructible sheaves and Lemma~\ref{i:tensor}. The same localizing argument therefore gives
the stated restriction of $j_!$.
\end{proof}

Comparison maps between coproduct-preserving exact functors on
$\Loc(\Dctf(-,\Lambda))$ can be checked on $\Dctf(-,\Lambda)$.
This applies to proper base change and to changes of compactification.

\begin{lemma}[Checking isomorphisms on constructible complexes]
\label{b:lem:localizing-comparison}
Let $\Lambda$ be admissible, let $F_1,F_2:\Loc(\Dctf(X,\Lambda))\to\Loc(\Dctf(Y,\Lambda))$ be exact functors
preserving direct sums, and let $\eta:F_1\to F_2$ be a natural
transformation of exact functors. If $\eta_K$ is an isomorphism for every
$K\in\Dctf(X,\Lambda)$, then it is an isomorphism for every
$K\in\Loc(\Dctf(X,\Lambda))$.
\end{lemma}

\begin{proof}
The objects $K$ for which $\eta_K$ is an isomorphism are closed under
shifts and cones by the five lemma, and under direct sums because $F_1$
and $F_2$ preserve them. They contain $\Dctf(X,\Lambda)$.
\end{proof}

Consider a commutative square of schemes with $f\colon X\to Y$,
$f'\colon X'\to Y'$, $g\colon Y'\to Y$, $g'\colon X'\to X$ and
$f\circ g'=g\circ f'$. Its base change map is
\[
 g^*Rf_*\longrightarrow Rf'_*g'^*
\]
of functors $D(X_{\proet},\Lambda)\to D(Y'_{\proet},\Lambda)$ obtained by applying the unit of $f'^*\dashv Rf'_*$, the isomorphism
$f'^*g^*\simeq g'^*f^*$, and the counit of $f^*\dashv Rf_*$. For a cartesian square with $f$ proper we show next that it is an
isomorphism on $\Loc(\Dctf(X,\Lambda))$; on $\Dctf(X,\Lambda)$ this is
Proposition~\ref{b:prop:proper-coefficients}. For integral
coefficients that proposition uses one more point. As functors
$\Dctf(Y,\mathcal O_E)\to\Dctf(Y',\mathcal O_E)$, the
inverse image used here agrees, after restriction of scalars to the
constant sheaf $\mathcal O_E$, with the completed inverse image of the
Stacks Project, by derived completeness and reduction modulo $\pi_E$
(Lemma~\ref{i:tensor}), so the proper base change theorem of the Stacks
Project for this completed inverse image (\stacks{09C9}) applies.

\begin{proposition}[Proper base change on the localizing categories]
\label{b:prop:proper-base-change-unbounded}
Let $\Lambda$ be admissible. Consider a cartesian square
\[
\begin{array}{ccc}
X'&\xrightarrow{g'}&X\\
\downarrow p'&&\downarrow p\\
Y'&\xrightarrow{g}&Y
\end{array}
\]
with $p$ proper. The base change map $g^*Rp_*K\to Rp'_*g'^*K$ in
$D(Y'_{\proet},\Lambda)$ is an isomorphism for every
$K\in\Loc(\Dctf(X,\Lambda))$.
\end{proposition}

\begin{proof}
Inverse image is exact, preserves direct sums, and sends
$\Dctf(-,\Lambda)$ into $\Dctf(-,\Lambda)$. For finite $A$ this follows
from flat constructible models, and for $\mathcal O_E$, $E$ and $\Qbar$
from pullback of the local coefficient models
(Lemma~\ref{i:tensor} in the integral case). Thus it preserves
$\Loc(\Dctf(-,\Lambda))$. By
Proposition~\ref{b:prop:proper-image-localizing}, both $g^*Rp_*$ and
$Rp'_*g'^*$ are exact functors
$\Loc(\Dctf(X,\Lambda))\to\Loc(\Dctf(Y',\Lambda))$ preserving direct sums.
The base change map is the natural transformation of functors
$D(X_{\proet},\Lambda)\to D(Y'_{\proet},\Lambda)$ defined above from the
adjunctions for inverse image and direct image. On $\Dctf(-,\Lambda)$ it is an isomorphism
by Proposition~\ref{ctf:finite-proper} for finite rings and by
Proposition~\ref{b:prop:proper-coefficients} for $\mathcal O_E$, $E$ and $\Qbar$.
Lemma~\ref{b:lem:localizing-comparison} extends it to
$\Loc(\Dctf(X,\Lambda))$.
\end{proof}

By Nagata's theorem (Theorem~\ref{app:nagata-theorem}), every separated
morphism of finite type $f:X\to Y$ with $Y$ qcqs has a
compactification: a factorization $f=p\circ j$ with
$j:X\to\overline X$ a quasi-compact open immersion and
$p:\overline X\to Y$ proper. In comparisons it
is convenient to replace $\overline X$ by the scheme-theoretic closure
of $X$, so that $X$ is dense; by
Lemma~\ref{app:nagata-density-normalization} this does not change
$Rp_*j_!$.

If $j:X\to Z$ is an open immersion, then $j_!\dashv j^*$ for
$j_!\colon D(X_{\proet},\Lambda)\to D(Z_{\proet},\Lambda)$ and
$j^*\colon D(Z_{\proet},\Lambda)\to D(X_{\proet},\Lambda)$, so a map out of
$j_!K$ in $D(Z_{\proet},\Lambda)$ is determined by its restriction to $X$. All comparison maps
between compactifications below restrict to the identity on $X$, and
this fixes them uniquely.

\begin{proposition}[Construction and independence of $Rf_!$]
\label{b:prop:lower-exceptional-image}
Let $\Lambda$ be admissible and let $f:X\to Y$ be separated of finite
type, with compactification
$f=p\circ j$ as above. The functor
\[
 Rf_!=Rp_*\circ j_!:\Loc(\Dctf(X,\Lambda))\to\Loc(\Dctf(Y,\Lambda))
\]
is exact, preserves direct sums, and maps $\Dctf(X,\Lambda)$ into
$\Dctf(Y,\Lambda)$. Two compactifications give canonically isomorphic
functors, and these isomorphisms satisfy the cocycle condition.
\end{proposition}

\begin{proof}
The properties of $Rp_*\circ j_!$ for a single compactification follow from
Proposition~\ref{b:prop:proper-image-localizing} and the preservation
of $\Dctf(-,\Lambda)$ by proper direct image
(Proposition~\ref{ctf:finite-proper} for finite rings and
Proposition~\ref{b:prop:proper-coefficients} for $\mathcal O_E$, $E$ and $\Qbar$). For independence it suffices to
compare a compactification with a refinement of it: after replacing
compactifications by dense ones, which does not change the functor
(Lemma~\ref{app:nagata-density-normalization}), any two have a common
refinement (Proposition~\ref{app:nagata-refinement}).

\emph{Refinements.} Let $j:X\to Z$ and $j':X\to Z'$ be dense
compactifications over $Y$, and let $h:Z'\to Z$ be a refinement
(Definition~\ref{app:nagata-definition}): a proper $Y$-morphism
with $h\circ j'=j$ and $h^{-1}(X)=X$. Restriction to the open subscheme
$X$ commutes with $Rh_*\colon D(Z'_{\proet},\Lambda)\to D(Z_{\proet},\Lambda)$
(restriction is exact and preserves K-injectives,
its left adjoint $j_!$ being exact), so $j^*Rh_*j'_!K\simeq K$ in
$D(X_{\proet},\Lambda)$. By adjunction this isomorphism
corresponds to a map
\begin{equation}
 e_{h,K}:j_!K\longrightarrow Rh_*j'_!K\quad\text{in }D(Z_{\proet},\Lambda)
 \label{b:eq:refinement-map}
\end{equation}
whose restriction to $X$ is the identity. We claim it is an isomorphism
for $K\in\Loc(\Dctf(X,\Lambda))$. By
Proposition~\ref{b:prop:proper-image-localizing} both sides are exact
functors of $K$ preserving direct sums, so by
Lemma~\ref{b:lem:localizing-comparison} it suffices to take
$K\in\Dctf(X,\Lambda)$. On $X$ the map is the identity. On the closed
complement $Z\setminus X$, proper base change on $\Dctf$
(Proposition~\ref{b:prop:proper-base-change-unbounded}, for the inclusion
of $Z\setminus X$ into $Z$) identifies, in
$D((Z\setminus X)_{\proet},\Lambda)$, the restriction
of $Rh_*j'_!K$ with the direct image of the restriction of $j'_!K$ to
$h^{-1}(Z\setminus X)=Z'\setminus X$, which is zero; the restriction of
$j_!K$ is zero as well. By the open--closed exact sequence, applied in
each degree (\stacks{09AH}), a map of complexes that is an
isomorphism on $X$ and on $Z\setminus X$ is an isomorphism. Applying $Rp_*$
and using $R(p\circ h)_*=Rp_*Rh_*$ gives an isomorphism between the two
candidate functors $Rf_!\colon\Loc(\Dctf(X,\Lambda))\to\Loc(\Dctf(Y,\Lambda))$.

If $Z''\xrightarrow{k}Z'\xrightarrow{h}Z$ are two successive refinements,
then $Rh_*(e_{k,K})\circ e_{h,K}$ and $e_{hk,K}$ are maps from $j_!K$ to
$R(hk)_*j''_!K$ that both restrict to the identity on $X$, so they are
equal by the observation before the proposition.

\emph{Common refinements.} Given dense compactifications $Z_1$ and
$Z_2$, Proposition~\ref{app:nagata-refinement} gives a dense
compactification $Z$ with proper maps $h_i:Z\to Z_i$ satisfying
$h_i^{-1}(X)=X$, and similarly for finitely many compactifications. The
refinement isomorphisms for $h_1$ and $h_2$ compose to an isomorphism
between the functors defined by $Z_1$ and $Z_2$. Transitivity of the
refinement maps on a common refinement shows that this comparison does
not depend on the chosen refinement and satisfies the cocycle
condition.
\end{proof}

\begin{proposition}[Composition and base change for $Rf_!$]
\label{b:prop:lower-composition}
For every admissible $\Lambda$ the following statements hold on
$\Loc(\Dctf(-,\Lambda))$ and restrict to $\Dctf(-,\Lambda)$.
\begin{enumerate}
\item For separated morphisms of finite type $X\xrightarrow fY
  \xrightarrow gZ$ there is a canonical isomorphism
  $Rg_!\circ Rf_!\simeq R(gf)_!$ of functors
  $\Loc(\Dctf(X,\Lambda))\to\Loc(\Dctf(Z,\Lambda))$. These isomorphisms are associative and
  compatible with the identity $R(\mathrm{id})_!=\mathrm{id}$.
\item Let $v:Y'\to Y$ be a morphism, $f':X'\to Y'$ the base change of $f$
  along $v$, and $u:X'\to X$ the projection. The base change map
  $v^*Rf_!\to Rf'_!u^*$, a natural transformation of functors
  $\Loc(\Dctf(X,\Lambda))\to\Loc(\Dctf(Y',\Lambda))$, is an isomorphism.
\end{enumerate}
\end{proposition}

\begin{proof}
\emph{The composition isomorphism.} Choose compactifications $f=p\circ j$
and $g=q\circ i$, with $j:X\to\overline X$ and $i:Y\to\overline Y$ open and
$p,q$ proper. In the composite $Rq_*i_!Rp_*j_!$ the middle factor $i_!Rp_*$
applies a proper direct image before an extension by zero; we
interchange the two.
Compactify the morphism $i\circ p:\overline X\to\overline Y$ as
$\overline X\xrightarrow{j'}W\xrightarrow{r}\overline Y$, with $j'$ open,
$r$ proper and $\overline X$ dense in $W$. Then $r^{-1}(Y)=\overline X$:
the open immersion $\overline X\to r^{-1}(Y)$ is a morphism from a proper
$Y$-scheme to a separated $Y$-scheme, hence proper, so its image is closed;
being also open and dense, it is everything.

For $A$ in $D(\overline X_{\proet},\Lambda)$, restriction to the open
$Y\subset\overline Y$ identifies $i^*Rr_*j'_!A$ with $Rp_*A$ in
$D(Y_{\proet},\Lambda)$, because
$r^{-1}(Y)=\overline X$. By the adjunction $i_!\dashv i^*$ this gives a map
\begin{equation}
 i_!Rp_*A\longrightarrow Rr_*j'_!A\quad\text{in }D(\overline Y_{\proet},\Lambda),
 \label{b:eq:open-proper-interchange}
\end{equation}
which is the unique map restricting to the identity of $Rp_*A$ over $Y$.
For $A\in\Dctf(\overline X,\Lambda)$ it is an isomorphism: over $Y$ it is
the identity, and over $\overline Y\setminus Y$ both sides vanish, the
right side by proper base change on $\Dctf$
(Proposition~\ref{b:prop:proper-base-change-unbounded}, for the inclusion
of $\overline Y\setminus Y$ into $\overline Y$) since $j'_!A$ vanishes on
$r^{-1}(\overline Y\setminus Y)=W\setminus\overline X$. Both sides preserve
direct sums on $\Loc(\Dctf(\overline X,\Lambda))$
(Proposition~\ref{b:prop:proper-image-localizing}), so
Lemma~\ref{b:lem:localizing-comparison} shows that
\eqref{b:eq:open-proper-interchange} is an isomorphism for all
$A\in\Loc(\Dctf(\overline X,\Lambda))$. Taking $A=j_!K$ gives, in $D(Z_{\proet},\Lambda)$,
\[
 Rq_*i_!Rp_*j_!K\simeq Rq_*Rr_*j'_!j_!K=R(qr)_*(j'j)_!K,
\]
and the right side is $R(gf)_!K$ computed with the compactification
$gf=(qr)\circ(j'j)$.

\emph{Canonicity and associativity.} Every map used is characterized by
its restriction to an open subscheme: the refinement maps
\eqref{b:eq:refinement-map} and the interchange maps
\eqref{b:eq:open-proper-interchange} restrict to identities. Replacing any
of the chosen compactifications by a refinement therefore commutes with
\eqref{b:eq:open-proper-interchange}, because both composites restrict to
the same identity and are then equal by adjunction. For three composable
morphisms, choose common refinements (by graph closure) of the finitely
many compactifications that occur in the binary composites. On these, each
of the two bracketings is built from composites of proper direct images,
composites of extensions by zero, and interchange maps. Composition of
proper direct images is associative, being composition of right adjoints
of the exact inverse images, and composition of extensions by zero is
associative by adjunction. Interchanging two adjacent open and proper
steps in either order gives the same map by the uniqueness of
\eqref{b:eq:open-proper-interchange}. Hence the two bracketings agree; the
same argument with four morphisms gives the pentagon, and the identity
constraints are the case where one compactification is the identity.
Transitivity of the refinement maps transports these equalities back to
the originally chosen compactifications.

\emph{Base change.} Pull back a compactification $f=p\circ j$ along $v$ to
a compactification $f'=p'\circ j'$, since open immersions and proper
morphisms are stable under base change. Extension by zero commutes with inverse
image: let $v':\overline X'\to\overline X$ be the projection; the functors
$v'^*j_!$ and $j'_!u^*$ from $D(X_{\proet},\Lambda)$ to
$D(\overline X'_{\proet},\Lambda)$ agree on the pulled-back open and are zero
on its closed complement, which characterizes them by \stacks{09AH}. Proper base change
on $\Loc(\Dctf(\overline X,\Lambda))$ is
Proposition~\ref{b:prop:proper-base-change-unbounded}. Composing the two
gives the isomorphism $v^*Rp_*j_!\simeq Rp'_*v'^*j_!\simeq Rp'_*j'_!u^*$ of
functors $\Loc(\Dctf(X,\Lambda))\to\Loc(\Dctf(Y',\Lambda))$. It is
compatible with change of compactification because the refinement maps are
characterized by their restriction to the open subscheme.
\end{proof}

\subsection{Compactness and proper finiteness for torsion coefficients}
\label{ctf:finite-foundations}

For Brown representability two facts are needed for each coefficient
ring: the constructible complexes of finite Tor amplitude are compact
generators of their localizing closure, and $Rf_!$ preserves this
closure and its direct sums. Here we prove both for torsion
coefficients. Let $A$ be a finite commutative ring annihilated by a
power of $\ell$. For discrete $A$, the condensed coefficient sheaf $A_X$
is the pullback of the constant coefficient sheaf on $X_{\et}$. Besides
the operations of \S\ref{b:sec:brown-construction}, the proofs use
sections with support $R\underline\Gamma_Z$ along closed immersions;
the exceptional inverse image of a general morphism is constructed from
these results in the next subsection.

\begin{lemma}[Flat models and generators]
\label{ctf:finite-models}
Every $C\in\Dctf(X,A)$ on $X_{\proet}$ is represented by a bounded complex of
flat constructible $A$-sheaves. A flat constructible sheaf has a finite
stratification on which it is locally constant with finite projective
$A$-module local models. Consequently $\Dctf(X,A)$ is generated, using
shifts, finite extensions and direct summands, by sheaves $a_!L$, where
$a:S\hookrightarrow X$ is locally closed and $L$ is locally constant
with finite projective local models.
\end{lemma}
\begin{proof}
By Corollary~\ref{pe:finite-tor-comparison}, $C=\nu_X^*C_{\et}$ with
$C_{\et}\in\Dctf(X_{\et},A)$, where
$\nu_X^*\colon D(X_{\et},A)\to D(X_{\proet},A)$ is the comparison inverse
image. The complex $C_{\et}$ has a bounded flat model by \stacks{0F4M},
Lemma~59.77.3, and $\nu_X^*$, which is exact on sheaves, pulls it back to a bounded flat
constructible model on $X_{\proet}$. Choose a common stratification on
which its finitely many terms are locally constant. A finite flat
module over the Noetherian ring $A$ is finite projective. For a flat
constructible sheaf $F$, the successive open--closed exact sequences
express $F$ as a finite extension of the sheaves $a_!L$ attached to its
strata, and the stupid filtration of the bounded flat complex expresses
$C$ as a finite extension of shifts of its \emph{terms}. Each of these
operations preserves finite Tor amplitude. The same argument gives the
corresponding statement on $X_{\et}$, which the compactness proof below
uses before passing to the pro-\'etale site.
\end{proof}

The lemma is phrased through flat models because the cohomology
sheaves of an object of $\Dctf$ need not have finite Tor amplitude.
For example, the complex
\[
 [A\xrightarrow{\ell}A],\qquad A=\mathbf Z/\ell^2\mathbf Z,
\]
is perfect, while its nonzero cohomology sheaves are copies
of $\mathbf F_\ell$, which have infinite projective dimension over $A$.
Thus $\Dctf$ is a thick triangulated category that is not the bounded
derived category of an abelian category of sheaves of finite Tor
amplitude.

\begin{proposition}[Compactness for finite coefficients]
\label{ctf:finite-compact}
Every object of $\Dctf(X,A)$ is compact in $\Loc(\Dctf(X,A))$,
and a small skeleton of $\Dctf(X,A)$ is a set of compact generators.
Moreover $\Loc(\Dctf(X,A))^c=\Dctf(X,A)$.
\end{proposition}
\begin{proof}
We first prove, on the \'etale site, compactness against arbitrary
objects of $D(X_{\et},A)$. By Lemma~\ref{ctf:finite-models} it
suffices to treat $a_!L$. Factor $a=i\circ u$ with
$i:Z\hookrightarrow X$ closed and $u:S\hookrightarrow Z$ open, and let
$j:X\setminus Z\hookrightarrow X$ be the complementary open immersion.
Since $L$ has finite projective stalks, it is dualizable with dual
$L^\vee=\underline{\operatorname{Hom}}_A(L,A_S)$, and the \'etale form of
Lemma~\ref{b:lem:formal-closed-support}(2) gives, for $K\in D(X_{\et},A)$,
\[
 \RHom_X(a_!L,K)
 =R\Gamma\bigl(S,L^\vee\otimes_A u^*R\underline\Gamma_ZK\bigr),
 \qquad R\underline\Gamma_ZK=i^*\operatorname{fib}\bigl(K\to Rj_*j^*K\bigr).
\]
Every functor on the right, as a functor between \'etale derived
categories, preserves arbitrary direct sums. For $Rj_*$
and $R\Gamma(S,-)$ this is \stacks{0F11}: a derived direct image of
finite cohomological dimension commutes with direct sums of unbounded
complexes. It applies to $A$-modules because derived direct
image and global sections commute with forgetting the coefficient ring
\stacks{0D6J}. The other functors are inverse images, tensor products
with finite projective sheaves, and a fiber. Taking $H^0$ gives
compactness of $a_!L$. Shifts, finite extensions and direct summands
preserve compactness, so Lemma~\ref{ctf:finite-models} extends it to
all of $\Dctf(X_{\et},A)$.

To pass to the pro-\'etale site, use
Corollary~\ref{pe:finite-tor-comparison} and
Proposition~\ref{pe:unbounded-classical}: the inverse image
$\nu_X^*\colon D(X_{\et},A)\to D(X_{\proet},A)$ is exact, fully
faithful and coproduct-preserving, so it identifies the localizing
closures of the corresponding finite-Tor constructible objects.
Compactness against all \'etale complexes therefore gives compactness
in $\Loc(\Dctf(X,A))$. The identification of this localizing closure
with the whole \'etale derived category is the next proposition.

Constructibility and finite Tor amplitude are preserved by direct
summands, so $\Dctf(X,A)$ is thick. It is essentially small,
being a full subcategory of the essentially small bounded constructible
category. A skeleton generates $\Loc(\Dctf(X,A))$ by definition, and the
usual compact-generation argument (\stacks{09SN} and the finite
factorization argument in \stacks{09SM}) identifies the compact objects
with the thick closure of this skeleton, which is $\Dctf(X,A)$.
\end{proof}

\begin{proposition}[Comparison with the \'etale derived category]
\label{ctf:finite-comparison}\label{ctf:eq:unbounded-comparison}
For a finite admissible ring $A$, the comparison functor
$\nu_X^*\colon D(X_{\et},A)\to D(X_{\proet},A)$ identifies
\[
 D(X_{\et},A)=\Loc(\Dctf(X_{\et},A))
       \simeq\Loc(\Dctf(X,A)).
\]
It preserves arbitrary sums and identifies ordinary direct image,
pullback and compactly supported direct image with their classical
calculations. In particular
$\nu_Y^*Rg_{\et,*}K\simeq Rg_{\proet,*}\nu_X^*K$ for every
$K\in D(X_{\et},A)$ and every finite-type $g$ in the standing range.
\end{proposition}
\begin{proof}
The bounded finite-Tor comparison is
Corollary~\ref{pe:finite-tor-comparison}, and its unbounded full
faithfulness, coproduct preservation and direct-image comparison are
Proposition~\ref{pe:unbounded-classical}. An exact fully faithful
coproduct-preserving functor identifies the localizing closures of
corresponding generators. The comparisons for pullback and extension by
zero then give $Rf_!$ through a compactification.

For the equality with the whole \'etale derived category, let
$u:U\to X$ range over a basis of quasi-compact separated \'etale
objects. The free representable module sheaves $u_!A_U$ are flat
constructible, hence belong to $\Dctf(X_{\et},A)$, and their shifts
detect zero objects of $D(X_{\et},A)$ through derived sections on $U$.
The \'etale part of the proof of Proposition~\ref{ctf:finite-compact}
shows that they are compact in all of $D(X_{\et},A)$: it uses local
support, continuity and cohomological dimension, which hold for them. By the standard
compact-generation theorem, their localizing closure is the entire
\'etale derived category.
\end{proof}

\begin{proposition}[Proper direct image for finite coefficients]
\label{ctf:finite-proper}
Let $p:X\to Y$ be proper. Then $Rp_*$ sends $\Dctf(X,A)$ to
$\Dctf(Y,A)$ and sends $\Loc(\Dctf(X,A))$ to $\Loc(\Dctf(Y,A))$,
preserving arbitrary direct sums on the latter categories. For a cartesian square as in
Proposition~\ref{b:prop:proper-base-change-unbounded}, the proper
base change map $g^*Rp_*\to Rp'_*g'^*$ is an isomorphism of functors
$\Loc(\Dctf(X,A))\to\Loc(\Dctf(Y',A))$. Extension by zero along an
open immersion preserves $\Dctf(-,A)$ and arbitrary sums, so a
compactification $f=p\circ j$ supplies the same properties for $Rp_*j_!$.
\end{proposition}
\begin{proof}
Work first on the \'etale site. Proper direct image of constructible
torsion sheaves is constructible (\stacks{0GL0}), and $Rp_*$ has finite
cohomological dimension (\stacks{0F10}); together they give
$Rp_*C\in\Dbc(Y_{\et},A)$. For finite Tor amplitude, suppose $C$
has Tor amplitude $[a,b]$, and choose a uniform bound $d$ for the
$\ell$-cohomological dimensions of the geometric fibers of $p$.
For any $A$-module $N$, proper base change and the constant-module
projection formula \stacks{0F12} in $D(Y_{\et},A)$ identify the stalk of
$(Rp_*C)\otimes_A^{\mathbf L}N_Y$ at $\bar y$ with
\[
 R\Gamma\bigl(X_{\bar y},
       C|_{X_{\bar y}}\otimes_A^{\mathbf L}N_{X_{\bar y}}\bigr),
\]
which has cohomology only in $[a,b+d]$. Testing against constant
modules suffices: a flat resolution of $Rp_*C$ restricts at every
geometric point to a flat resolution of the stalk, and tensoring with
an arbitrary sheaf is computed stalkwise. Thus
$Rp_*C\in\Dctf(Y_{\et},A)$.

On unbounded \'etale derived categories, proper direct image
$Rp_*\colon D(X_{\et},A)\to D(Y_{\et},A)$ preserves all direct sums by \stacks{0F0D}, and unbounded proper base change
follows from the same uniform cohomological bound and the bounded below
theorem (equivalently \stacks{0F0C}). The $A$-linear statements follow
by forgetting coefficients, and Proposition~\ref{ctf:finite-comparison}
transfers them to the pro-\'etale localizing categories. Since $Rp_*$
preserves the small subcategory and sums, it preserves its localizing
closure. Finally, $j_!$ is exact, is a left adjoint, and is stalkwise
extension by zero, so it preserves constructibility, finite Tor
amplitude and all sums.

The comparison of compactifications and the composition isomorphisms
(Propositions~\ref{b:prop:lower-exceptional-image}
and~\ref{b:prop:lower-composition}) need only the base change and
direct-sum assertions just proved.
\end{proof}

Propositions~\ref{ctf:finite-compact} and~\ref{ctf:finite-proper} are
the hypotheses on $A$ used in the construction of $Rf_!$ and in Brown
representability, so these constructions apply to $A$ as to the
integral and rational coefficient rings. We write $f_A^!$ for the right
adjoint of $Rf_!\colon\Loc(\Dctf(X,A))\to\Loc(\Dctf(Y,A))$ given by Brown
representability. That it preserves $\Dctf(-,A)$ is a geometric
statement, proved in Section~\ref{sec:restriction-chapter}.


\subsection{Brown representability for all coefficient rings}

Let $\Lambda$ be an admissible coefficient ring. By
Propositions~\ref{b:prop:lower-exceptional-image}
and~\ref{b:prop:lower-composition}, $Rf_!$ is defined on
$\Loc(\Dctf(-,\Lambda))$ independently of the compactification, with
composition and base change isomorphisms; it preserves arbitrary
direct sums and sends $\Dctf(-,\Lambda)$ into $\Dctf(-,\Lambda)$.
Brown representability further needs the generators to be compact.
This is proved separately for finite, integral and rational
coefficients.

For integral coefficients, compactness is proved in
Section~\ref{i:heart}, and only the structural, Ext-continuity and
compactness results of that section are used here. Their proofs use
local coefficient models, sections with closed support and torsion
cohomological bounds; since $f^!$ and open direct image play no role in
them, the forward reference is not circular.

\begin{lemma}[Hypotheses of Brown representability]
\label{b:lem:brown-hypotheses}\label{ctf:unified-compact}
For every admissible $\Lambda$, the category $\Loc(\Dctf(X,\Lambda))$
is locally small and triangulated, has arbitrary direct sums, and
has split idempotents. The thick category $\Dctf(X,\Lambda)$ is
essentially small, and a small skeleton is a set of compact generators
of $\Loc(\Dctf(X,\Lambda))$.
\end{lemma}
\begin{proof}
The ambient category is the derived category of a Grothendieck
abelian category of module sheaves, hence locally small, and so is its
full subcategory $\Loc(\Dctf(-,\Lambda))$. It has direct sums by
definition, so idempotents split (\stacks{05QW}). Constructibility and
finite Tor amplitude are preserved by shifts, cones and direct
summands. For essential smallness: the admissible constructible hearts
have a set of representatives, and bounded complexes with such
cohomology are obtained by finitely many extensions classified by
ambient Hom sets; the full subcategory with finite Tor amplitude is
contained in this set. The truncations used here are those of the
ambient category.

Compactness for finite $A$ is Proposition~\ref{ctf:finite-compact}.
For $E$ and $\Qbar$ it is
Theorem~\ref{b:thm:constructible-compactness}, and for $\mathcal O_E$
it is Proposition~\ref{i:compactness}; in the integral and rational cases
$\Dctf=\Dbc$. Finally, suppose $K\in\Loc(\Dctf(X,\Lambda))$ receives no
nonzero maps from shifts of objects of a small skeleton. The objects
$M$ with $\Hom(M[n],K)=0$ for all $n$ form a localizing subcategory
containing $\Dctf$, hence containing $K$. Thus $\mathrm{id}_K=0$, and
the skeleton generates.
\end{proof}

\begin{theorem}[Existence of the exceptional inverse image]
\label{b:thm:brown-adjoint}\label{ctf:unified-brown}
Let $\Lambda\in\{A,\Zl,\mathcal O_E,\Ql,E,\Qbar\}$ be any admissible
coefficient ring. For every separated finite-type morphism $f:X\to Y$,
the coproduct-preserving functor
$Rf_!:\Loc(\Dctf(X,\Lambda))\to\Loc(\Dctf(Y,\Lambda))$ has an exact
right adjoint
\[
 f^!:\Loc(\Dctf(Y,\Lambda))\longrightarrow\Loc(\Dctf(X,\Lambda)),
 \qquad
 \Hom(Rf_!M,N)\simeq\Hom(M,f^!N),
\]
for $M\in\Loc(\Dctf(X,\Lambda))$ and $N\in\Loc(\Dctf(Y,\Lambda))$.
It is unique up to unique isomorphism compatible with the adjunction
and preserves arbitrary direct sums.
\end{theorem}
\begin{proof}
For $N\in\Loc(\Dctf(Y,\Lambda))$, the functor
\[
 \Loc(\Dctf(X,\Lambda))^{\mathrm{op}}\longrightarrow\Ab,
 \qquad M\longmapsto\Hom(Rf_!M,N)
\]
is cohomological and sends direct sums to products. By
Lemma~\ref{b:lem:brown-hypotheses}, Brown representability
(\stacks{0A8F}) gives a representing object $f^!N$. The Yoneda lemma
makes these objects into a right adjoint, and a right adjoint of an
exact functor of triangulated categories is exact (\stacks{0A8G}).

For direct sums, test the canonical map
$\bigoplus_\alpha f^!N_\alpha\to f^!(\bigoplus_\alpha N_\alpha)$
in $\Loc(\Dctf(X,\Lambda))$ on a compact generator $G\in\Dctf(X,\Lambda)$ and its shifts.
Both $G$ and $Rf_!G\in\Dctf(Y,\Lambda)$ are compact, so adjunction
identifies the two Hom groups with
$\bigoplus_\alpha\Hom(Rf_!G[n],N_\alpha)$. The generators detect
isomorphisms, which proves the assertion.
\end{proof}


\subsection{Composition, open restriction and base change}
\label{n:compactifications}
Later sections use the compatibilities of $Rf_!$ in the following form. Any two compactifications of $f$ define
canonically isomorphic functors $Rp_*\circ j_!$, and these isomorphisms
satisfy the cocycle condition
(Proposition~\ref{b:prop:lower-exceptional-image}). They are built from
the refinement maps~\eqref{b:eq:refinement-map}, each of which is
characterized by restricting to the identity on $X$. For composable $f$
and $g$ there are associative isomorphisms
$R(gf)_!\simeq Rg_!\circ Rf_!$, satisfying the pentagon and compatible
with the identity (Proposition~\ref{b:prop:lower-composition}(1)). For a
cartesian square
\[
 \begin{array}{ccc}
 X'&\xrightarrow{g'}&X\\
 {\scriptstyle f'}\downarrow&&\downarrow{\scriptstyle f}\\
 Y'&\xrightarrow{g}&Y
 \end{array}
\]
the base change map $g^*Rf_!\to Rf'_!g'^*$ of functors
$\Loc(\Dctf(X,\Lambda))\to\Loc(\Dctf(Y',\Lambda))$ is an isomorphism that does not depend on the compactification
(Proposition~\ref{b:prop:lower-composition}(2)). Let $f'=p'\circ j'$ be
the pulled-back compactification and $\bar g\colon\overline X'\to\overline X$
the projection. The base change map is the composite of the isomorphism
$\bar g^*j_!\simeq j'_!g'^*$ of functors
$D(X_{\proet},\Lambda)\to D(\overline X'_{\proet},\Lambda)$ with the
proper base change isomorphism for $p$
(Proposition~\ref{b:prop:proper-base-change-unbounded}). The same
construction shows that these isomorphisms are compatible with
composing two cartesian squares. All of this restricts to
$\Dctf(-,\Lambda)$: inverse image and extension by zero
preserve this category, and proper direct image preserves it by
Proposition~\ref{ctf:finite-proper} for finite rings and
Proposition~\ref{b:prop:proper-coefficients} for $\mathcal O_E$, $E$ and $\Qbar$.



Whether $f^!$ preserves $\Dctf$ is a geometric question, answered in
Theorems~\ref{ctf:finite-restriction},~\ref{b:thm:restriction}
and~\ref{i:restriction}. Composition, open restriction and the
factorization through a compactification follow formally from the
corresponding properties of $Rf_!$ by passing to right adjoints. Base
change for $f^!$ requires a separate exchange analysis.

\begin{proposition}[Composition, open restriction and compactifications for $f^!$]
\label{b:prop:composition-locality}\label{n:shriek-compactification}
For every admissible $\Lambda$, the adjoints $f^!$ on
$\Loc(\Dctf(-,\Lambda))$ have the following canonical compatibilities.
\begin{enumerate}
\item For composable $X\xrightarrow fY\xrightarrow gZ$ there are canonical
  isomorphisms $(gf)^!\simeq f^!g^!$ of functors
  $\Loc(\Dctf(Z,\Lambda))\to\Loc(\Dctf(X,\Lambda))$, associative and
  compatible with the identity.
\item For an open immersion $j:U\to X$, $j^!\simeq j^*$ as functors
  $\Loc(\Dctf(X,\Lambda))\to\Loc(\Dctf(U,\Lambda))$.
\item Let $v:V\to Y$ be an open immersion, $X_V=f^{-1}(V)$ with inclusion
  $u:X_V\to X$, and $f_V:X_V\to V$ the restriction of $f$. Then
  $u^*f^!\simeq f_V^!v^*$ as functors
  $\Loc(\Dctf(Y,\Lambda))\to\Loc(\Dctf(X_V,\Lambda))$.
\item If $f=p\circ j$ is a compactification, then $f^!\simeq j^*\circ p^!$
  on $\Loc(\Dctf(Y,\Lambda))$.
\end{enumerate}
In particular, $f^!$ can be computed locally on $X$ and on $Y$ for open
coverings.
\end{proposition}

\begin{proof}
(1) For $A\in\Loc(\Dctf(X,\Lambda))$ and $B\in\Loc(\Dctf(Z,\Lambda))$,
\[
 \Hom(A,f^!g^!B)
 \simeq\Hom(Rf_!A,g^!B)
 \simeq\Hom(Rg_!Rf_!A,B)
 \simeq\Hom(R(gf)_!A,B),
\]
the last step by Proposition~\ref{b:prop:lower-composition}. The Yoneda
lemma gives $f^!g^!\simeq(gf)^!$. Passing to right adjoints transforms
the associativity and identity constraints of
Proposition~\ref{b:prop:lower-composition}(1) into those for $f^!$.

(2) For an open immersion, $Rj_!=j_!$ (take the trivial compactification
$p=\mathrm{id}$), and $j_!\dashv j^*$ on the localizing categories. By
uniqueness of right adjoints, $j^!\simeq j^*$.

(3) From $f\circ u=v\circ f_V$ and (1), (2) we get
$u^*f^!\simeq(fu)^!=(vf_V)^!\simeq f_V^!v^*$. Together with
$j^*f^!\simeq(fj)^!$ for an open $j:U\to X$, which is (1) and (2), this
gives the locality statement.

(4) On $\Loc(\Dctf(-,\Lambda))$, $j^*$ is right adjoint to $j_!$ and $p^!$
to $Rp_*=Rp_!$, so $j^*\circ p^!$ is right adjoint to $Rp_*\circ j_!=Rf_!$,
and so is $f^!$. Uniqueness of right adjoints gives the isomorphism.
\end{proof}

The proof uses only adjunctions on $\Loc(\Dctf(-,\Lambda))$. Once
Theorems~\ref{ctf:finite-restriction}, \ref{b:thm:restriction} and
\ref{i:restriction} show that $f^!$ preserves $\Dctf(-,\Lambda)$,
these isomorphisms restrict to $\Dctf$ and identify the adjunctions
between $\Dctf(X,\Lambda)$ and $\Dctf(Y,\Lambda)$.

\paragraph*{What is now available.} For every admissible $\Lambda$ and every
separated finite-type $f$, the functor $Rf_!$ on $\Loc(\Dctf(-,\Lambda))$
and its right adjoint $f^!$, compatible with composition and open
restriction. It remains to prove that $f^!$ preserves
$\Dctf(-,\Lambda)$.

\section{The exceptional inverse image preserves constructibility}
\label{sec:restriction-chapter}
The right adjoint $f^!$ of Theorem~\ref{b:thm:brown-adjoint} lives on the
category $\Loc(\Dctf(-,\Lambda))$, and its construction gives no
reason for it to preserve $\Dctf$. It can be computed in two cases. For
the projection $q:\mathbf P^1_Y\to Y$ it is $q^*(1)[2]$, by an explicit
adjunction $Rq_*\dashv q^*(1)[2]$. For a closed immersion it is
controlled by direct image along the complementary open immersion, once
we know that this direct image preserves bounded constructibility and
finite Tor amplitude. Both facts are proved first with finite
coefficients, by calculations on the \'etale site transported to
$X_{\proet}$. They give preservation of $\Dctf(-,A)$ directly, and the
result then lifts to integral and rational coefficients, where
$\Dctf=\Dbc$. The geometry thus comes from the classical \'etale theory,
applied to the adjoint already constructed on the pro-\'etale site.

\subsection{Torsion coefficients: the projective line and open direct image}
\label{b:sec:finite-geometric}

With torsion coefficients we need two facts: the adjunction
\(Rp_*\dashv p^*(1)[2]\) for the projective line
\(p:\mathbf P^1_S\to S\) (Proposition~\ref{b:prop:geo-p1-adjunction}),
and, over a field, the statement that direct image along an open
immersion preserves bounded constructible complexes and commutes with
base change (Corollary~\ref{b:cor:geo-open-finiteness}).

The second rests on the first. Since a smooth morphism is \'etale
locally an affine space, the projective-line adjunction shows that
\(f^!\) of a smooth \(f\) is \(f^*\) tensored with an invertible
complex. For a locally constant sheaf this duality turns \(Rf_*\) into
the dual of \(Rf_!\), whose constructibility is known, and so gives
generic base change for \(Rf_*\). The general case follows by induction
on the dimension of the generic fibre: projecting to a coordinate line
lowers that dimension, and the locus the projections miss is finite over
the base, where a proper direct image detects constructibility.

\paragraph*{Conventions for this subsection.}
Here $\Lambda$ is a finite commutative $\ell$-primary ring, and
$\Lambda_X=p_X^*\Lambda$ is the coefficient sheaf obtained from its
discrete condensed ring. Every geometric statement below is on the
pro-\'etale site. Let
\[
 D_{\mathrm{class}}(X_{\proet},\Lambda)\subset D(X_{\proet},\Lambda)
\]
be the essential image of $\nu_X^*:D(X_{\et},\Lambda)\to
D(X_{\proet},\Lambda)$; its bounded below part is the category
$D^+_{\mathrm{class}}$ of Corollary~\ref{pe:classical-derived}. In the
standing finite-dimensional range, Proposition~\ref{pe:unbounded-classical}
identifies it with the \'etale derived category, and
Proposition~\ref{ctf:finite-comparison} identifies it with the
localizing category $\Loc(\Dctf(X,\Lambda))$ on which $f^!$ was constructed. For finite
coefficients the objects of $\Dbc(X,\Lambda)$ are
classical by Theorem~\ref{pe:constructible-comparison}.

All schemes and base changes are of finite type over the standing field
$k$, except where a lemma states a more general hypothesis. The generic constructibility argument uses $\Lambda=\Fell$;
a finite filtration gives the assertion for any finite $\ell$-primary
ring. For smooth morphisms and the projective-line adjunction,
statements about arbitrary complexes concern $D_{\mathrm{class}}$, where
finite cohomological dimension permits unbounded comparison.

Each proof below calculates on the \'etale site and transports the
result to the pro-\'etale site. Theorem~\ref{pe:cohomology-comparison} identifies cohomology,
and Theorem~\ref{pe:direct-image-comparison} identifies the canonical
direct-image and base-change maps; its unbounded extension in
Proposition~\ref{pe:unbounded-classical} applies in the standing range.
Inverse image, tensor products, extension by zero and the localization
triangles are identified by the comparisons of
\S\ref{pe:functor-comparison-section}, and
Corollary~\ref{pe:finite-tor-comparison} controls the finite Tor
condition. These comparisons are compatible with units and counits, so an
\'etale calculation of a canonical map proves the corresponding
pro-\'etale assertion. In particular, for a separated finite-type morphism $f:X\to S$,
the functors $Rf_!\colon D_{\mathrm{class}}(X_{\proet},\Lambda)\to
D_{\mathrm{class}}(S_{\proet},\Lambda)$ and its right adjoint $f^!$
are those of Theorem~\ref{ctf:unified-brown} on
$\Loc(\Dctf(-,\Lambda))$, and transporting their adjunction through the
classical comparison permits the computations below. Accordingly, in a
statement $g^*$, $Rf_*$, $Rf_!$, $f^!$ and $\otimes^{\mathbf L}$ are the
functors between the categories $D_{\mathrm{class}}(-_{\proet},\Lambda)$
(or their bounded below parts), and in a proof that calculates on the
\'etale site they are the corresponding functors between the categories
$D(-_{\et},\Lambda)$.

The geometry rests on three theorems about the \'etale site.
\begin{itemize}
\item \emph{Proper base change}: for proper $p$, the exchange map
  $g^*Rp_*K\to Rp'_*g'^*K$ is an isomorphism.
  Theorem~\ref{pbc:torsion-theorem} proves this for bounded below
  torsion complexes. Its finite-coefficient
  unbounded extension is \stacks{0F0C}, and comparison transports it to
  the pro-\'etale exchange map
  $g^*Rp_*K\to Rp'_*g'^*K$ in $D(S'_{\proet},\Lambda)$ for
  $p:X\to S$ proper, a base change $g:S'\to S$ with projections
  $p':X'\to S'$, $g':X'\to X$, and
  $K\in D_{\mathrm{class}}(X_{\proet},\Lambda)$.
  The geometric stalk is the pro-\'etale cohomology of the fibre.
\item \emph{Finiteness}: after classical comparison, \stacks{0GL0}
  says that $Rf_!$ maps $\Dbc(X,\Lambda)$ into $\Dbc(S,\Lambda)$.
\item \emph{Cohomological dimension}: by \stacks{0F10}, the \'etale
  direct images along morphisms of finite type $k$-schemes have finite
  cohomological amplitude on torsion sheaves; comparison gives the same
  bound for
  $Rf_*,Rf_!\colon D^+_{\mathrm{class}}(X_{\proet},\Lambda)\to
  D^+_{\mathrm{class}}(S_{\proet},\Lambda)$ on classical pro-\'etale
  coefficient sheaves, so $\operatorname{cd}_\ell(f)<\infty$ in the
  sense of Definition~\ref{n:def-cd}.
\end{itemize}
The formal properties of $Rf_!$ (composition, base change, projection
formula) are cited where used.

This subsection is part of the proof that $f^!$ preserves bounded
constructible complexes, so the only property of $f^!$ used here is the
adjunction $Rf_!\dashv f^!$ on $\Loc(\Dctf(-,\Lambda))$. The smooth
calculation computes $f^!$ on coordinate charts and proves
the tensor and base-change formulas needed for finiteness of $Rj_*$;
the global orientation, trace and composition compatibility of smooth
purity are proved in Section~\ref{sec:smooth-purity}.

\subsubsection{The cycle class of a divisor and the projective line}

The adjunction \(Rp_*\dashv p^*(1)[2]\) for
\(p:\mathbf P^1_S\to S\) is built from two maps: a trace
\(Rp_*\Lambda(1)[2]\to\Lambda\), read off from the cohomology of
\(\mathbf P^1\), and the cycle class of the diagonal in
\(\mathbf P^1_S\times_S\mathbf P^1_S\). The trace is the counit and the
diagonal class is the unit. Both triangle identities reduce to one
geometric fact: the diagonal meets each fibre of either projection
\(\mathbf P^1\times\mathbf P^1\to\mathbf P^1\) in a single point, so
integrating its class along either factor gives \(1\). We first
construct cycle classes of effective Cartier divisors.

\begin{lemma}[The cycle class of a Cartier divisor]
\label{b:lem:geo-cartier}
Let \(i:D\hookrightarrow Z\) be an effective Cartier divisor on a
scheme \(Z\) on which \(\ell\) is invertible. There is a morphism
\[
 c_D:i_*\Lambda_D\longrightarrow\Lambda_Z(1)[2]\quad\text{in }D(Z_{\proet},\Lambda)
\]
whose composite with \(\Lambda_Z\to i_*\Lambda_D\) is the Kummer first
Chern class \(c_1(\mathcal O_Z(D))\). It is compatible with change of
coefficients and with every base change along which the pulled-back
divisor is still an effective Cartier divisor.
\end{lemma}

\begin{proof}
We construct the map in $D(Z_{\et},\Lambda)$ and apply
$\nu_Z^*\colon D(Z_{\et},\Lambda)\to D(Z_{\proet},\Lambda)$. Closed direct
image and inverse image comparison identify its source with
$i_*\Lambda_D$ on $Z_{\proet}$, and exactness and tensor compatibility
of $\nu_Z^*$ identify the target and all base-change maps; only
bounded below complexes occur.

The idea is that \(c_D\) is the boundary map of the extension of
\(i_*\mathbf Z_D\) by \(\mathbf G_m\) formed by meromorphic functions
with poles along \(D\), reduced modulo \(\ell^r\).

Let \(j:Z\setminus D\hookrightarrow Z\) be the open complement. Let
\(\mathcal M_D\subset j_*\mathbf G_m\) be the subsheaf which, on an open
where \(D\) has a regular equation \(t\), consists of the products
\(ut^{-a}\) with \(u\) a unit and \(a\) a locally constant integer on
\(D\); equivalently, it is generated by \(\mathbf G_m\) and \(t^{-1}\).
Replacing \(t\) by a unit multiple gives the same subsheaf, so
\(\mathcal M_D\) is globally defined. Sending \(t^{-1}\) to \(1\) gives
a short exact sequence of abelian sheaves on \(Z_{\et}\)
\[
 0\longrightarrow\mathbf G_m\longrightarrow \mathcal M_D
   \longrightarrow i_*\mathbf Z_D\longrightarrow0.       \tag{G1}
\]
Exactness is checked on geometric stalks. Away from \(D\) the quotient
is zero. At a point of \(D\), no nonzero power of a regular nonunit
\(t\) is a unit, so the exponent is well defined, and it is surjective
because exponents lift locally to powers of \(t^{-1}\).

The boundary of (G1) is a morphism \(i_*\mathbf Z_D\to\mathbf G_m[1]\)
in \(D(Z_{\et},\mathbf Z)\).
Apply \(-\otimes^{\mathbf L}_{\mathbf Z}\mathbf Z/\ell^r\mathbf Z\). The
source becomes \(i_*\mathbf Z/\ell^r\mathbf Z\). For the target, the
Kummer sequence \stacks{03PL} shows that the cone of multiplication by
\(\ell^r\) on \(\mathbf G_m\) has the single cohomology sheaf
\(\mu_{\ell^r}\) in degree \(-1\), so the target becomes
\(\mu_{\ell^r}[2]\). This is \(c_D\) for \(\mathbf Z/\ell^r\mathbf Z\).
For a finite \(\mathbf Z/\ell^r\mathbf Z\)-algebra \(\Lambda\), apply
\(-\otimes^{\mathbf L}_{\mathbf Z/\ell^r}\Lambda\); since
\(\mu_{\ell^r}\) is locally free of rank one over
\(\mathbf Z/\ell^r\mathbf Z\), this gives \(\Lambda_Z(1)\) whether or not
\(\Lambda\) is flat.

To identify the composite with \(c_1\), cover \(Z\) by opens with local
equations \(t_a\). The local lifts of the section \(1\) of
\(i_*\mathbf Z_D\) through (G1) are the \(t_a^{-1}\), and their ratios
are the transition functions of \(\mathcal O_Z(D)\) (this fixes the
sign convention). So the boundary of (G1) followed by the Kummer
boundary is the class of \(\mathcal O_Z(D)\). Changing a local equation
multiplies a lift by a unit and changes the cocycle by a coboundary.

For naturality, pull back the local equations, the units, and the
sequence (G1); the boundary maps are compatible. On the source one uses
the base-change identity \(g^*i_*=i'_*g_D^*\) of functors
\(\Mod(D_{\et},\Lambda)\to\Mod(Z'_{\et},\Lambda)\), for a base change
\(g:Z'\to Z\) with \(i':D'\hookrightarrow Z'\) and \(g_D:D'\to D\),
checked on geometric stalks.
\end{proof}

\begin{lemma}[Cohomology of the projective line and the trace]
\label{b:lem:geo-p1-cohomology}
Let \(S\) be a Noetherian scheme on which \(\ell\) is invertible and
\(p:\mathbf P^1_S\to S\) the projection. The unit
\(\Lambda_S\to Rp_*\Lambda_{\mathbf P^1_S}\) and the class
\(c_1(\mathcal O_{\mathbf P^1_S}(1))\) define an isomorphism
\[
 \Lambda_S\oplus\Lambda_S(-1)[-2]
       \ \xrightarrow{\ \sim\ }\ Rp_*\Lambda_{\mathbf P^1_S}
       \quad\text{in }D(S_{\proet},\Lambda),   \tag{G2}
\]
compatible with base change: for \(g:S'\to S\) with \(S'\) as in the
lemma, the base-change map \(g^*Rp_*\Lambda\to Rp'_*\Lambda\) in
\(D(S'_{\proet},\Lambda)\) carries the pullback of (G2) to (G2) for
\(S'\). Let
\[
 \operatorname{Tr}_p:Rp_*\Lambda_{\mathbf P^1_S}(1)[2]\longrightarrow\Lambda_S
                                                               \tag{G3}
\]
be the morphism in \(D(S_{\proet},\Lambda)\) given by the projection
onto the summand \(\Lambda_S\) of the twist of (G2).
Then \(\operatorname{Tr}_p\) is compatible with every Noetherian base
change.
\end{lemma}

\begin{proof}
Use the bounded-below cohomology and proper direct-image comparisons
of Theorems~\ref{pe:cohomology-comparison} and
\ref{pe:direct-image-comparison}. They reduce (G2) and (G3) to the
following calculation in \(D^+(S_{\et},\Lambda)\), also in the
Noetherian generality of the lemma.

Both components of (G2) are globally defined, the second by
Lemma~\ref{b:lem:geo-cartier} applied to the section at infinity. By
proper base change and since geometric stalks detect isomorphisms, it
suffices to check (G2) on a geometric fibre, first for
\(\Lambda=\mathbf Z/\ell^r\mathbf Z\). There the genus-zero case of
Proposition~\ref{app:curve-kummer} gives cohomology only in degrees
\(0\) and \(2\), with \(H^0=\mathbf Z/\ell^r\mathbf Z\) and
\(H^2=(\mathbf Z/\ell^r\mathbf Z)(-1)\), and its Kummer description sends
a line bundle to its degree modulo \(\ell^r\). Since
\(\mathcal O(1)\) has degree one, its class generates \(H^2\), and (G2)
is an isomorphism. The unit and the class of \(\mathcal O(1)\) commute
with base change, hence so do (G2), its inverse, and (G3). For a finite
\(\mathbf Z/\ell^r\mathbf Z\)-algebra \(\Lambda\), the case
\(\mathbf Z/\ell^r\mathbf Z\) and the projection formula \stacks{0GL5}
give the assertion.
\end{proof}

Over \(S=\operatorname{Spec}\bar k\), (G2) reads
\(R\Gamma(\mathbf P^1_{\bar k},\Lambda)\simeq\Lambda\oplus\Lambda(-1)[-2]\).
The shift \([-2]\) places the second summand in degree \(2\): it is
\(H^2\). The twist \((-1)\) records how Galois acts on \(H^2\): the
Kummer map identifies \(H^2\) with
\(\operatorname{Pic}(\mathbf P^1)\otimes\mu_{\ell^r}^{\otimes-1}\), that is, with
\(\Lambda(-1)\). Twisting by \((1)\) and shifting
by \([2]\) moves this summand to degree \(0\) with trivial twist, and
\(\operatorname{Tr}_p\) is the projection onto it; it sends the class
of a point to \(1\). The same calculation appears three more times. In
Proposition~\ref{b:prop:geo-p1-adjunction} it becomes the counit of
\(Rp_*\dashv p^*(1)[2]\). In \S\ref{b:sec:late-coefficients} it is
lifted to \(\mathcal O_E\). In \S\ref{sp:grp-trace} it is the smooth
trace of \(\mathbf P^1\), and the Lefschetz formula turns it into
\(\#\mathbf P^1(\mathbb F_q)=1+q\).

\begin{example}[Compact supports on the running example]
Let \(j:\mathbf A^1_{\bar k}\hookrightarrow\mathbf P^1_{\bar k}\) with
complement the point \(\infty\). By definition
\(R\Gamma_c(\mathbf A^1,\Lambda)=R\Gamma(\mathbf P^1,j_!\Lambda)\). The
localization triangle
\(j_!\Lambda\to\Lambda_{\mathbf P^1}\to i_{\infty*}\Lambda\) in
\(D((\mathbf P^1_{\bar k})_{\proet},\Lambda)\) and (G2)
give
\[
 R\Gamma_c(\mathbf A^1_{\bar k},\Lambda)
 \simeq\fib\bigl(\Lambda\oplus\Lambda(-1)[-2]\to\Lambda\bigr)
 \simeq\Lambda(-1)[-2],
\]
since restriction to \(\infty\) is an isomorphism on \(H^0\). So
\(H^2_c(\mathbf A^1)=\Lambda(-1)\) is the only compactly supported
cohomology, and the trace on it is the same \(\operatorname{Tr}_p\).
\end{example}

\paragraph*{Exercise.} Using the localization sequence of
Example~\ref{b:ex:running}, show that
\(H^1_c(\mathbf G_{m,\bar k},\Lambda)\simeq\Lambda\) and
\(H^2_c(\mathbf G_{m,\bar k},\Lambda)\simeq\Lambda(-1)\), and that these
are the only nonzero groups.

\begin{proposition}[The projective-line adjunction]
\label{b:prop:geo-p1-adjunction}
Let \(S\) and \(p:\mathbf P^1_S\to S\) be as in
Lemma~\ref{b:lem:geo-p1-cohomology}, and assume in addition that $S$ is
of finite type over the standing field $k$. The functors
\(Rp_*\colon D_{\mathrm{class}}((\mathbf P^1_S)_{\proet},\Lambda)\to
D_{\mathrm{class}}(S_{\proet},\Lambda)\) and
\(p^*(1)[2]\colon D_{\mathrm{class}}(S_{\proet},\Lambda)\to
D_{\mathrm{class}}((\mathbf P^1_S)_{\proet},\Lambda)\) form
an adjunction
\[
 Rp_* \ \dashv\ p^*(1)[2],
                                                               \tag{G4}
\]
with unit (G6) and counit (G7) below. Both are compatible with
base change in the standing range (as natural transformations of
functors between the categories over \(S'\), for \(g:S'\to S\)) and
with change of finite coefficients.
\end{proposition}

\begin{proof}
We calculate on $D((\mathbf P^1_S)_{\et},\Lambda)$ and
$D(S_{\et},\Lambda)$, using the classical equivalence of
Proposition~\ref{pe:unbounded-classical}; until the last paragraph,
\(Rp_*\), \(p^*\), \(Rp_{a*}\), \(p_a^*\), \(\Delta_*\) and
\(\otimes^{\mathbf L}\) are the functors between these \'etale derived
categories and that of \(\mathbf P^1_S\times_S\mathbf P^1_S\).
Transporting the resulting unit and counit gives (G4) on the stated
pro-\'etale categories.
The direct-image comparison identifies the canonical exchange and
projection maps used in this transport.

The counit will be the trace (G3), and the unit will be built from the
cycle class of the diagonal. The triangle identities then reduce to a
single fact: integrating the class of the diagonal along either factor
gives \(1\).

Let \(p_1,p_2:\mathbf P^1_S\times_S\mathbf P^1_S\to\mathbf P^1_S\) be
the projections and
\(\Delta:\mathbf P^1_S\hookrightarrow\mathbf P^1_S\times_S\mathbf P^1_S\)
the diagonal. It is cut out by \(x_0y_1-x_1y_0\), so it is a relative
effective Cartier divisor, and it stays one after every base change.
Lemma~\ref{b:lem:geo-cartier} gives its cycle class
\[
 \delta:\Delta_*\Lambda_{\mathbf P^1_S}\longrightarrow
   \Lambda_{\mathbf P^1_S\times_S\mathbf P^1_S}(1)[2]
   \quad\text{in }D((\mathbf P^1_S\times_S\mathbf P^1_S)_{\et},\Lambda).
\]
For \(a=1,2\), the projection \(p_a\) is a projective line over the
other factor, with trace \(\operatorname{Tr}_{p_a}\) given by (G3).
Integrating \(\delta\) along \(p_a\) gives an endomorphism in
\(D((\mathbf P^1_S)_{\et},\Lambda)\)
\[
 s_a:\Lambda_{\mathbf P^1_S}=Rp_{a*}\Delta_*\Lambda_{\mathbf P^1_S}
       \xrightarrow{Rp_{a*}\delta}
       Rp_{a*}\Lambda_{\mathbf P^1_S\times_S\mathbf P^1_S}(1)[2]
       \xrightarrow{\operatorname{Tr}_{p_a}}\Lambda_{\mathbf P^1_S}.                 \tag{G5}
\]

\emph{Claim: \(s_1=s_2=1\).} Both sides are endomorphisms of a sheaf
placed in degree zero, and
\[
 \Hom_{D((\mathbf P^1_S)_{\et},\Lambda)}
      (\Lambda_{\mathbf P^1_S},\Lambda_{\mathbf P^1_S})
   =\Hom_{\Mod((\mathbf P^1_S)_{\et},\Lambda)}
                   (\Lambda_{\mathbf P^1_S},\Lambda_{\mathbf P^1_S}),
\]
so \(s_a\) is a morphism of sheaves and is determined by its geometric
stalks. By proper base change, the stalk of (G5) at a geometric point
\(\bar x\) is the trace of the cycle class of the single point cut out
by the diagonal in the fibre \(\mathbf P^1_{\bar x}\). The line bundle of
that point has degree one, so by (G2) and (G3) its trace is \(1\).

\emph{The unit and counit.} For \(M\in D((\mathbf P^1_S)_{\et},\Lambda)\)
define the morphism in \(D((\mathbf P^1_S)_{\et},\Lambda)\)
\[
 \begin{split}
 \eta_M:M
 &=Rp_{2*}\Delta_*M\\
 &\longrightarrow Rp_{2*}\bigl(p_1^*M(1)[2]\bigr)
   \xrightarrow{\ \sim\ }p^*Rp_*M(1)[2] .
 \end{split}                                                    \tag{G6}
\]
Here the first arrow is \(Rp_{2*}\) applied to \(\delta\otimes p_1^*M\),
using the projection formula
\(\Delta_*M=\Delta_*\Lambda\otimes^{\mathbf L}p_1^*M\) in
\(D((\mathbf P^1_S\times_S\mathbf P^1_S)_{\et},\Lambda)\), and the second
is the proper base-change isomorphism \(p^*Rp_*\simeq Rp_{2*}p_1^*\) of
functors \(D((\mathbf P^1_S)_{\et},\Lambda)\to
D((\mathbf P^1_S)_{\et},\Lambda)\). For \(N\in D(S_{\et},\Lambda)\)
define the morphism in \(D(S_{\et},\Lambda)\)
\[
 \epsilon_N:Rp_*\bigl(p^*N(1)[2]\bigr)
   \xrightarrow{\ \sim\ }N\otimes_\Lambda^{\mathbf L}
                       Rp_*\Lambda_{\mathbf P^1_S}(1)[2]
   \xrightarrow{1\otimes \operatorname{Tr}_p}N,                        \tag{G7}
\]
the first arrow being the projection formula. These maps exist for
unbounded complexes by \stacks{0F0C} and \stacks{0GL5}.

\emph{First triangle identity.} We show
\(\epsilon_{Rp_*M}\circ Rp_*(\eta_M)=1_{Rp_*M}\). Using composition of
direct images, \(pp_1=pp_2\), proper base change and the projection
formula, this composite equals
\[
 \begin{split}
 Rp_*M
 &=R(pp_1)_*\Delta_*M\\
 &\longrightarrow R(pp_1)_*p_1^*M(1)[2]\\
 &\xrightarrow{\ \sim\ }
       Rp_*\bigl(M\otimes_\Lambda^{\mathbf L}
                       Rp_{1*}\Lambda_{\mathbf P^1_S\times_S\mathbf P^1_S}(1)[2]\bigr)\\
 &\xrightarrow{Rp_*(1_M\otimes \operatorname{Tr}_{p_1})}Rp_*M .
 \end{split}                                                    \tag{G8}
\]
In words: the unit (G6) integrates over the first coordinate with the
second fixed, and the counit (G7), after \(Rp_*\), integrates over the
second. Rewriting \(R(pp_2)_*\) as \(R(pp_1)_*\) and applying the
projection formula with coefficient \(M\) moves the integration into the
kernel factor, where the trace is \(\operatorname{Tr}_{p_1}\), the base
change of \(\operatorname{Tr}_p\).

This rewriting is an equality of maps because of how the canonical maps
are defined: the proper base-change
map is the adjoint transpose of the counit of \(g^*\dashv Rg_*\) for the
relevant pullback \(g\), and the projection map is the adjoint transpose
of tensoring with the same counit. Expanding these definitions along
either route uses this counit once in each coordinate, in the same
order once \(pp_1=pp_2\) is used. Hence the first arrow of (G8) is
\(1_M\otimes Rp_{1*}\delta\) under the projection identification, and
the composite is \(Rp_*(1_M\otimes s_1)=1_{Rp_*M}\) by the claim.

\emph{Second triangle identity.} Take \(M=p^*N(1)[2]\) in (G6). Then
\(p_1^*M\) and \(p_2^*M\) are both \((pp_1)^*N(1)[2]=(pp_2)^*N(1)[2]\),
since \(N\) and the Tate twist are pulled back from \(S\). The
projection formula therefore takes this common factor outside
\(Rp_{2*}\), and applying \(p^*(\epsilon_N)(1)[2]\) contracts the
remaining kernel by \(\operatorname{Tr}_{p_2}\):
\[
 p^*(\epsilon_N)(1)[2]\circ\eta_{p^*N(1)[2]}
      =1_{p^*N(1)[2]}\otimes
          (\operatorname{Tr}_{p_2}\circ Rp_{2*}\delta)
      =1_{p^*N(1)[2]}\otimes s_2=1_{p^*N(1)[2]}.                         \tag{G9}
\]
All shifts moved in (G8) and (G9) are even, so no Koszul signs appear.

Finally, the unit is built from \(\delta\) and proper base change, and
the counit from (G3) and the projection formula. All of these are
compatible with base change and change of coefficients, hence so are
\(\eta\) and \(\epsilon\), as natural transformations of functors on
\(D((\mathbf P^1_S)_{\et},\Lambda)\) and \(D(S_{\et},\Lambda)\); the
comparison transports this to the pro-\'etale categories of (G4).
\end{proof}

\subsubsection{The exceptional inverse image of a smooth morphism and base change}

A smooth morphism is \'etale locally an \'etale map followed by a
projection from affine space, so to compute \(f^!\) for smooth \(f\) we
first treat these two cases.

\begin{lemma}[Affine space and \'etale localization]
\label{b:lem:geo-affine-local}
Let \(S\) be of finite type over the standing field \(k\).
\begin{enumerate}
\item For the projection \(a:\mathbf A^d_S\to S\), the right adjoint
\(a^!\) of \(Ra_!\colon D_{\mathrm{class}}((\mathbf A^d_S)_{\proet},\Lambda)
\to D_{\mathrm{class}}(S_{\proet},\Lambda)\) is \(a^*(d)[2d]\), with unit and counit compatible
with base change.
\item Let \(f:X\to S\) be separated of finite type, let
\(N\in D_{\mathrm{class}}(S_{\proet},\Lambda)\), and let \(u:U\to X\) be separated, \'etale
and of finite type. There is a canonical identification
\[
 u^*f^!N=(fu)^!N\quad\text{in }D_{\mathrm{class}}(U_{\proet},\Lambda).
                                                               \tag{G10}
\]
It is compatible with units, counits, composition, and further
\'etale localization.
\end{enumerate}
\end{lemma}

\begin{proof}
By the classical comparison, the following \'etale computations
calculate the pro-\'etale adjoints together with their units and
counits.

(1) Let \(j:\mathbf A^1_S\hookrightarrow\mathbf P^1_S\) be the open
immersion, so that \(Ra_!=Rp_*j_!\) for \(d=1\). Since \(j_!\dashv j^*\),
(G4) gives \(a^!=j^*p^*(1)[2]=a^*(1)[2]\). For general \(d\), write
\(a\) as a composite of \(d\) affine-line projections and use
\stacks{0F7J}. The counit is the successive composite of the \(d\)
projective-line traces and the open-immersion counits. Associativity
\stacks{0F7K} and compatibility with base change \stacks{0F7L} show
that this is the counit for the canonical composition isomorphism, and
that it commutes with base change.

(2) Extension by zero
\(u_!\colon\Mod(U_{\et},\Lambda)\to\Mod(X_{\et},\Lambda)\) is exact and
left adjoint to \(u^*\) by its elementary sheaf-theoretic construction
(\stacks{03S2}), and
the adjunction extends to derived categories because both functors are
exact. For every \(A\in D(U_{\et},\Lambda)\),
\[
 \Hom(A,u^*f^!N)
  =\Hom(u_!A,f^!N)
  =\Hom(Rf_!u_!A,N)
  =\Hom(R(fu)_!A,N),
\]
with the Hom groups taken in \(D(U_{\et},\Lambda)\),
\(D(X_{\et},\Lambda)\) and \(D(S_{\et},\Lambda)\),
the last step by the composition isomorphism of \stacks{0F7J}. The
Yoneda lemma gives (G10), and since these bijections use exactly that
composition isomorphism, the counit of \(R(fu)_!\dashv(fu)^!\) is the
composite of the two counits. Repeating the argument for a further
\'etale map gives the compatibilities with composition and
localization.
\end{proof}

\begin{proposition}[The exceptional inverse image of a smooth morphism]
\label{b:prop:geo-smooth-tensor}
Let \(f:X\to S\) be separated, smooth and of finite type, of pure
relative dimension \(d\), and let
\(\omega_f=f^!\Lambda_S\in D_{\mathrm{class}}(X_{\proet},\Lambda)\) be its dualizing
complex. Then \(\omega_f\) is pro-\'etale locally isomorphic to
\(\Lambda_X(d)[2d]\), and the canonical map
\[
 f^*N\otimes_\Lambda^{\mathbf L}\omega_f
       \longrightarrow f^!N
       \quad\text{in }D_{\mathrm{class}}(X_{\proet},\Lambda)  \tag{G11}
\]
is an isomorphism for every \(N\in D_{\mathrm{class}}(S_{\proet},\Lambda)\). Moreover, let
\(g:S'\to S\) be a morphism, \(X'=X\times_SS'\), and
\(f':X'\to S'\), \(g':X'\to X\) the projections. Then the canonical
base-change map
\[
 g'^*f^!N\longrightarrow f'^!g^*N
       \quad\text{in }D_{\mathrm{class}}(X'_{\proet},\Lambda)  \tag{G12}
\]
is an isomorphism. It identifies \(g'^*\omega_f\) with
\(\omega_{f'}\) and is compatible with (G11).
\end{proposition}

\begin{proof}
Represent the classical objects by their \'etale antecedents, so that
the functors below are those between the categories
\(D(-_{\et},\Lambda)\).
Proposition~\ref{pe:unbounded-classical} and the functor comparisons
identify the adjunctions and canonical maps below with the pro-\'etale
ones. The coordinate
charts form a pro-\'etale covering, so the computed local
isomorphisms prove the pro-\'etale statement.

Both maps are defined globally as adjoint transposes, and we check
locally that they are isomorphisms. Write \(\epsilon_N:Rf_!f^!N\to N\),
a morphism in \(D(S_{\et},\Lambda)\), for the counit of
\(Rf_!\dashv f^!\) (for \(f=p\) this is (G7)). The map (G11) is the
transpose of the morphism in \(D(S_{\et},\Lambda)\)
\[
 Rf_!(f^*N\otimes^{\mathbf L}\omega_f)
    \xrightarrow{\ \sim\ }
 N\otimes^{\mathbf L}Rf_!\omega_f
    \xrightarrow{1\otimes\epsilon_{\Lambda_S}}N ,          \tag{G13}
\]
the first arrow being the projection formula. The map (G12) is the
transpose, for \(f'\), of the morphism in \(D(S'_{\et},\Lambda)\)
\[
 Rf'_!g'^*f^!N
   \xrightarrow{\ \sim\ }g^*Rf_!f^!N
   \xrightarrow{g^*\epsilon_N}g^*N,                       \tag{G14}
\]
where the first arrow is the inverse of the base-change isomorphism
for \(Rf_!\) \stacks{0F7L}.

\emph{Local charts.} By the local structure of smooth morphisms
\stacks{054L}, every point of \(X\) has an affine open neighbourhood
\(u:U\hookrightarrow X\) mapping into an affine open
\(v:V\hookrightarrow S\), together with a separated \'etale map
\(e:U\to\mathbf A^d_V\) such that \(fu=vae\), where
\(a:\mathbf A^d_V\to V\) is the projection. By
Lemma~\ref{b:lem:geo-affine-local} and uniqueness of right adjoints,
in \(D(U_{\et},\Lambda)\)
\[
 u^*f^!N=e^*a^*v^*N(d)[2d],\qquad
 u^*\omega_f=\Lambda_U(d)[2d].                            \tag{G15}
\]
In particular \(\omega_f\) has a single nonzero cohomology sheaf, in
degree \(-2d\), and it is locally free of rank one.

\emph{(G11) is an isomorphism.} Under (G15), source and target of (G11)
are both \(e^*a^*v^*N(d)[2d]\), and one must check that the map is the
obvious reassociation of tensor product and twist. Compare transposes:
both are the projection formula followed by the counit for \(v\), the
\(d\) affine-line counits, and the \'etale counit for \(e\). They agree
by the formula for the counit of a composite in
Lemma~\ref{b:lem:geo-affine-local}. So (G11) is an isomorphism on each
such \(U\), hence on \(X\).

\emph{(G12) is an isomorphism.} Consider first each type of chart
separately. For the projective line, the unit (G6) and counit (G7)
commute with base change, so transposing (G14) gives the ordinary
pullback identification \(g'^*p^*N(1)[2]=p'^*g^*N(1)[2]\). Restriction
to the affine line and composition give the same for affine space. For
a separated \'etale map, extension by zero is the direct sum over the
geometric sheets and commutes with base change, and its counit is the
sum of the evaluations; so here too the transpose of (G14) is the
ordinary inverse-image identification. To combine these, use that base
change for \(Rf_!\) is compatible with composition \stacks{0F7L} and
that the counit of a composite is the composite of the counits.
Substituting this into (G14) shows that the restriction of (G12) to
\(U\times_XX'\) is the composite of the \'etale, affine-space and
open-immersion cases just treated, so (G12) is an isomorphism on a
cover of \(X'\).

Applying the same calculation to (G13), base change respects its
projection map and its counit, and the resulting commutative square is
the compatibility of (G11) with (G12).
\end{proof}

\begin{lemma}[Smooth morphisms and locally constant sheaves]
\label{b:lem:geo-smooth-generic}
Let \(S\) be integral of finite type over the standing field \(k\), let
\(\ell\ne\operatorname{char}(k)\), and let \(f:X\to S\) be separated,
smooth and of finite type. For a finite locally constant
\(\Fell\)-sheaf \(L\) on \(X_{\proet}\), there is a nonempty open
\(S_0\subset S\) such that, for the restriction
\(f_{S_0}:X_{S_0}\to S_0\), the complex \(Rf_{S_0,*}L\) is bounded with
finite locally constant cohomology sheaves, and its formation commutes
with base change: for every finite-type $k$-scheme \(v:T\to S_0\), with
base change \(f_T:X_T\to T\) and \(v_X:X_T\to X_{S_0}\), the canonical
base-change map \(v^*Rf_{S_0,*}L\to Rf_{T,*}v_X^*L\) in
\(D(T_{\proet},\Fell)\) is an isomorphism.
\end{lemma}

\begin{proof}
Finite-presentation local-system comparison
(Proposition~\ref{pe:finite-local-systems}) gives an \'etale antecedent
of $L$. We make the following duality and constructibility computation
there, in the categories \(D(-_{\et},\Fell)\), and then use the direct-image and constructible comparisons
to obtain the assertion on $X_{\proet}$ and $S_{\proet}$.

The idea is to write \(L\) as the dual of a complex \(N\) relative to
\(\omega_f\). Then \(Rf_*L\) is the \(\Fell\)-linear dual of \(Rf_!N\)
(see (G16)). After shrinking \(S\), \(Rf_!N\) has locally constant
cohomology, and the dual of such a complex is again of this kind and
commutes with pullback. It remains to check that the identification
(G16) is compatible with base change.

\(X\) is the disjoint union of finitely many open pieces of pure
relative dimension, so we may assume \(f\) has pure relative dimension
\(d\). Let \(\omega_f=f^!\mathbf F_{\ell,S}\) as in
Proposition~\ref{b:prop:geo-smooth-tensor}, and let
\[
 N=\cHom_X(L,\mathbf F_{\ell,X})\otimes^{\mathbf L}\omega_f .
\]
By Proposition~\ref{b:prop:geo-smooth-tensor}, \(\omega_f\) is locally
\(\mathbf F_{\ell,X}(d)[2d]\), so \(N\) is bounded constructible and
\(\cHom_X(N,\omega_f)=L\).

\emph{Duality.} There is a canonical isomorphism
\[
 Rf_*L
   = \cHom_S(Rf_!N,\mathbf F_{\ell,S})
   \quad\text{in }D(S_{\et},\Fell).                          \tag{G16}
\]
To see this, write \(L=\cHom_X(N,f^!\mathbf F_{\ell,S})\).
Test the canonical pairing on shifts of $P=a_!\Fell$, where
$a:V\to S$ is quasi-compact, separated and \'etale. These free
representable module sheaves lie in $\Dctf$ and their shifts detect
isomorphisms in $D(S_{\et},\Fell)$. Since
$f^*P\otimes^{\mathbf L}N$ lies in $\Dctf(X,\Fell)$, tensor-Hom
adjunction, the Brown adjunction $Rf_!\dashv f^!$ between
$\Loc(\Dctf(X,\Fell))$ and $\Loc(\Dctf(S,\Fell))$, and the projection
formula give, with Hom groups in $D(S_{\et},\Fell)$,
\[
 \Hom(P,Rf_*\cHom(N,f^!\mathbf F_{\ell,S}))
 =\Hom(Rf_!(f^*P\otimes^{\mathbf L}N),\mathbf F_{\ell,S})
 =\Hom(P\otimes^{\mathbf L}Rf_!N,\mathbf F_{\ell,S}).
\]
The canonical map adjoint to the pairing (G17) below realizes these
bijections for all shifts of the generators, hence is an isomorphism.
This proves (G16) in $D(S_{\et},\Fell)$; the constructibility of
$Rf_*L$ will be deduced from it.

\emph{Shrinking the base.} By \stacks{0GL0} and \stacks{0G2A}, \(Rf_!N\)
is bounded constructible. Each of its finitely many nonzero cohomology
sheaves is locally constant on a nonempty open of \(S\); let \(S_0\) be
the intersection of these opens. Over the field \(\Fell\), a finite
locally constant sheaf is locally free of finite rank, its dual is
exact, and evaluation identifies it with its double dual. By induction
on the number of nonzero cohomology sheaves, using truncation triangles
(the relevant canonical maps form morphisms of these triangles), a
bounded complex with finite locally constant cohomology is dualizable,
its dual is again of this kind, and forming the dual commutes with
every pullback. By (G16), this proves the first assertion over \(S_0\).

\emph{Base change.} The isomorphism (G16) is adjoint to the pairing
in \(D(S_{\et},\Fell)\)
\[
 Rf_!N\otimes^{\mathbf L}Rf_*L
  \longrightarrow Rf_!(N\otimes^{\mathbf L}L)
  \longrightarrow Rf_!\omega_f
  \xrightarrow{\epsilon_{\mathbf F_{\ell,S}}}\mathbf F_{\ell,S},               \tag{G17}
\]
where the first arrow is the projection formula followed by the counit
\(f^*Rf_*L\to L\), and the second is evaluation. Let \(v:T\to S_0\) be a
finite-type $k$-scheme over \(S_0\), and let \(f_T:X_T\to T\) and
\(v_X:X_T\to X\) be the base changes. Base change for \(Rf_!\) and (G12)
identify \(v_X^*N\), \(v^*Rf_!N\) and \(v_X^*\omega_f\) with the
corresponding objects for \(f_T\). The base-change map
\(v^*Rf_*L\to Rf_{T,*}v_X^*L\) in \(D(T_{\et},\Fell)\) is characterized by the fact that it
carries the pulled-back counit \(f^*Rf_*L\to L\) to the counit for
\(f_T\). Hence the pullback of (G17) is the pairing for \(f_T\),
composed with this base-change map. Since \(v^*Rf_!N\) is dualizable
with dual \(v^*\cHom(Rf_!N,\mathbf F_{\ell,S})\), the base-change map
for \(Rf_*L\) is identified with the dual of the base-change
isomorphism for \(Rf_!N\). It is therefore an isomorphism.
\end{proof}

\subsubsection{Complements finite over the base and coordinate projections}

If a complex is known to be constructible away from a closed subset
that is finite over the base, a proper direct image detects whether it
is constructible everywhere.

\begin{lemma}[Detection by a proper direct image]
\label{b:lem:geo-finite-boundary}
Let \(b:Y\to S\) be proper, let \(u:G\hookrightarrow Y\) be an open
immersion with reduced closed complement \(i:Z\hookrightarrow Y\), and
suppose that \(bi:Z\to S\) is finite.
\begin{enumerate}
\item If \(K\in D^+_{\mathrm{class}}(Y_{\proet},\Fell)\), \(u^*K\in\Dbc(G,\Fell)\), and
\(Rb_*K\in\Dbc(S,\Fell)\), then \(K\in\Dbc(Y,\Fell)\).
\item If \(c:K\to K'\) is a morphism in \(D^+_{\mathrm{class}}(Y_{\proet},\Fell)\) such
that \(u^*c\) is an isomorphism in \(D^+(G_{\proet},\Fell)\) and
\(Rb_*c\) is an isomorphism in \(D^+(S_{\proet},\Fell)\), then \(c\) is an
isomorphism.
\end{enumerate}
\end{lemma}

\begin{proof}
By bounded-below direct-image and constructible comparison, it is
enough to prove these assertions for the \'etale antecedents of $K$
and $K'$. The localization triangles are preserved by the same
comparison. All functors and stalk calculations in this proof are
on the \'etale sites, in the categories \(D^+(-_{\et},\Fell)\) and
\(\Mod(-_{\et},\Fell)\).

We first treat a finite morphism \(q:Z\to S\). For a sheaf \(F\) of
\(\Fell\)-modules on \(Z\), \(q_*F\) has stalks
\[
 (q_*F)_{\bar s}
       =\prod_{\bar z\in |Z_{\bar s}|}F_{\bar z}
                                                               \tag{G18}
\]
\stacks{03QP}, so
\(q_*\colon\Mod(Z_{\et},\Fell)\to\Mod(S_{\et},\Fell)\) is exact and
conservative. Moreover, if
\(q_*F\) is constructible then so is \(F\): the counit
\(q^*q_*F\to F\) is surjective, since on stalks it is a projection from
the finite product (G18), its source is constructible, and a quotient of
a constructible sheaf is constructible \stacks{09BH}. Since \(q_*\) is
exact, \(H^n(Rq_*M)=q_*H^n(M)\) for every complex \(M\). Hence, by
conservativity and the sheaf case in each degree, \(Rq_*M\in\Dbc\)
implies \(M\in\Dbc\).

Now apply the open-closed exact sequence \stacks{095L} termwise to
\(K\); it gives a triangle in \(D^+(Y_{\et},\Fell)\)
\[
 u_!u^*K\longrightarrow K\longrightarrow i_*i^*K
       \longrightarrow u_!u^*K[1].                       \tag{G19}
\]
(1) The first term is bounded constructible, so its image under
\(Rb_*=Rb_!\) is bounded constructible by \stacks{0GL0} and
\stacks{0G2A}. Applying \(Rb_*\) to (G19) shows that
\(R(bi)_*i^*K\) is bounded constructible. By the previous paragraph
applied to \(q=bi\), \(i^*K\in\Dbc(Z,\Fell)\). Since \(i_*\) is exact
and preserves constructibility, (G19) gives \(K\in\Dbc(Y,\Fell)\).

(2) Let \(M\) be a cone of \(c\). Then \(u^*M=0\), so (G19) gives
\(M=i_*i^*M\), and \(0=Rb_*M=(bi)_*i^*M\). By exactness and
conservativity of \((bi)_*\), applied to cohomology sheaves, \(M=0\).
\end{proof}

The hypotheses of Lemma~\ref{b:lem:geo-finite-boundary} are preserved
by every base change in the standing range, since \(bi\) stays finite and \(b\)
stays proper; this is how the lemma is applied to base-change maps.
The next lemma produces such a finite complement: projections to
coordinate lines lower the dimension of the generic fibre, and the
subset they leave uncovered becomes finite over a smaller base.

\begin{lemma}[Coordinate projections]
\label{b:lem:geo-coordinates}
Let \(S\) be an integral Noetherian affine scheme with function field
\(k(S)\), and let \(Y\to S\) be projective. Choose a finite affine open
cover \(Y=\bigcup_\alpha Y_\alpha\) and closed immersions
\(Y_\alpha\hookrightarrow\mathbf A^{m_\alpha}_S\) with
\(m_\alpha\geq1\), and write \(t_{\alpha r}:Y_\alpha\to\mathbf A^1_S\)
(\(1\le r\le m_\alpha\)) for the coordinate functions.
\begin{enumerate}
\item If \(X\subset Y\) is open and \(\dim X_{k(S)}\leq n\), then for
every \(\alpha,r\) the generic fibre of
\(t_{\alpha r}:X\cap Y_\alpha\to\mathbf A^1_S\) has dimension at most
\(n-1\).
\item Suppose given a dense open \(V_{\alpha r}\subset\mathbf A^1_S\) for
each \(\alpha,r\), and let
\[
 G=\bigcup_{\alpha,r}
     t_{\alpha r}^{-1}(V_{\alpha r})\subset Y.
\]
After shrinking \(S\), the reduced closed complement \(Z=Y\setminus G\)
is finite over \(S\).
\end{enumerate}
\end{lemma}

\begin{proof}
(1) The dimension of an integral scheme of finite type over a field is
the transcendence degree of its function field. Consider an irreducible
component of \(X_{k(S)}\cap Y_\alpha\) with function field \(L\). If
\(t_{\alpha r}\) is algebraic over \(k(S)\) on it, the component does
not dominate \(\mathbf A^1_{k(S)}\) and does not meet the generic fibre
of \(t_{\alpha r}\). Otherwise \(t_{\alpha r}\) is transcendental and
\[
 \operatorname{trdeg}_{k(S)(t_{\alpha r})}L
      =\operatorname{trdeg}_{k(S)} L-1\leq n-1.
\]

(2) Let \(t\) be the coordinate of \(\mathbf A^1_S\). The ideal of the
closed complement of \(V_{\alpha r}\) contains a nonzero polynomial
\(P_{\alpha r}\in\mathcal O(S)[t]\). After inverting the leading
coefficients of these finitely many polynomials we may assume that all
of them are monic. The function \(P_{\alpha r}(t_{\alpha r})\) vanishes
at every point of the reduced affine scheme \(Z\cap Y_\alpha\), hence is
zero in its coordinate ring. That ring is therefore a quotient of
\[
 \mathcal O(S)[t_{\alpha1},\ldots,t_{\alpha m_\alpha}]/
       (P_{\alpha1}(t_{\alpha1}),\ldots,
                         P_{\alpha m_\alpha}(t_{\alpha m_\alpha})),
                                                               \tag{G20}
\]
which is generated as an \(\mathcal O(S)\)-module by the monomials whose
exponent in each variable is less than the degree of the corresponding
polynomial. Thus each \(Z\cap Y_\alpha\) is finite over \(S\), and
\(Z\to S\) has finite fibres. To conclude that \(Z\to S\) is finite we
use that \(Z\) is closed in the projective \(Y\): then \(Z\to S\) is
proper with finite fibres, hence finite \stacks{02LS}.
\end{proof}

\subsubsection{Generic constructibility of direct image along an open immersion}

\begin{theorem}[Generic constructibility and base change for \(Rj_*\)]
\label{b:thm:geo-generic-open}
Let \(S\) be integral of finite type over the standing field \(k\), let
\(\ell\ne\operatorname{char}(k)\), and let \(j:X\hookrightarrow Y\) be
an open immersion of \(S\)-schemes of finite type. For a constructible
\(\Fell\)-sheaf \(F\) on \(X_{\proet}\), there is a nonempty open
\(S_0\subset S\) such that
\[
 Rj_{S_0,*}(F|_{X_{S_0}})\in\Dbc(Y_{S_0},\Fell),
                                                               \tag{G21}
\]
and such that for every finite-type $k$-scheme \(T\to S_0\), with base change
\(j_T:X_T\hookrightarrow Y_T\), the canonical base-change map
\[
 (Rj_{S_0,*}F)|_{Y_T}
          \longrightarrow Rj_{T,*}(F|_{X_T})
          \quad\text{in }D^+((Y_T)_{\proet},\Fell)             \tag{G22}
\]
is an isomorphism.
\end{theorem}

\begin{proof}
We prove the assertion for the \'etale antecedent of $F$ and then
transport it by Theorems~\ref{pe:direct-image-comparison} and
\ref{pe:constructible-comparison}. These compare the canonical
base-change maps, so the induction below, carried out with the functors
between the categories \(D^+(-_{\et},\Fell)\), proves (G21) and (G22)
on the pro-\'etale site.

We prove (G21) and (G22) together by induction on
\(n=\dim X_{k(S)}\), allowing at each stage all integral bases \(S\) of
finite type over \(k\). In the inductive step, coordinate projections
give the assertions away from a part of \(Y\) finite over \(S\); when
\(X\) is smooth and \(F\) locally constant,
Lemma~\ref{b:lem:geo-smooth-generic} controls the proper direct image
to \(S\), which by Lemma~\ref{b:lem:geo-finite-boundary} fills in the
finite part; and a general \(F\) differs from such a sheaf by something
supported in smaller dimension.

\emph{Preliminary reductions.} The assertions extend from open to
locally closed immersions with the same source: factor such an
immersion as an open immersion into its scheme-theoretic closure
followed by a closed immersion, and note that closed direct image is
exact, preserves constructibility, and commutes with arbitrary pullback
\stacks{0959}. We will use the induction hypothesis in this form.
Replacing \(Y\) by the scheme-theoretic closure of \(X\), we may assume
\(X\) dense in \(Y\). We may also shrink \(S\) to be affine, and remove
the closures of the images of the irreducible components of \(Y\) that
do not dominate \(S\): there are finitely many, and none contains the
generic point of \(S\).

The assertions are local on \(Y\), and since we only use finite affine
covers, the finitely many resulting opens of \(S\) can be intersected.

\emph{Base case.} If the generic fibre \(X_{k(S)}\) is empty, so is
\(Y_{k(S)}\), by density. Then \(Y\) becomes empty after shrinking
\(S\): on a finite affine cover, a coordinate algebra that vanishes
over \(k(S)\) has its unit killed by a nonzero element of
\(\mathcal O(S)\). If \(n=0\), the dense open \(X_{k(S)}\) of the
zero-dimensional Noetherian space \(Y_{k(S)}\) is all of it; the same
argument applied to \(Y\setminus X\) shows that \(j\) becomes an
isomorphism after shrinking \(S\).

Now let \(n\geq1\) and assume the theorem, including (G22), for
generic fibres of dimension less than \(n\).

\emph{Step 1: for \(Y\) projective over \(S\), both assertions hold on
an open \(G\subset Y\) whose complement is finite over \(S\).} Choose
charts and coordinates as in Lemma~\ref{b:lem:geo-coordinates}. Via
\(t_{\alpha r}\), regard \(X\cap Y_\alpha\hookrightarrow Y_\alpha\) as an
open immersion of schemes over the integral base \(\mathbf A^1_S\),
which is of finite type over \(k\). By part (1) of that lemma its
generic fibre has dimension less than \(n\), so the induction
hypothesis gives a dense open \(V_{\alpha r}\subset\mathbf A^1_S\) over
which both assertions hold. They also hold after a base
change in the standing range \(T\to S\): apply the hypothesis to the inverse image of
\(V_{\alpha r}\) in \(\mathbf A^1_T\). Hence on the finite union \(G\)
of the opens \(t_{\alpha r}^{-1}(V_{\alpha r})\), the complex
\(Rj_*F\) is bounded constructible (there are finitely many pieces, so
its cohomology lies in a common finite range of degrees), and (G22) is
an isomorphism. By part (2) of the lemma, after shrinking \(S\) the
reduced complement \(Z=Y\setminus G\) is finite over \(S\).

\emph{Step 2: the theorem holds in dimension \(n\) when \(X\) is
smooth over \(S\) and \(F\) is locally constant.} First reduce to
\(Y\) projective. Restrict to a member of a finite affine cover of
\(Y\), embed it in affine space over the affine base, and take its
closure in projective space. This affine member, and hence \(X\), is
open in that closure. Replacing the closure by the closure of \(X\)
makes \(X\) dense without changing its generic fibre, and a result
there restricts to the affine member. So we may assume \(b:Y\to S\)
projective.

Step~1 gives \((Rj_*F)|_G\in\Dbc\) with \(Z\to S\) finite; also
\(Rj_*F\in D^+(Y_{\et},\Fell)\). Since \(bj:X\to S\) is smooth and \(F\)
is locally constant, Lemma~\ref{b:lem:geo-smooth-generic}, after a
further shrinking of \(S\), shows that \(Rb_*Rj_*F=R(bj)_*F\) lies in
\(\Dbc(S,\Fell)\) and that its formation commutes with every base change
in the standing range. Part~(1) of Lemma~\ref{b:lem:geo-finite-boundary} gives
\(Rj_*F\in\Dbc(Y,\Fell)\).

For a finite-type $k$-scheme \(T\to S\), let \(c\) be the map (G22). It is an
isomorphism on \(G_T\) by Step~1. By proper base change for \(b\) and
the compatibility of base-change maps with composition \stacks{0E46},
\(Rb_{T,*}c\) is identified with the base-change map for \(R(bj)_*F\),
which is an isomorphism. Part~(2) of
Lemma~\ref{b:lem:geo-finite-boundary}, over \(T\), shows that \(c\) is
an isomorphism. (A priori the target of \(c\) is only bounded below;
the isomorphism shows it is bounded constructible.)

\emph{Step 3: the general case in dimension \(n\).} We first make the
smooth locus large. By \stacks{0550}, after shrinking \(S\) there is a
finite universal homeomorphism \(S'\to S\) with \(S'\) integral, such
that \((X_{S'})_{\mathrm{red}}\to S'\) is flat of finite presentation
with geometrically reduced generic fibre. That generic fibre has a dense
smooth open \stacks{056V}, so by the fibrewise criterion \stacks{01V9}
its generic points lie in the smooth locus of
\((X_{S'})_{\mathrm{red}}\to S'\). After shrinking the base once more to
remove non-dominant components, this smooth locus contains a dense open
of \((X_{S'})_{\mathrm{red}}\). Since \(F\) is constructible, there is a
smaller dense open \(V\) on which the pullback of \(F\) is locally
constant. Its complement contains no generic point of a component of
the generic fibre, so its generic fibre has dimension less than \(n\).
Passing to reductions and along finite universal homeomorphisms does
not change \'etale topoi \stacks{03SI}. We may therefore work over
\(S'\), with all schemes replaced by their reductions, and descend to
\(S\) at the end.

Let \(v:V\hookrightarrow X\) be the open immersion and
\(i:X\setminus V\hookrightarrow X\) its reduced closed complement. By
Step~2, after shrinking the base, both \(v\) and \(jv\) have bounded
constructible direct images of \(v^*F\), and their base-change maps are
isomorphisms after every base change in the standing range. Complete the unit of
\(v^*\dashv Rv_*\) to a triangle in \(D^+(X_{\et},\Fell)\)
\[
 F\longrightarrow Rv_*v^*F\longrightarrow D
       \longrightarrow F[1].                            \tag{G23}
\]
The first two terms are bounded constructible, hence so is \(D\), by
\stacks{09BH} and the long exact cohomology sequence. The unit is an
isomorphism on \(V\), so \(v^*D=0\) and (G19) gives
\[
 D=i_*i^*D.                                              \tag{G24}
\]
The locally closed immersion \(ji:X\setminus V\to Y\) has generic fibre
of dimension less than \(n\). Applying the induction hypothesis to each
of the finitely many constructible cohomology sheaves of \(i^*D\), and
using their truncation triangles, we get
\[
 R(ji)_*i^*D=Rj_*D\in\Dbc(Y,\Fell),
\]
with base-change maps that are isomorphisms. Now apply \(Rj_*\) to
(G23). The middle term is \(R(jv)_*v^*F\), treated by Step~2, and the
third term was just treated. Hence \(Rj_*F\) is bounded constructible.

For base change along \(T\), the base-change map for \(v\) identifies
the pullback of the first arrow of (G23) with the unit for \(v_T\).
Since cones are not functorial, we fix the triangle for this unit to be
the pullback of (G23), transported through that identification. The
base-change map is a natural transformation
\((Rj_*-)|_{Y_T}\to Rj_{T,*}(-|_{X_T})\) of exact functors
\(D^+(X_{\et},\Fell)\to D^+((Y_T)_{\et},\Fell)\), so it
gives a morphism from the pullback of \(Rj_*\) of (G23) to \(Rj_{T,*}\)
of this triangle. It is an isomorphism on the second and third terms,
hence on the first, which is (G22).

\emph{Descent to \(S\).} A finite universal homeomorphism is a
homeomorphism of underlying spaces. If the argument over \(S'\) gives a
nonempty open \(U'\subset S'\), remove from \(S\) the closed image of
\(S'\setminus U'\). The remaining open \(U\) is nonempty and its inverse
image lies in \(U'\). Over \(U\), pullback along \(S'_U\to U\) is an
equivalence of \'etale topoi \stacks{03SI}; it preserves constructible
sheaves in both directions and commutes with direct images and their
derived functors. So it detects boundedness and constructibility, and
(G21) holds over \(U\). For a finite-type $k$-scheme \(T\to U\), the map
\(T\times_US'_U\to T\) is again a universal homeomorphism, and
compatibility of base-change maps with composition identifies the
pullback of (G22) to \(T\times_US'_U\) with the base-change map already
shown to be an isomorphism over \(S'\). Since the equivalence is
conservative, (G22) is an isomorphism over \(T\). This completes the
induction.
\end{proof}

Over a field the only nonempty open of \(\operatorname{Spec}k\) is
itself, so the theorem gives a statement with no shrinking.

\begin{corollary}[Open direct image over a field]
\label{b:cor:geo-open-finiteness}
\label{b:thm:finite-open}
Let \(j:X\hookrightarrow Y\) be an open immersion of schemes of finite
type over the standing field \(k\), with \(\ell\ne\operatorname{char}(k)\), and let
\(\Lambda\) be a finite \(\ell\)-primary ring. Then
\[
 Rj_*:\Dbc(X,\Lambda)\longrightarrow\Dbc(Y,\Lambda),
                                                               \tag{G25}
\]
and for every finite-type $k$-scheme \(T\to\operatorname{Spec}k\) and every
\(K\in\Dbc(X,\Lambda)\), the base-change map
\((Rj_*K)|_{Y_T}\to Rj_{T,*}(K|_{X_T})\) in \(D((Y_T)_{\proet},\Lambda)\)
is an isomorphism.
\end{corollary}

\begin{proof}
We argue on the \'etale site; bounded constructible and direct-image
comparison identify the result and its base-change maps with the
asserted pro-\'etale ones.

For a constructible \(\Fell\)-sheaf this is
Theorem~\ref{b:thm:geo-generic-open} with \(S=\operatorname{Spec}k\). A
constructible sheaf killed by \(\ell^r\) has a finite filtration by the
images of multiplication by powers of \(\ell\), with constructible
\(\Fell\)-sheaves as successive quotients; the long exact cohomology
sequences give (G25), and morphisms between the resulting triangles give
the base-change assertion. For a finite \(\ell\)-primary ring
\(\Lambda\), apply this to the underlying
\(\mathbf Z/\ell^r\mathbf Z\)-sheaf, carrying the \(\Lambda\)-action
along all maps. This is possible because \(Rj_*\colon D^+(X_{\et},\Lambda)\to
D^+(Y_{\et},\Lambda)\) commutes with the forgetful functors to
\(D^+(-_{\et},\mathbf Z/\ell^r\mathbf Z)\) \stacks{0D6J}, and because a \(\Lambda\)-sheaf
is constructible if and only if its underlying sheaf is: on a finite
stratification trivializing the underlying sheaf, the finitely many
endomorphisms giving the action become constant after a further finite
refinement. Finally, truncation triangles extend both conclusions from
sheaves to bounded constructible complexes.
\end{proof}

Over a field, $Rj_*$ also has bounded amplitude on classical torsion sheaves.

\begin{corollary}[Finite cohomological dimension of open direct image]
\label{b:cor:geo-open-finiteness-b}
In the situation of Corollary~\ref{b:cor:geo-open-finiteness},
$\operatorname{cd}_\ell(j)<\infty$ (Definition~\ref{n:def-cd}); that is,
$Rj_*\colon D^+(X_{\proet},\Lambda)\to D^+(Y_{\proet},\Lambda)$ has
finite cohomological amplitude on classical $\ell$-power torsion
sheaves.
\end{corollary}

\begin{proof}
For \(Rj_*\colon D^+(X_{\et},\Lambda)\to D^+(Y_{\et},\Lambda)\) this
is the bound of \stacks{0F10} recalled at the start of this subsection,
which holds over every field \(k\).
Proposition~\ref{pe:unbounded-classical} transfers it to classical
pro-\'etale sheaves.
\end{proof}


\subsection{Finite coefficients}
\label{ctf:finite-geometric-restriction}

Let $A$ be a finite commutative $\ell$-primary ring. Schemes and
morphisms of finite type over $k$ have finite $\ell$-cohomological
dimension $\operatorname{cd}_\ell$ (Definition~\ref{n:def-cd}), which is
what Theorem~\ref{ctf:unified-brown} requires, also for the smooth base
changes in its finiteness argument. We show that its right adjoint
$f^!$ preserves $\Dctf$.

Finite Tor amplitude is a condition on complexes: the cohomology
modules of an object of $\Dctf$ may have infinite projective dimension. The
Tor bounds below therefore come from models. By \stacks{0F4M}, every
object of $\Dctf(X,A)$ is represented by a bounded complex of
constructible flat $A$-sheaves, and on the strata of a finite
stratification such a model is locally a bounded complex of finite
projective $A$-modules.

\begin{proposition}[Open direct image preserves finite Tor amplitude]
\label{ctf:finite-open}
Let $j:U\hookrightarrow X$ be an open immersion. Then
$Rj_*\colon D(U_{\proet},A)\to D(X_{\proet},A)$ restricts to
\[
 Rj_*:\Dctf(U,A)\longrightarrow\Dctf(X,A).
\]
More precisely, if $K$ has Tor amplitude in $[a,b]$ and $d\ge
\operatorname{cd}_\ell(j)$, then $Rj_*K$ has Tor amplitude in $[a,b+d]$.
\end{proposition}

\begin{proof}
Bounded constructibility is Theorem~\ref{b:thm:finite-open}, and
$\operatorname{cd}_\ell(j)$ is finite by
Corollary~\ref{b:cor:geo-open-finiteness-b}. For the Tor amplitude we work on the \'etale site, where geometric
points are conservative. Write $K=\nu_U^*K_{\et}$ by
Corollary~\ref{pe:finite-tor-comparison}; the pro-\'etale pullback of
$Rj_{\et,*}K_{\et}\in D(X_{\et},A)$ is $Rj_*K$, computed by
$Rj_*\colon D(U_{\proet},A)\to D(X_{\proet},A)$, by
Proposition~\ref{pe:unbounded-classical}.

The main step is the scalar projection formula: for every $A$-module
$M$, of arbitrary projective dimension and viewed as a constant sheaf,
the projection morphism is an isomorphism in $D(X_{\proet},A)$
\begin{equation}\label{ctf:scalar-open-projection}
 (Rj_*K)\otimes_A^{\mathbf L}\underline M_X
   \xrightarrow{\ \sim\ }
 Rj_*\bigl(K\otimes_A^{\mathbf L}\underline M_U\bigr).
\end{equation}
Its \'etale version, the projection morphism for
$Rj_{\et,*}\colon D(U_{\et},A)\to D(X_{\et},A)$, follows by applying the projection formula of \stacks{0F12} to every
quasi-compact \'etale $V\to X$ and to $V\times_XU$, both of finite
$\ell$-cohomological dimension; the resulting isomorphisms identify the
derived sections of the two sides, compatibly with the projection
morphism. Since inverse image is exact and symmetric monoidal and
commutes with unbounded direct image, the pullback of the \'etale
isomorphism is \eqref{ctf:scalar-open-projection}.

The complex $K\otimes_A^{\mathbf L}\underline M_U$ has cohomology in
$[a,b]$, so the \'etale version of the right side of
\eqref{ctf:scalar-open-projection} lies in $[a,b+d]$. Taking geometric
stalks gives this Tor bound for every stalk of $Rj_{\et,*}K_{\et}$.
On the \'etale site, testing against constant $A$-modules suffices:
tensor product and cohomology commute with the conservative geometric
stalks, and the stalk of any \'etale module sheaf is an $A$-module.
Hence $Rj_{\et,*}K_{\et}$ has Tor amplitude in $[a,b+d]$.

To pass to $X_{\proet}$, pull back a bounded flat constructible model
of $Rj_{\et,*}K_{\et}$ in degrees $[a,b+d]$ (\stacks{0F4M}). This gives
the same interval against every pro-\'etale module sheaf, including
nonclassical ones, for which geometric stalks are not available as a
test.
\end{proof}

\begin{theorem}[$f^!$ with finite coefficients]
\label{ctf:finite-restriction}
For every separated finite-type morphism $f:X\to Y$, the right adjoint
$f^!\colon\Loc(\Dctf(Y,A))\to\Loc(\Dctf(X,A))$ of
Theorem~\ref{ctf:unified-brown} satisfies
\[
 f^!\bigl(\Dctf(Y,A)\bigr)\subseteq\Dctf(X,A).
\]
Consequently its adjunction restricts to
\[
 Rf_!:\Dctf(X,A)\ \rightleftarrows\ \Dctf(Y,A):f^!.
\]
For smooth $f$ of pure relative dimension $d$, the functor
$f^!\colon\Dctf(Y,A)\to\Dctf(X,A)$ is locally on $X$ isomorphic to
$f^*(-)(d)[2d]$.
\end{theorem}

The canonical global isomorphism for smooth $f$ is constructed in
Section~\ref{sec:smooth-purity}.

\begin{proof}
Locally on $X$ and $Y$, $f$ is a closed immersion followed by the
projection from an affine space, so we treat these two cases and then
glue.

Let $i:Z\hookrightarrow X$ be a closed immersion with open
complement $j:U\hookrightarrow X$, and let $R\underline\Gamma_Z$ be as in
Lemma~\ref{b:lem:formal-closed-support}, so that for
$K\in\Dctf(X,A)$ there is a triangle in $D(X_{\proet},A)$
\[
 i_*R\underline\Gamma_ZK\longrightarrow K\longrightarrow Rj_*j^*K
       \longrightarrow i_*R\underline\Gamma_ZK[1].
\]
By Proposition~\ref{ctf:finite-open} the two right-hand terms lie in
$\Dctf$, hence so does $R\underline\Gamma_ZK$; explicitly, a Tor interval
$[a,b]$ for $K$ gives the interval
$[a,b+\operatorname{cd}_\ell(j)+1]$ for $R\underline\Gamma_ZK$. Since
$R\underline\Gamma_Z\colon D(X_{\proet},A)\to D(Z_{\proet},A)$ is right
adjoint to $i_*$ (Lemma~\ref{b:lem:formal-closed-support}(1)) and
$\Loc(\Dctf(-,A))$ is a full subcategory of $D((-)_{\proet},A)$,
uniqueness of right adjoints identifies $R\underline\Gamma_ZK$ with
$i^!K$ in $\Loc(\Dctf(Z,A))$.

For the projection $a:\mathbf A^d_Y\to Y$, the projective-line
adjunction $Rq_*\dashv q^*(1)[2]$ of
Proposition~\ref{b:prop:geo-p1-adjunction} and Lemma~\ref{b:lem:geo-affine-local} give
$a^!=a^*(d)[2d]$ as functors
$\Loc(\Dctf(Y,A))\to\Loc(\Dctf(\mathbf A^d_Y,A))$. This functor preserves $\Dctf$, because pullback
preserves a bounded constructible flat model and the Tate twist is
locally free of rank one. More generally,
Proposition~\ref{b:prop:geo-smooth-tensor} proves the same for smooth
morphisms, whose dualizing complexes are locally an invertible sheaf
shifted by $2d$.

Finally, after restricting to affine opens in $X$ and $Y$, a
finite-type morphism factors as a closed immersion into an affine
space over $Y$ followed by its projection. Composition and open
locality of the adjunction identify the restriction of $f^!K$ with the
composite of the two functors just considered, so $f^!K$ is bounded
constructible of finite Tor amplitude locally on $X$. A finite affine
cover of the quasi-compact scheme $X$ gives one cohomological interval
and one Tor interval valid everywhere. Since $Rf_!$ preserves $\Dctf$ by
Proposition~\ref{ctf:finite-proper}, both sides of the adjunction lie
in $\Dctf$, and its Hom bijections restrict.
\end{proof}

\begin{corollary}[The six operations on $\Dctf(-,A)$]
\label{ctf:finite-operations}
For separated finite-type morphisms of the standing schemes, the
operations $f^*$, $Rf_*$, $Rf_!$, and $f^!$, together with
$\otimes_A^{\mathbf L}$ and $R\mathcal Hom_A$, preserve bounded
constructible complexes of finite Tor amplitude. Here $f^*$, $Rf_*$,
$\otimes_A^{\mathbf L}$ and $R\mathcal Hom_A$ are the functors on
$D((-)_{\proet},A)$, and $Rf_!$, $f^!$ are the functors on
$\Loc(\Dctf(-,A))$ of Theorem~\ref{ctf:unified-brown}.
\end{corollary}

\begin{proof}
Pullback, tensor product and extension by zero preserve bounded
constructible flat models. If $p$ is proper, proper finiteness gives
constructibility of $Rp_*K$, and the scalar projection formula together
with the finite cohomological dimension of $p$ gives its finite Tor
amplitude as in Proposition~\ref{ctf:finite-open}. For a
compactification $f=pj$ this proves the assertion for $Rf_!=Rp_*j_!$,
and $Rf_*=Rp_*Rj_*$ gives it for $Rf_*$. The case of $f^!$ is
Theorem~\ref{ctf:finite-restriction}.

For internal Hom we argue by d\'evissage along the strata on which the
first argument is perfect. Choose a finite stratification on which a
bounded constructible flat model of the first argument is locally a
finite complex of finite projective modules. On such a stratum a
complex $P$ is dualizable, so
\[
 R\mathcal Hom_A(P,L)=P^\vee\otimes_A^{\mathbf L}L
\]
lies in $\Dctf$ whenever $L$ does. To pass from the strata to $X$, use open
and closed localization triangles in the first argument together with
the identities
\[
\begin{aligned}
 R\mathcal Hom_A(j_!P,L)
   &\simeq Rj_*R\mathcal Hom_A(P,j^*L),\\
 R\mathcal Hom_A(i_*P,L)
   &\simeq i_*R\mathcal Hom_A(P,i^!L)
\end{aligned}
\]
in $D(X_{\proet},A)$, for an open immersion $j$ and a closed immersion $i$. They follow from
tensor-Hom adjunction and the adjunctions for $j_!$ and $i_*$, and the
preceding results keep both right sides in $\Dctf$. Induction on the
number of strata gives the result.
\end{proof}


\subsection{Open direct image and the projective line with \texorpdfstring{$\ell$}{l}-adic coefficients}
\label{b:sec:late-coefficients}

In this subsection and \S\ref{b:sec:restriction}, $\Lambda=E$ or
$\Qbar$, so $\Dctf(-,\Lambda)=\Dbc(-,\Lambda)$. Sheaves live on the
pro-\'etale site, and $f^!$ denotes the right adjoint on
$\Loc(\Dctf(-,\Lambda))=\Loc(\Dbc(-,\Lambda))$ from
Theorem~\ref{ctf:unified-brown}. Some intermediate results are stated
also for $\mathcal O_E$; they are used again in Section~\ref{i:heart}.
The two torsion results behind $f^!$ are the finiteness of $Rj_*$ for
an open immersion $j$ and the projective-line adjunction
(Proposition~\ref{b:prop:geo-p1-adjunction}). The first transfers to
$\ell$-adic coefficients by reduction modulo $\pi_E$, as proper direct
image did in Proposition~\ref{b:prop:proper-coefficients}. The second
needs more care. It is built from the cycle class of the diagonal,
which we must produce as a single morphism with
$\mathcal O_E$-coefficients satisfying the integration identities, and
its unit and counit live on $\Loc(\Dbc(-,\Lambda))$, where we need the
projection formula and proper base change.

\begin{proposition}[Direct image along an open immersion]
\label{b:prop:open-coefficients}
Let $j:U\hookrightarrow X$ be a quasi-compact open immersion. Then
$Rj_*\colon D(U_{\proet},\Lambda)\to D(X_{\proet},\Lambda)$ maps $\Dbc(U,\Lambda)$ into $\Dbc(X,\Lambda)$ for
$\Lambda=\mathcal O_E$, $E$ and $\Qbar$.
\end{proposition}
\begin{proof}
We follow Proposition~\ref{b:prop:proper-coefficients}, with
Theorem~\ref{b:thm:finite-open} in place of proper finiteness.

\emph{Integral coefficients.} Let $K\in\Dbc(U,\mathcal O_E)$. Then
$Rj_*K$ is derived complete by \stacks{099L}. Reduction modulo $\pi_E$
commutes with $Rj_*$, because $K\otimes^{\mathbf L}\mathcal O_{E,U}/\pi_E$
is the cone of multiplication by $\pi_E$ on $K$ and $Rj_*$ is exact; this
is the same observation as in \eqref{b:eq:proper-reduction}. The
reduction of $K$ is pulled back from a bounded constructible complex on
$U_{\et}$. By \stacks{09A6} its pro-\'etale direct image in
$D(X_{\proet},\mathcal O_E/\pi_E)$ is the pullback of its \'etale direct
image in $D(X_{\et},\mathcal O_E/\pi_E)$, which is bounded constructible by
Theorem~\ref{b:thm:finite-open}, since the coefficient ring
$\mathcal O_E/\pi_E$ is a finite field of characteristic $\ell$. Thus
$Rj_*K$ is derived complete with bounded constructible reduction, and
Lemma~\ref{b:lem:integral-models}(2) gives $Rj_*K\in\Dbc(X,\mathcal O_E)$.

\emph{Rational coefficients.} Inverting $\ell$ is the colimit of
$K\xrightarrow{\ell}K\xrightarrow{\ell}\cdots$, and extension of scalars
to $\Qbar$ is a filtered colimit over finite extensions of $E$. Both are
colimits of systems with a common lower bound, so
Lemma~\ref{b:lem:bounded-continuity} shows that $Rj_*$ commutes with
them. By the integral case, $Rj_*$ therefore maps the generators of
Lemma~\ref{b:lem:integral-models}(4) into $\Dbc(X,E)$, respectively
$\Dbc(X,\Qbar)$. Since these categories are triangulated
(Corollary~\ref{b:lem:coefficient-heart-b}), $Rj_*$ preserves
$\Dbc(-,E)$ and $\Dbc(-,\Qbar)$.
\end{proof}

We turn to the projective line. With torsion coefficients,
Proposition~\ref{b:prop:geo-p1-adjunction} builds the adjunction
$Rq_*\dashv q^*(1)[2]$ between $D_{\mathrm{class}}((\mathbf P^1_Y)_{\proet},\Lambda)$
and $D_{\mathrm{class}}(Y_{\proet},\Lambda)$ from two maps, the cycle class
$\delta:\Delta_*\Lambda\to\Lambda(1)[2]$ of the diagonal of
$\mathbf P^1_Y\times_Y\mathbf P^1_Y$ and the trace
$\operatorname{Tr}_q:Rq_*\Lambda(1)[2]\to\Lambda$, and from two
identities between them. For $\ell$-adic coefficients we need the same
data over $\mathcal O_E$. The classes with coefficients
$\mathbf Z/\ell^n$ form a family of morphisms in the separate derived
categories $D((\mathbf P^1_Y\times_Y\mathbf P^1_Y)_{\proet},\mathbf Z/\ell^n)$, compatible under reduction. Such a
family does not in general determine a morphism after taking the derived
inverse limit, because morphisms in a derived category are not computed
levelwise. We therefore construct the class once, as a morphism in the
derived category of inverse systems of sheaves on
$(\mathbf P^1_Y\times_Y\mathbf P^1_Y)_{\proet}$, and only then apply
$R\varprojlim$.

\begin{proposition}[The $\ell$-adic cycle class of the diagonal]
\label{b:prop:diagonal-tower}
Let $q:\mathbf P^1_Y\to Y$ be the projection and
$\Delta:\mathbf P^1_Y\hookrightarrow\mathbf P^1_Y\times_Y\mathbf P^1_Y$
the diagonal. There is a morphism
\[
 \delta_{\mathcal O_E}:\Delta_*\mathcal O_{E,\mathbf P^1_Y}
              \longrightarrow
              \mathcal O_{E,\mathbf P^1_Y\times_Y\mathbf P^1_Y}(1)[2]
\]
in $D((\mathbf P^1_Y\times_Y\mathbf P^1_Y)_{\proet},\mathcal O_E)$ whose
reductions are the cycle classes of $\Delta$ with torsion coefficients
(Lemma~\ref{b:lem:geo-cartier}), and a trace
$\operatorname{Tr}_q:Rq_*\mathcal O_{E,\mathbf P^1_Y}(1)[2]\to\mathcal O_{E,Y}$
in $D(Y_{\proet},\mathcal O_E)$ lifting the torsion traces of Lemma~\ref{b:lem:geo-p1-cohomology}. The two
integration identities of Proposition~\ref{b:prop:geo-p1-adjunction},
\[
 \operatorname{Tr}_{p_i}\circ Rp_{i*}(\delta_{\mathcal O_E})
   =1_{\mathcal O_{E,\mathbf P^1_Y}}\qquad(i=1,2),
\]
hold in $D((\mathbf P^1_Y)_{\proet},\mathcal O_E)$ for these morphisms,
with $Rp_{i*}\colon D((\mathbf P^1_Y\times_Y\mathbf P^1_Y)_{\proet},\mathcal O_E)\to D((\mathbf P^1_Y)_{\proet},\mathcal O_E)$, where $p_1,p_2$ are the two projections of
$\mathbf P^1_Y\times_Y\mathbf P^1_Y$, each a projective line over
$\mathbf P^1_Y$. Extending scalars gives $E$- and $\Qbar$-linear
morphisms $\delta_\Lambda$ and $\operatorname{Tr}_q$ with the same
identities. All of these are compatible with base change along
separated morphisms of finite type $Y'\to Y$ between schemes as in
Section~\ref{n:scope}.
\end{proposition}
\begin{proof}
All products are over $Y$. We build the torsion cycle classes for all
$n$ at once as one morphism of inverse systems, apply $R\varprojlim$,
and prove the identities already for inverse systems, where they can be
checked levelwise because the objects involved sit in degree $0$.

\emph{The extension defining the cycle class.} The diagonal is an
effective Cartier divisor, with equation $x_0y_1-x_1y_0$. The proof of
Lemma~\ref{b:lem:geo-cartier} gives the extension of \'etale sheaves
\[
 0\longrightarrow\mathbf G_m\longrightarrow \mathcal M_\Delta
       \longrightarrow\Delta_*\mathbf Z\longrightarrow0,
\]
where $\mathcal M_\Delta$ is generated by the units and the inverse of a
local equation $t$ of $\Delta$, and the map to $\Delta_*\mathbf Z$ sends
$t^{-1}$ to $+1$. We keep this sign convention; with it the class of a
section has degree $+1$, and both integration identities hold with
sign $+1$. Pull the extension back to the pro-\'etale site, keeping the
names $\mathbf G_m$ and $\mathcal M_\Delta$. As a morphism in
$D((\mathbf P^1_Y\times_Y\mathbf P^1_Y)_{\proet},\mathbf Z)$, the
extension is the diagram
\[
 \Delta_*\mathbf Z\ \xleftarrow{\ \sim\ }\
 [\mathbf G_m\longrightarrow \mathcal M_\Delta]\ \longrightarrow \mathbf G_m[1],
\]
with the two-term complex in degrees $-1$ and $0$. The left arrow is a
quasi-isomorphism by exactness of the extension (on the pro-\'etale
site, $\Delta_*$ of the pulled-back sheaf is the pullback of
$\Delta_*\mathbf Z$ by closed base change).

\emph{Functorial derived completion.} We reduce modulo all $\ell^n$ at
once. On $(\mathbf P^1_Y\times_Y\mathbf P^1_Y)_{\proet}$, the derived
tensor product with the inverse system of quotient rings,
\[
 K\longmapsto K\otimes^{\mathbf L}_{\mathbf Z}
                  (\underline{\mathbf Z/\ell^n})_{n\geq1},
\]
is a functor from $D((\mathbf P^1_Y\times_Y\mathbf P^1_Y)_{\proet},\mathbf Z)$ to the
derived category of inverse systems of $\underline{\mathbf Z_\ell}$-modules
on the same site.
It is computed by one K-flat resolution of $K$ tensored termwise with
that system, and $\mathbf Z_\ell$ acts on the $n$-th term through
$\mathbf Z/\ell^n$. Evaluation of systems is exact and detects
quasi-isomorphisms, so this is a derived functor, and it
sends the boundary morphism of the Cartier extension to a single
morphism of inverse systems. Composed with $R\varprojlim$, which takes values in
$D((\mathbf P^1_Y\times_Y\mathbf P^1_Y)_{\proet},\underline{\mathbf Z_\ell})$,
it is derived $\ell$-completion $K\mapsto\widehat K$ (\stacks{099M}), so the
limit morphism is the completion of an existing morphism.

Its source is the system $\Delta_*\mathbf Z/\ell^\bullet$: the formula
$(\Delta_*\mathbf Z)\otimes F\simeq\Delta_*\Delta^{-1}F$ holds because
both sides agree after restriction to the diagonal and to its open
complement, and these restrictions are jointly conservative
(\stacks{09AD}); the right side is exact in $F$. Its target is
computed by the Kummer sequence,
$\mathbf G_m\otimes^{\mathbf L}_{\mathbf Z}\mathbf Z/\ell^n\simeq\mu_{\ell^n}[1]$.
To identify the transition maps, use the resolution
$[\mathbf Z\xrightarrow{\ell^n}\mathbf Z]$ of $\mathbf Z/\ell^n$, with
transition $\ell$ in degree $-1$ and the identity in degree $0$; it
shows that the transition maps of the target are
$\zeta\mapsto\zeta^\ell$ on $\mu_{\ell^n}$. The morphism is
the image of the boundary morphism, and canonical truncations
identify its target with the system $\mathbf Z/\ell^\bullet(1)[2]$ as
linear objects. We have constructed one morphism
\[
 \Delta_*\mathbf Z/\ell^\bullet\longrightarrow \mathbf Z/\ell^\bullet(1)[2]
\]
in the derived category of inverse systems.

\emph{Passage to the limit.} Apply $R\varprojlim$ to this morphism.
Since the pro-\'etale site is replete and products of sheaves are exact
(Propositions~\ref{pe:replete} and~\ref{pe:products}),
\[
 R\varprojlim_n\Delta_*\mathbf Z/\ell^n=\Delta_*\Zl,\qquad
 R\varprojlim_n \mathbf Z/\ell^n(1)=\Zl(1).
\]
For the coefficient systems this can be seen directly: on a weakly
contractible affine, sections are exact and the transition maps are
surjective, so the Milnor complex computing $R\varprojlim$ has no
cohomology in degree one. The functor $\Delta_*$ is exact and preserves
limits. More generally, derived direct image commutes with
$R\varprojlim$ (\stacks{099L}): both are right adjoints of exact functors
(inverse image and the constant-system functor), so both preserve
K-injectives and their derived functors commute. Evaluation of an
inverse system at one level also preserves injectives, since its left
adjoint is exact, so derived direct image of a system is computed
levelwise.

The limit morphism is linear over the constant sheaf
$\underline{\mathbf Z_\ell}$, and we need it over the completed
coefficient sheaf $\mathbf Z_{\ell}$ of $\mathbf P^1_Y\times_Y\mathbf P^1_Y$.
Its source and target are derived $\ell$-complete, and
the map from $\underline{\mathbf Z_\ell}$ to it is flat and an
isomorphism modulo every $\ell^n$. Hence restriction of scalars is an
equivalence on derived complete objects (\stacks{099I}, as in
Lemma~\ref{b:lem:integral-models}(1)), and the morphism is, in a unique
way, a $\Zl$-linear morphism $\Delta_*\Zl\to\Zl(1)[2]$ over the
completed coefficient sheaf.

Here the Tate twist $(1)$ is the tensor product, termwise and over the
completed coefficient ring, with the invertible sheaf $\Zl(1)$; the
twist $(-1)$ uses its dual. The sheaf $\Zl(1)$ is free of rank one on
the pro-\'etale cover given by the system of primitive $\ell^n$-th roots
of unity, and its reduction modulo $\ell^n$ is $\mu_{\ell^n}$ with the
transition maps computed above. Hence the twist is exact, preserves
products, and commutes with $R\varprojlim$.

\emph{The trace and the integration identities.} Form the decomposition
of Lemma~\ref{b:lem:geo-p1-cohomology} in the derived category of
inverse systems on $Y_{\proet}$, from the unit and the class of a section:
\[
 \mathbf Z/\ell^\bullet\oplus \mathbf Z/\ell^\bullet(-1)[-2]
   \xrightarrow{\ \sim\ }Rq_*\mathbf Z/\ell^\bullet .
\]
It is an isomorphism because its cone is acyclic at every level, and
cohomology of systems is computed levelwise. Twisting and projecting to
the second summand through its inverse defines the trace on systems, and
$R\varprojlim$ gives $\operatorname{Tr}_q$. Derived direct image
preserves completeness, so its source and target are complete, and the
same equivalence makes it linear over the completed coefficient sheaf
on $Y$.

Each composite $\operatorname{Tr}_{p_i}\circ Rp_{i*}(\delta)$ is now an
endomorphism of the system $\mathbf Z/\ell^\bullet$ in the derived
category of inverse systems on $(\mathbf P^1_Y)_{\proet}$, which is concentrated in degree $0$. Endomorphisms of an
object in degree $0$ in a derived category are the endomorphisms in the
abelian category, here endomorphisms of inverse systems of sheaves,
that is, compatible families of endomorphisms of the levels
$\mathbf Z/\ell^n$. At each level the endomorphism is the identity by
Proposition~\ref{b:prop:geo-p1-adjunction}. Hence the identities hold in
the derived category of inverse systems, and applying $R\varprojlim$
gives them over $\Zl$. For morphisms between arbitrary objects of the
derived category of inverse systems, equality at every level would not
imply equality; this is why the identities are checked on objects in
degree $0$.

\emph{Coefficients and base change.} The ring $\mathcal O_E$ is finite
free over $\Zl$, and extension of scalars along it commutes with products
and $R\varprojlim$. Since $\ell^n\mathcal O_E=\pi_E^{en}\mathcal O_E$ up
to a unit, where $e$ is the ramification index, the systems of
reductions modulo $\ell^n$ and modulo $\pi_E^n$ are cofinal. This gives
$\delta_{\mathcal O_E}$, its trace and its identities. Inverting $\ell$
and the filtered colimit over finite extensions of $E$ give
$\delta_E$, $\delta_{\Qbar}$ and their traces; direct image commutes
with these colimits of bounded below objects
(Lemma~\ref{b:lem:bounded-continuity}), so the identities are preserved.

For a morphism $g:Y'\to Y$, the diagonal and the sections remain Cartier
divisors after base change, and the Cartier extension and the Kummer
maps are natural under pullback. The adjunction between the
constant-system functor and $R\varprojlim$ gives, for an inverse
system $K_\bullet$ of sheaves on $T_{\proet}$ with $T$ one of
$Y$, $\mathbf P^1_Y$, $\mathbf P^1_Y\times_Y\mathbf P^1_Y$, and with $g$
also denoting its base change $T\times_YY'\to T$, a natural comparison
\[
 g^{-1}R\varprojlim K_\bullet\longrightarrow
                          R\varprojlim g^{-1}K_\bullet,
\]
in $D((T\times_YY')_{\proet},\underline{\mathbf Z_\ell})$, which is an isomorphism for the systems used here, by
Lemma~\ref{b:lem:coefficient-pullback}, closed base change, and the
compatibility of the Tate twist with pullback. Naturality of this
comparison gives the compatibility of $\delta_\Lambda$ with base change,
and the decomposition of $Rq_*$ gives that of the trace.
\end{proof}

The triangle identities for the projective line use the projection
formula and proper base change on all of $\Loc(\Dbc(-,\Lambda))$. These
extend from $\Dbc$ by the usual localizing argument.

\begin{proposition}[Projection formula and proper base change with rational coefficients]
\label{b:prop:pf-bc-generated}
For $\Lambda=E$ or $\Qbar$, the category $\Loc(\Dctf(X,\Lambda))$ is closed
under $\otimes^{\mathbf L}$. For a proper $p:X\to Y$, the projection
formula map $(Rp_*A)\otimes^{\mathbf L}B\to Rp_*(A\otimes^{\mathbf L}p^*B)$
is an isomorphism in $\Loc(\Dctf(Y,\Lambda))$ for
$A\in\Loc(\Dctf(X,\Lambda))$ and $B\in\Loc(\Dctf(Y,\Lambda))$, and for a
cartesian square as in Proposition~\ref{b:prop:proper-base-change-unbounded}
the base change map $g^*Rp_*\to Rp'_*g'^*$ is an isomorphism of
functors $\Loc(\Dctf(X,\Lambda))\to\Loc(\Dctf(Y',\Lambda))$.
\end{proposition}
\begin{proof}
Each claim is proved by fixing all but one variable and observing that
the objects satisfying it form a localizing subcategory containing
$\Dbc$. Recall from \S\ref{b:sec:brown-construction} that $g^*$ and
$Rp_*$ for proper $p$ restrict to exact functors on
$\Loc(\Dbc(-,\Lambda))$ preserving direct sums.

\emph{Tensor products.} Fix $B\in\Dbc(X,\Lambda)$. Since
$-\otimes^{\mathbf L}B$ is exact and preserves direct sums, the objects
$A\in\Loc(\Dbc(X,\Lambda))$ with $A\otimes^{\mathbf L}B\in\Loc(\Dbc(X,\Lambda))$
form a localizing subcategory. It contains $\Dbc(X,\Lambda)$, because
$\Dbc(X,\Lambda)$ is closed under $\otimes^{\mathbf L}$: by
Corollary~\ref{b:lem:coefficient-heart-b} it suffices to treat
$a_!L\otimes^{\mathbf L}b_!M$ with $L$, $M$ lisse on strata $S$, $T$,
and the formula
$(a_!L)\otimes^{\mathbf L}N\simeq a_!(L\otimes^{\mathbf L}a^*N)$ from the
proof of Lemma~\ref{b:lem:integral-models}(3) identifies it with the
extension by zero of the lisse sheaf $L\otimes M$ on $S\cap T$, and the
subcategory is everything. Now fix $A\in\Loc(\Dbc(X,\Lambda))$ and repeat the
argument in $B$.

\emph{Projection formula.} Fix $B\in\Dbc(Y,\Lambda)$. Both sides of
\[
 (Rp_*A)\otimes^{\mathbf L}B\longrightarrow
                    Rp_*(A\otimes^{\mathbf L}p^*B)
 \quad\text{in }\Loc(\Dbc(Y,\Lambda))
\]
are exact in $A$ and preserve direct sums, and the map is an isomorphism
for $A\in\Dbc(X,\Lambda)$ by Proposition~\ref{b:prop:proper-coefficients}.
Lemma~\ref{b:lem:localizing-comparison} extends it to all $A$. Then fix
$A$ and argue in the same way in $B$. Proper base change is
Proposition~\ref{b:prop:proper-base-change-unbounded}. All these maps
are the canonical ones and satisfy the usual compatibilities with
composition.
\end{proof}

\begin{proposition}[The projective line]
\label{b:prop:p1-adjunction-rational}
For $\Lambda=E$ or $\Qbar$ and $q:\mathbf P^1_Y\to Y$, the functor
$q^*(1)[2]\colon\Loc(\Dctf(Y,\Lambda))\to\Loc(\Dctf(\mathbf P^1_Y,\Lambda))$
is right adjoint to
$Rq_*:\Loc(\Dctf(\mathbf P^1_Y,\Lambda))\to\Loc(\Dctf(Y,\Lambda))$. Hence
$q^!\simeq q^*(1)[2]$ as functors
$\Loc(\Dctf(Y,\Lambda))\to\Loc(\Dctf(\mathbf P^1_Y,\Lambda))$, where $q^!$
is the right adjoint of Theorem~\ref{b:thm:brown-adjoint}.
\end{proposition}
\begin{proof}
All products are over $Y$. Let $p_1,p_2:\mathbf P^1_Y\times\mathbf P^1_Y\to\mathbf P^1_Y$
be the projections and $\delta_\Lambda$, $\operatorname{Tr}$ the class
and traces of Proposition~\ref{b:prop:diagonal-tower}. The unit and
counit are defined as in Proposition~\ref{b:prop:geo-p1-adjunction}.
For $M\in \Loc(\Dbc(\mathbf P^1_Y,\Lambda))$ the unit is the morphism in
$\Loc(\Dbc(\mathbf P^1_Y,\Lambda))$
\[
 \eta_M:M=Rp_{2*}\Delta_*M
 \longrightarrow Rp_{2*}\bigl(p_1^*M(1)[2]\bigr)
 \simeq q^*Rq_*M(1)[2],
\]
where $Rp_{2*}$ is the direct image
$\Loc(\Dbc(\mathbf P^1_Y\times\mathbf P^1_Y,\Lambda))\to\Loc(\Dbc(\mathbf P^1_Y,\Lambda))$,
the first arrow is $Rp_{2*}$ of $1\otimes\delta_\Lambda$ under
$\Delta_*M\simeq p_1^*M\otimes^{\mathbf L}\Delta_*\Lambda_{\mathbf P^1_Y}$,
and the second is proper base change. The isomorphism
$\Delta_*M\simeq p_1^*M\otimes^{\mathbf L}\Delta_*\Lambda_{\mathbf P^1_Y}$
holds for every $M$: both sides restrict to $M$ on the diagonal and to
$0$ on its complement, and these two restrictions are jointly
conservative (\stacks{09AD}). For $K\in \Loc(\Dbc(Y,\Lambda))$ the counit
is the morphism in $\Loc(\Dbc(Y,\Lambda))$
\[
 \varepsilon_K:Rq_*q^*K(1)[2]
 \simeq K\otimes^{\mathbf L}Rq_*\Lambda_{\mathbf P^1_Y}(1)[2]
 \xrightarrow{1\otimes \operatorname{Tr}_q}K,
\]
using the projection formula.

The triangle identities are verified by the calculation in the proof of
Proposition~\ref{b:prop:geo-p1-adjunction}. That calculation uses only
the projection formula, proper base change and composition of derived
direct images, each applied to the canonical maps, and these hold on
$\Loc(\Dbc(-,\Lambda))$ by Proposition~\ref{b:prop:pf-bc-generated}.
It moves the factor $M$, respectively $q^*K$, outside the direct image by
the projection formula; the base change map, being defined by the
adjunction counit, cancels against restriction to the diagonal and gives
$Rp_{i*}\Delta_*=\mathrm{id}$. The result is the pair of equalities of
endomorphisms in $\Loc(\Dbc(Y,\Lambda))$ and
$\Loc(\Dbc(\mathbf P^1_Y,\Lambda))$
\[
\begin{aligned}
 \varepsilon_{Rq_*M}\circ Rq_*(\eta_M)
   &=Rq_*\bigl(1_M\otimes
                   (\operatorname{Tr}_{p_1}\circ Rp_{1*}\delta_\Lambda)\bigr),\\
 q^*(\varepsilon_K)(1)[2]\circ\eta_{q^*K(1)[2]}
   &=1_{q^*K(1)[2]}\otimes(\operatorname{Tr}_{p_2}\circ Rp_{2*}\delta_\Lambda),
\end{aligned}
\]
where the projection formula for $p_i$ identifies the morphism in
parentheses with an endomorphism of $\Lambda_{\mathbf P^1_Y}$. By the
integration identities of Proposition~\ref{b:prop:diagonal-tower} both
endomorphisms are the identity, so both triangle identities hold.
Uniqueness of right adjoints gives $q^!\simeq q^*(1)[2]$.
\end{proof}

\subsection{\texorpdfstring{$f^!$}{f!} preserves bounded constructible complexes}
\label{b:sec:restriction}

Locally on $X$, a morphism
$f:X\to Y$ is a closed immersion into an affine space $\mathbf A^d_Y$
followed by the projection, and $f^!$ is compatible with composition and
with restriction to open subsets
(Proposition~\ref{b:prop:composition-locality}). So it suffices to treat
the projection $\mathbf A^d_Y\to Y$ and closed immersions.

\begin{corollary}\label{b:cor:affine-space-rational}
For $\Lambda=E$ or $\Qbar$ and the projection
$p:\mathbf A^d_V\to V$, there is an isomorphism of functors
$\Loc(\Dctf(V,\Lambda))\to\Loc(\Dctf(\mathbf A^d_V,\Lambda))$
\[
 p^!\simeq p^*(d)[2d].
\]
In particular $p^!$ maps $\Dctf(V,\Lambda)$ into
$\Dctf(\mathbf A^d_V,\Lambda)$.
\end{corollary}
\begin{proof}
Let $d=1$, and write $p=q\circ j$ with $j:\mathbf A^1_V\hookrightarrow\mathbf P^1_V$
the standard open immersion and $q:\mathbf P^1_V\to V$ the projection.
Then $Rp_!=Rq_*\circ j_!$. On the categories $\Loc(\Dbc(-,\Lambda))$,
$j_!$ is left adjoint to $j^*$ (both preserve these categories), and
$Rq_*$ is left adjoint to $q^*(1)[2]$ by
Proposition~\ref{b:prop:p1-adjunction-rational}. Composing, $Rp_!$ is left
adjoint to $j^*q^*(1)[2]=p^*(1)[2]$, and uniqueness of right adjoints
gives $p^!\simeq p^*(1)[2]$. For general $d$, write $p$ as a composite of
$d$ projections $\mathbf A^{k}_V\to\mathbf A^{k-1}_V$ and use
$(gf)^!\simeq f^!g^!$ (Proposition~\ref{b:prop:composition-locality}).

The functor $p^*(d)[2d]$ preserves $\Dbc$: pullback and the Tate twist
send lisse sheaves on the strata to lisse sheaves on the pulled-back
strata, and shifts preserve boundedness.
\end{proof}

For a closed immersion $i:Z\hookrightarrow X$ with open complement
$j:U\hookrightarrow X$, the functor $R\underline\Gamma_Z$ of sections with
support in $Z$ is a functor $D(X_{\proet},\Lambda)\to D(Z_{\proet},\Lambda)$
right adjoint to $i_*\colon D(Z_{\proet},\Lambda)\to D(X_{\proet},\Lambda)$
(Lemma~\ref{b:lem:formal-closed-support}(1)), and it sits in the triangle
in $D(X_{\proet},\Lambda)$
\[
 i_*R\underline\Gamma_ZK\longrightarrow K\longrightarrow Rj_*j^*K
 \longrightarrow i_*R\underline\Gamma_ZK[1].
\]
Constructibility of $Rj_*$ (Proposition~\ref{b:prop:open-coefficients})
therefore makes $R\underline\Gamma_ZK$ constructible, and the next proposition identifies it with $i^!K$.

\begin{example}\label{b:ex:closed-support}
Let $X=\mathbf A^1_k$ with $k$ algebraically closed, let $Z$ be the origin
and $U=\mathbf G_{m,k}$ its complement. Then $Rj_*\Lambda_U$ has
$\mathcal H^0=\Lambda_X$ and $\mathcal H^1=i_*\Lambda(-1)$; the latter
comes from the Kummer covers $t\mapsto t^{\ell^n}$ of $\mathbf G_m$ near
$0$. The unit $\Lambda_X\to Rj_*\Lambda_U$ is an isomorphism on
$\mathcal H^0$, so the triangle above with $K=\Lambda_X$ gives
$i_*R\underline\Gamma_Z\Lambda_X\simeq\mathcal H^1(Rj_*\Lambda_U)[-2]$, that is
$i^!\Lambda_X\simeq\Lambda(-1)[-2]$. This is the answer one expects for a
point of codimension one, and it agrees with smooth purity
(Section~\ref{sec:smooth-purity}).
\end{example}

\begin{proposition}\label{b:prop:closed-adjoint-constructible}
Let $i:Z\hookrightarrow X$ be a closed immersion. For
$K\in\Dbc(X,\Lambda)$ there is an isomorphism
$i^!K\simeq R\underline\Gamma_ZK$ in $\Loc(\Dbc(Z,\Lambda))$, natural in $K$. In particular $i^!$
maps $\Dbc(X,\Lambda)$ into $\Dbc(Z,\Lambda)$.
\end{proposition}
\begin{proof}
Let $K\in\Dbc(X,\Lambda)$. The triangle above comes from the localization
sequence \stacks{09AH} and exactness of $i_*$ (\stacks{09BM}). Its third
term $Rj_*j^*K$ is bounded constructible by
Proposition~\ref{b:prop:open-coefficients}, so $i_*R\underline\Gamma_ZK$ is
bounded constructible. Applying $i^*$, and using $i^*i_*\simeq\mathrm{id}$
(\stacks{09AA}), gives $R\underline\Gamma_ZK\in\Dbc(Z,\Lambda)$.

Both $i^!K$ and $R\underline\Gamma_ZK$ now lie in $\Loc(\Dbc(Z,\Lambda))$.
For $A$ in this category, $Ri_!=i_*$ and
Theorem~\ref{b:thm:brown-adjoint} give $\Hom(A,i^!K)=\Hom(i_*A,K)$, with
Hom groups in $\Loc(\Dbc(Z,\Lambda))$ and $\Loc(\Dbc(X,\Lambda))$, and
Lemma~\ref{b:lem:formal-closed-support}(1) gives the same identity
$\Hom(A,R\underline\Gamma_ZK)=\Hom(i_*A,K)$ in $D(Z_{\proet},\Lambda)$ and
$D(X_{\proet},\Lambda)$, whose full subcategories these are. By the Yoneda lemma,
$i^!K\simeq R\underline\Gamma_ZK$, naturally in $K$.
\end{proof}

\begin{theorem}[$f^!$ with rational coefficients]\label{b:thm:restriction}
For $\Lambda=E$ or $\Qbar$ and every separated morphism of finite type
$f:X\to Y$, $f^!$ satisfies
\[
 f^!\bigl(\Dctf(Y,\Lambda)\bigr)\subset\Dctf(X,\Lambda).
\]
Its restriction is right adjoint to
$Rf_!:\Dctf(X,\Lambda)\to\Dctf(Y,\Lambda)$.
\end{theorem}
\begin{proof}
Let $K\in \Dbc(Y,\Lambda)$. Cover $Y$ by finitely many affine opens
$V_\beta$, and each $f^{-1}(V_\beta)$ by finitely many affine opens
$U_{\beta\alpha}$. Since $f$ is of finite type, each $U_{\beta\alpha}$
admits a closed immersion $i:U_{\beta\alpha}\hookrightarrow\mathbf A^d_{V_\beta}$
over $V_\beta$. With $p:\mathbf A^d_{V_\beta}\to V_\beta$ the projection
and $u:U_{\beta\alpha}\hookrightarrow X$, $v:V_\beta\hookrightarrow Y$ the
inclusions, we have $f\circ u=v\circ p\circ i$. By
Proposition~\ref{b:prop:composition-locality}, as functors
$\Loc(\Dbc(Y,\Lambda))\to\Loc(\Dbc(U_{\beta\alpha},\Lambda))$,
\[
 u^*f^!\simeq (fu)^!=(vpi)^!\simeq i^!p^!v^*.
\]
Now $v^*K$ is bounded constructible, so $p^!v^*K$ is bounded
constructible by Corollary~\ref{b:cor:affine-space-rational}, and
$u^*f^!K\simeq i^!p^!v^*K$ is bounded constructible by
Proposition~\ref{b:prop:closed-adjoint-constructible}.

Both conditions defining $\Dbc$ are local for a finite open cover. The
cohomology of $f^!K$ vanishes outside the union of the finitely many
bounded intervals found on the $U_{\beta\alpha}$. For constructibility,
each locally closed stratum of $U_{\beta\alpha}$ is locally closed in
$X$, and a common refinement of the finitely many resulting finite
partitions of $X$ into locally closed subsets is a stratification on
whose strata all cohomology sheaves of $f^!K$ are lisse. Hence
$f^!K\in \Dbc(X,\Lambda)$.

Finally, $Rf_!$ maps $\Dbc(X,\Lambda)$ into $\Dbc(Y,\Lambda)$
(Proposition~\ref{b:prop:lower-exceptional-image}), and
$\Dbc(-,\Lambda)\subset\Loc(\Dbc(-,\Lambda))$ is full. So for
$A\in \Dbc(X,\Lambda)$ the bijection of Theorem~\ref{b:thm:brown-adjoint}
reads
\[
 \Hom_{\Dbc(Y,\Lambda)}(Rf_!A,K)
   \simeq\Hom_{\Dbc(X,\Lambda)}(A,f^!K).\qedhere
\]
\end{proof}

\begin{remark}
The results used in this proof are proved before $f^!$ is known on
constructible complexes. Compactness (\S\ref{b:sec:compactness}) uses
$R\underline\Gamma_Z$ only through open--closed localization. The
projective-line adjunction is built from the Cartier extension, Kummer
theory, the cohomology of $\mathbf P^1$, and properties of $Rf_!$ and
$Rf_*$, and the torsion theorem for open immersions behind
Proposition~\ref{b:prop:open-coefficients} is built from that
adjunction with torsion coefficients.
\end{remark}

This proves Theorem~\ref{n:existence} for $\Lambda=E,\Qbar$: $f^!$
preserves $\Dbc$, and on bounded constructible complexes it is right
adjoint to $Rf_!$.

\section{Integral coefficients}
\label{i:heart}
The rational arguments used at several places that a constructible sheaf
is locally free on the strata of a stratification. With integral
coefficients a constructible sheaf may have torsion, such as
$\mathcal O_{E,X}/\pi_E$. What survives is that on each stratum a
constructible $\mathcal O_{E,X}$-module is locally the completion of a
finitely generated $\mathcal O_E$-module. Since $\mathcal O_E$ is a
discrete valuation ring, such a module has a free resolution of length
one, so tensoring with it has Tor amplitude in $[-1,0]$. This bound
adapts the rational compactness estimates to $\Lambda=\mathcal O_E$,
and the resulting estimates supply the integral case of
Theorem~\ref{ctf:unified-brown}. Some integral statements were proved
alongside the rational ones and are used below as they stand: proper direct image,
the projection formula and proper base change on $\Dbc(-,\mathcal O_E)$
(Proposition~\ref{b:prop:proper-coefficients}), open direct image
(Proposition~\ref{b:prop:open-coefficients}), and the $\mathcal O_E$-linear
cycle class of the diagonal (Proposition~\ref{b:prop:diagonal-tower}).

Write $\mathbf F_E=\mathcal O_E/\pi_E$ for the residue field and
$\mathbf F_{E,X}$ for the pullback of $\underline{\mathbf F_E}$ along
$p_X$. Here
$\Dctf(-,\mathcal O_E)=\Dbc(-,\mathcal O_E)$.
In this section $f^!$ is the right adjoint of
Theorem~\ref{ctf:unified-brown} for $\Lambda=\mathcal O_E$, on
$\Loc(\Dctf(-,\mathcal O_E))$ (Proposition~\ref{i:brown}). Reduction
modulo $\pi_E^n$ always means the derived tensor product
$K\otimes^{\mathbf L}_{\mathcal O_{E,X}}\mathcal O_{E,X}/\pi_E^n$, which is
the cone of multiplication by $\pi_E^n$ on $K$.

The category $\Loc(\Dbc(X,\mathcal O_E))$ is not a category of derived
complete objects. Its direct sums are the ordinary ones, and it contains
\[
 \mathcal O_{E,X}[1/\ell]
 =\operatorname{hocolim}\bigl(\mathcal O_{E,X}
   \xrightarrow{\ell}\mathcal O_{E,X}
   \xrightarrow{\ell}\cdots\bigr),
\]
whose derived $\ell$-completion is zero. Constructible integral objects
are complete, but a direct sum of them need not be. So the equivalence
on complete objects in Lemma~\ref{b:lem:integral-models}(1) does not
apply to the whole localizing category, and
Theorem~\ref{ctf:unified-brown} is applied to
$\Loc(\Dctf(-,\mathcal O_E))$ with its ordinary direct sums. The
structure and compactness results below verify its hypotheses for
$\mathcal O_E$, and the last subsection proves that
$f^!K$ is complete and constructible when
$K\in\Dctf(-,\mathcal O_E)$.

\subsection{Integral constructible complexes}

\begin{lemma}[Structure of integral constructible complexes]
\label{i:integral-heart}
\leavevmode
\begin{enumerate}
\item $\Dbc(X,\mathcal O_E)$ is a triangulated subcategory of
  $D(X_{\proet},\mathcal O_E)$, its objects are bounded, and it is stable
  under the standard truncations. The category $\Cons(X,\mathcal O_E)$ of
  constructible $\mathcal O_{E,X}$-modules (Definition~\ref{n:def-constructible})
  consists of the objects of $\Dbc(X,\mathcal O_E)$ concentrated in degree
  $0$, and is an essentially small abelian subcategory of
  $\Mod(X_{\proet},\mathcal O_E)$ closed under extensions.
\item For $C\in\Cons(X,\mathcal O_E)$ there is a finite stratification of
  $X$ such that on each stratum $C$ is pro-\'etale locally the completion
  of a finitely generated $\mathcal O_E$-module, possibly with torsion.
\item On each such stratum, $C$ is dualizable in the derived category,
  of Tor amplitude in $[-1,0]$, and its dual has Tor amplitude in
  $[0,1]$.
\end{enumerate}
\end{lemma}
\begin{proof}
By Lemma~\ref{b:lem:integral-models}(1),(2), restriction of scalars to
the constant sheaf $\underline{\mathcal O_E}$, which is flat, identifies
$\Dbc(X,\mathcal O_E)$ with the category $\Dc(X,\mathcal O_E)$ of the
Stacks Project (\stacks{099I}, \stacks{099L}). The two conditions of
\S\ref{n:constructible}, derived completeness and constructibility of the
reduction modulo $\pi_E$, are stable under cones and shifts, so
$\Dbc(X,\mathcal O_E)$ is triangulated, and its objects are bounded
because our schemes are quasi-compact (\stacks{09C2}).

The local structure is the key point, and it is \stacks{09C7}, whose
Noetherian hypothesis holds for our schemes (since $\mathcal O_E$ is
already complete, the completed coefficient ring of \stacks{09C4} is
$\mathcal O_E$ itself). The statement is: \emph{for
$K\in\Dbc(X,\mathcal O_E)$ there is a finite stratification of $X$ such
that on each stratum $K$ is pro-\'etale locally the completion of a
bounded complex $M^\bullet$ of finite free $\mathcal O_E$-modules, with
differentials given by matrices over $\mathcal O_E$.} The idea of the
proof is that the reduction $K/\pi_E$ is a classical constructible complex,
hence on a stratum is locally a perfect complex of
$\mathcal O_E/\pi_E$-modules; on a weakly contractible cover, where
cohomology vanishes in positive degrees, such a perfect complex and its
successive deformations to $\mathcal O_E/\pi_E^n$ can be lifted
compatibly, and the derived limit of the lifts is $K$ by completeness. As $\underline{\mathcal O_E}\to\mathcal O_{E,X}$ is flat,
$H^q(K)$ is then locally the completion of the finitely generated
$\mathcal O_E$-module $H^q(M^\bullet)$; this is (2) once we know that
$H^q(K)$ lies in $\Cons(X,\mathcal O_E)$.

First, $H^q(K)$ is derived complete: the
pro-\'etale site has enough weakly contractible affines
(\stacks{0990}), and \stacks{099Q}(5) then shows that cohomology sheaves
of derived complete complexes are derived complete. Second, a finitely
generated $\mathcal O_E$-module has a free resolution of length one, in
degrees $-1,0$. The reduction of $H^q(K)$ modulo $\pi_E$ therefore has
cohomology only in degrees $-1$ and $0$, and on the chosen strata these
cohomology sheaves are locally constant with finite stalks, hence come
from constructible sheaves on the \'etale site (\stacks{099Y},
\stacks{09AK}). Thus $H^q(K)\in\Dbc(X,\mathcal O_E)$. Every truncation
of $K$ is a finite extension of shifts of the $H^q(K)$, so it also lies
in $\Dbc(X,\mathcal O_E)$.

Stability under truncation gives (1) for the heart: the kernel and
cokernel of a map of constructible sheaves are cohomology sheaves of its
cone, and an extension of two constructible sheaves is the middle term
of a triangle with constructible ends. Essential smallness holds because
the data describing a constructible sheaf (a finite stratification,
finite presentations of $\mathcal O_E$-modules, sheaves on a fixed small
site, and gluing morphisms) form a set. For (3), the length one free
resolution of the local model gives the derived dual and the evaluation
and coevaluation maps locally, and the Tor amplitudes can be read off
from it.
\end{proof}

For the local model $\mathcal O_E/\pi_E$ the amplitudes in (3) can be seen
directly. The resolution
$0\to\mathcal O_E\xrightarrow{\pi_E}\mathcal O_E\to\mathcal O_E/\pi_E\to0$
shows that $M\otimes^{\mathbf L}\mathcal O_E/\pi_E$ is the cone of $\pi_E$
on $M$, which lives in degrees $[-1,0]$ for a module $M$. Applying
$\RHom(-,\mathcal O_E)$ to the same resolution gives the dual
$\mathcal O_E/\pi_E[-1]$, whose Tor amplitude is $[0,1]$.

Arguments by d\'evissage along a stratification need inverse image,
tensor products and extension by zero along strata to preserve integral
constructible complexes.

\begin{lemma}[Inverse image and tensor product]
\label{i:tensor}
Inverse image $g^*\colon D(Y_{\proet},\mathcal O_E)\to D(X_{\proet},\mathcal O_E)$
along a morphism $g:X\to Y$ preserves $\Dbc(-,\mathcal O_E)$, and on these
categories it agrees, after restriction of scalars to $\underline{\mathcal O_E}$, with
the completed inverse image of the Stacks Project. For
$K,L\in \Dbc(X,\mathcal O_E)$, the tensor product
$K\otimes_{\mathcal O_{E,X}}^{\mathbf L}L\in D(X_{\proet},\mathcal O_E)$, without completion, lies in
$\Dbc(X,\mathcal O_E)$. Extension by zero along a quasi-compact open
immersion and direct image along a closed immersion preserve
$\Dbc(-,\mathcal O_E)$.
\end{lemma}
\begin{proof}
\emph{Extension by zero and closed direct image.} For a quasi-compact
open immersion $j$, $j_!$ is exact and preserves all limits
(\stacks{09BP}). For a closed immersion $i$, the complementary open
immersion is quasi-compact since our schemes are Noetherian, and $i_*$
is exact (\stacks{09BM}) and preserves limits (\stacks{09AA}). Both
therefore preserve derived completeness, and being exact they commute
with reduction modulo $\pi_E^n$, where they preserve constructibility
with torsion coefficients. Hence they preserve $\Dbc(-,\mathcal O_E)$. Combining
this with open--closed localization along a stratification, every
$K\in\Dbc(X,\mathcal O_E)$ is built by finitely many extensions from
objects $a_!M$, where $a:S\hookrightarrow X$ is the inclusion of a
stratum and $M$ is pro-\'etale locally the completion of a bounded
complex of finite free $\mathcal O_E$-modules (Lemma~\ref{i:integral-heart}).

\emph{Inverse image.} For such a local model $M$, the inverse image is
the same complex of finite free $\mathcal O_E$-modules, completed over the
coefficient ring of the new scheme. It is derived complete and its
reduction is constructible. Since inverse image commutes with $a_!$ (base
change for open extension and for closed direct image), the conclusion
passes to $a_!M$ and then, through the finitely many localization
triangles, to all of $\Dbc(-,\mathcal O_E)$. The comparison map to the
completed inverse image is an isomorphism on the local models, hence on
all objects by the same triangles.

\emph{Tensor product.} For a locally closed immersion $a$,
\[
 (a_!M)\otimes^{\mathbf L} L\simeq a_!(M\otimes^{\mathbf L} a^*L);
\]
this holds for an open immersion and for a closed immersion, first
termwise and then on K-flat resolutions, and a locally closed immersion
is a composite of the two. It therefore suffices to show that
$M\otimes^{\mathbf L}a^*L\in\Dbc(S,\mathcal O_E)$. The local models of
$M$ provide a derived dual and evaluation and coevaluation maps locally;
these descend, and the triangle identities can be checked locally. Thus
$M$ is dualizable, and $M\otimes^{\mathbf L}-$, having adjoints on both
sides, preserves products and derived inverse limits. Since $a^*L$ is
derived complete, so is $M\otimes^{\mathbf L}a^*L$. Its reduction modulo
$\pi_E$ is the tensor product of the reductions, which is constructible.
By the first paragraph, $a_!$ preserves these properties, and finitely
many triangles in the first variable give the claim for all $K$ and $L$.
\end{proof}


\subsection{Compactness}
In the rational case, compactness came from reducing maps out of a
constructible sheaf to a finite range of target degrees
(Lemma~\ref{b:lem:finite-window}). That reduction rests on two facts
about Ext out of a constructible sheaf, continuity and a uniform bound,
and both hold with integral coefficients.

\begin{lemma}[Continuity and boundedness of Ext]
\label{i:ext}
\leavevmode
\begin{enumerate}
\item For $C\in \Dbc(X,\mathcal O_E)$, the functors
  $\Hom_{D(X_{\proet},\mathcal O_E)}(C,-[q])$ commute
  with filtered colimits of complexes with a common lower bound.
\item For each $C\in\Cons(X,\mathcal O_E)$ there is an integer $d_C$ such
  that
  \[
   \Ext^q_{\mathcal O_{E,X}}(C,N)=0
   \qquad\text{for all }q>d_C\text{ and all }N\in\Cons(X,\mathcal O_E).
  \]
\end{enumerate}
\end{lemma}
\begin{proof}
Both statements reduce, by the long exact Ext sequences of a finite
filtration of $C$, to $C=a_!L$, where $a=i\circ u:S\hookrightarrow X$ is
the inclusion of a stratum, with $u:S\hookrightarrow Z$ open and
$i:Z\hookrightarrow X$ closed, and $L$ is locally the completion of a
finitely generated $\mathcal O_E$-module. Let $j$ be the open immersion
complementary to $i$. Lemma~\ref{b:lem:formal-closed-support} holds for
$\mathcal O_{E,X}$-modules with the same proof. Tensor--Hom
adjunction and the adjunction between $i_*$ and $R\underline\Gamma_Z$
give
\[
 \RHom_X(a_!L,K)
  =R\Gamma\bigl(S,\cHom_{\mathcal O_{E,S}}(L,\mathcal O_{E,S})
     \otimes_{\mathcal O_{E,S}}^{\mathbf L} u^*R\underline\Gamma_ZK\bigr).
\]
By Lemma~\ref{i:integral-heart}(3), tensoring with the dual
$\cHom_{\mathcal O_{E,S}}(L,\mathcal O_{E,S})$ has Tor amplitude $[0,1]$
and commutes with filtered colimits.

(1) The proof of Lemma~\ref{b:lem:bounded-continuity} works for any sheaf
of rings. Apply it to the two terms of the fibre sequence defining
$R\underline\Gamma_Z$ (as in the proof of
Lemma~\ref{b:lem:constructible-ext-continuity}) and use that tensoring
with the dual has finite Tor amplitude. This gives continuity for
$C=a_!L$, and in particular for $\Hom$ and $\Ext^1$ between
constructible sheaves.

(2) Let $N\in\Cons(X,\mathcal O_E)$. Its reductions
$N\otimes^{\mathbf L}\mathcal O_{E,X}/\pi_E^n$ are constructible torsion
complexes in degrees $[-1,0]$, by the length one resolution, and
$N\simeq R\varprojlim_n(N\otimes^{\mathbf L}\mathcal O_{E,X}/\pi_E^n)$ by
derived completeness (\stacks{099Q}). The cohomological dimensions
$\operatorname{cd}_\ell(j)$ and $\operatorname{cd}_\ell(S)$ of
Definition~\ref{n:def-cd} are finite and depend only on $j$ and $S$. We
bound the degrees of the formula above, applied to each reduction of $N$:
\begin{itemize}
\item $R\underline\Gamma_Z$ of a torsion complex in degrees $[-1,0]$ has
  cohomology in degrees $\le\operatorname{cd}_\ell(j)+1$, by its fibre
  description;
\item tensoring with the dual of $L$ raises the upper bound by at most
  one;
\item $R\Gamma(S,-)$ raises it by at most $\operatorname{cd}_\ell(S)$.
\end{itemize}
Each reduction therefore contributes only in degrees
$\le\operatorname{cd}_\ell(j)+\operatorname{cd}_\ell(S)+2$. All the functors
involved ($R\underline\Gamma_Z$, open restriction, derived sections,
closed inverse image by \stacks{09AA}, and tensoring with the dualizable
dual of $L$) commute with derived inverse limits. The derived inverse
limit of a countable tower raises the upper bound by at most one (Milnor
sequence, \stacks{08TB}); this applies because the complexes and
transition maps can be taken in the abelian category of inverse systems
before deriving. Hence
\[
 \Ext^q_X(a_!L,N)=0\qquad(q>\operatorname{cd}_\ell(j)+\operatorname{cd}_\ell(S)+3).
\]
The bound depends only on $a$. For general $C$, take the maximum over
the finitely many pieces $a_!L$ of a filtration of $C$.
\end{proof}

For $\Lambda=\mathcal O_E$, the hypotheses of
Lemma~\ref{ctf:unified-compact} and of the construction of $Rf_!$ by
compactifications are compactness of $\Dctf(-,\mathcal O_E)$ and
preservation of direct sums by proper direct image and by extension by
zero. They follow from the lemma above.

\begin{proposition}[Compactness and proper direct image]
\label{i:compactness}
Every object of $\Dctf(X,\mathcal O_E)=\Dbc(X,\mathcal O_E)$
is compact in $\Loc(\Dctf(X,\mathcal O_E))$. For proper $p:X\to Y$,
$Rp_*\colon D(X_{\proet},\mathcal O_E)\to D(Y_{\proet},\mathcal O_E)$ maps $\Dctf(X,\mathcal O_E)$ into
$\Dctf(Y,\mathcal O_E)$, maps $\Loc(\Dctf(X,\mathcal O_E))$
into $\Loc(\Dctf(Y,\mathcal O_E))$ and preserves direct sums there.
\end{proposition}
\begin{proof}
\emph{Compactness.} The rational proof
(Theorem~\ref{b:thm:constructible-compactness}) used three facts about
$\Cons(X,\Lambda)$: it is an essentially small abelian category closed
under extensions, Ext out of a constructible sheaf commutes with
filtered colimits, and the uniform Ext bound. These are
Lemmas~\ref{i:integral-heart} and~\ref{i:ext} for $\mathcal O_E$. The
proofs of Proposition~\ref{b:prop:ind-heart} and
Corollary~\ref{b:prop:ind-heart-b} then give an exact fully faithful
realization of $\IndCons(X,\mathcal O_E)=\operatorname{Ind}(\Cons(X,\mathcal O_E))$
in $\Mod(X_{\proet},\mathcal O_E)$, closed under extensions and under
kernels and cokernels of its own morphisms, and the bound
$\Ext^q(C,B)=0$ in $\Mod(X_{\proet},\mathcal O_E)$ for $q>d_C$ and $B\in\IndCons(X,\mathcal O_E)$. Truncations are taken in the full
subcategory $\Dic(X,\mathcal O_E)$ of complexes with cohomology in
$\IndCons(X,\mathcal O_E)$, which is closed under direct sums and
truncations and contains $\Loc(\Dbc(X,\mathcal O_E))$; the latter need
not be stable under truncation.

The proof of Lemma~\ref{b:lem:finite-window} applies with these
hypotheses, since products of sheaves on the pro-\'etale site are exact
and its derived category is left complete
(Propositions~\ref{pe:products} and~\ref{pe:left-complete}). Its key step
is the lower tail. Continuity only covers systems with a common lower
bound, while objects of $\Loc(\Dbc(X,\mathcal O_E))$ may be unbounded
below. Along the tower of good truncations, uniform Ext vanishing makes
the relevant tower of Hom groups eventually constant, so the
$\varprojlim^1$ term in the Milnor sequence vanishes. As a result,
$\Hom(C,K[m])$, computed in $D(X_{\proet},\mathcal O_E)$, depends only on the degrees $[m-d_C,m]$ of $K$. On a
finite range of degrees, Ext commutes with direct sums by continuity,
and induction on the length of the range shows that $\Hom(C,-)$ commutes
with direct sums. Finite truncation of the source extends this to all of
$\Dbc(X,\mathcal O_E)$.

\emph{Proper direct image of $\Dbc$.} This is the integral case of
Proposition~\ref{b:prop:proper-coefficients}.

\emph{Direct sums.} Let $d$ bound the dimensions of the fibres of $p$. On
a weakly contractible affine $V\to Y$, proper direct image of a torsion
sheaf has cohomology in degrees $\le 2d$ (\stacks{095U}), and since
$\Gamma(V,-)$ is exact this bounds cohomology of sections. The reductions
of a constructible sheaf live in degrees $[-1,0]$, and one derived
inverse limit adds at most one degree, so $Rp_*$ of a constructible
sheaf has cohomology in degrees $[0,2d+1]$. By continuity for bounded
below systems, the same bound holds for objects of
$\IndCons(X,\mathcal O_E)$: for $N=\varinjlim_iN_i$ with
$N_i\in\Cons(X,\mathcal O_E)$, one has
$H^qRp_*N=\varinjlim_iH^qRp_*N_i$, and these lie in
$\IndCons(Y,\mathcal O_E)$ because each $H^qRp_*N_i$ is constructible by
the previous paragraph. Now the proofs of
Lemmas~\ref{b:lem:proper-unbounded-window}
and~\ref{b:lem:proper-unbounded-window-b} apply: in a fixed degree $q$,
$H^qRp_*K$ depends only on the degrees $[q-2d-1,q]$ of $K$, and on such
a finite range direct sums commute with $Rp_*$ by induction on the
length. This holds for arbitrary families in $\Dic(X,\mathcal O_E)$,
with no common lower bound. Finally, the objects of
$\Loc(\Dbc(X,\mathcal O_E))$ whose direct image lies in
$\Loc(\Dbc(Y,\mathcal O_E))$ form a localizing subcategory containing
$\Dbc(X,\mathcal O_E)$, hence everything.
\end{proof}

\begin{corollary}\label{i:compactness-b}
For a quasi-compact open immersion $j:U\hookrightarrow X$, extension by
zero $j_!\colon D(U_{\proet},\mathcal O_E)\to D(X_{\proet},\mathcal O_E)$
maps $\Dbc(U,\mathcal O_E)$ into $\Dbc(X,\mathcal O_E)$, maps
$\Loc(\Dbc(U,\mathcal O_E))$ into $\Loc(\Dbc(X,\mathcal O_E))$, and
preserves direct sums.
\end{corollary}
\begin{proof}
The functor $j_!$ is exact and preserves direct sums, being a left
adjoint, and it preserves $\Dbc(-,\mathcal O_E)$ by Lemma~\ref{i:tensor}.
The localizing argument at the end of the proof of
Proposition~\ref{i:compactness} gives the second assertion.
\end{proof}

\subsection{The exceptional inverse image with integral coefficients}

\begin{proposition}[Compactly supported direct image and its adjoint]
\label{i:brown}
Compactifications define a functor
$Rf_!:\Loc(\Dctf(X,\mathcal O_E))\to\Loc(\Dctf(Y,\mathcal O_E))$,
independent of the compactification, compatible with composition and
base change, preserving direct sums and sending
$\Dctf(X,\mathcal O_E)$ into $\Dctf(Y,\mathcal O_E)$. Its exact right adjoint
$f^!$ is the right adjoint of Theorem~\ref{ctf:unified-brown} for
$\Lambda=\mathcal O_E$.
\end{proposition}
\begin{proof}
This is the case $\Lambda=\mathcal O_E$ of
Theorem~\ref{ctf:unified-brown}, using
$\Dctf(-,\mathcal O_E)=\Dbc(-,\mathcal O_E)$. Its hypotheses for
$\mathcal O_E$ are Propositions~\ref{b:prop:proper-coefficients}
and~\ref{i:compactness} and Corollary~\ref{i:compactness-b}.
\end{proof}

To show that $f^!$ preserves $\Dbc(-,\mathcal O_E)$, we need the
projection formula and proper base change on the localizing categories
for the projective line, and the finiteness of open direct image for
closed immersions.

\begin{lemma}[Open direct image and projection formula]
\label{i:open-pf}
For a quasi-compact open immersion $j:U\hookrightarrow X$,
$Rj_*\colon D(U_{\proet},\mathcal O_E)\to D(X_{\proet},\mathcal O_E)$ maps
$\Dctf(U,\mathcal O_E)$ into $\Dctf(X,\mathcal O_E)$.
Proper base change and the proper projection formula, with the tensor
product $\otimes^{\mathbf L}_{\mathcal O_{E,X}}$, hold on
$\Loc(\Dctf(-,\mathcal O_E))$: the base change and projection formula maps
are isomorphisms in $\Loc(\Dctf(Y',\mathcal O_E))$, respectively
$\Loc(\Dctf(Y,\mathcal O_E))$.
\end{lemma}
\begin{proof}
The first assertion is the integral case of
Proposition~\ref{b:prop:open-coefficients}. In detail: $Rj_*K$ is derived
complete by \stacks{099L}, and
\[
 (Rj_*K)\otimes^{\mathbf L}_{\mathcal O_{E,X}}\mathbf F_{E,X}
       \simeq Rj_*(K\otimes^{\mathbf L}_{\mathcal O_{E,U}}\mathbf F_{E,U})
\]
in $D(X_{\proet},\mathcal O_E)$, because reduction is a cone. The right side is bounded constructible by
Theorem~\ref{b:thm:finite-open}, which applies to $\mathbf F_E$ since it
is a finite free extension of $\Fell$, and the residue criterion
\stacks{09C1} gives $Rj_*K\in\Dbc(X,\mathcal O_E)$.

Proper base change on $\Loc(\Dctf(-,\mathcal O_E))$ is the integral case
of Proposition~\ref{b:prop:proper-base-change-unbounded}, using the
estimates of Proposition~\ref{i:compactness}. For the projection
formula, let $p$ be proper,
$K\in \Dbc(X,\mathcal O_E)$ and $L\in \Dbc(Y,\mathcal O_E)$. Both sides of
\[
 Rp_*K\otimes^{\mathbf L}L\longrightarrow Rp_*(K\otimes^{\mathbf L}p^*L)
 \quad\text{in }D(Y_{\proet},\mathcal O_E)
\]
are derived complete, by Lemma~\ref{i:tensor} and proper finiteness. Its
reduction modulo $\pi_E$ is the torsion projection formula map in
$D(Y_{\proet},\mathbf F_E)$, an
isomorphism by \stacks{0F0G}. The cone is derived complete with zero
reduction modulo $\pi_E$, so it vanishes by derived Nakayama
(\stacks{099I}). The categories $\Loc(\Dbc(-,\mathcal O_E))$ are closed under
$\otimes^{\mathbf L}$ by Lemma~\ref{i:tensor} and the two-variable
localizing argument, and $Rp_*$ preserves direct sums there, so the
localizing argument of Proposition~\ref{b:prop:pf-bc-generated} extends the
projection formula to all of $\Loc(\Dbc(-,\mathcal O_E))$.
\end{proof}

\begin{theorem}[$f^!$ preserves integral constructible complexes]
\label{i:restriction}
For every separated morphism of finite type $f:X\to Y$, the
right adjoint $f^!$ maps $\Dctf(Y,\mathcal O_E)$ into
$\Dctf(X,\mathcal O_E)$. Its restriction is right adjoint to
$Rf_!:\Dctf(X,\mathcal O_E)\to\Dctf(Y,\mathcal O_E)$.
\end{theorem}
\begin{proof}
\emph{Projective line and affine space.}
Proposition~\ref{b:prop:diagonal-tower} constructs the cycle class
$\delta_{\mathcal O_E}$, the trace, and the two integration identities
directly over $\mathcal O_E$, as morphisms and equalities in
$D((-)_{\proet},\mathcal O_E)$. With the projection formula and proper base change of Lemma~\ref{i:open-pf}, the
unit, counit and triangle identities in the proof of
Proposition~\ref{b:prop:p1-adjunction-rational} work over $\mathcal O_E$.
Together with the argument of Corollary~\ref{b:cor:affine-space-rational},
this gives isomorphisms of functors from $\Loc(\Dbc(Y,\mathcal O_E))$ to
$\Loc(\Dbc(\mathbf P^1_Y,\mathcal O_E))$, respectively
$\Loc(\Dbc(\mathbf A^d_Y,\mathcal O_E))$,
\[
 q^!\simeq q^*(1)[2]\quad(q:\mathbf P^1_Y\to Y),\qquad
 p^!\simeq p^*(d)[2d]\quad(p:\mathbf A^d_Y\to Y).
\]
These functors preserve $\Dbc(-,\mathcal O_E)$ by Lemma~\ref{i:tensor}.

\emph{Closed immersions.} For a closed immersion $i:Z\hookrightarrow X$,
the proof of Proposition~\ref{b:prop:closed-adjoint-constructible}
applies with $\mathcal O_E$-coefficients, using Lemma~\ref{i:open-pf} for
the finiteness of $Rj_*$: for $K\in\Dbc(X,\mathcal O_E)$,
$i^!K\simeq R\underline\Gamma_ZK\in\Dbc(Z,\mathcal O_E)$.

\emph{General $f$.} With finite affine covers and closed immersions
$U\hookrightarrow\mathbf A^d_V$ as in the proof of
Theorem~\ref{b:thm:restriction}, composition of compactly supported direct
images and uniqueness of right adjoints give $u^*f^!\simeq i^!p^!v^*$ as functors
$\Loc(\Dbc(Y,\mathcal O_E))\to\Loc(\Dbc(U,\mathcal O_E))$,
so $u^*f^!K\in\Dbc(U,\mathcal O_E)$ for each open $U$ of the cover. The
conditions defining $\Dbc(-,\mathcal O_E)$ are local for a finite open
cover. Boundedness and constructibility of the reduction are local as in
the rational case. Derived completeness is local because restriction to
an open commutes with reduction modulo $\pi_E^n$ and with $R\varprojlim$,
and restrictions to the opens of a cover are jointly conservative. Hence
$f^!K\in\Dbc(X,\mathcal O_E)$. Since the inclusions
$\Dbc(-,\mathcal O_E)\subset\Loc(\Dbc(-,\mathcal O_E))$ are full, the
adjunction of Proposition~\ref{i:brown} restricts to $\Dbc$. This
completes the proof of Theorem~\ref{n:existence} for $\Lambda=\mathcal O_E$.
\end{proof}

Thus for $\mathcal O_E$-coefficients $Rf_!$ and its right adjoint $f^!$
restrict to $\Dbc$, the last coefficient case of
Theorem~\ref{n:existence}.

\section{The six functors on constructible complexes of finite Tor amplitude}
\label{n:proof}

\subsection{Proof of the main theorem}

\begin{proof}[Proof of Theorem~\ref{n:existence}]
The functor $Rf_!\colon\Loc(\Dctf(X,\Lambda))\to\Loc(\Dctf(Y,\Lambda))$
and its right adjoint $f^!$ are constructed for all $\Lambda$ at once in Theorem~\ref{ctf:unified-brown}.
Its compactness and direct-image hypotheses are verified for finite $A$
in Propositions~\ref{ctf:finite-compact} and~\ref{ctf:finite-proper},
for fields in Theorem~\ref{b:thm:constructible-compactness} and
Proposition~\ref{b:prop:proper-image-localizing}, and for
$\mathcal O_E$ in Proposition~\ref{i:compactness}.

That $f^!$ preserves $\Dctf(-,\Lambda)$ is
Theorem~\ref{ctf:finite-restriction} for finite $A$,
Theorem~\ref{b:thm:restriction} for $E,\Qbar$, and
Theorem~\ref{i:restriction} for $\mathcal O_E$; in the integral and
rational cases $\Dctf=\Dbc$. The field proof is defined over $E$: its
diagonal class and trace are constructed over $\mathcal O_E$ before
localization and finite-extension descent. This proves all parts of the
theorem for every admissible coefficient ring.
\end{proof}

\begin{proof}[Proof of Corollary~\ref{n:existence-b}]
For proper $p$ take the trivial compactification, so $Rp_!=Rp_*$ as functors
$\Dctf(X,\Lambda)\to\Dctf(Y,\Lambda)$, and restrict the adjunction $Rf_!\dashv f^!$ to $\Dctf(-,\Lambda)$.
\end{proof}

\subsection{The remaining operations}

The operations on module sheaves, together with the adjunction
$Rf_!\dashv f^!$ of Theorem~\ref{ctf:unified-brown}, restrict to the
categories $\Dctf(-,\Lambda)$. The proofs work with perfect local models
on strata: the cohomology sheaves of a complex of finite Tor amplitude
need not have finite Tor dimension, so a d\'evissage by cohomology
sheaves would leave $\Dctf$.

\begin{proposition}[The six functors on $\Dctf$]
\label{n:six-small}
For every admissible $\Lambda\in\{A,\Zl,\mathcal O_E,\Ql,E,\Qbar\}$, the
operations $f^*$, $Rf_*$, $Rf_!$, $f^!$, $\otimes^{\mathbf L}$ and
$\cHom$ preserve $\Dctf(-,\Lambda)$ for morphisms of the standing
schemes. Their adjunctions are the restrictions of the adjunctions
$f^*\dashv Rf_*$ and $-\otimes^{\mathbf L}K\dashv\cHom(K,-)$ on
$D(-_{\proet},\Lambda)$ or of the adjunction $Rf_!\dashv f^!$ on
$\Loc(\Dctf(-,\Lambda))$.
\end{proposition}

For integral and rational coefficients these categories are the full
categories $\Dbc$, so the proposition concerns all bounded constructible
complexes in those cases.
\begin{proof}
Pullback and tensor product preserve the perfect local models on
strata, and extension by zero preserves their finite Tor intervals.
For finite $A$ these statements can be checked on the bounded flat
constructible models of Lemma~\ref{ctf:finite-models}; for
$\mathcal O_E$, $E$ and $\Qbar$ they follow from Lemma~\ref{i:tensor} and
the local models of lisse sheaves.

Proper direct image preserves $\Dctf(-,\Lambda)$ by
Proposition~\ref{b:prop:proper-coefficients}. Open direct image does
so by Proposition~\ref{ctf:finite-open} for finite $A$ and
Proposition~\ref{b:prop:open-coefficients} for $\mathcal O_E$, $E$ and $\Qbar$.
A compactification $f=pj$ gives $Rf_*=Rp_*Rj_*$ and
$Rf_!=Rp_*j_!$, proving both direct-image assertions. Preservation
by $f^!$ is the restriction part of Theorem~\ref{n:existence}.

For internal Hom, a finite stratification and finite cones reduce the
first argument to $a_!P$, where $a:S\hookrightarrow X$ is locally
closed and $P$ is a locally constant perfect complex on $S$.
For finite $A$ this follows from the bounded flat model and its
stupid filtration; for $\mathcal O_E$, $E$ and $\Qbar$ it follows from
their constructible local models. Tensor-Hom adjunction and the projection formula for $a_!$ give
\[
 \cHom_X(a_!P,K)
   \simeq Ra_*\cHom_S(P,a^!K)
   \simeq Ra_*(P^\vee\otimes^{\mathbf L}a^!K)
   \quad\text{in }D(X_{\proet},\Lambda).
\]
Here $P^\vee=\cHom_S(P,\Lambda_S)$, and the second identity holds
because $P$ is dualizable. All terms on the right
belong to $\Dctf(-,\Lambda)$ by the operations already established.
Finite cones and shifts prove the assertion for every first argument.
The full inclusions then restrict the ambient Hom bijections to the
stated adjunctions.
\end{proof}

\subsection{Change of coefficients}
\label{n:change}

A change of coefficients is derived extension of scalars along a map
of condensed rings, pulled back to $X$. For a map $\Lambda\to\Lambda'$ write
\[
 K_{\Lambda'}=\Lambda'_X\otimes^{\mathbf L}_{\Lambda_X}K.
\]
This functor $D(X_{\proet},\Lambda)\to D(X_{\proet},\Lambda')$ is left
adjoint to the exact restriction of scalars
$D(X_{\proet},\Lambda')\to D(X_{\proet},\Lambda)$.
The coefficient maps used in these notes are:
\begin{enumerate}
\item any homomorphism $A\to B$ of finite admissible rings;
\item the maps $\mathcal O_E\to\mathcal O_F$ and $E\to F$ for a
  finite extension $F/E$ inside $\Qbar$;
\item localization $\mathcal O_E\to E$, inclusions $E\to\Qbar$,
  and maps $\mathcal O_E\to A$ to finite admissible rings;
\item composites of these maps.
\end{enumerate}
Taking $E=\Ql$ gives $\Zl\to\mathcal O_F$, $\Ql\to F$, $\Zl\to\Ql$,
$\Zl\to A$ and $\Ql\to\Qbar$, so these maps connect all the admissible
coefficient rings. A map from $\mathcal O_E$ to a finite ring is
continuous, and its target is killed by a power of $\pi_E$; the
finite-level reductions of $\mathcal O_E$ are of this kind. The map
$A\to B$ in (1) need not be flat.

\begin{proposition}[The six operations and coefficient comparison]
\label{n:scalar-six}
For the listed coefficient maps, extension of scalars sends
$\Dctf(X,\Lambda)$ into $\Dctf(X,\Lambda')$.
For $f\colon X\to Y$, the canonical coefficient comparison maps for
$f^*$, $Rf_*$, $Rf_!$, $f^!$, tensor product and internal Hom, such as
$(Rf_*K)_{\Lambda'}\to Rf_*(K_{\Lambda'})$ in $\Dctf(Y,\Lambda')$ for
$K\in\Dctf(X,\Lambda)$, are isomorphisms of functors from
$\Dctf(-,\Lambda)$ to $\Dctf(-,\Lambda')$. They respect composition and the units,
counits and evaluation maps defining the comparisons.
\end{proposition}
\begin{proof}
It suffices to treat the listed elementary maps. A perfect local model
on a stratum remains perfect after derived extension of scalars.
Extension of scalars also commutes with extension by zero and with
finite cones. This proves preservation of $\Dctf(-,\Lambda)$.

\emph{Pullback and extension by zero.}
Associativity of derived pullback for ringed topoi (\stacks{06YV})
gives $(f^*K)_{\Lambda'}\simeq f^*(K_{\Lambda'})$ in
$D(X_{\proet},\Lambda')$ for $K\in D(Y_{\proet},\Lambda)$.
The same assertion for $j_!$ follows from its stalkwise description
or from the equality of the corresponding right adjoints. These
identifications respect composition.

\emph{Ordinary direct image.}
The comparison
\[
 (Rf_*K)_{\Lambda'}\longrightarrow Rf_*(K_{\Lambda'})
 \quad\text{in }D(Y_{\proet},\Lambda')
\]
is formed from the unit, the pullback identification and the counit
of $f^*\dashv Rf_*$ on $D(-_{\proet},\Lambda')$. For finite $A\to B$, calculate on the \'etale
site using Proposition~\ref{ctf:finite-comparison}. On a
quasi-compact \'etale test object $V\to Y$, the scalar projection
formula (\stacks{0F12}) on both $V$ and $X_V=X\times_YV$ gives
\[
 R\Gamma(V,(Rf_*K)_B)
  \simeq R\Gamma(X_V,K)\otimes_A^{\mathbf L}B
  \simeq R\Gamma(X_V,K_B).
\]
Thus the comparison induces an isomorphism on derived sections over
every such $V$. The scalar projection formula needs only cohomological
bounds for $\ell$-primary coefficients, which hold for the standing
schemes, so it applies to every $B$, whatever its projective dimension
over $A$. These isomorphisms commute with restriction in $V$, so they
identify the canonical morphism in $D(Y_{\et},B)$, and hence, by
Proposition~\ref{ctf:finite-comparison}, in $D(Y_{\proet},B)$.

For $\mathcal O_E\to A$ with $A$ finite, $A$ has a finite free
$\mathcal O_E$-resolution of length at most one. Its condensed
version is exact (a continuous map to a finite quotient lifts on a
finite clopen partition), and pullback gives a finite free resolution
of $A_X$ over $\mathcal O_{E,X}$. Tensoring with $A_X$ is consequently
computed by finite sums and a cone, which commute with the exact
derived functor $Rf_*$. For $\mathcal O_E\to\mathcal O_F$ and
$E\to F$ the coefficient extension is finite free, so the assertion
follows from finite sums. For $\mathcal O_E\to E$, use the system
$K\xrightarrow{\ell}K\xrightarrow{\ell}\cdots$; for $E\to\Qbar$,
use the filtered system of finite field extensions. Since $K$ is
bounded and these systems have a common lower bound,
Lemma~\ref{b:lem:bounded-continuity} commutes their colimits with
$Rf_*$. This proves the direct-image comparison in all cases.

\emph{Compactly supported direct image.}
For $f=pj$, combine the preceding comparison for $Rp_*$ with the
one for $j_!$. Both commute with the refinement maps of the
compactification construction: they are formed from the same
pullback identifications and adjunction maps, and restrict to the
identity on the open subscheme. Thus
\[
 (Rf_!K)_{\Lambda'}\simeq Rf_!(K_{\Lambda'})
 \quad\text{in }\Dctf(Y,\Lambda')
\]
is independent of the compactification and respects composition.

\emph{Exceptional inverse image.}
For $K\in\Dctf(Y,\Lambda)$ define
\[
 \beta_{f,K}:(f^!K)_{\Lambda'}\longrightarrow f^!(K_{\Lambda'})
\]
in $\Dctf(X,\Lambda')$ as the transpose, under the adjunction
$Rf_!\dashv f^!$ on $\Loc(\Dctf(-,\Lambda'))$, of the morphism in
$\Dctf(Y,\Lambda')$
\[
 Rf_!((f^!K)_{\Lambda'})
 \xrightarrow{\sim}(Rf_!f^!K)_{\Lambda'}
 \xrightarrow{(\epsilon_K)_{\Lambda'}}K_{\Lambda'}.
\]
The first arrow is the comparison for $Rf_!$, which applies because
$f^!K\in\Dctf(X,\Lambda)$ by Theorem~\ref{n:existence}. The counit of a composite
is the composite of its counits, so these maps respect composition
and open restriction.

For a closed immersion $i$ with open complement $j$, use the triangle
in $\Dctf(Y,\Lambda)$
\[
 i_*i^!K\longrightarrow K\longrightarrow Rj_*j^*K.
\]
Extension of scalars commutes with $i_*$, $j^*$ and, by the preceding
argument, $Rj_*$. It therefore carries this triangle to the triangle
over $\Lambda'$. The fibre comparison is $i_*\beta_{i,K}$ because
the first map is the counit; hence $\beta_{i,K}$ is an isomorphism.
For $q:\mathbf P^1_Y\to Y$, the diagonal class and trace defining
the adjunction $Rq_*\dashv q^*(1)[2]$ between $\Dctf(\mathbf P^1_Y,\Lambda)$
and $\Dctf(Y,\Lambda)$ commute with all the listed coefficient maps:
the finite classes are constructed from Kummer and derived scalar
extension, the integral classes are their compatible limits, and
the field classes are their scalar extensions
(Propositions~\ref{b:prop:geo-p1-adjunction}
and~\ref{b:prop:diagonal-tower}). For $\mathcal O_E\to A$, the
perfectness of $A_X$ proved above makes its derived tensor commute
with the inverse limits defining the integral classes; the resulting
finite-level classes are exactly the finite coefficient classes.
Thus $\beta_{q,K}$ is the usual
pullback-and-twist identification. Open restriction and composition
give the affine-space case. A finite affine cover and local
factorization as a closed immersion into affine space prove the
assertion for every $f$.

\emph{Tensor product and internal Hom.}
Derived extension of scalars
$D(X_{\proet},\Lambda)\to D(X_{\proet},\Lambda')$ is symmetric monoidal,
by associativity of tensor product. For internal Hom, reduce the first argument by
perfect strata and finite cones to $a_!P$, as in the preceding
proposition. The identity
\[
 \cHom_X(a_!P,K)
    \simeq Ra_*(P^\vee\otimes^{\mathbf L}a^!K)
\]
commutes with extension of scalars by the comparisons for $Ra_*$,
$a^!$, and the dual of a perfect complex. These are the canonical
comparison maps, since all are defined by the same evaluation
pairing. Finite cones give the result for every first argument.
\end{proof}

Finite Tor amplitude is what makes the statement uniform in the
coefficients. On a geometric point $X$, the sheaf $\mathbf F_\ell$
belongs to $\Dbc(X,\mathbf Z/\ell^2)$ but not to
$\Dctf(X,\mathbf Z/\ell^2)$, and its extension of scalars
$\mathbf F_\ell\otimes^{\mathbf L}_{\mathbf Z/\ell^2}\mathbf F_\ell$ is
unbounded. So $\Dbc$ is not preserved by arbitrary finite-ring change of
coefficients.

For rational coefficients $\Dctf=\Dbc$. With $\Qbar$-coefficients,
Lemma~\ref{b:lem:coefficient-finite-descent} assembles $\Dctf(X,\Qbar)$
from the categories $\Dctf(X,E)$ for finite $E$; as a statement about
ordinary categories this reads as follows.

\begin{proposition}[$\Dctf(X,\Qbar)$ as a $2$-colimit]
\label{n:two-colimit}
Extension of scalars identifies
\[
 \Dctf(X,\Qbar)\simeq\mathop{\text{2-colim}}_{E\subset\Qbar}\Dctf(X,E)
\]
as ordinary categories: every object of $\Dctf(X,\Qbar)$ is isomorphic to
$K_{\Qbar}$ for some finite $E$ and some $K\in\Dctf(X,E)$, and for
$K,L\in\Dctf(X,E)$ extension of scalars induces a bijection
\[
 \varinjlim_{F\supset E}\Hom_{\Dctf(X,F)}(K_F,L_F)
       \xrightarrow{\sim}\Hom_{\Dctf(X,\Qbar)}(K_{\Qbar},L_{\Qbar}),
\]
the colimit running over finite extensions $F$ of $E$ in $\Qbar$.
\end{proposition}
\begin{proof}
The statement about objects is
Lemma~\ref{b:lem:coefficient-finite-descent}(3). Its proof descends the
lisse sheaves on the strata, together with their transition matrices,
and then the connecting morphisms of a finite Postnikov tower, so it
descends the complexes themselves. The bijection
on morphisms is Lemma~\ref{b:lem:coefficient-finite-descent}(2) with
$r=0$, since $\Dctf$ is a full subcategory; it rests on the continuity of
Ext out of a constructible complex
(Lemma~\ref{b:lem:constructible-ext-continuity}). Surjectivity says that
every morphism descends to a finite stage, and injectivity that two
morphisms with the same image agree at a later finite stage. Together
with descent of objects, these say exactly that the functor from the
category with objects the pairs $(E,K)$, $K\in\Dctf(X,E)$, and Hom sets
these colimits, to $\Dctf(X,\Qbar)$ is fully faithful and essentially
surjective.
\end{proof}

The proof of Proposition~\ref{n:two-colimit} uses only
Lemmas~\ref{b:lem:coefficient-finite-descent}
and~\ref{b:lem:constructible-ext-continuity}, which precede the uniform
Ext bound and Brown representability; so the proposition is available
in the proof of Theorem~\ref{n:existence}.

\section{Smooth purity}\label{sec:smooth-purity}

Let $f:X\to Y$ be smooth of pure relative dimension $d$. Smooth purity
says that the right adjoint of $Rf_!$ is inverse image followed by the
Tate twist $(d)$ and the shift $[2d]$. With $\Qbar$-coefficients it
supplies the trace maps of smooth morphisms used in the Weil
conjectures.

Throughout this section $f^!$ is the right adjoint
$\Loc(\Dctf(Y,\Lambda))\to\Loc(\Dctf(X,\Lambda))$ of
Theorem~\ref{ctf:unified-brown}, restricted to a functor
$\Dctf(Y,\Lambda)\to\Dctf(X,\Lambda)$ by Theorem~\ref{n:existence}. The projective line and affine space compute
this adjoint and normalize its unit and counit.

\subsection{Statement and set-up}
\label{sp:grp-setup}

\subsubsection{Coefficients}

Tate twists are those of Definition~\ref{def:sp-Tate-twist}, and
$\Lambda$ is any admissible coefficient ring:
a finite commutative ring killed by a power of $\ell$, $\mathcal O_E$,
$E$ or $\Qbar$. Every coefficient sheaf is pulled back from its
condensed ring, and a finite discrete ring gives the constant sheaf.
For the integral and field coefficients $\Dctf=\Dbc$; we retain the
notation $\Dbc$ in calculations confined to those cases. The towers of
\S\ref{sp:grp-integral} present derived complete objects with these
coefficients. The proof runs through finite rings $A$ with $\ell^rA=0$,
such as $\mathcal O_E/\ell^n$, and for these $\Dctf$ is the category of
Definition~\ref{def:sp-finite-category}; its Tor condition keeps derived
change of coefficients bounded, as the example after
Proposition~\ref{n:scalar-six} shows. On the pro-\'etale site $\Dctf(T,A)$ is identified with the pullback
of its \'etale counterpart by Corollary~\ref{pe:finite-tor-comparison}.

With integral coefficients no extra condition is needed. By
Lemma~\ref{i:integral-heart}, an object $K$ of $\Dbc(T,\mathcal O_E)$ is a
finite extension of shifted sheaves $a_!C$ with $C$ of Tor amplitude
$[-1,0]$ on a stratum, so $K$ has finite Tor amplitude and its reductions
$K/\ell^n$ lie in $\Dctf(T_{\proet},\mathcal O_E/\ell^n)$, with bounds
independent of $n$. Conversely, if $K$ is derived complete and $K/\ell$
lies in $\Dctf(T_{\proet},\mathcal O_E/\ell)$, then
$K/\pi_E=(K/\ell)\otimes^{\mathbf L}_{\mathcal O_E/\ell}\mathcal O_E/\pi_E$
is bounded constructible, so $K\in\Dbc(T,\mathcal O_E)$. Moreover inverse
image preserves $\Dbc(-,\mathcal O_E)$ and needs no completion
(Lemma~\ref{i:tensor}).

When coefficient rings are compared, a functor carries its ring as a
subscript, as in $Rf_{!,A}$ or $f^!_{\mathcal O_E}$; otherwise the
subscript is dropped. With coefficients $A$, $Rf_!=R\bar f_*\circ j_!$ as functors
$\Dctf(X_{\proet},A)\to\Dctf(Y_{\proet},A)$ for
a compactification $f=\bar f\circ j$, and on $\Dctf$ these functors and
their right adjoints are the pro-\'etale functors constructed in
Theorem~\ref{ctf:unified-brown} and restricted in
Theorem~\ref{ctf:finite-restriction}. The geometric identities of
\S\ref{b:sec:finite-geometric} give proper base change, the
projection formula and the localization triangle; their proofs may
calculate on the auxiliary \'etale sites via comparison. For \'etale
$e:U\to T$, extension by zero is exact on sheaves of $A$-modules on the
pro-\'etale sites, and $e_!\colon D(U_{\proet},A)\to D(T_{\proet},A)$ is
left adjoint to $e^*\colon D(T_{\proet},A)\to D(U_{\proet},A)$; on
geometric stalks its counit is the sum over the fibre.

\subsubsection{The purity statement}

\begin{theorem}[Smooth purity for all admissible coefficients]
\label{thm:sp-integral}
Let $\Lambda\in\{A,\Zl,\mathcal O_E,\Ql,E,\Qbar\}$ be any admissible
coefficient ring and let $f:X\to Y$ be smooth of pure relative dimension $d$.
There is an isomorphism of functors
$\Dctf(Y,\Lambda)\to\Dctf(X,\Lambda)$, given at $K\in\Dctf(Y,\Lambda)$
by a canonical isomorphism in $\Dctf(X,\Lambda)$
\begin{equation}\label{eq:sp-main}
 \gamma_{f,\Lambda}(K):
    f_\Lambda^{*}K(d)[2d]\xrightarrow{\sim}f_{\Lambda}^{!}K.
\end{equation}
Here $f^!_\Lambda$ is the restriction of the right adjoint from
Theorem~\ref{ctf:unified-brown}. The isomorphism is compatible with
change of coefficients, base change and composition, with the
normalization given by integration on the affine line.
\end{theorem}

For \'etale $f$ we have $d=0$, and the theorem says
$f_{\Lambda}^{!}\simeq f_\Lambda^{*}$ for every admissible $\Lambda$. For a smooth curve $C$
with structure map $a:C\to\operatorname{Spec}k$ it gives
$a_{\Zl}^!\mathbf Z_{\ell,\operatorname{Spec}k}\simeq\mathbf Z_{\ell,C}(1)[2]$,
a single sheaf in degree $-2$.

\begin{definition}[Relative dualizing complex]\label{def:sp-dualizing}
For smooth $f:X\to Y$ and any admissible $\Lambda$, the
\emph{relative dualizing complex} is
\begin{equation}\label{eq:sp-omega}
 \omega_{f,\Lambda}=f_\Lambda^{!}\Lambda_Y
       \in\Dctf(X,\Lambda).
\end{equation}
\end{definition}

For finite coefficients $\omega_{f,\Lambda}$ can be computed on the
pro-\'etale site or, equivalently, on the \'etale site
(Proposition~\ref{prop:sp-et-proet}).

\paragraph*{Plan of the proof.}
With finite coefficients, the right adjoint $f^!$ of Theorem~\ref{ctf:unified-brown} is computed
on coordinate charts by the explicit projective-line adjunction, whose unit
is the class of the diagonal and whose counit is integration along the
fibres. The projection formula reduces $f^!K$ to
$f^*K\otimes^{\mathbf L}\omega_f$, and deforming the diagonal to the
zero section of the relative tangent bundle $T_{X/Y}$ identifies
$\omega_f$ with $(\mathbf Z/\ell^n\mathbf Z)_X(d)[2d]$ by an isomorphism
defined globally on $X$. To pass to $\mathcal O_E$ these constructions must be compatible at
every finite level, so that they form a single inverse system (a
coefficient tower). Its derived limit computes the integral Brown
adjoint with its normalized purity map. For field coefficients, objects
and Hom groups are compared with integral models to transport this
normalization to the existing adjunction.

\subsection{Finite coefficients}
\label{sp:grp-finite}
We first write the geometric calculation with
$\mathbf Z/\ell^n\mathbf Z$ coefficients, for a fixed $n\geq1$.
Until \S\ref{sp:grp-integral} the subscript is omitted from functors,
dualizing complexes and purity maps in this calculation:
$f^*$, $Rf_!$, $f^!$, $\omega_f$, $\gamma_f$ refer to this cyclic
coefficient ring. Proposition~\ref{prop:sp-finite} records the result
for every finite admissible ring $A$: the same calculation extends
along $\mathbf Z/\ell^n\mathbf Z\to A$ when $\ell^nA=0$.

The functor $f^!$ and its local form for smooth $f$ were obtained in
\S\ref{b:sec:finite-geometric} with $\Lambda=\mathbf Z/\ell^n\mathbf Z$,
and the projective-line adjunction was proved in
Proposition~\ref{b:prop:geo-p1-adjunction} and lifted to $\ell$-adic
coefficients in Proposition~\ref{b:prop:diagonal-tower}. We record the
explicit unit and counit with their normalization, since the passage to
towers in \S\ref{sp:grp-integral} repeats them level by level. What is
new here is the coordinate-free normalization: integration on vector
bundles and the orientation of $\omega_f$ from the diagonal.

\subsubsection{Integration on affine space}

Chern classes come from the Kummer sequence on $T_{\et}$,
\begin{equation}\label{eq:sp-kummer}
 1\longrightarrow\mu_{\ell^n}\longrightarrow\mathbb G_m
     \xrightarrow{\ (\cdot)^{\ell^n}\ }\mathbb G_m
     \longrightarrow1,
\end{equation}
whose boundary $H^1(T_{\et},\mathbb G_m)\to
H^2(T_{\et},(\mathbf Z/\ell^n\mathbf Z)_T(1))$ sends a line bundle $L$
to $c_1(L)$. Let $q:\mathbf P^1_S\to S$ be the projection. By
Lemma~\ref{b:lem:geo-p1-cohomology}, the unit and
$c_1(\mathcal O_{\mathbf P^1_S}(1))$ give an isomorphism in
$D(S_{\proet},\mathbf Z/\ell^n\mathbf Z)$, compatible with base change,
\begin{equation}\label{eq:sp-p1-splitting}
 Rq_*(\mathbf Z/\ell^n\mathbf Z)_{\mathbf P^1_S}
       \simeq(\mathbf Z/\ell^n\mathbf Z)_{S}\oplus(\mathbf Z/\ell^n\mathbf Z)_{S}(-1)[-2],
\end{equation}
and the \emph{trace map}
$\operatorname{Tr}_q:Rq_*(\mathbf Z/\ell^n\mathbf Z)_{\mathbf P^1_S}(1)[2]
\to(\mathbf Z/\ell^n\mathbf Z)_{S}$, a morphism in
$D(S_{\proet},\mathbf Z/\ell^n\mathbf Z)$, is the projection onto the twisted
second summand. It sends $c_1(\mathcal O_{\mathbf P^1_{\bar s}}(1))$ on
each geometric fibre to $1$. Let $\Delta$ be the diagonal of
$\mathbf P^1_S\times_S\mathbf P^1_S$, cut out by $x_0y_1-x_1y_0$, and
$p_1,p_2$ the projections. Lemma~\ref{b:lem:geo-cartier} gives its cycle
class, a morphism in
$D((\mathbf P^1_S\times_S\mathbf P^1_S)_{\proet},\mathbf Z/\ell^n\mathbf Z)$,
\[
 \delta:\Delta_*(\mathbf Z/\ell^n\mathbf Z)_{\mathbf P^1_S}
      \longrightarrow(\mathbf Z/\ell^n\mathbf Z)_{\mathbf P^1_S\times_S\mathbf P^1_S}(1)[2],
\]
the Kummer boundary of $\mathcal O(\Delta)$ with its canonical section,
compatible with every base change on $S$. With
$\operatorname{Tr}_{p_i}$ the pullback of $\operatorname{Tr}_q$, the
diagonal integrates to $1$ along either factor:
\begin{equation}\label{eq:sp-diagonal-integrals}
 \operatorname{Tr}_{p_i}\circ Rp_{i*}\delta
      =\operatorname{id}_{(\mathbf Z/\ell^n\mathbf Z)_{\mathbf P^1_S}}
 \qquad(i=1,2),
\end{equation}
since on a geometric fibre of $p_i$ the diagonal is one rational point,
whose line bundle has degree one (this is (G5) of
Proposition~\ref{b:prop:geo-p1-adjunction}).

\begin{lemma}[The projective-line adjunction]\label{lem:sp-p1}
The functor
$Rq_*\colon\Dctf(\mathbf P^1_{S,\proet},\mathbf Z/\ell^n\mathbf Z)\to
\Dctf(S_{\proet},\mathbf Z/\ell^n\mathbf Z)$ is left adjoint to $K\mapsto q^{*}K(1)[2]$. The
counit is integration along the fibres,
\[
 \varepsilon_K:Rq_*\bigl(q^{*}K(1)[2]\bigr)
   \simeq K\otimes^{\mathbf L} Rq_*(\mathbf Z/\ell^n\mathbf Z)_{\mathbf P^1_S}(1)[2]
   \xrightarrow{\ \operatorname{id}_K\otimes\operatorname{Tr}_q\ } K,
\]
and the unit is the class of the diagonal,
\[
 \eta_M:\ M\simeq Rp_{2*}\Delta_*M
  \xrightarrow{\ Rp_{2*}(\delta\otimes p_1^*M)\ } Rp_{2*}p_{1}^{*}M(1)[2]
  \simeq q^{*}Rq_*M(1)[2],
\]
using the projection formulas and proper base change. Both commute with
base change on $S$.
\end{lemma}
\begin{proof}
This is the pro-\'etale adjunction of
Proposition~\ref{b:prop:geo-p1-adjunction},
with unit (G6) and counit (G7) as displayed; the triangle identities
reduce to \eqref{eq:sp-diagonal-integrals}.
\end{proof}

Let $j:\mathbf A^1_S\hookrightarrow\mathbf P^1_S$ be the complement of
the section at infinity $i_\infty$, and $p=q\circ j$. Composing
Lemma~\ref{lem:sp-p1} with $j_!\dashv j^*$ gives
$Rp_!\dashv p^*(-)(1)[2]$, with counit the open-immersion counit
followed by $\varepsilon_K$. At the constant sheaf, $Rq_*$ of the
localization triangle for $i_\infty$ shows that
$Rp_!(\mathbf Z/\ell^n\mathbf Z)_{\mathbf A^1_S}\simeq
(\mathbf Z/\ell^n\mathbf Z)_S(-1)[-2]$ is the second summand of
\eqref{eq:sp-p1-splitting}, because restriction to infinity kills
$c_1(\mathcal O(1))$. So the counit at the constant sheaf, the
\emph{trace map of the affine line}
$\operatorname{Tr}_p:Rp_!(\mathbf Z/\ell^n\mathbf Z)_{\mathbf A^1_S}(1)[2]
\to(\mathbf Z/\ell^n\mathbf Z)_S$, is an isomorphism in
$D(S_{\proet},\mathbf Z/\ell^n\mathbf Z)$. Over a geometric
point this is $H^2_c(\mathbf A^1,(\mathbf Z/\ell^n\mathbf Z)(1))\simeq
\mathbf Z/\ell^n\mathbf Z$, generated by the class $\kappa$ that maps to
the class of a point in $H^2(\mathbf P^1,(\mathbf Z/\ell^n\mathbf Z)(1))$.

For affine space, let $p_j:\mathbf A^j_S\to S$ be the projection and
$r_j:\mathbf A^j_S\to\mathbf A^{j-1}_S$ forget the last coordinate.
Define $\operatorname{Tr}_{p_j}$ inductively as the composite
\[
 \begin{aligned}
 Rp_{j!}(\mathbf Z/\ell^n\mathbf Z)_{\mathbf A^j_S}(j)[2j]
 &\simeq Rp_{j-1!}\bigl(Rr_{j!}(\mathbf Z/\ell^n\mathbf Z)_{\mathbf A^j_S}(1)[2]\bigr)(j-1)[2j-2]\\
 &\xrightarrow{\ \operatorname{Tr}_{p_{j-1}}\circ Rp_{j-1!}(\operatorname{Tr}_{r_j})\ }
 (\mathbf Z/\ell^n\mathbf Z)_S,
 \end{aligned}
\]
using $Rp_{j!}=Rp_{j-1!}Rr_{j!}$ and the projection formula. For $j=d$
this is the \emph{integration isomorphism} (the identity when $d=0$)
\begin{equation}\label{eq:sp-affine-integration}
 \operatorname{Tr}_{p_d}:Rp_{d!}(\mathbf Z/\ell^n\mathbf Z)_{\mathbf A^d_S}(d)[2d]
       \xrightarrow{\sim}(\mathbf Z/\ell^n\mathbf Z)_{S}.
\end{equation}
By Lemma~\ref{b:lem:geo-affine-local}(1), $Rp_{d!}\dashv
p_d^*(-)(d)[2d]$, with counit the composite of the $d$ affine-line
counits, which at the constant sheaf is
\eqref{eq:sp-affine-integration}; unit and counit commute with base
change. On a geometric fibre the external product
$\kappa\boxtimes\cdots\boxtimes\kappa\in
H_c^{2d}(\mathbf A^d_{\bar s},(\mathbf Z/\ell^n\mathbf Z)(d))$
integrates to $1$, since by the projection formula integrating along the
last coordinate replaces the last factor by its integral $1$.

\begin{definition}[Representing objects]\label{def:sp-representing-pair}
Let $f$ be smooth, and work in $D(Y,\mathbf Z/\ell^n\mathbf Z)$ and
$D(X,\mathbf Z/\ell^n\mathbf Z)$ on either the \'etale or the
pro-\'etale sites. For $K$ on $Y$, a \emph{representing pair} for $K$
is an object $H$ on $X$ with a morphism $\varepsilon:Rf_!H\to K$ such
that $b\mapsto\varepsilon\circ Rf_!(b)$ is a bijection
$\operatorname{Hom}(M,H)\to\operatorname{Hom}(Rf_!M,K)$ for all
$M\in\Dctf(X,\mathbf Z/\ell^n\mathbf Z)$.
\end{definition}

The finite geometric computations of
\S\ref{b:sec:finite-geometric} are stated on the pro-\'etale site
and proved by auxiliary \'etale calculations. We recall the comparison
used in those calculations, including its compatibility with units
and counits.

\begin{proposition}[\'Etale and pro-\'etale finite coefficients]
\label{prop:sp-et-proet}
For finite admissible $A$, the category
$\Dctf(T,A)\subset D(T_{\proet},A)$ is equivalent to the category
$\Dctf(T_{\et},A)$ of constructible complexes of finite Tor amplitude on
the \'etale site. The equivalence respects pullback, qcqs derived
direct image and extension by zero, and hence $Rf_!$.
On the categories $\Loc(\Dctf(-,A))$ these comparisons are compatible with
the adjunction $Rf_!\dashv f^!$ of Theorem~\ref{ctf:unified-brown}.
\end{proposition}
\begin{proof}
The first assertion is Corollary~\ref{pe:finite-tor-comparison} and
Theorem~\ref{pe:direct-image-comparison}, with
\eqref{pe:eq:jshriek-comparison}. The localizing extension is
Proposition~\ref{ctf:finite-comparison}.
\end{proof}

So an adjunction between $Rf_!$ and some functor $G$ on the \'etale
categories $\Dctf(-_{\et},A)$ transfers to the pro-\'etale categories
$\Dctf(-_{\proet},A)$. The Hom groups agree by full faithfulness, and
$Rf_!$ corresponds to $Rf_!$.

\begin{lemma}[Local computation of the finite Brown adjoint]
\label{lem:sp-finite-existence}
Let $f:X\to Y$ be smooth of pure relative dimension $d$.
On $\Dctf(-_{\proet},\mathbf Z/\ell^n\mathbf Z)$, the Brown
adjunction $Rf_!\dashv f^!$ has the following local form.
On an open $U\subset X$ with
$f|_U=p_d\circ e$, where $e:U\to\mathbf A^d_Y$ is \'etale,
\[
 f^!K|_U\simeq e^*p_{d}^*K(d)[2d]
 \quad\text{in }\Dctf(U_{\proet},\mathbf Z/\ell^n\mathbf Z),
\]
with counit the sum over the fibres of $e$ followed by
\eqref{eq:sp-affine-integration} tensored with $K$.
\end{lemma}
\begin{proof}
The adjunction is that of
Theorem~\ref{ctf:unified-brown} and restricts to $\Dctf$ by
Theorem~\ref{ctf:finite-restriction}. By (G10) of
Lemma~\ref{b:lem:geo-affine-local}, its restriction to $U$ is $e^!p_d^!$,
which is $e^*p_d^*(-)(d)[2d]$ by the same lemma, with the composite
counit. Proposition~\ref{prop:sp-et-proet} transfers the adjunction to the
pro-\'etale sites.
\end{proof}

These chart counits agree with the Brown counit on overlaps by
uniqueness of the adjunction comparison. For later use with towers,
recall the open-cover triangle in $\Dctf(X_{\proet},\mathbf Z/\ell^n\mathbf Z)$
for $X=U\cup V$ and $N\in\Dctf(X_{\proet},\mathbf Z/\ell^n\mathbf Z)$:
\begin{equation}\label{eq:sp-open-cover}
 j_{U\cap V!}(N|_{U\cap V})
 \xrightarrow{\ (1,-1)\ }
 j_{U!}(N|_U)\oplus j_{V!}(N|_V)
 \xrightarrow{\ (1,1)\ }N\longrightarrow{} .
\end{equation}
Its Hom sequences verify the representing property of a glued
chart model, and the counit-preserving comparison identifies that model
with the already constructed $f^!$. This gluing is used below to realize
the compatible finite counits as a single map of towers.

On a vector bundle, \eqref{eq:sp-affine-integration} gives an
integration map after choosing a local frame. To get a global map one
must show that it does not depend on the frame.

\begin{proposition}[Vector-bundle purity]\label{prop:sp-vector}
Let $\pi:V\to S$ be a vector bundle of rank $d$.
For $K\in\Dctf(S_{\proet},\mathbf Z/\ell^n\mathbf Z)$, integration induces an
isomorphism in $\Dctf(V_{\proet},\mathbf Z/\ell^n\mathbf Z)$
\[
 \pi^{*}K(d)[2d]\xrightarrow{\sim}\pi^{!}K.
\]
\end{proposition}
\begin{proof}
\emph{Local integration maps.} Over an open $U\subset S$ on which $V$
is trivial, a frame $\tau:V\times_S U\xrightarrow{\sim}\mathbf A^d_U$
transports \eqref{eq:sp-affine-integration} to an isomorphism in
$D(U_{\proet},\mathbf Z/\ell^n\mathbf Z)$
\[
 \operatorname{Tr}_{\tau}:
 R(\pi|_{V\times_S U})_!(\mathbf Z/\ell^n\mathbf Z)_{V\times_S U}(d)[2d]
       \xrightarrow{\sim}(\mathbf Z/\ell^n\mathbf Z)_U.
\]
These maps commute with base change in $U$: the construction uses
only $c_1(\mathcal O_{\mathbf P^1_U}(1))$, the splitting
\eqref{eq:sp-p1-splitting} and localization at infinity, and pullback
preserves the Kummer class while proper base change and localization
identify the pulled-back constructions with those over the new base.

\emph{Independence of the frame.} The key point is that a linear change
of coordinates acts trivially on integration, because $\GL_d$ has
connected geometric fibres. Let
\[
 a:\mathbf A^d_{\GL_{d,S}}\longrightarrow\mathbf A^d_{\GL_{d,S}},
       \qquad (g,v)\longmapsto(g,gv),
\]
be the universal linear change of coordinates, an automorphism over
$\GL_{d,S}$. It induces an automorphism in
$D(\GL_{d,S,\proet},\mathbf Z/\ell^n\mathbf Z)$ of the compactly supported
direct image of $(\mathbf Z/\ell^n\mathbf Z)_{\mathbf A^d_{\GL_{d,S}}}(d)[2d]$
to $\GL_{d,S}$, which integration identifies with the constant sheaf
$(\mathbf Z/\ell^n\mathbf Z)_{\GL_{d,S}}$ in degree zero. So this
automorphism is multiplication by a section
\[
 u_a\in H^0\bigl((\GL_{d,S})_{\et},
                    (\mathbf Z/\ell^n\mathbf Z)_{\GL_{d,S}}^{\times}\bigr),
\]
a locally constant function with values in
$(\mathbf Z/\ell^n\mathbf Z)^{\times}$. For a geometric point
$\bar s\to S$, the fibre $\GL_{d,\bar s}$ is connected (for $d\geq1$ it
is a nonempty open subset of affine space of $d\times d$ matrices, for
$d=0$ a point), so $u_a$ is constant on it; at the identity matrix $a$
is the identity, so the constant is $1$. Hence $u_a=1$.

If $\tau,\tau'$ are two frames over $U$, then $\tau'\circ\tau^{-1}$ is
the pullback of $a$ along the transition matrix $U\to\GL_{d,S}$. By
base change, $u_a=1$ gives $\operatorname{Tr}_{\tau'}=\operatorname{Tr}_{\tau}$.
Since source and target are, locally, sheaves in degree zero, the
local maps glue over a trivializing cover to an isomorphism in
$D(S_{\proet},\mathbf Z/\ell^n\mathbf Z)$
\[
 \operatorname{Tr}_{\pi}:
 R\pi_!(\mathbf Z/\ell^n\mathbf Z)_V(d)[2d]
       \xrightarrow{\sim}(\mathbf Z/\ell^n\mathbf Z)_S.
\]

\emph{The purity map.} For general $K$, the projection formula and
$\operatorname{Tr}_\pi$ give
\[
 R\pi_!\bigl(\pi^{*}K(d)[2d]\bigr)
 \simeq K\otimes^{\mathbf L} R\pi_!(\mathbf Z/\ell^n\mathbf Z)_V(d)[2d]
 \xrightarrow{\ \mathrm{id}_K\otimes\operatorname{Tr}_{\pi}\ }K.
\]
Its adjoint under $R\pi_!\dashv\pi^{!}$
(Lemma~\ref{lem:sp-finite-existence}) is the asserted map. Over each
trivializing open it is the affine-space adjunction isomorphism
following \eqref{eq:sp-affine-integration}, so its cone vanishes on an
open cover and the map is an isomorphism.
\end{proof}

The Kummer classes, projective-line integration and vector-bundle
integration all commute with base change on $S$, by the same arguments.

\subsubsection{The tensor formula and local coordinates}

Purity has two parts. The local part says that
$f^!K\simeq f^*K\otimes^{\mathbf L}\omega_f$ and that $\omega_f$ is
isomorphic to $(\mathbf Z/\ell^n\mathbf Z)(d)[2d]$ on each coordinate
chart; this, with base change, was proved in
\S\ref{b:sec:finite-geometric}. It leaves open whether the chart
isomorphisms agree on overlaps. The global part is an orientation of
$\omega_f$ defined without charts. We construct it from the diagonal in
the next subsubsection and check that it restricts to the chart
isomorphisms.

Let $f:X\to Y$ be smooth of pure relative dimension $d$ and
$K\in\Dctf(Y_{\proet},\mathbf Z/\ell^n\mathbf Z)$. The projection
formula and the counit $Rf_!\omega_{f}\to(\mathbf Z/\ell^n\mathbf Z)_{Y}$
give $Rf_!\bigl(f^{*}K\otimes^{\mathbf L}\omega_{f}\bigr)
\simeq K\otimes^{\mathbf L} Rf_!\omega_{f}\to K$, and its adjoint is
the \emph{tensor map}, a natural transformation of functors
$\Dctf(Y_{\proet},\mathbf Z/\ell^n\mathbf Z)\to
\Dctf(X_{\proet},\mathbf Z/\ell^n\mathbf Z)$
\begin{equation}\label{eq:sp-tensor-map}
 f^{*}K\otimes^{\mathbf L}_{(\mathbf Z/\ell^n\mathbf Z)_{X}}\omega_{f}
      \longrightarrow f^{!}K.
\end{equation}

\begin{proposition}[The tensor formula]\label{prop:sp-tensor}
For $K\in\Dctf(Y_{\proet},\mathbf Z/\ell^n\mathbf Z)$, the tensor map
\eqref{eq:sp-tensor-map} is an isomorphism. On a chart $f|_U=p_d\circ e$
as in Lemma~\ref{lem:sp-finite-existence},
$\omega_f|_U\simeq(\mathbf Z/\ell^n\mathbf Z)_U(d)[2d]$, so
$\omega_f[-2d]$ is a locally constant sheaf of free
$\mathbf Z/\ell^n\mathbf Z$-modules of rank one.
\end{proposition}
\begin{proof}
This is (G11) and (G15) of
Proposition~\ref{b:prop:geo-smooth-tensor}.
\end{proof}

The orientation compares $\omega$ on the fibres of a family of
sections, and this comparison rests on base change. Consider a Cartesian square
\begin{equation}\label{eq:sp-cartesian}
\begin{CD}
 X' @>{u}>> X\\
 @V{f'}VV @VV{f}V\\
 Y' @>{v}>> Y,
\end{CD}
\end{equation}
with $f$ smooth of pure relative dimension $d$. For
$K\in\Dctf(Y_{\proet},\mathbf Z/\ell^n\mathbf Z)$, pulling back the
counit $Rf_!f^{!}K\to K$ along $v$ and using compact-support base change
gives a morphism $R(f')_!u^{*}f^{!}K\to v^{*}K$ in
$\Dctf(Y'_{\proet},\mathbf Z/\ell^n\mathbf Z)$. Its adjoint under
$R(f')_!\dashv(f')^!$ is the \emph{exceptional base-change map}
$u^*f^!K\to(f')^!v^*K$, a morphism in
$\Dctf(X'_{\proet},\mathbf Z/\ell^n\mathbf Z)$.

\begin{proposition}[Exceptional base change for a smooth map]
\label{prop:sp-exceptional-bc}
For the square \eqref{eq:sp-cartesian} and
$K\in\Dctf(Y_{\proet},\mathbf Z/\ell^n\mathbf Z)$, the exceptional
base-change map is an isomorphism, compatible with the tensor maps for
$f$ and $f'$:
\begin{equation}\label{eq:sp-exceptional-bc}
 u^{*}f^{!}K
       \xrightarrow{\sim}(f')^{!}v^{*}K.
\end{equation}
\end{proposition}
\begin{proof}
This is the pro-\'etale comparison (G12) of
Proposition~\ref{b:prop:geo-smooth-tensor}: on a chart both sides are
$(f'\circ a')^*v^*K(d)[2d]$ and the map is the identity, because the
\'etale counit and the affine-space counit commute with base change.
\end{proof}

\subsubsection{The orientation from the diagonal}

By the tensor formula, the finite-coefficient proof is reduced to one
global isomorphism $\omega_f\simeq(\mathbf Z/\ell^n\mathbf Z)_X(d)[2d]$.
It is obtained from the diagonal
$\Delta:X\to X\times_YX$, which is a section of the smooth projection
$p_1$. Deformation to the normal bundle moves $\Delta$ to the zero
section of the relative tangent bundle $T_{X/Y}$, and there
vector-bundle purity (Proposition~\ref{prop:sp-vector}) provides the
isomorphism.

We first treat an arbitrary section. Let $g:V\to S$ be smooth of pure
relative dimension $d$ and let $s:S\to V$ be a section of $g$ that is a
closed immersion. Then $s$ is a regular immersion of codimension $d$.
Let $N_s$ be its normal bundle, with projection $\pi:N_s\to S$ and zero
section $z:S\to N_s$.

The \emph{deformation to the normal bundle} is the open subscheme
$\operatorname{Def}_s(V)$ of the blowup of $V\times\mathbf A^1$ along
$s(S)\times\{0\}$ obtained by removing the strict transform of
$V\times\{0\}$. It comes with a smooth morphism
$\widetilde g:\operatorname{Def}_s(V)\to\mathbf A^1_S$ and a section
$\widetilde s:\mathbf A^1_S\to\operatorname{Def}_s(V)$ that is a closed
immersion. Over $t=1$ the pair $(\widetilde g,\widetilde s)$ restricts
to $(g,s)$, and over $t=0$ to $(\pi,z)$. In local terms: if $s(S)$ is
cut out by relative coordinates $x_1,\ldots,x_d$, then
$\operatorname{Def}_s(V)$ has coordinates $u_1,\ldots,u_d$ with
$x_i=tu_i$, and at $t=0$ the $u_i$ become linear coordinates on $N_s$.

\begin{lemma}[A section and its normal bundle]\label{lem:sp-dnc}
There is a canonical isomorphism in
$\Dctf(S_{\proet},\mathbf Z/\ell^n\mathbf Z)$
\begin{equation}\label{eq:sp-dnc}
 s^{*}\omega_{g}\xrightarrow{\sim}z^{*}\omega_{\pi}.
\end{equation}
It is obtained by restricting to $t=0$ the unique isomorphism
$\widetilde s^{*}\omega_{\widetilde g}\to a^{*}s^{*}\omega_{g}$ in
$\Dctf(\mathbf A^1_{S,\proet},\mathbf Z/\ell^n\mathbf Z)$ that is the identity at $t=1$, and inverting; here
$a:\mathbf A^1_S\to S$ is the projection.
\end{lemma}

The fibres of $\widetilde s^{*}\omega_{\widetilde g}$ at $t=1$ and $t=0$
are the two sides of \eqref{eq:sp-dnc}. The proof uses coordinates to
show that this shifted rank-one sheaf is pulled back from $S$, \'etale
locally on $S$. The geometric fibres of $\mathbf A^1_S\to S$ are
connected, so a comparison fixed at $t=1$ is unique, and the local
comparisons glue. Coordinates are needed for the first step: in
positive characteristic a locally constant sheaf on $\mathbf A^1_S$ need
not be pulled back from $S$, and here it is because the deformation is
locally a product.

\begin{proof}
By Proposition~\ref{prop:sp-exceptional-bc}, applied to the inclusions
of the fibers $t=1$ and $t=0$, the restrictions of
$\widetilde s^{*}\omega_{\widetilde g}$ to $t=1$ and $t=0$ are
$s^{*}\omega_{g}$ and $z^{*}\omega_{\pi}$. By Proposition~\ref{prop:sp-tensor}, both
$\widetilde s^{*}\omega_{\widetilde g}[-2d]$ and
$s^{*}\omega_{g}[-2d]$ are locally constant sheaves, free of rank one
over $\mathbf Z/\ell^n\mathbf Z$. Morphisms between two such complexes
are therefore morphisms of sheaves, and they glue.

\emph{Local construction.} After passing to an \'etale cover
$\{S_i\to S\}$ we may choose, for the base-changed section
$s_i:S_i\to V\times_SS_i$, a neighborhood $U_i$ of $s_i(S_i)$ and an
\'etale map $e_i:U_i\to\mathbf A^d_{S_i}$ sending $s_i$ to the zero
section. The inverse image of the zero section under $e_i$ is
separated and \'etale over $S_i$, so $s_i(S_i)$ is open and closed in
it. Its complement is closed in $U_i$, and we remove it from $U_i$.
After this, the inverse image of the zero section is exactly
$s_i(S_i)$, so $e_i$ is an \'etale map of pairs whose square of centers
is Cartesian. Blowing up and removing the strict transform commute
with \'etale base change, so $e_i$ induces an \'etale map from a
neighborhood of the deformed section in $\operatorname{Def}_s(V)$ to the
deformation space of the zero section of $\mathbf A^d_{S_i}$. The latter
is $\mathbf A^d_{\mathbf A^1_{S_i}}$, mapping to $\mathbf A^d_{S_i}$ by
$x_i=tu_i$.

On $\mathbf A^d_{\mathbf A^1_{S_i}}$ the affine-space adjunction
identifies the dualizing complex with the constant sheaf
$(\mathbf Z/\ell^n\mathbf Z)(d)[2d]$, and so does the chart $e_i$ for
$s^{*}\omega_{g}$ over $S_i$. Composing these two identifications
gives an isomorphism over $\mathbf A^1_{S_i}$
\[
 \widetilde s^{*}\omega_{\widetilde g}\xrightarrow{\sim}
 a^{*}s^{*}\omega_{g}.
\]
At $t=1$ it restricts to an automorphism of $s^{*}\omega_{g}$ over
$S_i$. Composing with the pullback along $a$ of the inverse of this
automorphism, we may assume that the local isomorphism is the identity
at $t=1$.

\emph{Gluing.} Over $S_i\times_SS_j$ the two normalized local
isomorphisms differ by an automorphism of the rank-one locally constant
sheaf $a^{*}s^{*}\omega_g[-2d]$, that is, by a section of a locally
constant sheaf of units of $\mathbf Z/\ell^n\mathbf Z$. The locus where
this section equals $1$ is open and closed. On each geometric fiber
$\mathbf A^1_{\bar b}$, which is connected, it contains the point $t=1$,
so it is the whole fiber. Hence the section is $1$ at every geometric
point, and the local isomorphisms agree on overlaps. They glue to a
unique isomorphism over $\mathbf A^1_S$ that is the identity at $t=1$.

Restricting it to $t=0$ gives an isomorphism
$z^{*}\omega_{\pi}\simeq s^{*}\omega_g$, whose inverse is
\eqref{eq:sp-dnc}.
\end{proof}

The construction commutes with base change along any $S'\to S$:
pullback preserves the local product model and the normalization at
$t=1$, so the uniqueness statement in the lemma identifies the
pulled-back isomorphism with the one constructed over $S'$.

We apply the lemma to the diagonal. Let $p_1,p_2:X\times_YX\to X$ be
the projections. Since $f$ is separated, $\Delta$ is a closed immersion,
and it is a section of the smooth morphism $p_1$. Its conormal sheaf is
$\Omega^1_{X/Y}$, so its normal bundle is the relative tangent bundle
$T_{X/Y}$; write $\pi:T_{X/Y}\to X$ for its projection and
$z:X\to T_{X/Y}$ for its zero section.

\begin{proposition}[The relative dualizing complex]
\label{prop:sp-orientation}
Let $f:X\to Y$ be smooth of pure relative dimension $d$. There is a
canonical isomorphism in $\Dctf(X_{\proet},\mathbf Z/\ell^n\mathbf Z)$
\begin{equation}\label{eq:sp-orientation}
 \omega_{f}\xrightarrow{\sim}(\mathbf Z/\ell^n\mathbf Z)_{X}(d)[2d].
\end{equation}
\end{proposition}
\begin{proof}
The projection $p_1$ is the base change of $f$ along $f$, so
Proposition~\ref{prop:sp-exceptional-bc} gives
$\omega_{p_1}\simeq p_2^{*}\omega_f$ in
$\Dctf((X\times_YX)_{\proet},\mathbf Z/\ell^n\mathbf Z)$. Since $p_2\circ\Delta$ is the
identity of $X$, this gives the first isomorphism in
\begin{align*}
 \omega_{f}
  &\simeq L\Delta^{*}\omega_{p_1}
  \xrightarrow{\sim}z^{*}\omega_{\pi}
  \xrightarrow{\sim}z^{*}(\mathbf Z/\ell^n\mathbf Z)_{T_{X/Y}}(d)[2d]
  \simeq(\mathbf Z/\ell^n\mathbf Z)_{X}(d)[2d].
\end{align*}
The second isomorphism is Lemma~\ref{lem:sp-dnc} for the section
$\Delta$ of $p_1$, and the third is the inverse of the vector-bundle
purity isomorphism of Proposition~\ref{prop:sp-vector} at the constant
sheaf.
\end{proof}

The orientation \eqref{eq:sp-orientation} is defined without
coordinates. We check that on a coordinate chart it agrees with the
identification given by affine-space integration
\eqref{eq:sp-affine-integration}; this is what makes the purity map
below the expected one, and it is used again for the smooth trace.

\emph{Vector bundles.} Let $\pi:V\to S$ be a vector bundle of rank $d$.
Here the diagonal orientation should be the vector-bundle orientation of
Proposition~\ref{prop:sp-vector}. Regard $V\times_SV$ as a vector bundle
over the first factor. The automorphism $(v,w)\mapsto(v,w-v)$ of
$V\times_SV$ over the first factor takes the diagonal to the zero
section. It preserves integration in the second factor: letting the
translation vector vary over $V$, its effect on the top compactly
supported direct image, identified with the constant sheaf by
integration, is multiplication by a locally constant unit on $V$. This
unit is $1$ at the zero vector, where the translation is the identity,
and each geometric fiber of $V\to S$ is connected, so it is $1$
everywhere. Next, the deformation to the normal bundle of the zero
section of a vector bundle is canonically the product with
$\mathbf A^1$: in a local frame the map to the original bundle is
$x=tu$, and these descriptions agree under linear changes of frame.
Vector-bundle integration in this product family is the pullback of the
original integration, so the comparison of Lemma~\ref{lem:sp-dnc},
being the identity at $t=1$, is the identity under these integrations,
in particular at $t=0$. Together with the translation, this identifies
the diagonal orientation for $\pi$ with the vector-bundle orientation.

\emph{Coordinate charts.} Let $U\subset X$ be open with an \'etale map
$e:U\to\mathbf A^d_Y$ over $Y$. The \'etale map
$e\times e:U\times_YU\to\mathbf A^d_Y\times_Y\mathbf A^d_Y$ carries the
diagonal of $U$ into the diagonal of $\mathbf A^d_Y$. The inverse image
of the latter is $U\times_{\mathbf A^d_Y}U$, in which the diagonal of
$U$ is open and closed because $e$ is separated and \'etale. Removing
the other components, as in the proof of Lemma~\ref{lem:sp-dnc}, gives
a Cartesian square of centers. The induced map of normal bundles is the
differential $T_{U/Y}\simeq e^*T_{\mathbf A^d_Y/Y}$ (pullback of vector
bundles). Both ingredients of the diagonal orientation are therefore
pulled back from the affine model. For the dualizing complex this holds
because, by composing adjunctions, the exceptional inverse image along
an \'etale map is its ordinary inverse image. For the deformation
comparison, the same applies to the induced \'etale map of deformation
spaces, the pulled-back comparison is the identity at $t=1$, and the
uniqueness in Lemma~\ref{lem:sp-dnc} identifies it with the comparison
for $U$. Vector-bundle integration is preserved by the linear
isomorphism of normal bundles, so by the vector-bundle case the diagonal
orientation on $U$ is the affine-space orientation given by its
coordinate adjunction.

Combining the orientation with the tensor formula gives purity with
finite coefficients.

\begin{proposition}[Finite-coefficient smooth purity]
\label{prop:sp-finite}
Let $A$ be any finite admissible coefficient ring, and let
$f:X\to Y$ be smooth of pure relative dimension $d$.
For $K\in\Dctf(Y_{\proet},A)$ there is a canonical
isomorphism in $\Dctf(X_{\proet},A)$, natural in $K$,
\begin{equation}\label{eq:sp-finite}
 f_A^{*}K(d)[2d]
   \xrightarrow{\ \gamma_{f,A}(K)\ }f_A^{!}K.
\end{equation}
\end{proposition}
\begin{proof}
Choose $n$ with $\ell^nA=0$. The calculation above works with $A$
throughout: extend the Kummer and diagonal classes along
$\mathbf Z/\ell^n\mathbf Z\to A$; the projective-line splitting has
two free summands, and its integration and triangle identities retain
the same normalization. The local adjunctions therefore compute the
Brown adjoint over $A$, and the vector-bundle and deformation
arguments give $\omega_{f,A}\simeq A_X(d)[2d]$. Define $\gamma_{f,A}(K)$
as the composite
\[
 f_A^{*}K\otimes^{\mathbf L}A_X(d)[2d]
   \xrightarrow{\sim}f_A^{*}K\otimes^{\mathbf L}\omega_{f,A}
   \xrightarrow{\sim}f_A^{!}K.
\]
These are the orientation and tensor-map isomorphisms just proved.
The construction is natural in $K$ and uses only free rank-one Tate
sheaves, so it applies to every finite admissible $A$.
Lemma~\ref{lem:sp-reduction} verifies the compatibility of the
normalization for different choices of $n$ and maps of finite rings.
\end{proof}

\subsection{From finite to integral coefficients}
\label{sp:grp-integral}

\subsubsection{Change of finite coefficients}

The integral orientation and counit are obtained at finite levels
and passed to the limit, so these levels must be compatible under reduction
$\mathcal O_E/\ell^n\to\mathcal O_E/\ell^m$; the passage to $\Qbar$ also
needs the maps $\mathcal O_E/\ell^n\to\mathcal O_F/\ell^n$ for
$E\subset F$. We therefore treat an arbitrary map of finite coefficient
rings. Here $\mathcal O_E/\ell^n$ stands for
$\mathcal O_E/\ell^n\mathcal O_E$.

For a finite commutative ring $A$ with $\ell^rA=0$ for some $r$,
Proposition~\ref{prop:sp-finite} computes the right
adjoint $f_A^!\colon\Dctf(Y_{\proet},A)\to\Dctf(X_{\proet},A)$ of
$Rf_{!,A}$ and gives, for $K\in\Dctf(Y_{\proet},A)$,
the purity isomorphism
$\gamma_{f,A}(K):f_A^*K(d)[2d]\xrightarrow{\sim}f_A^!K$.

Let $f:X\to Y$ be smooth of pure relative dimension $d$, and fix a map
$A\to B$ of such rings.

\begin{definition}[Extension of coefficients]
\label{def:sp-scalar-extension}
For $K\in\Dctf(T_{\proet},A)$ we write
$K_B=K\otimes^{\mathbf L}_{A_T}B_T\in\Dctf(T_{\proet},B)$, and for a
morphism $u:K\to L$ we write
$u_B=u\otimes^{\mathbf L}_{A_T}\mathrm{id}_{B_T}$; this is a functor
$\Dctf(T_{\proet},A)\to\Dctf(T_{\proet},B)$. For a further map
$B\to C$, associativity of the tensor product gives
$(K_B)_C\simeq K_C$, compatibly with morphisms.
\end{definition}

\begin{lemma}[Coefficient comparison]
\label{lem:sp-reduction}
For $M\in\Dctf(X_{\proet},A)$ and $K\in\Dctf(Y_{\proet},A)$ there are
isomorphisms, natural in $M$ and $K$,
\[
\begin{aligned}
 b(M)&:(Rf_{!,A}M)_B\xrightarrow{\sim} Rf_{!,B}(M_B)
 &&\text{in }\Dctf(Y_{\proet},B),\\
 c_{A,B}(K)&:(f_A^!K)_B\xrightarrow{\sim}f_B^!(K_B)
 &&\text{in }\Dctf(X_{\proet},B).
\end{aligned}
\]
The map $c_{A,B}(K)$ is determined by the counit identity
\begin{equation}\label{eq:sp-cAB-counit}
 \varepsilon_B(K_B)\circ Rf_{!,B}\bigl(c_{A,B}(K)\bigr)
 =(\varepsilon_A(K))_B\circ b(f_A^!K)^{-1},
\end{equation}
where $\varepsilon_A:Rf_{!,A}f_A^!\to\mathrm{id}$, a natural
transformation of endofunctors of $\Dctf(Y_{\proet},A)$, and
$\varepsilon_B$ are the counits of Lemma~\ref{lem:sp-finite-existence}.
\end{lemma}
\begin{proof}
For $b$, use one compactification $f=\bar f\circ j$ for both rings.
Extension by zero commutes with $-\otimes^{\mathbf L}_AB$. For
$R\bar f_*$, take a bounded-above resolution of $B$ whose terms are
finite free $A$-modules; it may be infinite to the left. Represent $M$
by the bounded flat constructible model of
Lemma~\ref{ctf:finite-models}. The comparison is an isomorphism on
each free term, and the cohomological bound
$\operatorname{cd}_\ell(\bar f)$ gives a finite window in every
total degree. As in the scalar projection argument of
Proposition~\ref{n:scalar-six}, this applies to every $B$, whatever its
projective dimension over $A$. Associativity gives compatibility
with a second map $B\to C$.

Define $c_{A,B}(K)$ as the transpose of the right side of
\eqref{eq:sp-cAB-counit}. Whether it is an isomorphism is local, so we
may check it on the coordinate charts $f|_U=p_d\circ e$ of
Lemma~\ref{lem:sp-finite-existence}. There the counit is built from
summation over the geometric fibres of $e$, the Kummer class of
$\mathcal O_{\mathbf P^1}(1)$, the diagonal class and the splitting
\eqref{eq:sp-p1-splitting} into two free summands of rank one, each of
which commutes with extension of coefficients (for the projective line
this is the coefficient compatibility in
Proposition~\ref{b:prop:geo-p1-adjunction}). Hence on a chart
$c_{A,B}(K)$ is the transpose of the canonical isomorphism
$(e^*p_d^*K(d)[2d])_B\simeq e^*p_d^*(K_B)(d)[2d]$.
\end{proof}

Transposing \eqref{eq:sp-cAB-counit} gives the further compatibilities
used below. The triangle identities give compatibility with units: for
$M\in\Dctf(X_{\proet},A)$, in $\Dctf(X_{\proet},B)$,
\[
 f_B^!\bigl(b(M)\bigr)\circ c_{A,B}(Rf_{!,A}M)\circ (\eta_A(M))_B
       =\eta_B(M_B),
\]
and for $A\to B\to C$ the composite of the comparisons for $A\to B$ and
$B\to C$ has the same transpose as $c_{A,C}$, hence equals it. Finally,
for $K\in\Dctf(Y_{\proet},A)$, in $\Dctf(X_{\proet},B)$,
\begin{equation}\label{eq:sp-scalar-gamma}
 c_{A,B}(K)\circ (\gamma_{f,A}(K))_B
       =\gamma_{f,B}(K_B),
\end{equation}
the sources being identified by
$(f_A^*K(d)[2d])_B\simeq f_B^*(K_B)(d)[2d]$. Indeed $\gamma_{f,A}$ is
the tensor map, built from the projection formula and the counit,
composed with the inverse of the orientation. The tensor map is
compatible by the lemma. The orientation comes from deformation to the
normal bundle, whose geometry does not depend on the coefficients; its
rank-one sheaves are locally free, so their derived extensions are
ordinary ones even when $A\to B$ is not flat, and the comparisons at
$t=1$ are the chart comparisons just treated. The uniqueness in
Lemma~\ref{lem:sp-dnc} then gives agreement at $t=0$.

The cases used below are reduction
$\mathcal O_E/\ell^n\to\mathcal O_E/\ell^m$ for $m\leq n$, and
$\mathcal O_E/\ell^n\to\mathcal O_F/\ell^n$ for finite extensions
$E\subset F$ of $\Ql$; the latter is finite free, also when $F/E$ is
ramified.

\subsubsection{Coefficient towers}

An integral constructible complex $M$ is derived complete,
$M\simeq R\varprojlim_n M/\ell^n$, so one would like to apply the finite
construction to all reductions $M/\ell^n$ and take $R\varprojlim$. The
finite-level units, counits and transition maps live in different
derived categories, so we work instead with inverse systems of sheaves
(towers), in whose derived category a compatible family of finite-level
constructions is a single morphism. Normalized constructible towers are
equivalent to integral constructible complexes
\eqref{eq:sp-normalized-equivalence}; we repeat the finite construction
on towers and transport it through this equivalence.

Fix a finite extension $E/\Ql$. The ideals $(\ell^n)$ and $(\pi_E^n)$
of $\mathcal O_E$ are cofinal, so we use powers of $\ell$. By
Lemma~\ref{n:coefficient-comparison},
$\mathcal O_{E,T}=\varprojlim_n(\mathcal O_E/\ell^n)_T$, and we write
$\mathcal O_{E,T}/\ell^n$ for $(\mathcal O_E/\ell^n)_T$ and
$M/\ell^n=M\otimes^{\mathbf L}_{\mathcal O_{E,T}}\mathcal O_{E,T}/\ell^n$.

\begin{definition}[Coefficient towers]\label{def:sp-coefficient-towers}
The category $\operatorname{Tow}(T,\mathcal O_E)$ has as objects the
systems $C_\bullet=(C_n,t_n)_{n\geq1}$, where $C_n$ is a sheaf of
$\mathcal O_{E,T}/\ell^n$-modules on $T_{\proet}$ and
$t_n:C_{n+1}\to C_n$ is linear over
$\mathcal O_{E,T}/\ell^{n+1}\to\mathcal O_{E,T}/\ell^n$. Morphisms are
families of module maps commuting with the $t_n$. We write
$\mathcal O_{E,T}/\ell^\bullet=(\mathcal O_{E,T}/\ell^n)_n$ for the
tower of coefficient rings.
\end{definition}

This is the category of $\mathcal O_{E,T}/\ell^\bullet$-modules in the
topos of inverse systems of sheaves on $T_{\proet}$ (compare
Bhatt--Scholze, Section~3.5), a Grothendieck abelian category in which
kernels, cokernels and filtered colimits are computed level by level.
We use its ordinary derived category. Since towers and their maps are
built before deriving, derived limits of compatible families of finite
constructions are functorial. Evaluation at level $n$ is exact; for
$C_\bullet\in D(\operatorname{Tow}(T,\mathcal O_E))$ we write $C_n$ for
its image in $D(T_{\proet},\mathcal O_E/\ell^n)$. The transition maps
induce, for $m\leq n$, \emph{reduction maps}
$C_n\otimes^{\mathbf L}_{\mathcal O_{E,T}/\ell^n}\mathcal O_{E,T}/\ell^m\to C_m$
in $D(T_{\proet},\mathcal O_E/\ell^m)$.

\begin{definition}[Constructible and normalized towers]
\label{def:sp-normalized-towers}
Let $\Dctf(T_{\proet},\mathcal O_E/\ell^\bullet)$ be the full
subcategory of $D(\operatorname{Tow}(T,\mathcal O_E))$ of objects
$C_\bullet$ such that each $C_n$ lies in
$\Dctf(T_{\proet},\mathcal O_E/\ell^n)$, with cohomological and Tor
amplitude bounds independent of $n$. Such a tower is
\emph{normalized} if all its reduction maps ($m\leq n$) are
isomorphisms. The full subcategory of normalized towers is denoted
$\Dctf(T_{\proet},\mathcal O_E/\ell^\bullet)_{\mathrm{norm}}$.
\end{definition}

For any $M\in D(T_{\proet},\mathcal O_E)$ the reductions $M/\ell^n$
form a normalized tower. Conversely:

\begin{lemma}[Reconstruction of a normalized tower]
\label{lem:sp-tower}
Let $C_\bullet\in\Dctf(T_{\proet},\mathcal O_E/\ell^\bullet)$ be
normalized. Then $R\varprojlim_n C_n$ lies in
$\Dbc(T,\mathcal O_E)$, and for every $m$ the natural map
\[
 \bigl(R\varprojlim_n C_n\bigr)\otimes^{\mathbf L}_{\mathcal O_{E,T}}
 \mathcal O_{E,T}/\ell^m\longrightarrow C_m
\]
is an isomorphism in $D(T_{\proet},\mathcal O_E/\ell^m)$.
\end{lemma}
\begin{proof}
Let $C=R\varprojlim_nC_n$ (\stacks{0940}). Reduction modulo $\ell^m$ is
the cone of $\ell^m$, so it commutes with $R\varprojlim$. Fix $m$; the
terms with $n<m$ do not affect the derived limit. For $n\geq m$,
$\mathcal O_{E,T}/\ell^n\otimes^{\mathbf L}_{\mathcal O_{E,T}}
\mathcal O_{E,T}/\ell^m$ has $H^0=H^{-1}=\mathcal O_{E,T}/\ell^m$.
Its truncation triangle and normalization give triangles
$C_m[1]\to C_n\otimes^{\mathbf L}_{\mathcal O_{E,T}}\mathcal O_{E,T}/\ell^m\to C_m$,
in which the transition maps on the left term are multiplication by
powers of $\ell$, hence zero along the indices $m,2m,3m,\dots$; so the
left term has vanishing $R\varprojlim$, and the reduction formula
follows. The complex $C$ is derived complete as a limit of complexes
killed by powers of $\ell$, and it is bounded since $T_{\proet}$ is
replete. Since $C/\ell\simeq C_1$ lies in
$\Dctf(T_{\proet},\mathcal O_E/\ell)$, $C$ lies in $\Dbc(T,\mathcal O_E)$
by \S\ref{sp:grp-setup}. The same calculation of normalized towers is
in Bhatt--Scholze, Lemmas~3.5.4--3.5.7, and \stacks{099P}.
\end{proof}

The reduction functor
$\operatorname{red}_T:D(T_{\proet},\mathcal O_E)\to
D(\operatorname{Tow}(T,\mathcal O_E))$, $M\mapsto(M/\ell^n)_n$, is
left adjoint to
$R\varprojlim\colon D(\operatorname{Tow}(T,\mathcal O_E))\to
D(T_{\proet},\mathcal O_E)$ (Bhatt--Scholze, Lemma~3.5.7). The unit
is an isomorphism on derived complete complexes, and by
Lemma~\ref{lem:sp-tower} the counit is an isomorphism on normalized
towers. So the adjunction restricts to mutually inverse equivalences
\begin{equation}\label{eq:sp-normalized-equivalence}
 \operatorname{red}_T:\Dbc(T,\mathcal O_E)
       \ \rightleftarrows\
       \Dctf(T_{\proet},\mathcal O_E/\ell^\bullet)_{\mathrm{norm}}
       :R\varprojlim_n .
\end{equation}

\paragraph{Operations on towers.}
For $g:T\to S$ define
$g_\bullet^*\colon\operatorname{Tow}(S,\mathcal O_E)\to\operatorname{Tow}(T,\mathcal O_E)$
and
$g_{*,\bullet}\colon\operatorname{Tow}(T,\mathcal O_E)\to\operatorname{Tow}(S,\mathcal O_E)$
levelwise. The first is exact and induces
$g_\bullet^*\colon D(\operatorname{Tow}(S,\mathcal O_E))\to D(\operatorname{Tow}(T,\mathcal O_E))$;
let
$Rg_{*,\bullet}\colon D^+(\operatorname{Tow}(T,\mathcal O_E))\to D^+(\operatorname{Tow}(S,\mathcal O_E))$
be the derived functor of the second. For an open immersion $j$ let
$j_{!,\bullet}$ be levelwise extension by zero, take tensor products
over $\mathcal O_{E,T}/\ell^\bullet$ in $D(\operatorname{Tow}(T,\mathcal O_E))$,
and twist levelwise.

\begin{lemma}[Levels of the tower direct image]
\label{lem:sp-tower-levels}
For $g:T\to S$ and $C_\bullet\in D^+(\operatorname{Tow}(T,\mathcal O_E))$,
\[
 \bigl(Rg_{*,\bullet}C_\bullet\bigr)_n\simeq Rg_{*,\mathcal O_E/\ell^n}C_n
 \quad\text{in }D^+(S_{\proet},\mathcal O_E/\ell^n).
\]
\end{lemma}
\begin{proof}
Evaluation at $n$ has an exact right adjoint $r_n$: for an
$\mathcal O_{E,T}/\ell^n$-module sheaf $J$,
\[
 (r_nJ)_r=\begin{cases}
 J,&r\geq n,\quad\text{with scalars restricted along }
     \mathcal O_{E,T}/\ell^r\to\mathcal O_{E,T}/\ell^n,\\
 0,&r<n,
 \end{cases}
\]
with identity transition maps from level $n$ on and zero maps below it.
It is right adjoint to evaluation because a map of towers
$C_\bullet\to r_nJ$ is determined by its component $C_n\to J$.

Embed each $C_n$ into an injective $\mathcal O_{E,T}/\ell^n$-module
sheaf $J_n$. By
adjunction this gives a monomorphism
$C_\bullet\hookrightarrow\prod_{n\geq1}r_nJ_n$; at level $r$ only the
factors with $n\leq r$ are nonzero, and one of them is the embedding into
$J_r$. If $C_\bullet$ is injective the monomorphism splits, so $C_r$ is a
direct summand of a finite product of the sheaves $J_n$ with $n\leq r$.

An injective module sheaf over any of these rings is acyclic for
sections over every object of $T_{\proet}$. This is a property of the
underlying abelian sheaf, so restriction of scalars does not affect it,
and finite products and direct summands preserve it. Hence every
component of an injective tower is $g_*$-acyclic. Since evaluation is
exact, evaluating a bounded below injective resolution of $C_\bullet$
computes $Rg_{*,\mathcal O_E/\ell^n}C_n$.
\end{proof}

For a compactification $f=\bar f\circ j$ put
$Rf_{!,\bullet}=R\bar f_{*,\bullet}\circ j_{!,\bullet}\colon
D^+(\operatorname{Tow}(X,\mathcal O_E))\to D^+(\operatorname{Tow}(Y,\mathcal O_E))$, so
$(Rf_{!,\bullet}C_\bullet)_n\simeq Rf_{!,\mathcal O_E/\ell^n}C_n$ in
$D^+(Y_{\proet},\mathcal O_E/\ell^n)$. Since
flatness and tensor products of towers are levelwise, the projection
formula, proper base change and localization triangles hold for towers:
the tower maps evaluate to the finite ones, and a map of towers that is
an isomorphism at every level is an isomorphism.

\begin{lemma}[Direct image preserves constructible towers]
\label{lem:sp-tower-bound}
For $C_\bullet\in\Dctf(X_{\proet},\mathcal O_E/\ell^\bullet)$, one has
$Rf_{!,\bullet}C_\bullet\in\Dctf(Y_{\proet},\mathcal O_E/\ell^\bullet)$.
\end{lemma}
\begin{proof}
Each level $Rf_{!}C_n$ lies in $\Dctf(Y_{\proet},\mathcal O_E/\ell^n)$
by Proposition~\ref{n:six-small}.
For uniformity, compactly supported base change and the projection
formula compute
$(Rf_{!}C_n)_{\bar y}\otimes^{\mathbf L}N$ as
$R\Gamma_c(X_{\bar y},C_n|_{X_{\bar y}}\otimes^{\mathbf L}N)$, and
$R\Gamma_c$ on fibres of dimension $\leq d$ has amplitude $[0,2d]$
(\stacks{0F10}); so a common Tor interval $[a,b]$ of the $C_n$ gives
the common interval $[a,b+2d]$.
\end{proof}

\begin{proposition}[The tower model of the smooth adjunction]
\label{prop:sp-tower-geometric}
Let $f:X\to Y$ be smooth of pure relative dimension $d$. There is an
adjunction
\begin{equation}\label{eq:sp-tower-adjunction}
 Rf_{!,\bullet}\dashv f_\bullet^!
       \quad\text{between }\Dctf(X_{\proet},\mathcal O_E/\ell^\bullet)
       \text{ and }\Dctf(Y_{\proet},\mathcal O_E/\ell^\bullet).
\end{equation}
On a coordinate chart $f|_U=p_d\circ e$ the right adjoint is
$e_\bullet^*p_{d,\bullet}^*(-)(d)[2d]$. At each level $n$ the unit and
counit are those of the finite adjunction; in particular
$(f_\bullet^!C_\bullet)_n\simeq f_{\mathcal O_E/\ell^n}^!C_n$ in
$\Dctf(X_{\proet},\mathcal O_E/\ell^n)$, and the
reduction maps of $f_\bullet^!C_\bullet$ are induced by those of
$C_\bullet$ and the coefficient comparisons of
Lemma~\ref{lem:sp-reduction}.
\end{proposition}
\begin{proof}
We repeat the finite construction in the ringed topos of towers on
$T_{\proet}$, with coefficient ring $\mathcal O_{E,T}/\ell^\bullet$,
and check each identity level by level.

\emph{Classes and integration.} In the abelian category of towers of
abelian sheaves there is an exact Kummer sequence
\[
 0\longrightarrow(\mu_{\ell^n})_n
 \longrightarrow(\mathbb G_m,(\cdot)^\ell)
 \xrightarrow{\ ((\cdot)^{\ell^n})_n\ }
 (\mathbb G_m,\mathrm{id})\longrightarrow0,
\]
with the $\ell$-power map as transition on $\mu_{\ell^n}$. A line bundle
$L$ gives a class in $H^1$ of the last tower, and its boundary is a class
in $H^2((\mu_{\ell^n})_n)$, that is, a morphism
$(\mathbf Z/\ell^n)_n\to(\mu_{\ell^n})_n[2]$ in the derived category of
$(\mathbf Z/\ell^n)_n$-modules. Extending coefficients to $\mathcal O_{E,T}/\ell^\bullet$ gives
$c_1(L):\mathcal O_{E,T}/\ell^\bullet\to (\mathcal O_{E,T}/\ell^\bullet)(1)[2]$. The Kummer sequence is used only in the
abelian category where it is exact, and linearity over $\mathcal O_{E,T}/\ell^\bullet$ is added
afterwards. The version with supports, applied to a Cartier divisor with
its canonical section, gives the diagonal class, a morphism in
$D(\operatorname{Tow}(\mathbf P^1_S\times_S\mathbf P^1_S,\mathcal O_E))$,
$\delta_\bullet:\Delta_*\mathcal O_{E,\mathbf P^1_S}/\ell^\bullet\to
(\mathcal O_{E,\mathbf P^1_S\times_S\mathbf P^1_S}/\ell^\bullet)(1)[2]$.

The unit and $c_1(\mathcal O(1))$ define a morphism in
$D(\operatorname{Tow}(S,\mathcal O_E))$
\[
 \mathcal O_{E,S}/\ell^\bullet\oplus (\mathcal O_{E,S}/\ell^\bullet)(-1)[-2]\longrightarrow
 Rq_{*,\bullet}\mathcal O_{E,\mathbf P^1_S}/\ell^\bullet.
\]
It is an isomorphism, since by Lemma~\ref{lem:sp-tower-levels} it
evaluates to \eqref{eq:sp-p1-splitting} at each level, and evaluations
jointly detect isomorphisms. Let $\operatorname{Tr}_{q,\bullet}$ be the
projection onto the twisted second summand.

\emph{Unit and counit on the projective line.} The counit is the
projection formula followed by $\mathrm{id}\otimes\operatorname{Tr}_{q,\bullet}$,
and the unit is given by the diagonal class:
\[
 \begin{aligned}
 Rq_{*,\bullet}(q_\bullet^*K_\bullet(1)[2])
 &\simeq K_\bullet\otimes^{\mathbf L}
 Rq_{*,\bullet}(\mathcal O_{E,\mathbf P^1_S}/\ell^\bullet)(1)[2]
 \longrightarrow K_\bullet,\\
 M_\bullet&\simeq Rp_{2*,\bullet}\Delta_*M_\bullet
 \xrightarrow{\ Rp_{2*,\bullet}(\delta_\bullet\otimes p_1^*M_\bullet)\ }
 Rp_{2*,\bullet}p_1^*M_\bullet(1)[2]
 \simeq q_\bullet^*Rq_{*,\bullet}M_\bullet(1)[2].
 \end{aligned}
\]
As in Lemma~\ref{lem:sp-p1}, proper base change and the projection
formulas in this topos reduce the triangle identities to
\[
 \operatorname{Tr}_{p_i,\bullet}\circ Rp_{i*,\bullet}\delta_\bullet
 =\mathrm{id}_{\mathcal O_{E,\mathbf P^1_S}/\ell^\bullet}\qquad(i=1,2).
\]
Both sides are morphisms between tower sheaves in degree zero, and they
agree at each level by \eqref{eq:sp-diagonal-integrals}; so they agree.
This gives the projective-line adjunction in $D(\operatorname{Tow})$.

\emph{Charts and gluing.} Restrict to $\mathbf A^1_S$ by extension by
zero, iterate to $\mathbf A^d_S$, and compose with the \'etale adjunction
$e_{!,\bullet}\dashv e_\bullet^*$. The local counits are the iterated
integrations and fibre sums of the finite construction, now as maps of
towers. For each $K_\bullet$, glue the local representing pairs by the
open-cover cone and the Hom sequences of \eqref{eq:sp-open-cover}, as
after Lemma~\ref{lem:sp-finite-existence}: compatible local counits lift
by exactness of the Hom sequence, and comparing Hom sequences shows that
the glued pair represents
$M_\bullet\mapsto\operatorname{Hom}(Rf_{!,\bullet}M_\bullet,K_\bullet)$.
Its restrictions to charts have the stated form, so the cohomological
and Tor bounds are uniform. Since a representing pair with its counit is
unique up to unique isomorphism, this defines the functor and its action
on maps without a functorial choice of cones.

Evaluation at level $n$ preserves these constructions and gives the
finite local counits, so by the same argument the evaluated pair
represents the finite functor. Uniqueness identifies it with
$f_{\mathcal O_E/\ell^n}^!C_n$, compatibly with counits and hence with
units. The coefficient comparisons have the same local counits, so
uniqueness also identifies the transition maps with those of
Lemma~\ref{lem:sp-reduction}.
\end{proof}

The orientation also passes to towers. By the chart form in
Proposition~\ref{prop:sp-tower-geometric}, $\omega_{f,\bullet}[-2d]$ is a
rank-one tower sheaf. Vector-bundle integration does not depend on
frames, by the universal $\GL_d$ calculation of
Proposition~\ref{prop:sp-vector}: the possible discrepancy is a map of
tower sheaves in degree zero, so it can be checked level by level. For
the deformation family, the coordinate construction of
Lemma~\ref{lem:sp-dnc} gives comparisons of tower sheaves normalized at
$t=1$; uniqueness glues them and determines them at $t=0$. So the
diagonal gives an orientation in $D(\operatorname{Tow})$ whose levels
are the finite orientations. Its inverse, composed with the tensor map
built from the tower counit and the projection formula, is the tower
purity isomorphism in $\Dctf(X_{\proet},\mathcal O_E/\ell^\bullet)$, natural in
$K_\bullet\in\Dctf(Y_{\proet},\mathcal O_E/\ell^\bullet)$,
\begin{equation}\label{eq:sp-tower-gamma}
 \gamma_{f,\bullet}(K_\bullet):
 f_\bullet^*K_\bullet(d)[2d]\xrightarrow{\sim}f_\bullet^!K_\bullet ,
\end{equation}
with level $n$ equal to $\gamma_{f,\mathcal O_E/\ell^n}(K_n)$. Both
functors in \eqref{eq:sp-tower-adjunction} preserve normalized towers:
$Rf_{!,\bullet}$ by the comparison $b$ of Lemma~\ref{lem:sp-reduction},
and $f_\bullet^!$ by \eqref{eq:sp-scalar-gamma} and
\eqref{eq:sp-tower-gamma}.

\paragraph{Integral coefficients.}
For $M\in\Dbc(X,\mathcal O_E)$ the tower
$\operatorname{red}_XM$ is normalized with the Tor bounds of $M$, hence
so is $Rf_{!,\bullet}\operatorname{red}_XM$. The existing
compact-support direct image has the canonical tower description, an
isomorphism in $\Dbc(Y,\mathcal O_E)$,
\begin{equation}\label{eq:sp-integral-F}
 Rf_{!,\mathcal O_E}M\simeq
 R\varprojlim_n Rf_{!,\mathcal O_E/\ell^n}(M/\ell^n).
\end{equation}
By Lemma~\ref{lem:sp-tower} it lies in $\Dbc(Y,\mathcal O_E)$
and, in $D(Y_{\proet},\mathcal O_E/\ell^n)$,
\begin{equation}\label{eq:sp-F-reduction}
 (Rf_{!,\mathcal O_E}M)\otimes^{\mathbf L}_{\mathcal O_{E,Y}}\mathcal O_{E,Y}/\ell^n
      \simeq Rf_{!,\mathcal O_E/\ell^n}(M/\ell^n).
\end{equation}
Independence of the compactification holds in the tower category,
hence after $R\varprojlim$. This is the compactly supported direct image
of \S\ref{b:sec:brown-construction}: $R\bar f_*$ and $j_!$ commute with
reduction modulo $\ell^n$ and with $R\varprojlim$ (\stacks{09BP}), and
$M\simeq R\varprojlim_nM/\ell^n$, so
$Rf_{!,\mathcal O_E}M\simeq R\bar f_*j_!M$. Likewise
$f_{\mathcal O_E}^{*}K$ has reductions
$f_{\mathcal O_E/\ell^n}^*(K/\ell^n)$. Since tensor products of objects
of $\Dbc(-,\mathcal O_E)$ stay in this category (Lemma~\ref{i:tensor}) with reductions the tensor products of reductions,
$R\varprojlim$ of the tower projection formula gives, in
$\Dbc(Y,\mathcal O_E)$,
\[
 Rf_{!,\mathcal O_E}\bigl(M\otimes^{\mathbf L} f_{\mathcal O_E}^{*}K\bigr)
 \simeq Rf_{!,\mathcal O_E}M\otimes^{\mathbf L} K.
\]

For $K\in\Dbc(Y,\mathcal O_E)$ the tower
$f_\bullet^!\operatorname{red}_YK$ is normalized. Its derived limit defines the functor
$H_{f,\mathcal O_E}\colon\Dbc(Y,\mathcal O_E)\to\Dbc(X,\mathcal O_E)$,
\begin{equation}\label{eq:sp-integral-H}
 H_{f,\mathcal O_E}(K):=
 R\varprojlim_n f_{\mathcal O_E/\ell^n}^!(K/\ell^n)
 \in\Dbc(X,\mathcal O_E).
\end{equation}
The next proposition identifies this model, with its counit, with
the integral Brown adjoint. Twisting by the rank-one Tate sheaf and
shifting commute with $R\varprojlim$, so the limit of
\eqref{eq:sp-tower-gamma}, followed by this identification, gives
an isomorphism in $\Dbc(X,\mathcal O_E)$
\begin{equation}\label{eq:sp-integral-gamma-O}
 \gamma_{f,\mathcal O_E}(K):
 f_{\mathcal O_E}^{*}K(d)[2d]
      \xrightarrow{\sim}f_{\mathcal O_E}^!K.
\end{equation}

\begin{proposition}[The normalized model of the integral adjunction]
\label{prop:sp-integral-adjunction}
The tower unit and counit make $H_{f,\mathcal O_E}$ a right adjoint
of $Rf_{!,\mathcal O_E}$ on $\Dbc(-,\mathcal O_E)$. Consequently there is a
unique counit-preserving isomorphism
$H_{f,\mathcal O_E}\simeq f_{\mathcal O_E}^!$ of functors
$\Dbc(Y,\mathcal O_E)\to\Dbc(X,\mathcal O_E)$, where $f_{\mathcal O_E}^!$
is the right adjoint of Theorem~\ref{ctf:unified-brown}.
\end{proposition}
\begin{proof}
By the tower comparisons and \eqref{eq:sp-normalized-equivalence},
$\operatorname{red}_Y Rf_{!,\mathcal O_E}M\simeq Rf_{!,\bullet}\operatorname{red}_X M$
and $\operatorname{red}_X H_{f,\mathcal O_E}(K)\simeq f_\bullet^!\operatorname{red}_Y K$,
isomorphisms in $\Dctf(Y_{\proet},\mathcal O_E/\ell^\bullet)$ and
$\Dctf(X_{\proet},\mathcal O_E/\ell^\bullet)$.
Hence the equivalence \eqref{eq:sp-normalized-equivalence} and the tower
adjunction \eqref{eq:sp-tower-adjunction} give natural bijections
\[
 \begin{aligned}
 \operatorname{Hom}(Rf_{!,\mathcal O_E}M,K)
 &\simeq\operatorname{Hom}(Rf_{!,\bullet}\operatorname{red}_X M,\operatorname{red}_Y K)\\
 &\simeq\operatorname{Hom}(\operatorname{red}_X M,f_\bullet^!\operatorname{red}_Y K)
 \simeq\operatorname{Hom}(M,H_{f,\mathcal O_E}(K)).
 \end{aligned}
\]
Here the outer Hom groups are taken in $\Dbc(Y,\mathcal O_E)$ and
$\Dbc(X,\mathcal O_E)$, and the middle ones in the tower categories.
Uniqueness of the right adjoint identifies this adjunction with the
restriction of Theorem~\ref{ctf:unified-brown}, preserving its unit and
counit.
\end{proof}

The unit and counit are therefore $R\varprojlim$ of the tower unit and
counit, and their reductions are the finite units and counits.

The isomorphism \eqref{eq:sp-integral-gamma-O} proves the
$\mathcal O_E$ case of Theorem~\ref{thm:sp-integral}, and with $E=\Ql$
the $\Zl$ case; the finite case is Proposition~\ref{prop:sp-finite}. The field cases
follow next, and the remaining compatibilities are proved in
\S\ref{sp:grp-trace}.

\subsection{Coefficient fields}
\label{sp:grp-fields}

\subsubsection{Uniqueness of right adjoints}

Each coefficient case produces an explicit right adjoint of $Rf_!$,
with its own unit and counit. It is identified with the Brown adjoint
$f^!_\Lambda$ by the following elementary fact, which was used for
$\mathcal O_E$ in Proposition~\ref{prop:sp-integral-adjunction} and is
used for $E$ and $\Qbar$ below.

\begin{proposition}[Comparison with any other right adjoint]
\label{prop:sp-unique-adjoint}
Let $\Lambda$ be any admissible coefficient ring, let
$f_\Lambda^!$ be the right adjoint of
$Rf_{!,\Lambda}:\Dctf(X,\Lambda)\to\Dctf(Y,\Lambda)$
with its counit $\varepsilon_{f,\Lambda}$, and let
$J:\Dctf(Y,\Lambda)\to\Dctf(X,\Lambda)$ be any other
right adjoint of $Rf_{!,\Lambda}$, with counit $\varepsilon^J$. There is a unique natural isomorphism
$\gamma:f_\Lambda^!\xrightarrow{\sim}J$ of functors
$\Dctf(Y,\Lambda)\to\Dctf(X,\Lambda)$ compatible with the counits, that
is, with $\varepsilon^J(K)\circ Rf_{!,\Lambda}(\gamma(K))
=\varepsilon_{f,\Lambda}(K)$ in $\Dctf(Y,\Lambda)$.
\end{proposition}
\begin{proof}
Let $\gamma(K)$ be the $J$-transpose of $\varepsilon_{f,\Lambda}(K)$ and
$\delta(K)$ the $f_\Lambda^!$-transpose of $\varepsilon^J(K)$; these are
natural, and the displayed identity holds by definition. Then
$\delta(K)\circ\gamma(K)$ and the identity both have
$f^!_\Lambda$-transpose $\varepsilon_{f,\Lambda}(K)$, so they are equal,
and symmetrically $\gamma(K)\circ\delta(K)=\mathrm{id}$. A map compatible
with the counits has $J$-transpose $\varepsilon_{f,\Lambda}(K)$, hence
equals $\gamma(K)$.
\end{proof}

Any right adjoint $J$ therefore inherits the coefficient comparisons
of $f_\Lambda^!$. For a coefficient extension $\Lambda\to\Lambda'$ (with
subscripts as in Definition~\ref{def:sp-scalar-extension}) and a right
adjoint $J'$ of $Rf_{!,\Lambda'}$, define
$c_J(K):(JK)_{\Lambda'}\to J'(K_{\Lambda'})$ in $\Dctf(X,\Lambda')$ as the transpose of
$(\varepsilon^J(K))_{\Lambda'}$ composed with the compatibility of
$Rf_!$ with extension of scalars, and let $c_f$ be the map defined in this way
for $f_\Lambda^!$. Let $\gamma'\colon f^!_{\Lambda'}\xrightarrow{\sim}J'$ be the
isomorphism of Proposition~\ref{prop:sp-unique-adjoint} over $\Lambda'$.
Comparing transposes with the counit compatibility
of $\gamma$ and $\gamma'$ gives
\[
 c_J(K)=\gamma'(K_{\Lambda'})\circ c_f(K)\circ(\gamma(K)^{-1})_{\Lambda'},
\]
so $c_J$ is an isomorphism, and it satisfies the composition identity
because $c_f$ does.

\subsubsection{Finite extensions of $\Ql$}

Fix a finite extension $E/\Ql$. Smooth purity with coefficients $E$
follows from the integral case by inverting $\ell$: every object of $\Dbc(T,E)$ is $K_0[1/\ell]$ for some
$K_0\in\Dbc(T,\mathcal O_E)$, and Hom groups are the integral ones with
$\ell$ inverted (\S\ref{n:ambient}):
\begin{equation}\label{eq:sp-localized-Hom}
 \Hom(K_0,L_0[q])\otimes_{\mathcal O_E}E
 \xrightarrow{\ \sim\ }
 \Hom_{\Dbc(T,E)}\bigl(K_0[1/\ell],L_0[1/\ell][q]\bigr).
\end{equation}
Indeed, $L_0[1/\ell]$ is the colimit of $L_0\xrightarrow{\ell}L_0\xrightarrow{\ell}\cdots$,
and Hom out of $K_0$ commutes with this colimit by
Lemma~\ref{i:ext}(1); essential surjectivity is
Lemma~\ref{b:lem:integral-models}(4).

The functors $f^*\colon\Dbc(Y,E)\to\Dbc(X,E)$ and
$Rf_!\colon\Dbc(X,E)\to\Dbc(Y,E)$ are those of
\S\ref{n:ambient} and \S\ref{b:sec:brown-construction}. On integral models
they are the integral functors with $\ell$ inverted:
$f^*(K_0[1/\ell])=(f^*K_0)[1/\ell]$ is clear, and
$Rf_!(M_0[1/\ell])\simeq(Rf_{!,\mathcal O_E}M_0)[1/\ell]$ in $\Dbc(Y,E)$ because
$Rf_{!,\mathcal O_E}M_0=R\bar f_*j_!M_0$ by \eqref{eq:sp-integral-F} and
$R\bar f_*$ commutes with the colimit along $\ell$
(Lemma~\ref{b:lem:bounded-continuity}).

\begin{proposition}[Smooth purity with coefficients $E$]
\label{cor:sp-rational}
Let $f:X\to Y$ be smooth of pure relative dimension $d$, and $E/\Ql$
finite. Then, as functors between $\Dbc(X,E)$ and $\Dbc(Y,E)$,
\[
 Rf_{!,E}\dashv\bigl(K\longmapsto f_E^*K(d)[2d]\bigr).
\]
\end{proposition}
\begin{proof}
Choose integral models $M_0,K_0$. By the comparisons above and
\eqref{eq:sp-integral-gamma-O}, the two functors are the scalar
extensions of $Rf_{!,\mathcal O_E}$ and $f^!_{\mathcal O_E}$, so
\eqref{eq:sp-localized-Hom} twice and the $\mathcal O_E$-linear
adjunction of Proposition~\ref{prop:sp-integral-adjunction} give
\[
 \begin{aligned}
 \operatorname{Hom}(Rf_{!,E}M,K)
 &\simeq\operatorname{Hom}(Rf_{!,\mathcal O_E}M_0,K_0)\otimes E\\
 &\simeq\operatorname{Hom}(M_0,f_{\mathcal O_E}^!K_0)\otimes E
 \simeq\operatorname{Hom}(M,f_E^*K(d)[2d]),
 \end{aligned}
\]
with Hom groups taken in $\Dbc(Y,E)$, $\Dbc(Y,\mathcal O_E)$,
$\Dbc(X,\mathcal O_E)$ and $\Dbc(X,E)$ respectively. The unit and counit are the scalar extensions of the integral ones,
transported along $\gamma_{f,\mathcal O_E}$, and the triangle identities
are images of the integral ones.
\end{proof}

Proposition~\ref{prop:sp-unique-adjoint} identifies this adjunction
with the Brown adjunction over $E$ and gives the counit-preserving purity
map $\gamma_{f,E}:f_E^*(-)(d)[2d]\xrightarrow{\sim}f_E^!$ of functors
$\Dbc(Y,E)\to\Dbc(X,E)$.
We keep writing $f_E^!$, $\eta_{f,E}$ and $\varepsilon_{f,E}$ for the
Brown adjoint and its unit and counit; through $\gamma_{f,E}$ they are
normalized by integral models.

\begin{lemma}[Compatibility with finite extensions of $E$]
\label{lem:sp-field-extension}
Let $E\subset F$ be finite extensions of $\Ql$. For $M\in\Dbc(X,E)$ and
$K\in\Dbc(Y,E)$ there are natural isomorphisms, in $\Dbc(Y,F)$ and
$\Dbc(X,F)$ respectively,
\[
\begin{aligned}
 (Rf_{!,E}M)\otimes^{\mathbf L}_{E_Y}F_Y
       &\simeq Rf_{!,F}(M\otimes^{\mathbf L}_{E_X}F_X),\\
 (f_E^*K(d)[2d])\otimes^{\mathbf L}_{E_X}F_X
       &\simeq f_F^*(K\otimes^{\mathbf L}_{E_Y}F_Y)(d)[2d],
\end{aligned}
\]
which carry the unit and counit of Proposition~\ref{cor:sp-rational} for
$E$ to those for $F$. For a further extension $F\subset L$, the
comparison for $E\subset L$ is the composite of those for $E\subset F$
and $F\subset L$.
\end{lemma}
\begin{proof}
Since $\mathcal O_F$ is finite free over $\mathcal O_E$,
$-\otimes\mathcal O_{F,T}$ commutes with $R\varprojlim$ and reduces at
level $n$ to extension along $\mathcal O_E/\ell^n\to\mathcal O_F/\ell^n$.
There Lemma~\ref{lem:sp-reduction} gives the comparisons with their
compatibilities with units, counits, purity maps and composition; they
commute with the reductions $n\to m$ by the composition statement, so
they are maps of the towers defining \eqref{eq:sp-integral-F} and
\eqref{eq:sp-integral-H}. Apply $R\varprojlim$ and invert $\ell$.
\end{proof}

\subsubsection{Algebraically closed coefficients}

With coefficients $\Qbar$ the functors are
$f_{\Qbar}^*=f^{-1}\colon\Dbc(Y,\Qbar)\to\Dbc(X,\Qbar)$ and
$Rf_{!,\Qbar}=R\bar f_{*,\Qbar}\circ j_{!,\Qbar}\colon\Dbc(X,\Qbar)\to\Dbc(Y,\Qbar)$. By
Proposition~\ref{n:two-colimit}, $\Dbc(T,\Qbar)$ is the $2$-colimit of the
categories $\Dbc(T,E)$; in particular, for $K,L\in\Dbc(T,E)$,
\begin{equation}\label{eq:sp-qbar-Hom}
 \varinjlim_{F\supset E}\Hom(K_F,L_F[q])
 \xrightarrow{\sim}\Hom(K_{\Qbar},L_{\Qbar}[q]),
\end{equation}
with Hom groups in $\Dbc(T,F)$ and $\Dbc(T,\Qbar)$. Hence any finite
diagram of $\Qbar$-complexes, with finitely many relations, is defined
over one finite extension of $\Ql$. By \S\ref{n:change},
$Rf_{!,\Qbar}$ and $f^*_{\Qbar}$ commute with extension of scalars from
$E$, and $Rf_{!,\Qbar}$ is independent of the compactification by
Proposition~\ref{b:prop:lower-exceptional-image}.

\begin{theorem}[Smooth purity with algebraically closed coefficients]
\label{thm:sp-qbar}
Let $f:X\to Y$ be smooth of pure relative dimension $d$. Then, as
functors between $\Dbc(X,\Qbar)$ and $\Dbc(Y,\Qbar)$,
\[
 Rf_{!,\Qbar}\dashv
       \bigl(K\longmapsto f_{\Qbar}^*K(d)[2d]\bigr).
\]
\end{theorem}
\begin{proof}
Choose models $M_E,K_E$ over one finite $E$. By the compatibility with
extension of scalars, \eqref{eq:sp-qbar-Hom} and
Proposition~\ref{cor:sp-rational} over each $F\supset E$,
\[
 \begin{aligned}
 \operatorname{Hom}(Rf_{!,\Qbar}M,K)
 &\simeq\varinjlim_{F}\operatorname{Hom}\bigl(Rf_{!,F}M_F,K_F\bigr)\\
 &\simeq\varinjlim_{F}\operatorname{Hom}\bigl(M_F,f_F^*K_F(d)[2d]\bigr)
 \simeq\operatorname{Hom}(M,f_{\Qbar}^*K(d)[2d]),
 \end{aligned}
\]
with Hom groups in $\Dbc(Y,\Qbar)$, $\Dbc(Y,F)$, $\Dbc(X,F)$ and
$\Dbc(X,\Qbar)$.
The middle bijections are given by $\eta_{f,F}$ and $\varepsilon_{f,F}$,
which are compatible with the transition maps by
Lemma~\ref{lem:sp-field-extension}. The bijection is independent of
the models and natural, because the isomorphisms, morphisms and
relations involved are defined over one finite $F$, where the
adjunction is natural.
\end{proof}

Proposition~\ref{prop:sp-unique-adjoint}, applied to the Brown adjoint,
gives the counit-preserving isomorphism $\gamma_{f,\Qbar}:f_{\Qbar}^*(-)(d)[2d]\xrightarrow{\sim}
f_{\Qbar}^!$ of functors $\Dbc(Y,\Qbar)\to\Dbc(X,\Qbar)$. Its unit $\eta_{f,\Qbar}$ and counit
$\varepsilon_{f,\Qbar}$ are identified with those transported from
finite extensions of $\Ql$, so the smooth trace with coefficients
$\Qbar$ is normalized by integration at finite level.

\subsection{Base change, composition and the trace}
\label{sp:grp-trace}

\subsubsection{Base change and composition}

The smooth trace is built from the purity map, and its transitivity and
Frobenius compatibility rest on the compatibility of purity with base
change and with composition. The geometric content of both lies at
finite level, where it comes from the construction of
Section~\ref{sec:smooth-purity}. The passage to $\Zl$ goes through
towers, and for it the finite isomorphisms must commute with change of
finite coefficients.

Fix a Cartesian square
\[
\begin{CD}
 X' @>{u}>> X\\
 @V{f'}VV @VV{f}V\\
 Y' @>{v}>> Y,
\end{CD}
\]
with $f$ smooth of pure relative dimension $d$, and a map $A\to B$ of
finite commutative rings killed by a power of $\ell$; a subscript $B$
denotes extension of scalars, as in
Definition~\ref{def:sp-scalar-extension}.

\paragraph{Change of coefficients and towers.}
For $M\in\Dctf(X_{\proet},A)$ and $K\in\Dctf(Y_{\proet},A)$ let
$b^f(M):(Rf_{!,A}M)_B\xrightarrow{\sim}Rf_{!,B}M_B$ in
$\Dctf(Y_{\proet},B)$ and
$c^f_{A,B}(K):(f_A^!K)_B\xrightarrow{\sim}f_B^!K_B$ in
$\Dctf(X_{\proet},B)$ be the coefficient
comparisons of Lemma~\ref{lem:sp-reduction}. With
$u_A^*\colon\Dctf(X_{\proet},A)\to\Dctf(X'_{\proet},A)$ and
$v_A^*\colon\Dctf(Y_{\proet},A)\to\Dctf(Y'_{\proet},A)$, let
$\beta_{f,A}(M):v_A^*Rf_{!,A}M\xrightarrow{\sim}R(f')_{!,A}u_A^*M$ in
$\Dctf(Y'_{\proet},A)$ be
compact-support base change, and let
$e_{f,A}(K):u_A^*f_A^!K\xrightarrow{\sim}(f')_A^!v_A^*K$ in
$\Dctf(X'_{\proet},A)$ be the
exceptional base change isomorphism of
Proposition~\ref{prop:sp-exceptional-bc}. Base change commutes with
change of coefficients:
\begin{equation}\label{eq:sp-mixed-direct-bc}
 \beta_{f,B}(M_B)\circ v_B^*(b^f(M))
 =b^{f'}(u_A^*M)\circ\bigl(\beta_{f,A}(M)\bigr)_B
 \quad\text{in }\Dctf(Y'_{\proet},B),
\end{equation}
\begin{equation}\label{eq:sp-mixed-exceptional-bc}
 c^{f'}_{A,B}(v_A^*K)\circ(e_{f,A}(K))_B
 =e_{f,B}(K_B)\circ u_B^*c^f_{A,B}(K)
 \quad\text{in }\Dctf(X'_{\proet},B).
\end{equation}
For \eqref{eq:sp-mixed-direct-bc} choose a compactification
$f=\bar f\circ j$ and pull it back to $Y'$. For the open immersion $j$
both sides are the evident identification of extensions by zero. For the
proper map $\bar f$, both $b^{\bar f}$ and proper base change are
transposes of the counit of $L\bar f^*\dashv R\bar f_*$ on the derived
categories of $A$-module sheaves on the pro-\'etale sites, after extension
of scalars and after pullback respectively, and both sides of
\eqref{eq:sp-mixed-direct-bc} transpose to the pulled-back counit with
scalars extended. Since $c^f$ and $e_{f,A}$ are defined by transposing
the counit composed with $b^f$, respectively with $\beta_{f,A}^{-1}$,
transposing \eqref{eq:sp-mixed-exceptional-bc} under
$R(f')_{!,B}\dashv(f')_B^!$, between $\Dctf(X'_{\proet},B)$ and
$\Dctf(Y'_{\proet},B)$, reduces it to \eqref{eq:sp-mixed-direct-bc}.

In the tower category of \eqref{eq:sp-tower-adjunction}, for a finite
extension $E/\Ql$, extension by zero and proper base change give
$\beta_{f,\bullet}:v_\bullet^*Rf_{!,\bullet}\to R(f')_{!,\bullet}u_\bullet^*$,
a natural transformation of functors
$\Dctf(X_{\proet},\mathcal O_E/\ell^\bullet)\to\Dctf(Y'_{\proet},\mathcal O_E/\ell^\bullet)$,
and transposing $v_\bullet^*(\varepsilon_{f,\bullet})\circ\beta_{f,\bullet}^{-1}$
under the tower adjunction for $f'$ gives
$e_{f,\bullet}:u_\bullet^*f_\bullet^!\to(f')_\bullet^!v_\bullet^*$, of functors
$\Dctf(Y_{\proet},\mathcal O_E/\ell^\bullet)\to\Dctf(X'_{\proet},\mathcal O_E/\ell^\bullet)$. At
level $n$ these are $\beta_{f,\mathcal O_E/\ell^n}$ and
\eqref{eq:sp-exceptional-bc}, so they are isomorphisms, and by
\eqref{eq:sp-mixed-exceptional-bc} $e_{f,\bullet}$ is compatible with
the reduction maps, hence preserves normalized towers. For $E=\Ql$ and
$K\in\Dbc(Y,\Zl)$, the reduction formulas for inverse
image and for $f^!_{\Zl}$ write both sides below as $R\varprojlim_n$ of
the finite-level terms for $K_n=K\otimes^{\mathbf L}_{\mathbf Z_{\ell,Y}}
(\mathbf Z/\ell^n\mathbf Z)_Y$, and $R\varprojlim$ of $e_{f,\bullet}$ is the
integral exceptional base change isomorphism
\begin{equation}\label{eq:sp-integral-bc}
 u_{\Zl}^{*}f_{\Zl}^{!}K
       \xrightarrow{\sim}(f')_{\Zl}^{!}v_{\Zl}^{*}K
       \quad\text{in }\Dbc(X',\Zl).
\end{equation}

\begin{proposition}[Base change of purity]\label{prop:sp-purity-bc}
Let $\Lambda$ be any admissible coefficient ring and let
$K\in\Dctf(Y,\Lambda)$. Identify the sources by
$u_\Lambda^{*}f_\Lambda^{*}\simeq(f')_\Lambda^{*}v_\Lambda^{*}$, as functors
$\Dctf(Y,\Lambda)\to\Dctf(X',\Lambda)$,
and the targets by exceptional base change. Then
\[
 u_\Lambda^{*}\gamma_{f,\Lambda}(K)
 =\gamma_{f',\Lambda}(v_\Lambda^{*}K)
 \quad\text{in }\Dctf(X',\Lambda).
\]
\end{proposition}
\begin{proof}
For a finite ring $A$ the following orientation and counit
calculation is made directly with $A$-coefficients. For $\mathcal O_E$
we make it for the tower purity map $\gamma_{f,\bullet}$ of
\eqref{eq:sp-tower-gamma} and apply $R\varprojlim$. Inverting $\ell$
and descending to a common finite extension then give $E$ and $\Qbar$.
We write out the tower calculation; dropping the index $\bullet$ gives
the finite one. The map $\gamma_{f,\bullet}$ is the composite of
the orientation
$(\mathcal O_{E,X}/\ell^\bullet)(d)[2d]\xrightarrow{\sim}\omega_{f,\bullet}$
and the tensor map
$f_\bullet^*K_\bullet\otimes\omega_{f,\bullet}\to f_\bullet^!K_\bullet$,
both morphisms in $\Dctf(X_{\proet},\mathcal O_E/\ell^\bullet)$,
and we treat the two factors separately.

\emph{Orientation.} The diagonal, its normal bundle, the deformation to
the normal bundle and the integration on vector bundles all commute
with base change; on each coordinate chart this is the finite-level
geometric calculation, at every level $n$. After removing the shift
$[2d]$, the two orientations to be compared are maps between rank-one
tower sheaves in degree zero. Equality of maps of sheaves can be tested
at each level, so the orientations agree in $\Dctf(X'_{\proet},\mathcal O_E/\ell^\bullet)$.

\emph{Tensor map.} Transpose both tensor maps under the tower adjunction
for $f'$, between $\Dctf(X'_{\proet},\mathcal O_E/\ell^\bullet)$ and
$\Dctf(Y'_{\proet},\mathcal O_E/\ell^\bullet)$. By the definition of $e_{f,\bullet}$, the transpose of the
pulled-back tensor map is the pulled-back counit for $f$, preceded by
the projection-formula map in $\Dctf(Y'_{\proet},\mathcal O_E/\ell^\bullet)$; by naturality of the projection formula
under pullback, the transpose of the tensor map for $f'$ is the same
composite. This is the counit calculation of
Proposition~\ref{prop:sp-exceptional-bc}, carried out with towers. So
the transposes, and hence the tensor maps, are equal.
\end{proof}

\paragraph{Composition.}
Let $X\xrightarrow{f}Y\xrightarrow{g}Z$ be smooth of pure relative
dimensions $d$ and $e$. Composing the finite adjunctions identifies
$(g\circ f)_A^!$ with $f_A^!g_A^!$ as functors
$\Dctf(Z_{\proet},A)\to\Dctf(X_{\proet},A)$, compatibly with
$R(g\circ f)_{!,A}\simeq Rg_{!,A}Rf_{!,A}$ as functors
$\Dctf(X_{\proet},A)\to\Dctf(Z_{\proet},A)$, and the counit of the
composite adjunction is the composite of the two counits. The
coefficient comparisons compose: for $M\in\Dctf(X_{\proet},A)$ and
$K\in\Dctf(Z_{\proet},A)$,
\[
 b^{g\circ f}(M)=Rg_{!,B}(b^f(M))\circ b^g(Rf_{!,A}M),
 \qquad
 c^{g\circ f}_{A,B}(K)=f_B^!c^g_{A,B}(K)\circ c^f_{A,B}(g_A^!K),
\]
identities of morphisms in $\Dctf(Z_{\proet},B)$ and
$\Dctf(X_{\proet},B)$ respectively. The first holds on each free term of a resolution of $B$ over $A$ by
naturality of the composition isomorphism, and the totalization in
Lemma~\ref{lem:sp-reduction} gives it for arbitrary $A\to B$; the second
follows by transposing under the composite adjunction. The tower
versions of these composition isomorphisms, on the categories
$\Dctf(T_{\proet},\mathcal O_E/\ell^\bullet)$ for $T=X,Y,Z$, evaluate at each level to the
finite ones, so by these identities they preserve normalized towers, and
for $E=\Ql$ applying $R\varprojlim$ gives the integral composition
isomorphisms for compact-support direct images
$\Dbc(X,\Zl)\to\Dbc(Z,\Zl)$ and
\[
 (g\circ f)_{\Zl}^{!}\simeq f_{\Zl}^{!}g_{\Zl}^{!}
   :\Dbc(Z,\Zl)\longrightarrow\Dbc(X,\Zl).
\]

\begin{proposition}[Composition of purity]\label{prop:sp-composition}
For every admissible $\Lambda$ and $K\in\Dctf(Z,\Lambda)$,
the purity map $\gamma_{g\circ f,\Lambda}(K)$ is the composite
\begin{align*}
 (g\circ f)^{*}K(d+e)[2d+2e]
   &\simeq f^{*}\bigl(g^{*}K(e)[2e]\bigr)(d)[2d]\\
   &\xrightarrow{\ f^{*}(\gamma_g(K))(d)[2d]\ }
       f^{*}\bigl(g^!K\bigr)(d)[2d]\\
   &\xrightarrow{\ \gamma_f(g^!K)\ }
       f^!g^!K
   \simeq(g\circ f)^!K
\end{align*}
in $\Dctf(X,\Lambda)$, where the coefficient $\Lambda$ is omitted
from the displayed composite.
\end{proposition}
\begin{proof}
As for base change, we compare orientations and tensor maps, first
over any finite admissible ring and then for towers.
$R\varprojlim$ gives the integral identity, and inversion of $\ell$
and descent to a finite extension of $\Ql$ give the field identities.

\emph{Orientation, finite coefficients.} Choose coordinates for $g$, and
coordinates for $f$ over such a chart. Over the resulting tower of
affine-space projections, both orientations are adjoint to iterated
integration along affine lines, in the same order: first the $d$
coordinates of $f$, then the $e$ coordinates of $g$. The integration on
$\mathbf A^{d+e}$ is defined by exactly this composite (the recursive
construction preceding \eqref{eq:sp-affine-integration}), and the
diagonal orientation agrees with this normalization on charts, so the
two orientations agree locally. After removing the common twist and
shift they are maps of locally constant rank-one sheaves, so local
agreement implies equality.

\emph{Tensor maps, finite coefficients.} Transpose both tensor maps
under the composite adjunction between $\Dctf(X,\Lambda)$ and
$\Dctf(Z,\Lambda)$. The transpose of the iterated tensor map
applies the projection formula for $f$ and the counit for $f$, then the
projection formula for $g$ and the counit for $g$. The transpose of the
tensor map for $g\circ f$ applies the projection formula for $g\circ f$
and its counit. The two projection-formula isomorphisms agree by
associativity of tensor products, and the counit of the composite
adjunction is the composite of the two counits. Hence the transposes,
and with them the tensor maps, agree.

\emph{Towers.} In the tower category
$\Dctf(X_{\proet},\mathcal O_E/\ell^\bullet)$, with its purity maps and
its composite adjunction, the two arguments apply as follows: the
orientations are maps of rank-one tower sheaves, whose equality can be
tested at each level, and the tensor maps are compared by the same
transposition.
\end{proof}

\subsubsection{The smooth trace}

The smooth trace is the counit of the purity adjunction, rewritten by
means of the purity map so that its source is the inverse image twisted
and shifted. For $X$ smooth of dimension $d$ over algebraically closed
$k$ it is the trace map
$H_c^{2d}(X_{\proet},\mathbf Z_{\ell,X}(d))\to\Zl$ used in the Weil
conjectures.

Let $f:X\to Y$ be smooth of pure relative dimension $d$, and let
$\Lambda$ be any admissible coefficient ring. We omit $\Lambda$ from
the notation and let $\varepsilon_f(K):Rf_!f^!K\to K$ be the counit
of the adjunction $Rf_!\dashv f^!$, restricted to $\Dctf(X,\Lambda)$
and $\Dctf(Y,\Lambda)$; it is a morphism in $\Dctf(Y,\Lambda)$.

\begin{definition}[Smooth trace]\label{def:sp-trace}
For $K\in\Dctf(Y,\Lambda)$, the \emph{smooth trace} is the composite in $\Dctf(Y,\Lambda)$
\begin{equation}\label{eq:sp-trace}
 \operatorname{Tr}_{f}(K):\quad
 Rf_!\bigl(f^{*}K(d)[2d]\bigr)
    \xrightarrow{\ Rf_!\gamma_{f}(K)\ }
                   Rf_!f^!K
    \xrightarrow{\ \varepsilon_{f}(K)\ }K.
\end{equation}
\end{definition}

For the projection $\mathbf A^d_Y\to Y$ and $K=\Lambda_Y$ this
is the corresponding coefficient version of the integration map
\eqref{eq:sp-affine-integration}, which on each geometric fibre sends
the class of the origin to $1$. For a finite \'etale morphism, where
$d=0$, it is summation over the fibres.

\begin{proposition}[Transitivity of the smooth trace]
\label{prop:sp-trace-composition}
Let $X\xrightarrow{f}Y\xrightarrow{g}Z$ be smooth of pure relative
dimensions $d$ and $e$, and $K\in\Dctf(Z,\Lambda)$ for any admissible
$\Lambda$. Then
$\operatorname{Tr}_{g\circ f}(K)$ is the composite in $\Dctf(Z,\Lambda)$
\[
\begin{gathered}
 R(g\circ f)_!
       \bigl((g\circ f)^{*}K(d+e)[2d+2e]\bigr)
 \simeq Rg_!Rf_!
   \bigl(f^{*}
     (g^{*}K(e)[2e])(d)[2d]\bigr)\\
 \xrightarrow{\ Rg_!
   (\operatorname{Tr}_{f}(g^{*}K(e)[2e]))\ }
 Rg_!\bigl(g^{*}K(e)[2e]\bigr)
 \xrightarrow{\ \operatorname{Tr}_{g}(K)\ }K.
\end{gathered}
\]
\end{proposition}
\begin{proof}
Expand both traces by \eqref{eq:sp-trace}. By naturality of $\gamma_f$,
the map $\gamma_g(K):g^{*}K(e)[2e]\to g^!K$ can be moved
past the purity map of the first trace. The composite then consists of
the two purity maps, which together are $\gamma_{g\circ f}(K)$ by
Proposition~\ref{prop:sp-composition}, followed by
\[
 Rg_!Rf_!f^!g^!K
 \xrightarrow{\ Rg_!(\varepsilon_{f}(g^!K))\ }
 Rg_!g^!K
 \xrightarrow{\ \varepsilon_{g}(K)\ }K.
\]
This is the counit of the composite adjunction
$R(g\circ f)_!\dashv(g\circ f)^!$ between $\Dctf(X,\Lambda)$ and
$\Dctf(Z,\Lambda)$, so the composite is purity for
$g\circ f$ followed by its counit, that is,
$\operatorname{Tr}_{g\circ f}(K)$.
\end{proof}

\paragraph{Change of coefficients.}
For finite rings, the compatibility of purity maps and traces with
change of coefficients is \eqref{eq:sp-scalar-gamma} together with the
counit identity \eqref{eq:sp-cAB-counit}. The limit over towers,
inversion of $\ell$ and descent to a finite extension of $\Ql$ carry
these identities to $\mathcal O_E$, $E$ and $\Qbar$. Hence purity and
trace satisfy the base change and composition laws above, with the same
normalization, for every admissible $\Lambda$; for $\mathcal O_E$, $E$
and $\Qbar$ one may read $\Dbc$ for $\Dctf$ throughout.

\begin{corollary}[The dualizing complex of a smooth scheme]
\label{cor:sp-all-coefficients}
Let $a:X\to\operatorname{Spec}k$ be smooth of pure dimension $d$, and
$\Lambda$ any admissible coefficient ring. There is a natural isomorphism
$a_\Lambda^!\Lambda_{\operatorname{Spec}k}\simeq\Lambda_X(d)[2d]$ in
$\Dctf(X,\Lambda)$.
\end{corollary}
\begin{proof}
Apply smooth purity with coefficients $\Lambda$ to
$K=\Lambda_{\operatorname{Spec}k}$.
\end{proof}

Over algebraically closed $k$, taking global sections turns the counit
into a map
$R\Gamma_c(X_{\proet},\Lambda_X(d))[2d]\to\Lambda$ in $D^b(\Lambda)$,
normalized by finite-level integration, and in degree zero into
$H_c^{2d}(X_{\proet},\Lambda_X(d))\to\Lambda$.
For $X=\mathbf P^1_k$ it sends $c_1(\mathcal O(1))$ to $1$.
Over a finite field the counit is first a map on
$(\operatorname{Spec}k)_{\proet}$; the geometric cohomology trace is
obtained after extending to an algebraic closure, with the Frobenius
action described next.

\subsubsection{Weil structures}

In the Weil conjectures the trace is applied to cohomology with its
Frobenius action, so it must commute with Frobenius. For this
subsubsection, fix $k_0=\mathbb F_q$ and its algebraic closure $k$.
Suppose that $f$ is obtained from a smooth morphism
$f_0:X_0\to Y_0$ over $k_0$, by extension to $k$. The arithmetic Frobenius is
$\sigma_{\mathrm{ar}}\in\operatorname{Gal}(k/k_0)$, $a\mapsto a^q$, and
the geometric Frobenius is its inverse
$\sigma=\sigma_{\mathrm{ar}}^{-1}$. For a scheme $T$ of the form
$T_0\times_{k_0}k$, with $T_0$ over $k_0$, let
\[
 \operatorname{Fr}_T:T\longrightarrow T
\]
be the semilinear automorphism induced by $\sigma$; it is a morphism of
schemes covering the automorphism $\operatorname{Spec}(\sigma)$ of
$\operatorname{Spec}k$.

We use $\operatorname{Fr}_T$ in all formulas below. Two other
Frobenius maps are common: the absolute $q$-power Frobenius
$F_T^{\mathrm{abs}}$, and the $k$-morphism $F_{T_0/k_0}\times_{k_0}k$
obtained by base change from the relative $q$-power Frobenius over $k_0$.
They are related by
\[
 F_{T_0/k_0}\times_{k_0}k
 =F_T^{\mathrm{abs}}\circ\operatorname{Fr}_T .
\]
Absolute Frobenius is a universal homeomorphism, so it induces an
equivalence of \'etale topoi. Through this relation, the Weil
structures below correspond to the usual formulation with relative
Frobenius. Tate
twists are always identified by the Kummer maps below, so geometric
Frobenius acts on $\Lambda(1)$ by $q^{-1}$.

For every admissible coefficient ring $\Lambda$, pullback of module
sheaves along $\operatorname{Fr}_T$ is exact and induces a functor
\[
 \operatorname{Fr}_T^*:
 \Dctf(T,\Lambda)\longrightarrow
 \Dctf(T,\Lambda).
\]
For $\Lambda=\mathcal O_E$ this pullback preserves derived completeness,
so
$\operatorname{Fr}_T^*
 \simeq (\operatorname{Fr}_T)_{\mathcal O_E}^{*}$ as functors on
$\Dctf(T,\mathcal O_E)$.
For $\Lambda=E$ or $\Qbar$ it is $L(\operatorname{Fr}_T)_\Lambda^*$.

Purity and trace commute with Frobenius pullback as natural
transformations, before any Weil structure is chosen. The square with
horizontal arrows $\operatorname{Fr}_X,\operatorname{Fr}_Y$ and vertical
arrows $f,f$ is Cartesian, so for $K\in\Dctf(Y,\Lambda)$ and
$M\in\Dctf(X,\Lambda)$ the base-change comparisons identify
\[
 \operatorname{Fr}_X^*f_\Lambda^!K
       \simeq f_\Lambda^!\operatorname{Fr}_Y^*K,
 \qquad
 \operatorname{Fr}_Y^*Rf_{!,\Lambda}M
       \simeq Rf_{!,\Lambda}\operatorname{Fr}_X^*M
\]
in $\Dctf(X,\Lambda)$ and $\Dctf(Y,\Lambda)$ respectively. With these identifications, and the Kummer identification
$\operatorname{Fr}_T^*\Lambda_T(1)\simeq\Lambda_T(1)$ for the twists,
we have
\[
 \operatorname{Fr}_X^*\gamma_{f,\Lambda}(K)
       =\gamma_{f,\Lambda}(\operatorname{Fr}_Y^*K),
 \qquad
 \operatorname{Fr}_Y^*\operatorname{Tr}_{f,\Lambda}(K)
       =\operatorname{Tr}_{f,\Lambda}(\operatorname{Fr}_Y^*K),
\]
identities of morphisms in $\Dctf(X,\Lambda)$ and $\Dctf(Y,\Lambda)$
respectively. The first identity is base change of purity
(Proposition~\ref{prop:sp-purity-bc}) for this Cartesian square. At
finite level the orientation comparison is obtained by pulling back
the diagonal, its deformation family and vector-bundle integration,
and the Kummer classes have the same pullback. The second follows from
the first and the defining counit identity for exceptional base change,
which says that after compact-support base change the pulled-back
counit is the counit preceded by the exceptional comparison. Both pass
to all admissible coefficients by the change of coefficients above.

\begin{definition}[Weil structure]\label{def:sp-Weil-structure}
For $K\in\Dctf(Y,\Lambda)$, a \emph{Weil structure}
is an isomorphism in $\Dctf(Y,\Lambda)$
\[
 \phi_K:\operatorname{Fr}_Y^*K\xrightarrow{\sim}K.
\]
A \emph{Weil complex} is a pair $(K,\phi_K)$ with this structure.
A morphism $(K,\phi_K)\to(L,\phi_L)$ is a morphism
$h:K\to L$ in $\Dctf(Y,\Lambda)$ satisfying
\[
 h\circ\phi_K=\phi_L\circ\operatorname{Fr}_Y^*(h).
\]
The same definition applies on $X$ or on any $T=T_0\times_{k_0}k$.
\end{definition}

A Weil structure $\phi_K$ gives the following structure on
$f_\Lambda^!K$, an isomorphism in $\Dctf(X,\Lambda)$:
\[
 \operatorname{Fr}_X^*f_\Lambda^!K
   \simeq f_\Lambda^!\operatorname{Fr}_Y^*K
   \xrightarrow{\ f_\Lambda^!(\phi_K)\ }f_\Lambda^!K.
\]
On $f_\Lambda^*K(d)[2d]$ the structure is the pullback of $\phi_K$
combined with the Kummer identification, for every admissible
$\Lambda$. If
$(M,\phi_M)$ is a Weil complex on $X$, the structure on its
compact-support direct image is the isomorphism in $\Dctf(Y,\Lambda)$
\[
 \operatorname{Fr}_Y^*Rf_{!,\Lambda}M
   \simeq Rf_{!,\Lambda}\operatorname{Fr}_X^*M
   \xrightarrow{\ Rf_{!,\Lambda}(\phi_M)\ }Rf_{!,\Lambda}M.
\]
By naturality and the two identities above, purity and trace respect
these structures; this is their Frobenius equivariance for a Weil
complex $K$.

For a Weil complex $(K,\phi_K)$ on $T$, the induced geometric
Frobenius operator on compactly supported cohomology is the
endomorphism in $D^b(\Lambda)$
\[
\begin{gathered}
 \operatorname{Frob}_{T,K}:R\Gamma_c(T_{\proet},K)
   \xrightarrow{\ \text{pullback along }\operatorname{Fr}_T\ }
 R\Gamma_c(T_{\proet},\operatorname{Fr}_T^*K)\\
   \xrightarrow{\ R\Gamma_c(T_{\proet},\phi_K)\ }
 R\Gamma_c(T_{\proet},K).
\end{gathered}
\]
The constant coefficient sheaf has a canonical Weil structure, so the
smooth dualizing identification and the constant-coefficient trace are
canonically equivariant. For $r\in\mathbf Z$, write
$\Lambda(r)=\Gamma((\operatorname{Spec}k)_{\proet},
\Lambda_{\operatorname{Spec}k}(r))$ for the Tate coefficient
module with its induced Frobenius operator. Geometric Frobenius acts on
$\Lambda(1)$ by $q^{-1}$ and on $\Lambda(-d)$ by $q^d$, and arithmetic
Frobenius acts by the inverse scalars. Untwisted, the trace is
therefore an equivariant map
\[
 H_c^{2d}(X_{\proet},\Lambda_X)\longrightarrow\Lambda(-d).
\]
For $X=\mathbf P^1_k$ this says that geometric Frobenius acts on
$H^2(\mathbf P^1_k,\Lambda)$ by $q$ and on $H^0$ by $1$, so that the
Lefschetz trace formula with characteristic-zero coefficients gives
the integer equality $\#\mathbf P^1(\mathbb F_q)=1+q$. With finite
coefficients the trace identity is interpreted in $\Lambda$.

\subsubsection{Comparison with the exceptional inverse image}

For every admissible $\Lambda$, Theorem~\ref{ctf:unified-brown}
constructs $f^!$, and Theorem~\ref{n:existence}(4) restricts it to
$\Dctf$. The finite calculations, their towers and the field
comparisons above compute this functor, with its counit. This
identification completes the proof of Theorem~\ref{thm:sp-integral},
and it makes the trace and its Frobenius equivariance statements about
the six operations.

\begin{proposition}[Smooth purity for $f^!$]
\label{prop:sp-comparison-brown}
Let $f:X\to Y$ be smooth of pure relative dimension $d$ and let
$\Lambda$ be any admissible coefficient ring. There is a unique
isomorphism of functors $\Dctf(Y,\Lambda)\to\Dctf(X,\Lambda)$, with
components for $K\in\Dctf(Y,\Lambda)$
\[
 f_\Lambda^*K(d)[2d]\xrightarrow{\sim}f^!K
\]
compatible with the counits, where $f^!$ is the functor of
Theorem~\ref{ctf:unified-brown} restricted to $\Dctf(Y,\Lambda)$.
\end{proposition}
\begin{proof}
Both adjunctions live on $\Dctf(X,\Lambda)$ and $\Dctf(Y,\Lambda)$, and in
both $Rf_!$ is the composite
$R\bar f_*\circ j_!$ of ordinary derived functors on module sheaves for a
compactification $f=\bar f\circ j$, independent of the compactification
(\S\ref{n:compactifications}). Finite purity (Proposition~\ref{prop:sp-finite}), the integral
adjunction (Proposition~\ref{prop:sp-integral-adjunction}) and its
rational extensions (Proposition~\ref{cor:sp-rational} and
Theorem~\ref{thm:sp-qbar}) show, in the respective
coefficient cases, that $K\mapsto f_\Lambda^*K(d)[2d]$ is
right adjoint to $Rf_{!,\Lambda}$. By Theorem~\ref{ctf:unified-brown}
and Theorem~\ref{n:existence}(4), the restriction of $f^!$ is also
right adjoint to $Rf_{!,\Lambda}$ on $\Dctf$.
Proposition~\ref{prop:sp-unique-adjoint} gives the unique
counit-preserving isomorphism from the explicit adjoint
$K\mapsto f_\Lambda^*K(d)[2d]$ to $f_\Lambda^!$.
\end{proof}

For every admissible $\Lambda$ the exceptional inverse image of a
smooth morphism is therefore the twisted and shifted inverse image,
normalized by integration on the affine line, and the smooth trace
\eqref{eq:sp-trace} is the counit of $Rf_!\dashv f^!$ in this form,
compatible with base change, composition, change of coefficients and
Frobenius.

\clearpage
\appendix
\noindent The appendices prove the geometric results used in the main
text: Nagata's compactification theorem and Chow's lemma
(Appendix~\ref{app:compactification}), the Picard scheme and the
Jacobian of a curve (Appendix~\ref{app:picard}), and the cohomology of
curves (Appendix~\ref{app:curves}). They use only scheme theory and
coherent cohomology, together with the comparison results of
Section~\ref{sec:proetale-foundations} where pro-\'etale cohomology
appears. A detailed account of these constructions is
in~\cite{BlueprintGeometry}.

\section{Compactification and projective domination}
\label{app:compactification}

The direct image with proper support is defined through
compactifications and their common refinements, and proper base change
is reduced to projective morphisms by Chow's lemma. This appendix proves
both results with schemes, schematic closures, blow-ups and valuation
rings; the only result of the main text it uses is
Proposition~\ref{pe:open-closed}, in
Lemma~\ref{app:nagata-density-normalization}. Two deep algebraic results
are quoted with their precise hypotheses: flattening by blow-up, through
its consequence elimination of indeterminacy
(Theorem~\ref{app:nagata-elimination}), and Noetherian approximation.

The proof of Nagata's theorem runs as follows. Blow-ups make an open
schematically dense and eliminate indeterminacy. A further blow-up
separates the closures of two disjoint boundary pieces, and this is
what allows two compactifiable opens to be glued to a compactification
of their union. Induction over an affine cover proves the theorem over
a Noetherian base, and Noetherian approximation gives the general case.

\subsection{Compactifications and refinements}

\begin{definition}[Compactification and refinement]
\label{app:nagata-definition}
Let $f:X\to S$ be a morphism. A compactification is a factorization
\[
 X\xrightarrow{j}\overline X\xrightarrow{p}S
\]
with $j$ a quasi-compact open immersion and $p$ proper. It is
\emph{dense} if $j$ is scheme-theoretically dense: the map
$\mathcal O_{\overline X}\to j_*\mathcal O_X$ is injective.
A refinement of compactifications is an $S$-morphism
$h:\overline X'\to\overline X$ with $hj'=j$ and
$h^{-1}(X)=X$, under the specified identifications with $X$.
\end{definition}

\begin{theorem}[Nagata compactification]
\label{app:nagata-theorem}
Let $S$ be quasi-compact and quasi-separated. Every separated morphism
$f:X\to S$ of finite type has a dense compactification.
\end{theorem}

The theorem is \stacks{0F41}; its proof is given in
Section~\ref{app:nagata-global-proof}. The theorem holds for an
arbitrary qcqs base $S$ and any $f$ of finite type, and this generality
is used in the proof of proper base change over arbitrary qcqs bases;
for schemes of finite type over $k$ the Noetherian case of the proof
suffices.

Density is a condition on structure sheaves. A nilpotent section
supported on the boundary is invisible on $X$ but obstructs schematic
density. Dense compactifications have unique morphisms between them and
admit common refinements (Proposition~\ref{app:nagata-refinement}), and
passing to a dense compactification does not change the direct image
with proper support (Lemma~\ref{app:nagata-density-normalization}).

\begin{lemma}[Schematic closure and uniqueness]
\label{app:nagata-closure}
Let $i:U\to T$ be a quasi-compact immersion, with $T$ quasi-separated.
Its scheme-theoretic image $\overline U\subseteq T$ is a closed
subscheme, $U\to\overline U$ is an open immersion, and $U$ is
scheme-theoretically dense in $\overline U$.
If $U\subseteq T$ is a scheme-theoretically dense open and $Y\to S$
is separated, two $S$-maps $T\to Y$ agreeing on $U$ agree on $T$.
\end{lemma}

\begin{proof}
The ideal of the image is
$\ker(\mathcal O_T\to i_*\mathcal O_U)$, which is quasi-coherent
because $i$ is quasi-compact and quasi-separated. On an open of $T$
where the immersion is closed, the image restricts to $U$ itself.
Hence $U$ is open in its image. By construction a section of
$\mathcal O_{\overline U}$ vanishing on $U$ is zero; this is schematic
density. These statements are the image and closure lemmas
\stacks{01R8} and \stacks{01RG}.
For uniqueness, the equalizer of the two maps is a closed subscheme
of $T$, since it is the inverse image of the diagonal of $Y/S$.
Its ideal vanishes on $U$ and thus vanishes on $T$.
\end{proof}

\begin{lemma}[Passing to the schematic closure]
\label{app:nagata-density-normalization}
Let $X\xrightarrow{j}\overline X\xrightarrow{p}S$ be any
compactification with $\overline X$ Noetherian, and let
$i:Z\hookrightarrow\overline X$ be the
schematic closure of $X$, with $j_Z:X\hookrightarrow Z$ the induced
open immersion. For a condensed coefficient ring $\Lambda$, use its
pulled-back coefficient sheaves on these schemes. For every
$K\in D(X_{\proet},\Lambda_X)$ there are canonical identifications,
in $D(\overline X_{\proet},\Lambda_{\overline X})$ and
$D(S_{\proet},\Lambda_S)$ respectively,
\[
 j_!K\simeq i_*j_{Z,!}K,\qquad
 Rp_*j_!K\simeq R(p\circ i)_*j_{Z,!}K.
\]
\end{lemma}

Thus, for schemes of finite type over $k$, replacing a compactification
by the schematic closure of $X$ does not change the expression
$Rp_*j_!$ defining $Rf_!$; that expression needs only an open immersion
followed by a proper morphism.

\begin{proof}
The open and closed immersion operations on the pro-\'etale site
are exact by Proposition~\ref{pe:open-closed}. Its quasi-compact-open hypothesis applies to the
complement of $Z$ here because $\overline X$ is Noetherian, and
therefore ensures exactness of $i_*$. Extension by zero along $j$
is supported on $X$, which
lies in $Z$. Its inverse image on $Z$ is extension by zero along
$j_Z$. Equivalently, the natural map
$j_!K\to i_*i^*j_!K=i_*j_{Z,!}K$ is the identity on the open $X$
and both sides restrict to zero on its closed complement. The
open--closed localization sequence therefore makes it an
isomorphism. Since the immersion operations are exact, the argument
applies degreewise to unbounded complexes.
Since $i_*$ is exact, $Ri_*=i_*$. The derived direct-image
composition isomorphism $Rp_*Ri_*\simeq R(p\circ i)_*$ proves the
second identification, as a composite of derived right adjoints.
\end{proof}

\begin{remark}
The same normalization proof works for qcqs $\overline X$ whenever
$\overline X\setminus Z$ is quasi-compact. For an arbitrary
non-Noetherian compactification this condition can fail, and then
Proposition~\ref{pe:open-closed} gives no exactness statement for the
closed pro-\'etale pushforward. The schematic closure and dense
refinements exist in the generality of Theorem~\ref{app:nagata-theorem}.
\end{remark}

\begin{proposition}[Common dense refinements]
\label{app:nagata-refinement}
For the morphism of Theorem~\ref{app:nagata-theorem}, any finite family
of dense compactifications has a common dense refinement. Morphisms
between dense compactifications restricting to the identity on $X$
are unique, proper, and automatically satisfy $h^{-1}(X)=X$.
The category of dense compactifications is therefore cofiltered.
\end{proposition}

\begin{proof}
For two compactifications take the schematic closure $Z$ of the
diagonal image of $X$ in $\overline X_1\times_S\overline X_2$.
The map from $X$ is a quasi-compact immersion: it is the graph over
the open $X\subseteq\overline X_1$, and the second factor is
separated over $S$. Lemma~\ref{app:nagata-closure} makes $X$ an open
subscheme, scheme-theoretically dense in $Z$. The product is proper
over $S$, so $Z$ and its projections are proper.

Over $X\times_S\overline X_2$ the schematic closure is exactly the
graph of $X\to\overline X_2$: that graph is closed and already contains
the entire restriction of the original immersion. Thus the first
projection has inverse image of $X$ equal to $X$ and is the identity
there. The second projection has the same property. Induction gives
the finite-family assertion.

For a map $h:\overline X'\to\overline X$, properness follows by
factoring through its closed graph in
$\overline X'\times_S\overline X$. Put $V=h^{-1}(X)$. The inclusion
$X\to V$ is a section of the separated map $V\to X$ and is therefore
closed, as well as open. Since $X$ is schematically dense in
$\overline X'$, it is schematically dense in $V$; a closed subscheme
containing it must be $V$. Hence $V=X$. Uniqueness is
Lemma~\ref{app:nagata-closure}, so parallel morphisms are equal, and
the category is nonempty by Nagata's theorem.
\end{proof}

\begin{remark}[Base change of a compactification]
\label{app:nagata-base-change}
For every $S'\to S$, the base change of a compactification is an open
immersion followed by a proper morphism, and this suffices to define
the compactly supported direct image of $X\times_SS'\to S'$. When the
base-changed ambient scheme is Noetherian, as for schemes of finite
type over $k$, Lemma~\ref{app:nagata-density-normalization} identifies
it with the one defined by the schematic closure. Schematic density
and schematic closure commute with flat base change, but density can
be lost under an arbitrary base change. In that case the schematic
closure of $X\times_SS'$ in $\overline X\times_SS'$ is a dense
compactification, and a base-changed common refinement can be replaced
by its closure in the same way. The comparisons in the main text are
made on a further common refinement, where they are canonical by
uniqueness.
\end{remark}

\subsection{Blow-ups, strict transforms and elimination}
\label{app:nagata-blowups}

\begin{definition}[Admissible blow-up]
\label{app:nagata-admissible}
Let $T$ be quasi-compact and quasi-separated and $U\subseteq T$ a
quasi-compact open. For a finite-type quasi-coherent ideal
$\mathcal I\subseteq\mathcal O_T$ with
$\mathcal I|_U=\mathcal O_U$, the blow-up
\[
 b:\operatorname{Bl}_{\mathcal I}(T)=
 \operatorname{Proj}_T\Bigl(\bigoplus_{m\geq0}\mathcal I^m\Bigr)
 \longrightarrow T
\]
is called $U$-admissible. Its center is disjoint from $U$.
\end{definition}

On an affine $\operatorname{Spec}(R)$, write
$I=(a_1,\ldots,a_r)$. The charts of the blow-up have rings
\[
 R[I/a_i]=\bigl((\bigoplus_{m\geq0}I^mt^m)_{a_it}\bigr)_0.
\]
This is the degree-zero part of a localization of the Rees algebra,
with all its relations; the polynomial ring with only the relations
$a_iT_j=a_j$ can have additional $a_i$-torsion. In the chart $R[I/a_i]$
the element $a_i$ is a nonzerodivisor. It generates the transformed ideal $I$.
These charts explain why the center becomes an effective Cartier
divisor and why the blow-up is an isomorphism away from the center.
Finite generation of $I$ also realizes the affine-base blow-up as a
closed subscheme of projective space. Thus $b$ is proper. These
constructions glue under restriction to opens.

\begin{lemma}[Making an open schematically dense]
\label{app:nagata-boundary-divisor}
There is a $U$-admissible blow-up $T'\to T$ such that its inverse
image of $U$ is schematically dense in $T'$ and its complement is
the support of an effective Cartier divisor.
\end{lemma}

\begin{proof}
Choose a finite-type quasi-coherent ideal $\mathcal I$ with
$V(\mathcal I)=T\setminus U$; such an ideal exists for a
quasi-compact open in a qcqs scheme (\stacks{01PH}). Blow up this
ideal. The transformed ideal defines the exceptional Cartier divisor.
On its affine charts the local defining element $a$ is a
nonzerodivisor, so the map from the chart ring to its localization at
$a$ is injective. The complement of the divisor is consequently
schematically dense. It is precisely the inverse image of $U$.
This density persists under subsequent admissible blow-ups. Locally,
if $t$ defines the boundary divisor, multiplication by $t$ is injective
on each power of a blow-up ideal, since those powers are submodules of
$\mathcal O_T$. It is therefore injective on the Rees algebra and
its degree-zero localized charts. The pulled-back $t$ still defines
a Cartier divisor with complement the inverse image of $U$.
\end{proof}

For a morphism $W\to T$ and a blow-up $T'\to T$, the strict
transform of $W$ is the closed subscheme of $W\times_TT'$ obtained
by removing the sections of its structure sheaf supported on the
inverse image of the exceptional divisor. In a chart with equation
$a$, its ring is
\[
 B/\{b\in B:\ a^nb=0\text{ for some }n\geq0\}.
\]
If $U$ is the complement of the center, this is exactly the schematic
closure of $W_U$ in $W\times_TT'$; this closure is what occurs in
elimination of indeterminacy.
The strict transform is a blow-up of $W$ in the pulled-back ideal
(\stacks{080E}), so it is proper over $W$.

\begin{lemma}[Finite refinements of blow-ups]
\label{app:nagata-blowup-refinements}
Finitely many $U$-admissible blow-ups of $T$ admit a common
$U$-admissible blow-up. A composite of two $U$-admissible blow-ups is
again such a blow-up. A $U$-admissible blow-up on a quasi-compact
open $V\subseteq T$ containing $U$ extends to a $U$-admissible
blow-up of $T$.
\end{lemma}

\begin{proof}
For ideals $\mathcal I,\mathcal J$, blowing up their product factors
through both blow-ups. Indeed, after the product becomes an
invertible ideal, each factor is invertible, and the universal
property applies. More explicitly, first blowing up $\mathcal I$
and then the pulled-back $\mathcal J$ gives the blow-up of
$\mathcal I\mathcal J$ (\stacks{080A}). Induction proves the first
assertion. For a center on the first blow-up which is not itself
pulled back from $T$, the composition assertion is the finite-type
ideal calculation in \stacks{080B}.
The extension assertion uses extension of finite-type
quasi-coherent ideals from a quasi-compact open, with a center
still outside $U$; the precise admissible extension result is
\stacks{080M}. These two results turn finitely many local
constructions into one global modification preserving $U$.
\end{proof}

Blow-ups commute with flat base change (\stacks{0805}), because flat
pullback preserves the inclusions $\mathcal I^m\subseteq\mathcal O_T$.
Under a nonflat base change these maps can acquire kernels, and the
Rees algebra of the image ideal need not give the base-changed blow-up.

\begin{theorem}[Elimination of indeterminacy]
\label{app:nagata-elimination}
Let $T$ be qcqs, let $U\subseteq T$ be quasi-compact open, and let
$W\to T$ be separated and of finite type, with $W_U\to U$ an
isomorphism. There is a $U$-admissible blow-up $T'\to T$ whose
strict transform $W'\to T'$ is an open immersion. The modification
can first be chosen so that $U$ is schematically dense in $T'$.
\end{theorem}

This is \stacks{081S}. Its main ingredient is flattening by blow-up
(\stacks{081R}): if $W\to T$ is of finite type and quasi-separated and
is flat and locally of finite presentation over $U$, then its strict
transform under a suitable $U$-admissible blow-up is flat and of
finite presentation. After Lemma~\ref{app:nagata-boundary-divisor},
this strict transform contains the inverse image of $U$ schematically
densely and is the identity there, so it is an open immersion.

Conrad's account of Deligne's proof \cite[Theorem~2.4]{ConradNagata}
reaches the same conclusions through partial sections. For nested
quasi-compact opens $U\subseteq V$ and a section over $U$ whose graph
is already closed above $V$, it resolves the section while preserving
all of $V$. Its stages are: clearing denominators for an affine-space
target over an affine base, gluing after extending finitely many
centers, separating the unwanted boundary for nested opens, and
reducing a general separated target to a quasi-affine one by Zariski's
Main Theorem.

\begin{corollary}[Extension to a proper target]
\label{app:nagata-proper-extension}
Let $T$ be qcqs, $U\subseteq T$ a dense quasi-compact open, and
$Y\to S$ proper. Every $S$-map $a:U\to Y$ extends to an $S$-map
$T'\to Y$ after a $U$-admissible blow-up $T'\to T$.
\end{corollary}

\begin{proof}
Take the schematic closure $G$ of the graph of $a$ in $T\times_SY$.
The projection $G\to T$ is proper, since $Y\to S$ is proper, and
is an isomorphism over $U$. First make $U$ schematically dense by
Lemma~\ref{app:nagata-boundary-divisor}, and replace $G$ by its strict
transform. Theorem~\ref{app:nagata-elimination} then makes the strict
transform of $G$ an open subscheme of a further blow-up $T'$.
It is also proper over $T'$, being closed in the proper base change
of $G$. Its image is thus closed as well as open and contains the
schematically dense $U$. It is all of $T'$. The strict transform is
therefore isomorphic to $T'$, and its map to $Y$ is the desired
extension. Lemma~\ref{app:nagata-blowup-refinements} combines the
successive modifications into a single admissible blow-up.
\end{proof}

\begin{lemma}[Graph extension from a compactification]
\label{app:nagata-graph-extension}
Let $S$ be qcqs, $M\to S$ separated of finite type, and
$V\subseteq M$ a quasi-compact open with a compactification
$V\subseteq Y$ over $S$. After a $V$-admissible blow-up of $Y$
there is an open $V\subseteq V'\subseteq Y$ and a proper map
$\psi:V'\to M$ with $\psi^{-1}(V)=V$ and $\psi|_V=\mathrm{id}_V$.
\end{lemma}

\begin{proof}
Form the schematic closure $G$ of the graph of $V\to M$ in
$Y\times_SM$. Its projection to $Y$ is separated of finite type
and is an isomorphism over $V$. By
Theorem~\ref{app:nagata-elimination} its strict transform is an open
$V'\subseteq Y'$ after an admissible blow-up. Formation of the
graph closure under this modification gives that strict transform:
both are the closure of the same graph outside the exceptional
divisor (\stacks{0F3Y}). On the other hand $G'$ is closed in
$Y'\times_SM$, whose projection to $M$ is proper. Thus
$V'\simeq G'\to M$ is proper. Its inverse image of $V$ is the
closed graph over $V$, namely $V$ itself. This proves the lemma;
it is also \stacks{0F3Z}.
\end{proof}

\subsection{Separating the boundaries}
\label{app:nagata-boundaries}

Two proper compactifications of two opens cannot in general be glued
along their overlap: the diagonal of the glued scheme can fail to be
closed, because a specialization can approach incompatible boundary
points on the two sides. A blow-up that separates the closures of the
boundaries removes this obstruction.

\begin{lemma}[Separation of closures]
\label{app:nagata-separation}
Let $T$ be qcqs and $V\subseteq T$ quasi-compact open.
If $F_1,F_2\subseteq V$ are disjoint constructible closed subsets,
there is a $V$-admissible blow-up of $T$ on which their closures are
disjoint. Equivalently, for finitely presented closed subschemes of
$T$ disjoint over $V$, an admissible blow-up separates their strict
transforms.
\end{lemma}

This is \stacks{0F3V}. Its proof extends finite-type ideals from $V$,
separates the two pieces on their union by the chart calculation of
\stacks{080P}, and extends the center to $T$; this is where finite
presentation of the closed subschemes is needed. In the Noetherian
applications below every closed subset is constructible and its
reduced ideal is of finite type.

\begin{lemma}[An isomorphism near a boundary closure]
\label{app:nagata-near-closure}
Let $q:P\to T$ be proper between qcqs schemes, let $V\subseteq T$
be quasi-compact open, and let $F\subseteq V$ be closed. Suppose
$q$ is an isomorphism over a neighbourhood of $F$ in $V$.
After a $V$-admissible blow-up the strict transform of $q$ is an
isomorphism over a neighbourhood of the closure of $F$.
\end{lemma}

\begin{proof}
Inside $V$, separate $F$ from the closed complement of the maximal
open on which $q$ is an isomorphism. Since $V$ is spectral, both
closed subsets can be enlarged to disjoint constructible closed
subsets. Apply Lemma~\ref{app:nagata-separation} so that their closures
in $T$ are disjoint. On the complement $T_0$ of the second closure,
the restriction of $q$ is an isomorphism over an open $V_0$ containing
$F$. Apply Theorem~\ref{app:nagata-elimination} on $T_0$. Its strict
transform is both open and proper over the modification, so it is
an isomorphism there after the initial schematic-density modification.
Extend this blow-up from $T_0$ to $T$, preserving $V$, by the
extension part of Lemma~\ref{app:nagata-blowup-refinements}. This provides an isomorphism on
an open containing the closure of $F$. The precise argument, including
the choice of constructible enlargements, is \stacks{0F3W}.
\end{proof}

\begin{lemma}[Separated common modification of two opens]
\label{app:nagata-common-modification}
Let $S$ be qcqs and $U\subseteq T_i$ ($i=1,2$) be open immersions
between separated finite-type $S$-schemes. There are $U$-admissible
blow-ups $T_i'\to T_i$ and open immersions $T_i'\to T$ agreeing
on $U$, with $T$ separated and of finite type over $S$.
If $U$ is dense in both $T_i$, it is dense in $T$.
\end{lemma}

\begin{proof}
Take the schematic closure $G$ of $U$ in $T_1\times_ST_2$.
Apply Theorem~\ref{app:nagata-elimination} to its two projections.
The strict transforms then become opens in suitable blow-ups of
the two sides. Their further common refinement is obtained by
blowing up the product of the two pulled-back centers on $G$.
Extend these further centers from the resulting opens to each side
using Lemma~\ref{app:nagata-blowup-refinements}. The final comparison
scheme $G'$ is now an open in both $T_i'$.

Its map to $T_1'\times_ST_2'$ is a closed immersion: it is an
immersion, and is proper because $G'\to G$ is proper,
$G\subseteq T_1\times_ST_2$ is closed, and the modified product
is separated over the original product. Glue the two $T_i'$ along
$G'$. On products of these two opens the diagonal is either their
own closed diagonal or the just established closed graph $G'$.
Thus the glued scheme is separated. Finite type follows from the
finite open cover. Density follows from density on each side;
Lemma~\ref{app:nagata-boundary-divisor} can ensure schematic density
before performing further modifications. This construction is
\stacks{0F3X}.
\end{proof}

\begin{proposition}[The two-open compactification construction]
\label{app:nagata-two-open}
Let $S$ be Noetherian and $M\to S$ separated of finite type.
Suppose $M=M_1\cup M_2$, each $M_i$ has a compactification over
$S$, and $D=M_1\cap M_2$ is dense in $M$. Then $M$ has a
compactification over $S$.
\end{proposition}

\begin{proof}
Choose dense compactifications $M_i\subseteq X_i$. By
Lemma~\ref{app:nagata-graph-extension}, after admissible blow-ups
there are opens $M_i\subseteq V_i\subseteq X_i$ and proper maps
$\psi_i:V_i\to M$ which are isomorphisms above $M_i$.
Put, for $\{i,j\}=\{1,2\}$,
\[
 Z_i=M\setminus M_j,\qquad
 Z_{i,r}=\psi_i^{-1}(Z_r),\qquad
 B_{i,r}=\overline{Z_{i,r}}\subseteq X_i.
\]
The subsets $Z_{i,1},Z_{i,2}$ are disjoint and closed in $V_i$.
Lemma~\ref{app:nagata-separation} lets us arrange
$B_{i,1}\cap B_{i,2}=\varnothing$. Subsequent $V_i$-admissible
blow-ups preserve this disjointness.

Let $X_{12}$ be the schematic closure of
$V_1\times_MV_2$ in $X_1\times_SX_2$, with proper projections
$p_i$. Above $V_1$ the first projection is the base change of
$\psi_2$, and is an isomorphism near $Z_{1,2}$.
Lemma~\ref{app:nagata-near-closure} makes it an isomorphism near
$B_{1,2}$. Replace $X_{12}$ by its strict transform, equivalently
by the new graph closure. Since $Z_{1,2}$ is dense in $B_{1,2}$,
its inverse image is dense in $p_1^{-1}(B_{1,2})$. The second
projection maps that inverse image into $Z_{2,2}$, and therefore
maps all of $p_1^{-1}(B_{1,2})$ into $B_{2,2}$. Consequently
\begin{equation}
 X_{12}\cap(B_{1,2}\times_SB_{2,1})=\varnothing.
 \label{app:nagata-key-disjointness}
\end{equation}
This disjointness is what makes the valuative criterion work below.

Glue
\[
 W_i=M\coprod_{M_i}(X_i\setminus B_{i,j}).
\]
These schemes are separated over $S$. To check the mixed diagonal,
suppose a point of the graph of $M_i$ specializes in
$M\times_S(X_i\setminus B_{i,j})$. Properness of $\psi_i$
lifts this specialization to $V_i$. If its image in $M$ lay in
$M\setminus M_i=Z_j$, the lifted point would lie in
$Z_{i,j}\subseteq B_{i,j}$, contradicting the second coordinate.
The specialization therefore remains in the graph. That graph is
an immersion with constructible image; in the Noetherian situation
specialization-closed constructible subsets are closed. The graph
is thus a closed immersion, proving separatedness.

For a valuation ring $R$ with fraction field $K$, an $S$-map
$\operatorname{Spec}(K)\to D$ extends to both proper $X_i$.
The pair of extensions factors through $X_{12}$: the graph
closure contains its generic point, and
$\operatorname{Spec}(K)$ is schematically dense in
$\operatorname{Spec}(R)$. By
\eqref{app:nagata-key-disjointness}, the closed point cannot land
in both forbidden boundaries $B_{1,2}$ and $B_{2,1}$.
For some $i$ the extension therefore factors through
$X_i\setminus B_{i,j}\subseteq W_i$; an open containing the closed
point of a local scheme contains the whole scheme.

Apply Lemma~\ref{app:nagata-common-modification} to $M\subseteq W_i$.
It gives a separated finite-type $W$, covered by admissible
modifications $W_i'$. The valuation lift into $W_i$ lifts into
$W_i'$ by properness of the modification, hence into $W$.
The open $D$ is dense in $W$, and separatedness supplies uniqueness
of valuation lifts. The dense-source valuative criterion
(\stacks{0894}) now says $W\to S$ is universally closed.
Together with finite type and separatedness, this makes $W$ proper.
Its open $M$ yields the desired compactification; taking the
schematic closure of $M$ makes it dense. The construction is also
recorded in \stacks{0F40}.
\end{proof}

\subsection{Proof of Nagata's theorem}
\label{app:nagata-global-proof}

\begin{lemma}[Enough compactifiable affine opens]
\label{app:nagata-affine-opens}
For a finite-type scheme $X$ over a Noetherian scheme $S$ there is
a finite affine cover $X=\bigcup_{r=1}^m U_r$ whose total
intersection is dense in $X$. If $X\to S$ is separated, each
$U_r\to S$ has a projective compactification.
\end{lemma}

\begin{proof}
There are finitely many irreducible components of $X$.
For any point of $X$, the affine-neighbourhood lemma
\stacks{01ZX} supplies an affine open containing that point and all
the generic points of those components. Quasi-compactness selects
finitely many such opens. Their total intersection still contains
all generic points and is therefore dense.
Each $U_r$ is affine and its map to $S$ is of finite type.
The relative affine-space immersion lemma \stacks{04II} embeds it
as a locally closed subscheme of some $\mathbf A^{n_r}_S$.
Its schematic closure in $\mathbf P^{n_r}_S$ is proper, and
Lemma~\ref{app:nagata-closure} makes $U_r$ an open in that closure.
\end{proof}

\begin{proof}[Proof of Theorem~\ref{app:nagata-theorem}]
If $S$ is Noetherian, apply Lemma~\ref{app:nagata-affine-opens} and
induct on the unions $U_1\cup\cdots\cup U_r$.
At each step the overlap with $U_{r+1}$ contains the dense total
intersection. Proposition~\ref{app:nagata-two-open} gives a
compactification of the next union. The last union is $X$.

For a general qcqs $S$, first suppose $X\to S$ is separated and
of finite presentation. Absolute Noetherian approximation gives
\[
 S=\varprojlim_i S_i
\]
with $S_i$ Noetherian and affine transition maps (\stacks{01ZA}).
Finite-presentation descent gives a model $X_i\to S_i$, and at a
later stage it is separated (\stacks{01ZM}, \stacks{01ZQ}).
Compactify that model by the Noetherian result and base change
its compactification to $S$. Open immersions and proper morphisms
are stable under base change. Taking the schematic closure of $X$
then ensures density.

Finally let $X\to S$ be only separated of finite type.
The finite-presentation envelope theorem \stacks{01ZJ} provides a
closed immersion $X\hookrightarrow X_0$ with $X_0\to S$
separated of finite presentation. Compactify $X_0$ by the case
already proved, and take the schematic closure of $X$ in that
compactification. Restricted to the open $X_0$, this closure is
exactly the closed subscheme $X$. Thus $X$ is open and schematically
dense in the closure, which is proper over $S$.
\end{proof}

Both approximation results rest on an elementary affine fact: a
finitely presented algebra, its morphisms, and finitely many gluing
identities involve finitely many coefficients and equations, and so
are defined at a finite stage. The envelope step is needed because in
a finite-type algebra $R[x_1,\ldots,x_n]/I$ the ideal $I$ need not be
finitely generated; the finite-type scheme is therefore first embedded
as a closed subscheme of a finitely presented one, which descends.

\subsection{Chow's lemma and projective domination}

\begin{theorem}[Projective domination of a proper morphism]
\label{app:chow-projective}
Let $f:X\to S$ be proper with $S$ qcqs. There are a proper
surjection $\pi:\widetilde X\to X$ and a closed immersion
$i_S:\widetilde X\hookrightarrow\mathbf P^n_S$ such that
$f\pi=q_Si_S$. Moreover
\[
 i_X=(\pi,i_S):\widetilde X\hookrightarrow
 X\times_S\mathbf P^n_S=\mathbf P^n_X
\]
is a closed immersion. Both $\pi$ and $f\pi$ are projective.
\end{theorem}

The theorem holds over every qcqs base, in particular over every affine
one, with no Noetherian or finite-presentation assumption on $f$.

\begin{proof}
The version of Chow's lemma in \stacks{0202} gives a proper
surjection $\pi:\widetilde X\to X$ and an immersion
$i_S:\widetilde X\to\mathbf P^n_S$ for every separated finite-type
$X\to S$ with $S$ qcqs. Here $f\pi$ is proper. A map from a proper
$S$-scheme to a separated $S$-scheme is proper: its graph is closed
and the projection of its product to the target is proper. Therefore
$i_S$ is proper. A proper immersion is a closed immersion.

For the second assertion factor $i_X$ as the graph of $\pi$
followed by the base change of $i_S$:
\[
 \widetilde X\longrightarrow\widetilde X\times_SX
 \longrightarrow\mathbf P^n_S\times_SX.
\]
The first arrow is closed because $X\to S$ is separated, and the
second is closed because $i_S$ is. Their composite is $i_X$.
\end{proof}

The form of Chow's lemma used in the proof follows from the Noetherian
one, even when $X\to S$ is not of finite presentation. Embed $X$ as a
closed subscheme of a separated finitely presented $X_0$ by
\stacks{01ZJ}. Descend $X_0$ to a Noetherian stage, apply the
Noetherian form of Chow's lemma there, and base change its proper
surjective cover to $S$. Pull that cover back to the closed subscheme
$X$. The resulting source is closed in the source over $X_0$, so its
map to projective space is still an immersion, and properness and
surjectivity of its map to $X$ follow by base change.

Nagata's theorem and Chow's lemma play different roles. Nagata embeds
$X$ as an open subscheme of a proper scheme, keeping every point of
$X$; it supplies the compactifications and refinements used to define
$Rf_!$. Chow replaces a proper $X$ by a projective scheme mapping
properly and surjectively onto it; it supplies the projective
morphisms $\pi$ and $f\pi$ in the reduction of proper base change.

\section{Picard schemes and Jacobians}
\label{app:picard}

The cohomology of curves in Appendix~\ref{app:curves} uses the Picard
scheme and the Jacobian. They are constructed here from coherent
cohomology and descent. An open subscheme of a Hilbert scheme of
divisors represents an open subfunctor of the Picard functor, its
translates glue to the Picard scheme, and the degree-zero component is
the Jacobian, an abelian variety whose $n$-torsion for $n$ invertible
in $k$ is computed in Section~\ref{app:jacobian-multiplication}.
Coherent cohomology means cohomology of modules over the structure
sheaf, computed by affine \v Cech covers; the fppf topology is used
only to define the Picard functor.

\subsection{The Picard functor and rigidification}
\label{app:picard-functor}

\begin{definition}[Relative Picard functor]
For a morphism of schemes $f:C\to S$, let
\[
 \operatorname{Pic}_{C/S}
   =\bigl(T\longmapsto
       \operatorname{Pic}(C_T)/f_T^*\operatorname{Pic}(T)\bigr)^{\mathrm{fppf}}.
\]
Here $C_T=C\times_ST$ and the superscript denotes sheafification on
the big fppf site of $S$. Equivalently one may sheafify
$T\mapsto\operatorname{Pic}(C_T)$, since an invertible sheaf on $T$
becomes trivial on a Zariski cover. When this functor is represented by
a scheme, we call the representing scheme the Picard scheme.
\end{definition}

Sheafification is needed: invertible sheaves satisfy effective
descent, but automorphisms can prevent their isomorphism classes from
forming a sheaf, so neither the displayed quotient nor
$T\mapsto\operatorname{Pic}(C_T)$ need be an fppf sheaf
(\stacks{0B9L}). In the presence of a section, the next lemma
describes the sheafification directly.

\begin{lemma}[Rigidification along a section]
\label{app:picard-rigidification}
Suppose that $f$ has a section $\sigma:S\to C$ and that
\[
 \mathcal O_T\longrightarrow f_{T,*}\mathcal O_{C_T}
 \quad\text{is an isomorphism for every }T\to S.
 \tag{\ensuremath{\ast}}
\]
Then $\operatorname{Pic}_{C/S}(T)$ is naturally the set of isomorphism
classes of pairs $(\mathcal L,\alpha)$, with $\mathcal L$ invertible on
$C_T$ and $\alpha:\sigma_T^*\mathcal L\simeq\mathcal O_T$.
Every such pair has only the identity automorphism. There is a split
exact sequence
\[
 0\longrightarrow\operatorname{Pic}(T)
   \xrightarrow{f_T^*}\operatorname{Pic}(C_T)
   \longrightarrow\operatorname{Pic}_{C/S}(T)
   \longrightarrow0,
\]
whose retraction is $\sigma_T^*$. Normalization is given by
\[
 \mathcal L\longmapsto
 \mathcal L\otimes f_T^*(\sigma_T^*\mathcal L)^{-1}.
\]
\end{lemma}

\begin{proof}
An automorphism of an invertible sheaf is multiplication by a global
unit. Condition $(\ast)$ identifies that unit with a unit on $T$; a
rigidification forces it to be $1$. Given matching rigidified classes
on an fppf cover, choose local isomorphisms and rescale them to preserve
the rigidifications. Uniqueness forces the cocycle condition, so
effective descent for invertible sheaves glues the pairs. The
rigidified functor is therefore a sheaf. Normalization identifies it
with the quotient functor, while $\sigma_T^*f_T^*=\mathrm{id}$ gives
the splitting. This is the argument of \stacks{0B9N}.
\end{proof}

For a smooth proper geometrically connected curve $C$ over a field
$k$, one has $H^0(C,\mathcal O_C)=k$. Flat coherent base change gives
$(\ast)$ for $C\to\operatorname{Spec}k$, since every $k$-scheme is
flat over $k$. If $C$ has a $k$-point $P$, the preceding lemma applies
with $\sigma=P$; see \stacks{0B9U}.

\subsection{Riemann--Roch and coherent base change}
\label{app:picard-coherent-tools}

Fix a smooth projective geometrically connected curve $C/k$, and set
$g=\dim_kH^1(C,\mathcal O_C)$. Smooth proper geometrically connected
curves over a field are projective (\stacks{0A26}), so this includes
the proper curves of the main text.

\begin{theorem}[Riemann--Roch and coherent Serre duality]
\label{app:picard-riemann-roch}
For an invertible sheaf $\mathcal L$ on $C$,
\[
 \chi(C,\mathcal L)=\deg\mathcal L+1-g,
 \qquad
 H^1(C,\mathcal L)^\vee
     \simeq H^0(C,\omega_{C/k}\otimes\mathcal L^{-1}),
 \qquad \deg\omega_{C/k}=2g-2.
\]
These identifications commute with field extension. In particular, if
$\deg\mathcal L>2g-2$, then
\[
 H^1(C,\mathcal L)=0,
 \qquad h^0(C,\mathcal L)=\deg\mathcal L+1-g.
\]
\end{theorem}

This is coherent duality and Riemann--Roch for smooth curves
(\stacks{0B5B}); the duality is given by the residue trace on
differentials. For the displayed
vanishing, the dual sheaf has negative degree. A nonzero section of an
invertible sheaf on an integral curve has an effective zero divisor,
so a negative-degree sheaf has no nonzero section; Riemann--Roch then
gives the dimension formula.

\begin{lemma}[Coherent base change for invertible sheaves]
\label{app:picard-coherent-base-change}
Let $T$ be a locally Noetherian $k$-scheme and let $\mathcal L$ be
invertible on $C_T$. Then $Rf_{T,*}\mathcal L$ is locally represented
by a bounded complex of finite locally free $\mathcal O_T$-modules.
For any $a:T'\to T$, its derived pullback is canonically
$Rf_{T',*}\mathcal L_{T'}$. The locus where this complex is an
invertible module in degree $0$ is open; there $f_{T,*}\mathcal L$ is
invertible, its formation commutes with arbitrary base change, and
$R^1f_{T,*}\mathcal L=0$.
\end{lemma}

This is proper coherent base change for a flat finitely presented
module, together with the openness criterion for a perfect complex;
see \stacks{0B91} and \stacks{0B9S}. One way to obtain the complex is
to use a projective embedding, sufficiently large twists and a finite
affine \v Cech complex; Serre vanishing supplies finite locally free
models. After such a model has been obtained, openness is a matrix
calculation: near a point where the fibre has a single cohomology group
of rank $1$, cancel the differentials that are invertible on the
fibre. The remaining complex is a line bundle in degree $0$.

For an arbitrary test scheme $T$, the same open locus is obtained
locally by descending the finitely presented line bundle to a
finitely generated $k$-algebra and pulling back this calculation.
The derived base-change isomorphism makes the result independent of
the chosen model, so the construction below applies to all test
schemes.

\subsection{Hilbert schemes of points and an open chart}
\label{app:picard-hilbert}

For $d\geq0$, let $H_d=\operatorname{Hilb}^d_{C/k}$ be the functor of
closed subschemes $D\subset C_T$ which are finite locally free of
rank $d$ over $T$.

\begin{theorem}[Hilbert schemes of points on a smooth curve]
\label{app:picard-hilbert-theorem}
The functor $H_d$ is represented by a smooth projective geometrically
integral $k$-scheme of dimension $d$, with $H_0=\operatorname{Spec}k$.
Its objects are precisely the relative effective Cartier divisors
which are finite locally free of degree $d$. It has a universal
divisor $D_d\subset C\times_kH_d$ and an associated invertible sheaf
$\mathcal O(D_d)$.
\end{theorem}

Representability is Grothendieck's theorem that the Hilbert functor of
a projective scheme with fixed Hilbert polynomial is represented by a
projective scheme, applied with constant polynomial $d$;
\stacks{0B9A} and \stacks{0B9I} give the construction for curves. The
remaining properties can be seen directly.
On a smooth fibre, local rings at closed points are discrete valuation
rings, so a zero-dimensional closed subscheme is an effective
Cartier divisor. The flatness criterion for Cartier divisors gives
the relative assertion. At such a divisor, the Hilbert tangent space
is $\Hom(\mathcal I_D,\mathcal O_D)$, of dimension $d$, and the
obstruction group $\Ext^1(\mathcal I_D,\mathcal O_D)$ vanishes because
$\mathcal I_D$ is invertible and $D$ is finite. This proves smoothness.
The addition map $C^d\to H_d$ is surjective on geometric points:
an effective divisor is a sum of points with multiplicities.
The image of the geometrically integral $C^d$ is therefore
geometrically irreducible. Smoothness makes $H_d$ geometrically
integral.

Assume now that $C(k)$ contains a point $P$, and describe
$\operatorname{Pic}_{C/k}$ by rigidified invertible sheaves as in
Lemma~\ref{app:picard-rigidification}. Define a subfunctor
$F\subset\operatorname{Pic}_{C/k}$ by requiring
\[
 Rf_{T,*}\mathcal L\simeq\mathcal N[0]
 \quad\text{for some invertible }\mathcal N\text{ on }T.
\]
Here $\mathcal N$ may be any invertible sheaf on $T$, trivial or not.
Lemma~\ref{app:picard-coherent-base-change} shows that
$F\subset\operatorname{Pic}_{C/k}$ is representable by open immersions.

\begin{proposition}[Divisors of degree $g$ as an open chart]
\label{app:picard-abel-chart}
The functor $F$ is represented by an open subscheme $W\subset H_g$.
The morphism $W\to\operatorname{Pic}_{C/k}$ associates to a divisor its rigidified
invertible sheaf. Consequently, $F$ is smooth of dimension $g$.
\end{proposition}

\begin{proof}
On $C\times H_g$, normalize the universal invertible sheaf as
\[
 \mathcal L_g=\mathcal O(D_g)\otimes
 f_{H_g}^*(P_{H_g}^*\mathcal O(D_g))^{-1}.
\]
Let $W$ be the open locus where $Rf_{H_g,*}\mathcal L_g$ is a line
bundle in degree $0$. This gives a morphism $W\to F$.

Conversely, take a rigidified $\mathcal L$ in $F(T)$ and put
$\mathcal N=f_{T,*}\mathcal L$. The evaluation map
$f_T^*\mathcal N\to\mathcal L$ is a section of
$\mathcal L\otimes f_T^*\mathcal N^{-1}$. On each geometric fibre
it is a nonzero section of the unique one-dimensional space of
sections. Its zero scheme $D$ is thus a relative effective Cartier
divisor. Riemann--Roch gives
\[
 \deg\mathcal L_t
     =\chi(\mathcal L_t)-(1-g)=1-(1-g)=g.
\]
The divisor is proper with finite fibres and flat over $T$, hence
finite locally free of degree $g$. It defines $T\to H_g$.
Since $\mathcal O(D)\simeq\mathcal L\otimes f_T^*\mathcal N^{-1}$,
normalizing $\mathcal O(D)$ along $P_T$ recovers $\mathcal L$.
Multiplication of a section by a base unit does not change its zero
divisor; the two constructions are consequently inverse and commute
with base change. This proves $W\simeq F$; compare \stacks{0B9X}.
\end{proof}

\begin{lemma}[Translating into the chart]
\label{app:picard-twist}
Suppose that $k$ is algebraically closed. For every field extension
$K/k$ and every invertible sheaf $\mathcal L$ on $C_K$, there is an
invertible sheaf $\mathcal M_0$ on $C$ such that
\[
 h^0(C_K,\mathcal L\otimes\mathcal M_{0,K})=1,
 \qquad H^1(C_K,\mathcal L\otimes\mathcal M_{0,K})=0.
\]
\end{lemma}

\begin{proof}
Choose an ample $\mathcal M$ on $C$. For $r\gg0$, Serre vanishing
and Riemann--Roch make $\mathcal L\otimes\mathcal M_K^{\otimes r}$
have zero $H^1$ and a positive-dimensional space of sections.
Whenever its $h^0$ is larger than $1$, choose a nonzero section and
a point $x\in C(k)$ where that section is nonzero. Such a point
exists: the section has a finite zero divisor on $C_K$, while
$C(k)$ is infinite and its distinct points remain distinct on $C_K$.
The sequence
\[
 0\longrightarrow\mathcal F(-x_K)
  \longrightarrow\mathcal F
  \longrightarrow\mathcal F|_{x_K}\longrightarrow0
\]
shows that subtraction of $x_K$ lowers $h^0$ by exactly $1$ and
preserves $H^1=0$. Continue until $h^0=1$. All subtracted points
come from $k$, so the resulting twist has the form
$\mathcal M_0=\mathcal M^{\otimes r}(-x_1-\cdots-x_s)$ on $C$.
Normalization along $P$ does not change the fibre cohomology
dimensions. This is the twist argument of \stacks{0B9Y}.
\end{proof}

\subsection{The Picard scheme and the Jacobian}

\begin{theorem}[The Picard scheme and Jacobian of a geometric curve]
\label{app:picard-curve}
Let $k$ be algebraically closed and let $C/k$ be a smooth proper
connected curve of genus $g$. The fppf Picard functor is represented
by a smooth separated group scheme, locally of finite type over $k$,
with a decomposition
\[
 \operatorname{Pic}_{C/k}
   =\coprod_{d\in\mathbf Z}\operatorname{Pic}^d_{C/k}.
\]
Each component is a smooth proper geometrically connected variety of
dimension $g$, parametrizing invertible sheaves of degree $d$.
The identity component
$J(C)=\operatorname{Pic}^0_{C/k}$ is an abelian variety, and
\[
 J(C)(k)=\operatorname{Pic}^0(C),
 \qquad T_0J(C)=H^1(C,\mathcal O_C).
\]
The formation of this representing functor, its degree components
and $J(C)$ commutes with extension of the ground field.
\end{theorem}

\begin{proof}
Choose $P\in C(k)$. The preceding construction gives a representable
open subfunctor $F\subset\operatorname{Pic}_{C/k}$. For a fixed class
$M\in\operatorname{Pic}_{C/k}(k)$,
$M^{-1}F$ is another representable open subfunctor. Given a class
on a test scheme $T$, apply Lemma~\ref{app:picard-twist} to its fibre
over each residue field. The open loci where the corresponding
translated class lies in $F$ cover $T$. These charts therefore
cover $\operatorname{Pic}_{C/k}$. Their intersections are open in either chart, and their
transition isomorphisms satisfy the cocycle condition by functoriality.
Gluing constructs a scheme representing $\operatorname{Pic}_{C/k}$. This is the
representability criterion of \stacks{0B9Q}; the resulting curve
theorem is \stacks{0B9Z}.

Fibrewise Euler characteristic is locally constant for a proper flat
family of line bundles. Riemann--Roch thus makes degree locally
constant and defines the displayed open-and-closed decomposition.
The charts in $H_g$ and their translates show smoothness and
dimension $g$. The unit is a closed immersion for this locally
finite type group scheme over a field; translating the unit shows
that its diagonal is closed, so the scheme is separated.

The universal divisor defines the Abel map
\[
 H_g\longrightarrow\operatorname{Pic}^g_{C/k},
 \qquad D\longmapsto\mathcal O_C(D).
\]
On every geometric fibre, Riemann--Roch gives $h^0(\mathcal L)\geq1$
for a degree-$g$ line bundle, so this map is surjective. Its target is
quasi-compact, hence of finite type because the charts are locally of
finite type. A separated finite type target that is the surjective
image of the proper $H_g$ is proper. The image is also geometrically
irreducible. Translation
by $\mathcal O_C((d-g)P)$ identifies this target with every degree
component. In particular the degree-zero component is smooth,
proper, connected and commutative, hence an abelian variety;
see \stacks{0BA0}.

Finally, first-order deformations of the trivial line bundle have
transition functions $1+\epsilon a_{ij}$, whose cocycles modulo
coboundaries give $H^1(C,\mathcal O_C)$. Rigidification removes
scalar automorphisms. This gives the tangent-space assertion.
Picard sheafification and its universal functorial description commute
with base change, so the representing scheme does too. Degree and
the zero component have the same compatibility.
\end{proof}

\begin{corollary}[The degree sequence]
\label{app:picard-degree}
For $C/k$ as in Theorem~\ref{app:picard-curve}, there is an exact
sequence
\[
 0\longrightarrow J(C)(k)
   \longrightarrow\operatorname{Pic}(C)
   \xrightarrow{\deg}\mathbf Z\longrightarrow0.
\]
A choice $P\in C(k)$ splits it by
$d\mapsto\mathcal O_C(dP)$. The kernel and the degree map are
intrinsic; the splitting depends on $P$.
\end{corollary}

\begin{proof}
Degree is additive under tensor product. Its zero subgroup is
$J(C)(k)$ by the theorem, and $\deg\mathcal O_C(P)=1$.
\end{proof}

\begin{remark}[Ground fields and relative representability]
\label{app:picard-relative-scope}
For a smooth proper geometrically connected curve $C/k$ over an
arbitrary field, its degree-zero fppf Picard functor is represented
by an abelian variety $J/k$, and $J_K$ is the Jacobian of $C_K$.
This is the classical Jacobian descent theorem, stated precisely
in \cite[Theorem~1.1 and Proposition~1.9]{MilneJV}. If $C(k)$ is
nonempty, rigidification gives $J(k)=\operatorname{Pic}^0(C)$.
Without a rational point this identification need not hold; the
applications use only geometric points.

Over a general base, a separate relative theorem says that
$\operatorname{Pic}_{C/S}$ is an algebraic space if $C\to S$ is
proper, flat and of finite presentation and satisfies $(\ast)$
universally; see \stacks{0D2C}. A smooth proper curve with geometrically
connected fibres satisfies this structure-sheaf condition. The curve
calculations use only the field case proved above, where the Picard
functor is a scheme.
\end{remark}

\subsection{Torsion points of the Jacobian}
\label{app:jacobian-multiplication}

\begin{proposition}[Multiplication on an abelian variety]
\label{app:abelian-multiplication}
Let $B/k$ be an abelian variety of dimension $g$, and let $n>0$ be
invertible in $k$. Multiplication $[n]:B\to B$ is a finite \'etale
surjection of degree $n^{2g}$. In particular $B[n]=\ker[n]$ is a
finite \'etale group scheme of rank $n^{2g}$.
\end{proposition}

\begin{proof}
The coherent theorem of the cube implies, for a symmetric ample
invertible sheaf $\mathcal M$,
\[
 [n]^*\mathcal M\simeq\mathcal M^{\otimes n^2}.
\]
Such a sheaf exists because abelian varieties are projective:
symmetrize an ample sheaf by tensoring it with its pullback by $[-1]$.
If the kernel had positive dimension, restriction of $[n]^*\mathcal M$
to a positive-dimensional component would be both trivial and ample,
which is impossible. All fibres are translates of the kernel.
Properness therefore makes $[n]$ finite, and its closed image has
dimension $g$, hence is all of $B$. Its tangent map at the identity
is $n\cdot\mathrm{id}$, an isomorphism; translations give the same
at every point, so the morphism is \'etale. Finally the intersection
formula for a finite map gives
\[
 \deg[n]\,(c_1(\mathcal M)^g)
    =(c_1([n]^*\mathcal M)^g)
    =n^{2g}(c_1(\mathcal M)^g).
\]
The ample intersection number is positive, proving the degree
formula. See \stacks{0BFG}, \stacks{03RP}, and
\cite[Theorem~8.2]{MilneAV}.
\end{proof}

\begin{corollary}[Jacobian torsion and geometric extension]
\label{app:jacobian-torsion}
Let $k$ be algebraically closed, let $C/k$ be a smooth proper
connected curve of genus $g$, and let $n>0$ be invertible in $k$.
Then
\[
 J(C)[n](k)\simeq(\mathbf Z/n\mathbf Z)^{2g}.
\]
This isomorphism is not canonical. For every algebraically closed
extension $K/k$, the natural map
\[
 J(C)[n](k)\longrightarrow J(C_K)[n](K)
\]
is a canonical isomorphism. Multiplication by $n$ on $J(C)(k)$ is
surjective, and consequently
\[
 \operatorname{Pic}(C)[n]=J(C)[n](k),
 \qquad
 \operatorname{Pic}(C)/n\operatorname{Pic}(C)
      \xrightarrow[\deg]{\ \sim\ }\mathbf Z/n\mathbf Z.
\]
\end{corollary}

\begin{proof}
A finite \'etale scheme of rank $r$ over algebraically closed $k$
is a disjoint union of $r$ copies of $\operatorname{Spec}k$.
This description is stable under base extension, so the natural map
on $k$-points is a bijection.
Apply it to the kernel, using the base-change compatibility of the
Jacobian and of $[n]$.

For a prime $q\mid n$ and $q^a\mid n$, the group
$J[q^a](k)$ has order $q^{2ga}$, is killed by $q^a$, and its subgroup
killed by $q$ has order $q^{2g}$. The classification of finite
abelian groups therefore gives $2g$ cyclic factors, each of exponent
at most $q^a$. Their orders have product $q^{2ga}$, so every factor
has order $q^a$. Combining the prime-power parts gives the stated
$\mathbf Z/n\mathbf Z$-module structure.

Every fibre of the finite surjective $[n]$ over a $k$-point has a
$k$-point, proving divisibility. In the degree sequence of
Corollary~\ref{app:picard-degree}, the group $\mathbf Z$ has no
$n$-torsion and $J(C)(k)/n=0$. This proves the two assertions for
$\operatorname{Pic}(C)$.
\end{proof}

For $\ell$ invertible in $k$, the inverse system of these kernels
also gives $T_\ell J\simeq\Zl^{2g}$. To see this, choose a basis of
$J[\ell](k)$ and lift it successively under the surjective maps
$[\ell]:J[\ell^{a+1}](k)\to J[\ell^a](k)$. The lifted elements form
a basis at every level by reduction modulo $\ell$, hence a basis of
the inverse limit. These are the integral ranks in the curve
calculations; continuity of cohomology and change of coefficients
pass from them to $\mathcal O_E$, $E$, $\Ql$ and $\Qbar$.

The Kummer and Brauer-group calculations for curves in
Appendix~\ref{app:curves} use from this appendix the Picard functor,
its degree kernel, and the base-change behaviour of its finite torsion.

\section{Cohomology of curves}
\label{app:curves}

The proof of proper base change rests on the cohomology of curves
over an algebraically closed field (Theorem~\ref{pbc:curve-input}).
This appendix proves those facts.  It uses the Picard scheme and the
Jacobian of Appendix~\ref{app:picard}, the cohomology comparison
Theorem~\ref{pe:cohomology-comparison}, and the stalk calculation for
finite morphisms of Lemma~\ref{pbc:finite}; nothing else from the main
text enters.  Both the primes invertible in $k$ and the characteristic
prime are needed, since proper base change is proved for all torsion
coefficients.

Throughout this appendix, cohomology means cohomology on $C_{\proet}$,
and the sheaves are \emph{classical}.  By
Theorem~\ref{pe:classical-sheaves} a classical sheaf is the inverse
image under the comparison morphism $\nu_C$ of a sheaf on $C_{\et}$,
and Theorem~\ref{pe:cohomology-comparison} identifies its cohomology,
together with restriction and pullback maps, with \'etale cohomology.
All calculations below are therefore carried out on $C_{\et}$ and
transported back by $\nu_C^*$.  The classical hypothesis is used
exactly here: the assertions below concern classical sheaves only.

\subsection{Statement}

\begin{theorem}[Cohomology of curves over an algebraically closed field]
\label{app:curve-package}
Let $k$ be algebraically closed, $C$ a separated scheme of finite type
over $k$ of dimension at most one, and $F$ a classical torsion abelian
sheaf on $C_{\proet}$.  The following statements hold.
\begin{enumerate}[label=$(C_{\arabic*})$,leftmargin=3.5em]
\item $H^q(C,F)=0$ for $q>2$.
\item If $C$ is affine, $H^q(C,F)=0$ for $q>1$.
\item If $\operatorname{char}k=p>0$ and $F$ is $p$-primary, then
  $H^q(C,F)=0$ for $q>1$.
\item If $F$ is constructible with finite stalks and its torsion orders
  are prime to $\operatorname{char}k$, all its cohomology groups are finite.
\item If $C$ is proper and $F$ is constructible with finite stalks,
  all its cohomology groups are finite.
\item For an extension $k'/k$ of algebraically closed fields, the natural
  pullback maps $H^q(C,F)\to H^q(C_{k'},F|_{C_{k'}})$ are isomorphisms
  if $F$ is prime to the characteristic.
\item If $C$ is proper, these pullback maps are isomorphisms for every
  classical torsion $F$.
\item For every open $U\subseteq C$, the map
  $H^2(C,F)\to H^2(U,F|_U)$ is surjective.
\end{enumerate}
In characteristic zero every torsion order is allowed in $(C_4)$ and
$(C_6)$.
\end{theorem}

This is the pro-\'etale form of \stacks{03SC}.  The proof has three
stages.  First, for constant coefficients on a smooth connected curve
the cohomology is computed explicitly: by the Kummer sequence and the
Jacobian for primes invertible in $k$, and by the Artin--Schreier
sequence and coherent cohomology for the characteristic prime.  These
calculations carry the content of the theorem.  Second,
normalization and removal of finitely many points change cohomology
only by sheaves supported on finitely many points, which affect degrees
zero and one in a controlled way.  Third, a local system with finite
monodromy becomes a successive extension of constant sheaves on a
finite \'etale cover of degree prime to the torsion order, and the
trace splits off the original sheaf.  A d\'evissage and a passage to
filtered colimits then give arbitrary torsion sheaves.

\subsection{The multiplicative group and the Kummer sequence}

On a smooth connected curve $C$ over an algebraically closed field let
$\mathcal G_m^{\mathrm{cl}}=\nu_C^*\mathbf G_{m,C_{\et}}$ be the
classical multiplicative sheaf on $C_{\proet}$.

\begin{lemma}[Cohomology of the multiplicative group]
\label{app:curve-gm}
For a smooth separated curve $C$ of finite type over an algebraically
closed field, one has
\[
 H^1(C,\mathcal G_m^{\mathrm{cl}})=\operatorname{Pic}(C),
 \qquad
 H^q(C,\mathcal G_m^{\mathrm{cl}})=0\quad(q\geq2).
\]
The identifications are natural under restriction and extension of
the algebraically closed ground field.
\end{lemma}
\begin{proof}
We compute on $C_{\et}$.  Descent for invertible modules identifies
$\mathbf G_m$-torsors with line bundles, which gives the assertion in
degree one.  In degrees at least two we use the theorem that
$H^q(C_{\et},\mathbf G_m)=0$ for $q\geq2$ on a smooth curve over an
algebraically closed field (\stacks{03RM}).  Its proof combines the
divisor sequence of rational functions and their discrete valuations
with the Galois cohomology of the function field $K=k(C)$; the
geometric input is Tsen's theorem, which we recall.

Tsen's theorem says that $K$ is $C_1$: a homogeneous form of degree
$d$ in $m>d$ variables has a nonzero zero.  On a smooth proper model,
bound the poles of its coefficients by a divisor $D$ and let the
unknowns range over $H^0(C,\cO_C(NH))$ for an ample divisor $H$.  The
value of the form then lies in $H^0(C,\cO_C(dNH+D))$.  By
Riemann--Roch there are $mN\deg H+O(1)$ unknown coefficients and
$dN\deg H+O(1)$ homogeneous equations.  For large $N$ the first number
is larger, so over the algebraically closed field $k$ the equations
have a nonzero solution (\stacks{03RD}).

The reduced norm form of a division algebra of degree $d>1$ over $K$
has degree $d$ in $d^2>d$ variables and does not vanish on a nonzero
element.  Tsen's theorem therefore gives $\operatorname{Br}(K)=0$
(\stacks{03RF}).  The vanishing of the Brauer group is one ingredient;
\stacks{03RM} combines it with the discrete-valuation calculation and
the cohomology of the function field to obtain vanishing in every
degree $q\geq2$.

Finally Theorem~\ref{pe:cohomology-comparison} transports these
identifications to $C_{\proet}$.  The comparison maps, the torsor
identifications and the restriction maps are functorial, which gives
the stated naturality.
\end{proof}

\begin{proposition}[Smooth proper curves with invertible torsion]
\label{app:curve-kummer}
Let $C$ be smooth, proper and connected over an algebraically closed
field $k$, with genus $g$, and let $n\geq1$ be invertible in $k$.
Write $J=\operatorname{Pic}^0_{C/k}$.  There are canonical
identifications
\[
 \begin{aligned}
 H^0(C,\mu_n)&=\mu_n(k),&
 H^1(C,\mu_n)&=J[n](k),\\
 H^2(C,\mu_n)&=\Z/n,&
 H^q(C,\mu_n)&=0\quad(q>2).
 \end{aligned}
\]
The degree-two isomorphism sends the Kummer boundary of a line bundle
$L$ to $\deg L\bmod n$.  These identifications commute with extension
of algebraically closed fields.
\end{proposition}
\begin{proof}
Pull back the Kummer sequence of $C_{\et}$ by $\nu_C^*$; it is the exact
sequence of sheaves on $C_{\proet}$
\[
 0\longrightarrow\mu_n\longrightarrow\mathcal G_m^{\mathrm{cl}}
 \xrightarrow{(\cdot)^n}\mathcal G_m^{\mathrm{cl}}
 \longrightarrow0.
\]
Exactness is local: for a unit $a$, the equation $T^n=a$ defines an
\'etale cover, since $n$ and $T$ are invertible.  Global units on the
proper connected curve are $k^\times$, an $n$-divisible group.
Lemma~\ref{app:curve-gm} therefore turns the long exact sequence into
\[
 H^1(C,\mu_n)=\operatorname{Pic}(C)[n],\qquad
 H^2(C,\mu_n)=\operatorname{Pic}(C)/n,
\]
with all higher groups zero.

The degree sequence of Corollary~\ref{app:picard-degree} is
\[
 0\longrightarrow J(k)\longrightarrow\operatorname{Pic}(C)
 \xrightarrow{\deg}\Z\longrightarrow0.
\]
Multiplication by $n$ on $J(k)$ is surjective, with kernel $J[n](k)$,
and $J[n]$ is finite \'etale of rank $n^{2g}$
(Corollary~\ref{app:jacobian-torsion}).  The snake lemma gives the two
identifications asserted above, and $J[n](k)$ is noncanonically
isomorphic to $(\Z/n)^{2g}$.  This is the calculation of \stacks{03RQ}.

Under an algebraically closed extension $k'/k$, $J$ and $J[n]$ base
change to the corresponding objects for $C_{k'}$.  The map
$J[n](k)\to J[n](k')$ is bijective because a finite \'etale scheme over
an algebraically closed field is a finite disjoint union of copies of
its spectrum; the same holds for roots of unity.  Pullback of a line
bundle preserves degree.  Hence the pullback maps on cohomology
themselves are isomorphisms.
\end{proof}

\begin{remark}[Twists and the projective line]
\label{app:curve-twists}
The canonical formula is $H^2(C,\mu_n)=\Z/n$, with
$\mu_n=(\Z/n)(1)$.  Consequently
\[
 H^2(C,\Z/n)=(\Z/n)(-1).
\]
A primitive $n$th root of unity trivializes this twist; the
identifications above are functorial without such a choice.  For
$C=\PP^1_k$ the Jacobian is zero, so the cohomology with constant
coefficients is concentrated in degrees zero and two, and the Kummer
class of $\cO(1)$ maps to $1$ in $H^2(\PP^1_k,\mu_n)$.  This is the
normalization of the geometric fibre used in
Section~\ref{pbc:projective-line-section}.
\end{remark}

\subsection{Punctured curves with invertible torsion}

\begin{proposition}[Cohomology of a smooth affine curve]
\label{app:curve-affine-kummer}
Let $\overline C$ be a smooth proper connected curve of genus $g$
over an algebraically closed field.  Let
$U=\overline C\setminus D$, where $D$ consists of $r\geq1$ closed
points, and let $n$ be invertible.  There is a natural exact sequence
\begin{equation}
 0\longrightarrow H^1(\overline C,\mu_n)
 \longrightarrow H^1(U,\mu_n)
 \xrightarrow{\partial}(\Z/n)^D
 \xrightarrow{\sum}\Z/n\longrightarrow0.
 \label{app:curve-residue-sequence}
\end{equation}
Moreover $H^0(U,\mu_n)=\mu_n(k)$ and $H^q(U,\mu_n)=0$ for $q>1$.
The group $H^1(U,\mu_n)$ is noncanonically $(\Z/n)^{2g+r-1}$, and
the pullback maps on all these groups are isomorphisms under
algebraically closed extension.
\end{proposition}
\begin{proof}
Every smooth connected nonproper curve is affine and arises in this
way: take a projective model of its function field and normalize it.
Over a perfect field regular curves are smooth, and the complement is
a finite nonempty set of points.

Divisors on the regular curve give
\[
 \operatorname{Pic}(U)=\operatorname{Pic}(\overline C)/
 \langle\cO_{\overline C}(x):x\in D\rangle.
\]
Since one of these line bundles has degree one, $J(k)$ surjects onto
$\operatorname{Pic}(U)$, so $\operatorname{Pic}(U)$ is $n$-divisible.
The Kummer sequence and Lemma~\ref{app:curve-gm} now give the vanishing
above degree one.

The sequence \eqref{app:curve-residue-sequence} comes from line
bundles.  A $\mu_n$-torsor on $U$ is a pair
$(L,\alpha\colon L^{\otimes n}\simeq\cO_U)$.  Extend $L$ to a line
bundle $\overline L$ on $\overline C$ and let $s=\alpha^{-1}(1)$ be the
corresponding rational section of $\overline L^{\otimes n}$.  Its
divisor is supported on $D$.  Its orders modulo $n$ do not depend on
the extension, and they define $\partial(L,\alpha)$.  Their sum is zero
modulo $n$, since the divisor has degree divisible by $n$.  If all
orders are divisible by $n$, twist $\overline L$ by the negative of
that divisor divided by $n$; the trivialization then extends without
zero or pole.  This proves exactness at $H^1(U,\mu_n)$.  For
injectivity, a trivialization on $U$ of a torsor on $\overline C$
gives a rational section of its line bundle whose $n$th power has no
zero or pole on $\overline C$.  The divisor of the section is
therefore zero, and it trivializes the torsor on $\overline C$.

For surjectivity of $\partial$ onto the kernel of $\sum$, lift a
residue vector with sum zero modulo $n$ to an integral divisor $B$
supported on $D$ whose degree is divisible by $n$.  There is a line
bundle $\overline L$ with $\overline L^{\otimes n}\simeq\cO_{\overline C}(B)$:
choose the degree first and then use divisibility of $J(k)$ to correct
the degree-zero difference.  On $U$ this is a Kummer torsor with the
prescribed residues.  The summation map is surjective since $D$ is
nonempty.  This proves the sequence; see also \stacks{03RR}.

The left term and the kernel of summation are free $\Z/n$-modules of
ranks $2g$ and $r-1$, so the resulting short exact sequence splits
noncanonically.  For naturality, the compactification, boundary
points, line bundles and valuation orders all base change.  The
leftmost pullback map is an isomorphism by
Proposition~\ref{app:curve-kummer}, and the residue groups and their
summation maps are unchanged.  The five lemma applied to the two
exact sequences shows that the middle pullback map is an isomorphism.
In particular no descent of line bundles from $U_{k'}$ to $U$ is
needed.
\end{proof}

For the constant sheaf $\mathbf F_r$ with $r\ne\operatorname{char}k$,
the two preceding propositions give $(C_1)$--$(C_8)$ on smooth curves.
A primitive $r$th root of unity identifies $\mathbf F_r$ with $\mu_r$;
keeping the twist makes the pullback maps independent of that choice.
For $(C_8)$, an open in a smooth connected proper curve is either the
whole curve or an affine curve, with $H^2=0$; an open in a smooth affine
curve is affine.

\subsection{The Artin--Schreier sequence}

\begin{proposition}[Characteristic torsion on a smooth proper curve]
\label{app:curve-artin-schreier}
Let $C$ be a smooth proper connected curve of genus $g$ over an
algebraically closed field of characteristic $p>0$.  Then
\[
 H^0(C,\mathbf F_p)=\mathbf F_p,\qquad
 H^1(C,\mathbf F_p)=
 \ker(F-1:H^1(C,\cO_C)\to H^1(C,\cO_C)),
\]
and $H^q(C,\mathbf F_p)=0$ for $q>1$.  The group $H^1(C,\mathbf F_p)$
is finite, of $\mathbf F_p$-dimension at most $g$, and its pullback map
is an isomorphism under algebraically closed extension.
\end{proposition}
\begin{proof}
On $C_{\et}$ the Artin--Schreier sequence
\[
 0\longrightarrow\mathbf F_p\longrightarrow\mathbf G_a
 \xrightarrow{F-1}\mathbf G_a\longrightarrow0
\]
is exact; surjectivity is local because $T^p-T-a$ has derivative $-1$.
The \'etale cohomology of $\mathbf G_a$ is the coherent cohomology of
$\cO_C$, by faithfully flat descent on affine opens and the coherent
\v Cech resolution.  Coherent cohomology vanishes above dimension one,
and properness makes $H^1(C,\cO_C)$ a finite-dimensional $k$-vector
space.

The cohomology sequence is then controlled by a fact of linear
algebra: if $F$ is a Frobenius-linear endomorphism of a
finite-dimensional $k$-vector space $V$, then $F-1$ is surjective and
its kernel is the set of $k$-points of a finite \'etale scheme of
order $p^f$ with $f\leq\dim_k V$.  To see this, reduce along a finite
filtration to the part where $F$ is nilpotent and the part where $F$ is
bijective.  On the nilpotent part $F-1$ is an additive polynomial
automorphism.  On a bijective part of dimension $f$, choose a basis,
invert the matrix of $F$, and write the fibre equations as
\[
 x_i^p-\sum_j b_{ij}x_j=c_i\qquad(1\leq i\leq f),
\]
where $(b_{ij})$ is invertible.  The coordinate algebra is free of
dimension $p^f$, with basis the monomials having each exponent less
than $p$.  Its Jacobian matrix is invertible, so it is finite \'etale
and has $p^f$ points over the algebraically closed field
(\stacks{0A3L}).

Apply this to $H^1(C,\cO_C)$ and to $H^0(C,\cO_C)=k$.  The long exact
Artin--Schreier sequence gives the asserted formulas and vanishing.
Under $k'/k$, flat base change for coherent cohomology identifies the
vector space and its Frobenius-linear map with their scalar
extensions.  The kernel is a finite \'etale scheme, whose $k$-points
and $k'$-points agree, so the pullback map itself is an isomorphism
(\stacks{0HAN}, \stacks{0A3P}).  Theorem~\ref{pe:cohomology-comparison}
returns the statement to $C_{\proet}$.
\end{proof}

\begin{lemma}[Characteristic torsion on an affine curve]
\label{app:curve-affine-artin-schreier}
If $U$ is affine over an algebraically closed field of characteristic
$p$, then, for the classical constant sheaf $\mathbf F_p$,
\[
 H^1(U,\mathbf F_p)=
 \Gamma(U,\cO_U)/(F-1)\Gamma(U,\cO_U),\qquad
 H^q(U,\mathbf F_p)=0\quad(q>1).
\]
\end{lemma}
\begin{proof}
Quasi-coherent cohomology on an affine scheme vanishes in positive
degrees, so the long exact Artin--Schreier sequence on $U_{\et}$ gives
both claims.  Theorem~\ref{pe:cohomology-comparison} transports them to
$U_{\proet}$.
\end{proof}

\begin{example}[The affine line in characteristic $p$]
\label{app:curve-wild-example}
For $U=\A^1_k$ in characteristic $p$ one obtains
\[
 H^1(U,\mathbf F_p)=k[t]/\{h^p-h:h\in k[t]\}.
\]
Every class has a unique representative with zero constant term
and all positive exponents prime to $p$.  Exponents divisible by
$p$ are removed using $p$th roots in the perfect field; the constant
term is removed by surjectivity of $a\mapsto a^p-a$ on $k$.
For uniqueness, a nonconstant polynomial $h^p-h$ has leading
degree divisible by $p$.  Hence the quotient is infinite: the classes
of $t^m$ for positive $m$ prime to $p$ are $\mathbf F_p$-linearly
independent.  For an algebraically closed extension $k\subsetneq k'$
and $a\in k'\setminus k$, the class of $at$ is not in the image of the
quotient over $k$.  Thus finiteness and field-extension invariance in
degree one fail for characteristic torsion on an affine curve, which
is why $(C_5)$ and $(C_7)$ require properness.
\end{example}

For constant $\mathbf F_p$ on smooth curves, the proper and affine
calculations prove $(C_1)$, $(C_2)$, $(C_3)$, $(C_5)$ and $(C_7)$ in
their respective scopes, and $(C_8)$ follows from $H^2=0$.  The
statements $(C_4)$ and $(C_6)$ concern torsion prime to the
characteristic and say nothing in this case.

\subsection{Normalization and removal of finitely many points}

Sheaves supported on finitely many points have no higher cohomology,
and this lets one pass between a curve, its normalization and a dense
open subset at the cost of an error in degrees zero and one.

\begin{lemma}[Finite-support sheaves and finite maps]
\label{app:curve-finite-support}
A classical sheaf supported on finitely many closed points of a
curve over an algebraically closed field has no positive cohomology.
Its sections are the direct sum of its stalks, unchanged under
algebraically closed extension.  For a finite morphism $h\colon Y\to C$,
the direct image $h_*$ of classical sheaves is exact, commutes with
arbitrary base change, and satisfies $H^q(C,h_*G)=H^q(Y,G)$.
\end{lemma}
\begin{proof}
Work first on the \'etale sites.  Over a strict localization of $C$,
a finite scheme is a finite disjoint union of strictly henselian local
schemes.  Every \'etale cover of such a local scheme has a section over
its closed geometric point, so global sections are exact there.  This
gives vanishing of the higher direct images under $h$ and the usual
stalk formula, which in turn gives finite base change (\stacks{03QP},
\stacks{0959}; this is the calculation of Lemma~\ref{pbc:finite}).  A
finite set of algebraically closed points has exact sections, which
gives the first assertion.  Comparison of classical direct images and
Theorem~\ref{pe:cohomology-comparison} transport these statements to
the pro-\'etale site.
\end{proof}

\begin{lemma}[Normalization of a constant sheaf]
\label{app:curve-normalization}
Let $C$ be a reduced curve of finite type over an algebraically closed
field and $a\colon C^{\mathrm{nor}}\to C$ its normalization.  For a
finite abelian group $M$, the natural map
\[
 \underline M_C\longrightarrow a_*\underline M_{C^{\mathrm{nor}}}
\]
is injective with cokernel supported on finitely many points; at a
point with $b$ geometric branches its cokernel is $M^b/\Delta M$.
These cokernels and the resulting exact sequences commute with
algebraically closed field extension.
\end{lemma}
\begin{proof}
Normalization is finite for reduced schemes of finite type over a
field.  On each one-dimensional component its normal model is regular,
hence smooth over the perfect field; isolated components of dimension
zero are algebraically closed points.  The normalization is an
isomorphism over a dense open of each component, so its
nonisomorphism locus is finite.

By the stalk calculation of Lemma~\ref{app:curve-finite-support}, the
displayed map at a geometric point is the diagonal $M\to M^b$, with
$b\geq1$.  Thus its kernel is zero and its cokernel has the asserted
form.  More generally, any map between the two constructible sheaves
which is an isomorphism over their dense regular open has kernel and
cokernel supported in the finite complement.

The normalization remains smooth after algebraically closed extension
and is a finite birational normal model of the reduced base change;
hence it is its normalization.  Finite base change identifies the
fibres and their branch sets, so the diagonal maps and their cokernels
are the scalar-extended ones.
\end{proof}

From the normalization sequence and
Lemma~\ref{app:curve-finite-support}, one gets isomorphisms
$H^q(C,\underline M)\simeq H^q(C^{\mathrm{nor}},\underline M)$
for $q\geq2$, and a natural exact sequence controlling degrees zero
and one by a finite group.  Since the normalization of an affine or
proper curve is again affine or proper, all assertions of
Theorem~\ref{app:curve-package} for constant $\mathbf F_r$ follow from
the smooth calculations.  For $(C_8)$ apply this comparison to $C$ and
to the inverse image of an open subset; the restriction square then
reduces surjectivity to the smooth components.

\begin{lemma}[Removing finitely many points]
\label{app:curve-dense-open}
Let $j\colon U\hookrightarrow C$ be a dense open in a reduced curve and
$i\colon Z\hookrightarrow C$ its finite complementary closed subscheme.
For a classical constructible sheaf $F$ there is an exact sequence
\[
 0\longrightarrow j_!j^*F\longrightarrow F
 \longrightarrow i_*i^*F\longrightarrow0.
\]
Consequently the assertions of Theorem~\ref{app:curve-package} for $F$
and for $j_!j^*F$ are equivalent, each within its stated scope.
\end{lemma}
\begin{proof}
The exact sequence is the open--closed sequence of classical sheaves;
it can be checked on geometric stalks on $C_{\et}$ and then pulled
back.  The last term has finite sections and no positive cohomology,
so the long exact sequence gives
\[
 \begin{split}
 0\to H^0(C,j_!j^*F)\to H^0(C,F)\to H^0(Z,i^*F)\to{}\qquad\qquad\\
 H^1(C,j_!j^*F)\to H^1(C,F)\to0,
 \end{split}
\]
and isomorphisms in degrees at least two.  The sequence is natural
under field extension and its point term is unchanged, so the five
lemma gives the equivalence of the invariance assertions.  For $(C_8)$
apply the same identifications in degree two to each open subset of
$C$.  See also \stacks{0GJA}.
\end{proof}

In particular $j_!\underline{\mathbf F_r}$ satisfies the assertions of
Theorem~\ref{app:curve-package} when $j$ is dense, by
Lemma~\ref{app:curve-normalization} and the constant calculation.  For
an arbitrary open, apply this on the reduced union of the components it
meets, which is a closed subscheme, and use exactness of finite direct
image.  Alternatively, shrink to the dense regular open of the
normalization and use Lemma~\ref{app:curve-finite-support}; this is the
route of \stacks{0A3Q}.

\subsection{Finite monodromy and the trace}

A local system with finite coefficients has finite monodromy, and
passing to the cover defined by a Sylow subgroup makes it unipotent.

\begin{lemma}[Reduction to a unipotent local system]
\label{app:curve-sylow-trace}
Let $U$ be connected and $G$ a finite classical local system of
$\mathbf F_r$-vector spaces.  There is a finite \'etale cover
$v\colon V\to U$ of degree $d$ prime to $r$ such that $v^*G$ has a
finite filtration with successive quotients the constant
sheaf $\mathbf F_r$.  Moreover $G$ is a direct summand of
$v_*v^*G$.
\end{lemma}
\begin{proof}
By Proposition~\ref{pe:finite-local-systems}, $G$ is the pullback of a
finite \'etale local system, which corresponds to a finite continuous
representation of the fundamental group.  Let $H$ be the finite image
of this representation, choose an $r$-Sylow subgroup $P$, and take the
cover associated to its inverse image in the fundamental group.  Its
degree is $[H:P]$, prime to $r$, and the restricted action factors
through $P$.  The augmentation ideal of $\mathbf F_r[P]$ is nilpotent,
so the simple modules of the finite $r$-group $P$ are trivial of
dimension one, and a composition series gives the required filtration
(\stacks{0A3R}).

The unit $G\to v_*v^*G$ is the diagonal on a fibre, and the trace map
$\operatorname{Tr}\colon v_*v^*G\to G$ adds the $d$ entries of a fibre.
Hence the composite is multiplication by $d$, which is invertible on an
$\mathbf F_r$-sheaf, and $d^{-1}\operatorname{Tr}$ is a retraction.
\end{proof}

To apply this to $j_!G$ on a reduced curve $C$, take for $U$ a dense
smooth open on which $G$ is locally constant.  Normalizing $C$ in the
finite extensions of function fields defined by $V$ gives a finite
morphism $\bar v\colon Y\to C$ with $V$ the inverse image of $U$, and an
open immersion $j'\colon V\hookrightarrow Y$.  The cartesian square
formed by these maps and the stalk formula for finite morphisms give
\[
 j_!v_*v^*G=\bar v_*j'_!v^*G,
\]
which is also the finite open--closed exchange of \stacks{0GJ1}.  If
$C$ is proper, $Y$ is proper; if $C$ is affine, $Y$ is affine.

The filtration of $v^*G$, exactness of $j'_!$, and the calculation for
constant sheaves extended by zero therefore give the assertions of
Theorem~\ref{app:curve-package} for $j'_!v^*G$.  Exactness of
$\bar v_*$ passes them to $j_!v_*v^*G$, and the trace splitting passes
them to its direct summand $j_!G$.  For invariance under field
extension, the splitting maps commute with ground-field extension, so
the isomorphism for the larger sheaf is block diagonal for the induced
decomposition.  This is the argument of \stacks{03SD}.

If $C$ is not equidimensional, first choose a dense open which is
smooth and is the disjoint union of the regular loci of the
one-dimensional components and the isolated points.  The isolated
points have exact sections.  Treat each connected one-dimensional piece
as above; finite direct sums preserve all the assertions.

\subsection{Proof of the theorem}

The assertions of Theorem~\ref{app:curve-package} pass along short
exact sequences.  Let $0\to F_1\to F\to F_2\to0$ be exact, with the
assertions known for $F_1$ and $F_2$.  Vanishing and finiteness follow
from the long exact sequence, and invariance under field extension
from the map between the two long exact sequences and the five lemma.
For $(C_8)$, take $b\in H^2(U,F)$.  Lift its image in $H^2(U,F_2)$ to
$c_2\in H^2(C,F_2)$.  The obstruction to lifting $c_2$ to $H^2(C,F)$
lies in $H^3(C,F_1)$, which vanishes by $(C_1)$.  After choosing a lift
$c$, the difference $b-c|_U$ comes from $H^2(U,F_1)$, and $(C_8)$ for
$F_1$ lifts this difference.  Adding that lift to $c$ gives the
required global class (\stacks{0A5C}).

\begin{proof}[Proof of Theorem~\ref{app:curve-package}]
By Theorem~\ref{pe:cohomology-comparison} we may work on $C_{\et}$.  A
universal homeomorphism does not change the \'etale topos, so we may
replace $C$ by $C_{\mathrm{red}}$.  A constructible torsion sheaf has a
finite exponent, because it has finitely many isomorphism types of
finite stalks.  It is the direct sum of its finitely many primary
parts, and a finite filtration of an $r$-primary part reduces to
constructible $\mathbf F_r$-sheaves.  By the stability under extensions
just proved, it suffices to treat such sheaves.

Choose a dense smooth open $j\colon U\hookrightarrow C$ on which the
sheaf is finite locally constant.  Lemma~\ref{app:curve-dense-open}
reduces to its extension by zero.  On the connected components of $U$,
Lemma~\ref{app:curve-sylow-trace} and finite direct image reduce to
constant $\mathbf F_r$ extended by zero on curves finite over $C$.
Normalization and Lemma~\ref{app:curve-finite-support} reduce these
sheaves to the constant calculations on smooth curves: the Kummer
sequence for $r\ne\operatorname{char}k$ and the Artin--Schreier
sequence for $r=\operatorname{char}k$.  All reductions preserve the
hypothesis that $C$ is proper or affine.  This proves the theorem for
constructible torsion sheaves.

An arbitrary torsion sheaf on the qcqs scheme $C$ is a filtered
colimit of constructible torsion sheaves, and cohomology commutes with
filtered colimits (\stacks{03SA}, \stacks{03Q5}).  For torsion prime to
the characteristic or $p$-primary torsion, the constructible sheaves
may be taken with the same kind of torsion, by splitting off primary
parts before passing to the colimit.  Pullback and restriction also
commute with filtered colimits, so vanishing and the pullback
isomorphisms pass to the colimit.  Surjectivity in $(C_8)$ passes to
it because a representative of a class already occurs at one stage.
Finiteness is asserted only for constructible sheaves.  Finally
Theorem~\ref{pe:cohomology-comparison} returns these conclusions and
their maps to $C_{\proet}$.
\end{proof}

\subsection{Separably closed fields}

\begin{corollary}[Proper curves over separably closed fields]
\label{app:curve-invariance}
Let $k\subseteq k'$ be separably closed fields, let $C$ be proper
of dimension at most one over $k$, and let $F$ be a classical
torsion abelian sheaf on $C_{\proet}$.  Its natural pullback map
\[
 H^q(C,F)\longrightarrow H^q(C_{k'},F|_{C_{k'}})
\]
is an isomorphism in every degree.
\end{corollary}

Here, as in $(C_7)$, the torsion orders of $F$ are arbitrary.

\begin{proof}
Choose compatible algebraic closures $\bar k\subseteq\bar k'$.  The
extensions $\bar k/k$ and $\bar k'/k'$ are purely inseparable, so the
base-change morphisms $C_{\bar k}\to C$ and $C_{\bar k'}\to C_{k'}$ are
universal homeomorphisms and induce equivalences of \'etale topoi
(\stacks{03SI}).  By Theorem~\ref{pe:cohomology-comparison}, the
vertical maps in
\[
 \begin{CD}
 H^q(C,F) @>>> H^q(C_{k'},F|)\\
 @V{\sim}VV @VV{\sim}V\\
 H^q(C_{\bar k},F|) @>>> H^q(C_{\bar k'},F|)
 \end{CD}
\]
are isomorphisms.  The bottom map is an isomorphism by $(C_7)$ of
Theorem~\ref{app:curve-package}.  All maps are pullback maps, so the
square commutes and the top map is an isomorphism (\stacks{0A5E}).
\end{proof}

In the main text, Theorem~\ref{pbc:curve-input} and
Lemma~\ref{pbc:curve-field-extension} restate
Theorem~\ref{app:curve-package} and Corollary~\ref{app:curve-invariance}
for the comparison of closed fibres, and the cohomology of $\PP^1$ over
an algebraically closed field, with its normalized Kummer class, comes
from Proposition~\ref{app:curve-kummer}.  The results thus enter the
proof of proper base change in the following order:
\[
 \begin{gathered}
 \text{Picard geometry, Tsen and coherent cohomology}\\
 \Downarrow\\
 \text{smooth constant curve calculations}\\
 \Downarrow\\
 \text{normalization, finite trace and torsion d\'evissage}\\
 \Downarrow\\
 \text{curve invariance and the closed-fibre calculation}\\
 \Downarrow\\
 \text{proper base change, then the relative six-functor results}.
 \end{gathered}
\]

\phantomsection
\addcontentsline{toc}{section}{References}
\begin{thebibliography}{Milne AV}
\bibitem[BS]{BS}
B.~Bhatt and P.~Scholze,
\emph{The pro-\'etale topology for schemes},
Ast\'erisque \textbf{369} (2015), 99--201.
\href{https://people.mpim-bonn.mpg.de/scholze/proetale.pdf}{Author's manuscript}.

\bibitem[Eke]{Ekedahl}
T.~Ekedahl, \emph{On the adic formalism},
in \emph{The Grothendieck Festschrift}, Vol.~II,
Progress in Mathematics \textbf{87}, Birkh\"auser, 1990, 197--218.

\bibitem[SGA4]{SGA4}
M.~Artin, A.~Grothendieck and J.-L.~Verdier (eds.),
\emph{Th\'eorie des topos et cohomologie \'etale des sch\'emas},
SGA~4, Lecture Notes in Mathematics \textbf{269}, \textbf{270},
\textbf{305}, Springer, 1972--1973.

\bibitem[SGA5]{SGA5}
A.~Grothendieck (dir.),
\emph{Cohomologie $\ell$-adique et fonctions $L$},
SGA~5, Lecture Notes in Mathematics \textbf{589}, Springer, 1977.

\bibitem[Proet]{ProetaleBlueprint}
The \texttt{chrisflav/proetale} contributors, \emph{Pro-\'etale cohomology},
online notes, revision \texttt{c815a076}.
\href{https://github.com/chrisflav/proetale/tree/c815a07678fc5937809e6d4c8af78a107b064056/blueprint/ProetaleBlueprint/Chapters}{Source}.

\bibitem[BP]{BlueprintPBC}
Formal Arithmetic Geometry contributors, \emph{The Weil conjectures},
online notes, section \emph{Proper Base Change}, revision \texttt{f93e404}.
\href{https://github.com/formal-arithmetic-geometry/weil-conjectures-blueprint/tree/f93e4044b15d3e66893be38fa472491956ffaf83/blueprint/sec-proper-base-change}{Source}.

\bibitem[GeoBP]{BlueprintGeometry}
Formal Arithmetic Geometry contributors, \emph{The Weil conjectures},
online notes, chapters \emph{Nagata Compactification}, \emph{Picard Framework},
and \emph{Curve Estimates}, revision \texttt{280c963}.
\href{https://github.com/formal-arithmetic-geometry/weil-conjectures-blueprint/tree/280c9634916fe1c920d86c56068cd0c41f2007bf/blueprint}{Source}.

\bibitem[Con]{ConradNagata}
B.~Conrad, \emph{Deligne's notes on Nagata compactifications},
Journal of the Ramanujan Mathematical Society \textbf{22} (2007), 205--257.
\href{https://math.stanford.edu/~conrad/papers/nagatafinal.pdf}{Author's manuscript}.

\bibitem[Milne AV]{MilneAV}
J.~S.~Milne, \emph{Abelian Varieties}, in \emph{Arithmetic Geometry},
1986; corrected online version.
\url{https://www.jmilne.org/math/xnotes/AVs.pdf}.

\bibitem[Milne JV]{MilneJV}
J.~S.~Milne, \emph{Jacobian Varieties}, in \emph{Arithmetic Geometry},
1986; corrected online version.
\url{https://www.jmilne.org/math/xnotes/JVs.pdf}.

\bibitem[Stacks]{Stacks}
The Stacks Project Authors, \emph{The Stacks Project},
\url{https://stacks.math.columbia.edu}.
Individual results are cited by their stable tags in the text.
\end{thebibliography}

\end{document}

